
\documentclass[11pt]{article}
\usepackage[margin=0.82in]{geometry}
\usepackage{amsmath,amssymb,amsthm,mathtools}
\usepackage{booktabs,longtable,array,tabularx,multirow}
\usepackage{enumitem}
\usepackage{xcolor}
\usepackage{url}
\usepackage{microtype}
\usepackage{titlesec}
\usepackage{fancyhdr}
\setlist[itemize]{leftmargin=1.4em,itemsep=0.25em,topsep=0.25em}
\setlist[enumerate]{leftmargin=1.7em,itemsep=0.25em,topsep=0.25em}
\newtheorem{theorem}{Theorem}[section]
\newtheorem{lemma}[theorem]{Lemma}
\newtheorem{proposition}[theorem]{Proposition}
\newtheorem{corollary}[theorem]{Corollary}
\newtheorem{definition}[theorem]{Definition}
\newtheorem{remark}[theorem]{Remark}
\newtheorem{audit}[theorem]{Audit}
\newcommand{\QQ}{\mathbb{Q}}
\newcommand{\CC}{\mathbb{C}}
\newcommand{\CH}{\operatorname{CH}}
\newcommand{\Hdg}{\operatorname{Hdg}}
\newcommand{\clQ}{\operatorname{cl}_{\QQ}}
\newcommand{\FX}{F_X(\alpha)}
\newcommand{\HC}{\operatorname{HC}}
\newcommand{\EndpointUse}{\operatorname{EndpointUse}}
\newcommand{\OfficialCounterexampleForce}{\operatorname{OfficialCounterexampleForce}}
\newcommand{\FixedDomain}{\operatorname{FixedDomain}}
\newcommand{\Aplus}{A^+}
\newcommand{\FibDisc}{D_{\mathrm{fib}}}
\newcommand{\NoFibDisc}{N_{\mathrm{fib}}}
\newcommand{\EmptySep}{\operatorname{EmptyFibSep}}
\newcommand{\ReadAlg}{\operatorname{ReadAlg}}
\newcommand{\StandHdg}{\operatorname{StandHdg}}
\newcommand{\Bridge}{B^{\mathrm{cyc}}_X(\alpha)}
\newcommand{\AASC}{\operatorname{AASC}}
\setlength{\headheight}{14pt}
\pagestyle{fancy}
\fancyhf{}
\lhead{Hodge by Empty Cycle-Fiber Endpoint Exclusion}
\rhead{\thepage}
\titleformat{\section}{\large\bfseries}{\thesection}{0.7em}{}
\titleformat{\subsection}{\normalsize\bfseries}{\thesubsection}{0.6em}{}
\sloppy
\emergencystretch=3em
\title{\textbf{The Hodge Conjecture by Exclusion of the Empty Cycle-Fiber Endpoint}\\[0.25em]
\large An AASC Endpoint-Rigidity Proof for Rational Hodge Classes via Cycle-Fiber Occupation}
\author{Amos Jay Maley}
\date{June 2026\\Full audit-grade endpoint-exclusion rebuild}
\begin{document}
\maketitle
\setcounter{tocdepth}{1}


\begin{abstract}
This manuscript gives an AASC endpoint-exclusion proof of the Hodge Conjecture on the fixed Hodge/cycle carrier. The proof does not construct algebraic cycles by a new algebraic-geometric method and does not redefine Hodge classes, Chow groups, rational coefficients, the cycle-class map, or algebraic cycles. The official endpoint is the classical assertion that every rational Hodge class on a smooth projective complex variety has nonempty rational cycle-class fiber. The proof has one nonclassical dependency: official endpoint-defeating counterexample force is theorem-bearing endpoint use. Once that dependency is granted, the remaining argument is target instantiation plus route exhaustion on the fixed Hodge/cycle carrier.

For a fixed smooth projective complex variety $X$, codimension $p$, and rational Hodge class $\alpha\in \Hdg^p(X)$, the positive endpoint image is the occupation of the cycle-class fiber
\[
F_X(\alpha)=\{Z\in \CH^p(X)_\QQ: \clQ(Z)=\alpha\}.
\]
The positive branch is $F_X(\alpha)\neq\varnothing$. The bare negative branch is $F_X(\alpha)=\varnothing$. Empty-fiber content is lawful as raw mathematical content, as obstruction or nonmembership content, and as ordinary negative theoremhood. It becomes AASC-relevant only when used to defeat the official Hodge endpoint on the same fixed Hodge/cycle carrier. Under such use, empty-fiber counterforce is endpoint use; endpoint use instantiates the kernel of reference, standing, admissibility, and irreversibility; the kernel forces the $A^+$ consequence layer; and $A^+$ excludes independent same-domain endpoint-status discriminators.

The native empty-fiber counterforce route is exhausted into bridge-neutralized status, support-only status, bookkeeping, carrier/domain shift, or the independent cycle-fiber discriminator $D_{\mathrm{fib}}(X,p,\alpha)$. In the live exact-complement reductio subproof, the first four routes are unavailable: bridge-neutralization contradicts the live empty-fiber assumption, support-only and bookkeeping lack endpoint counterforce, and domain shift leaves the fixed Hodge endpoint. Thus the local negative countercase induces $D_{\mathrm{fib}}(X,p,\alpha)$, while endpoint use supplies $N_{\mathrm{fib}}(X,p,\alpha)$ and hence $\neg D_{\mathrm{fib}}(X,p,\alpha)$. The contradiction discharges the local empty-fiber assumption. Since $X,p,\alpha$ were arbitrary, every rational Hodge class has nonempty cycle-class fiber. This is the Hodge Conjecture.
\end{abstract}
\tableofcontents
\newpage

\section{Mathematical-status lock, proof-class lock, and front-loaded dependency lock}



\par\noindent\textbf{Mode declaration.} This manuscript is written in AASC endpoint-exclusion mode. AASC is used as a mathematical constraint formalism: it classifies fixed-domain endpoint use, report standing, admissibility, lawful redescription, and forbidden same-domain discriminator import. It is not an analytic cycle-construction method.

\par\noindent\textbf{Proof-class lock.} The proof class is kernel-forced empty cycle-fiber endpoint-countercase exclusion. It is not a Lefschetz theorem, not a motivic construction, not an Abel-Jacobi or normal-function argument, not a proof from the standard conjectures, not a deformation or spread argument, and not a computation. Classical Hodge-theoretic objects enter to identify the official carrier and endpoint. The branch-exclusion mechanism is AASC-internal.

\par\noindent\textbf{Front-loaded dependency lock.} This manuscript has one nonclassical dependency: official endpoint-defeating counterexample force is theorem-bearing endpoint use. A bare object, formula, theorem, or local assumption with empty-fiber content is not endpoint use merely by being written. But when that content is used to defeat the official Hodge endpoint on the same fixed carrier, the mathematical act fixes target, carrier, branch, reportability, downstream reuse, branch exclusion, act-time finality, and lawful-redescription-stable endpoint role. That act is endpoint use. Once this dependency is granted, the remaining proof is target instantiation plus route exhaustion.

\par\noindent\textbf{Immediate denial fork.} A same-mode objection must therefore do one of the following: deny that official counterexample force is endpoint use; deny that endpoint use instantiates the kernel and the $A^+$ consequence layer; deny that algebraic readout is exactly cycle-fiber occupation; deny native empty-fiber counterforce exhaustion; deny the no-independent-discriminator closure; or produce a fifth endpoint role for same-carrier empty-fiber counterforce.

\section{Read this first: endpoint proof spine}



For fixed $X,p,\alpha$, write
\[
\StandHdg(X,p,\alpha)\Longleftrightarrow \alpha\in\Hdg^p(X),\qquad
\ReadAlg(X,p,\alpha)\Longleftrightarrow F_X(\alpha)\neq\varnothing.
\]
The local proof opens the exact-complement countercase $[F_X(\alpha)=\varnothing]_i$ against the target $F_X(\alpha)\neq\varnothing$. The proof route is
\[
\begin{aligned}
[F_X(\alpha)=\varnothing]_i&\Rightarrow \operatorname{LocalEmptyFiberCountercase}(X,p,\alpha),\\
&\Rightarrow \operatorname{EndpointCounterforce}_{\HC}(X,p,\alpha),\\
&\Rightarrow \OfficialCounterexampleForce_R(F_X(\alpha)=\varnothing),\\
&\Rightarrow \EndpointUse_R(F_X(\alpha)=\varnothing),\\
&\Rightarrow K_R\Rightarrow A^+_R.
\end{aligned}
\]
Under the official Hodge route, empty-fiber endpoint counterforce is native empty-fiber counterforce. The native route is exhausted:
\[
\begin{aligned}
\operatorname{NativeEmptyFiberCounterforce}(X,p,\alpha)\Rightarrow{}&
\operatorname{BridgeNeutralized}
\vee \operatorname{SupportOnly}\vee \operatorname{Bookkeeping}\\
&{}\vee \operatorname{DomainShift}\vee D_{\mathrm{fib}}(X,p,\alpha).
\end{aligned}
\]
In the live subproof, bridge-neutralization contradicts the live empty-fiber assumption; support-only and bookkeeping lack endpoint counterforce; domain shift is not the fixed Hodge endpoint. Hence $D_{\mathrm{fib}}(X,p,\alpha)$. But endpoint use gives $N_{\mathrm{fib}}(X,p,\alpha)$, and by definition
\[
N_{\mathrm{fib}}(X,p,\alpha)\Rightarrow \neg D_{\mathrm{fib}}(X,p,\alpha).
\]
So the exact-complement subproof yields contradiction and discharges the sole object-level assumption $[F_X(\alpha)=\varnothing]_i$. Therefore $F_X(\alpha)\neq\varnothing$. Universal generalization gives the Hodge Conjecture.

\section{Official Hodge endpoint and target instantiation}


\subsection{Classical endpoint retained}


\begin{definition}[Cycle-fiber carrier]
Let $X$ be a smooth projective complex variety and $p$ a codimension. Define $\Hdg^p(X)=H^{2p}(X,\QQ)\cap H^{p,p}(X)$, let $\clQ:\CH^p(X)_\QQ\to H^{2p}(X,\QQ)$ be the rational cycle-class map, and let $F_X(\alpha)=\{Z\in\CH^p(X)_\QQ:\clQ(Z)=\alpha\}$.
\end{definition}


\begin{definition}[Positive and negative branch slots]
The positive local branch is $F_X(\alpha)\neq\varnothing$. The bare local negative branch is $F_X(\alpha)=\varnothing$. The global endpoint is $\forall X,p,\alpha(\alpha\in\Hdg^p(X)\Rightarrow F_X(\alpha)\neq\varnothing)$.
\end{definition}


\begin{theorem}[Official endpoint identity]
The universal nonempty cycle-fiber statement is the classical Hodge endpoint in rational cycle-class form.
\end{theorem}
\begin{proof}
The condition $F_X(\alpha)\neq\varnothing$ is precisely the existence of $Z\in\CH^p(X)_\QQ$ with $\clQ(Z)=\alpha$, equivalently membership of $\alpha$ in the image of the rational cycle-class map. Universalizing over smooth projective complex varieties, codimensions, and rational Hodge classes gives the usual Hodge assertion.
\end{proof}
\par\noindent\textbf{Audit note.} This theorem identifies the endpoint; it does not assert the positive branch.


\begin{definition}[Official Hodge endpoint context]
$C_{\HC}(R,X,p,\alpha)$ is the neutral context object that records the official Hodge carrier and endpoint slots: the smooth projective complex variety $X$, codimension $p$, rational class $\alpha$, the Hodge-standing predicate, the rational Chow group $\CH^p(X)_\QQ$, the rational cycle-class map $\clQ$, the fiber $F_X(\alpha)$, the positive endpoint slot $F_X(\alpha)\neq\varnothing$, the negative slot $F_X(\alpha)=\varnothing$, the empty-fiber separator slot, and the official Hodge route.  It is a carrier-and-slot object, not a truth or endpoint-standing assertion.
\end{definition}


\begin{theorem}[Context non-smuggling]
The Hodge endpoint context fixes the carrier and branch slots but asserts neither nonempty nor empty fiber, neither $A^+$ nor $D_{\mathrm{fib}}$.
\end{theorem}
\begin{proof}
A context object declares what is being evaluated: $X,p,\alpha$, the Hodge-standing predicate, the Chow carrier, the rational cycle-class map, the fiber, and the two endpoint slots. None of these declarations occupies either slot. Occupation or non-occupation enters only through proof roles.
\end{proof}
\par\noindent\textbf{Audit note.} This is the first anti-circularity lock.


\begin{theorem}[Target-instantiation theorem]
Under the official Hodge route, algebraic readout is exactly nonempty cycle-fiber occupation, and its exact complement is empty-fiber non-occupation.
\end{theorem}
\begin{proof}
The official route reads Hodge standing through the rational cycle-class map into the rational Chow fiber. Thus positive endpoint role is $F_X(\alpha)\neq\varnothing$, and the exact complement is $F_X(\alpha)=\varnothing$. This fixes roles, not truth.
\end{proof}
\par\noindent\textbf{Audit note.} This theorem is the Hodge analogue of endpoint-image projection.


\section{Cycle-class carrier and model adequacy}


\begin{lemma}[Smooth projective carrier exactness]
The carrier is the classical smooth projective complex variety carrier.
\end{lemma}
\begin{proof}
This is a model-adequacy clause. It records the official carrier and the corresponding readout role. It supplies no positive endpoint occupant.
\end{proof}
\par\noindent\textbf{Audit note.} If denied, the referee must provide a different official Hodge carrier without changing the problem.


\begin{lemma}[Hodge-standing exactness]
$\StandHdg(X,p,\alpha)$ is exactly rational $(p,p)$ classhood.
\end{lemma}
\begin{proof}
This is a model-adequacy clause. It records the official carrier and the corresponding readout role. It supplies no positive endpoint occupant.
\end{proof}
\par\noindent\textbf{Audit note.} If denied, the referee must provide a different official Hodge carrier without changing the problem.


\begin{lemma}[Chow carrier exactness]
The readout carrier is $\CH^p(X)_\QQ$, not a surrogate unless a bridge resets the target.
\end{lemma}
\begin{proof}
This is a model-adequacy clause. It records the official carrier and the corresponding readout role. It supplies no positive endpoint occupant.
\end{proof}
\par\noindent\textbf{Audit note.} If denied, the referee must provide a different official Hodge carrier without changing the problem.


\begin{lemma}[Cycle-class exactness]
The readout map is the rational cycle-class map $\clQ$.
\end{lemma}
\begin{proof}
This is a model-adequacy clause. It records the official carrier and the corresponding readout role. It supplies no positive endpoint occupant.
\end{proof}
\par\noindent\textbf{Audit note.} If denied, the referee must provide a different official Hodge carrier without changing the problem.


\begin{lemma}[Fiber exactness]
$F_X(\alpha)$ is exactly the set of rational Chow classes mapping to $\alpha$.
\end{lemma}
\begin{proof}
This is a model-adequacy clause. It records the official carrier and the corresponding readout role. It supplies no positive endpoint occupant.
\end{proof}
\par\noindent\textbf{Audit note.} If denied, the referee must provide a different official Hodge carrier without changing the problem.


\begin{lemma}[Readout exactness]
$\ReadAlg(X,p,\alpha)\leftrightarrow F_X(\alpha)\neq\varnothing$.
\end{lemma}
\begin{proof}
This is a model-adequacy clause. It records the official carrier and the corresponding readout role. It supplies no positive endpoint occupant.
\end{proof}
\par\noindent\textbf{Audit note.} If denied, the referee must provide a different official Hodge carrier without changing the problem.


\begin{lemma}[Empty-fiber exactness]
$\neg\ReadAlg(X,p,\alpha)\leftrightarrow F_X(\alpha)=\varnothing$.
\end{lemma}
\begin{proof}
This is a model-adequacy clause. It records the official carrier and the corresponding readout role. It supplies no positive endpoint occupant.
\end{proof}
\par\noindent\textbf{Audit note.} If denied, the referee must provide a different official Hodge carrier without changing the problem.


\begin{lemma}[No algebraicity smuggling]
The carrier does not include nonempty-fiber occupation by definition.
\end{lemma}
\begin{proof}
This is a model-adequacy clause. It records the official carrier and the corresponding readout role. It supplies no positive endpoint occupant.
\end{proof}
\par\noindent\textbf{Audit note.} If denied, the referee must provide a different official Hodge carrier without changing the problem.


\begin{lemma}[Surrogate separation]
Motives, absolute Hodge classes, Hodge loci, Abel-Jacobi data, and known cases are support or bridge programs until they supply the Chow fiber.
\end{lemma}
\begin{proof}
This is a model-adequacy clause. It records the official carrier and the corresponding readout role. It supplies no positive endpoint occupant.
\end{proof}
\par\noindent\textbf{Audit note.} If denied, the referee must provide a different official Hodge carrier without changing the problem.


\begin{theorem}[Model adequacy assembly]
The neutral Hodge context and fixed-domain preservation imply the context-bound Hodge/cycle model adequacy packet.
\end{theorem}
\begin{proof}
The context fixes $X,p,\alpha$, the Hodge-standing predicate, the Chow group, the rational cycle-class map, and the cycle fiber. Fixed-domain preservation prevents silent changes of coefficients, equivalence relation, target class, or readout slot. Thus the adequacy clauses are in force.
\end{proof}
\par\noindent\textbf{Audit note.} This supplies target instantiation, not a cycle construction.


\section{Kernel, A+, ATS, and UEAP replay for Hodge endpoint use}


\begin{theorem}[Target adequacy forces the kernel roles]
Target determinacy, step evaluability, act-time finality, and same-regime fidelity force reference, standing, admissibility, and irreversibility for Hodge endpoint acts.
\end{theorem}
\begin{proof}
Reference fixes the object of evaluation; standing fixes endpoint usability; admissibility fixes the boundary between licensed continuation and carrier shift; irreversibility blocks later repair as the same act. Without these roles, the alleged Hodge endpoint act becomes raw notation or a different target.
\end{proof}
\par\noindent\textbf{Audit note.} This is the local kernel replay.


\begin{theorem}[Endpoint use forces A+]
Endpoint use on the fixed Hodge/cycle carrier triggers the $A^+$ consequence layer.
\end{theorem}
\begin{proof}
Endpoint use is governed construction. The fixed-domain kernel consequences supply AMetric no-selector discipline, no carrier transfer, no repair, report preservation, unique admissible interior, standing conservation, and no independent same-domain endpoint discriminator.
\end{proof}
\par\noindent\textbf{Audit note.} A+ is downstream consequence machinery.


\begin{theorem}[ATS locking]
An endpoint act preserving the Hodge route must preserve $X,p,\alpha$, Hodge standing, Chow carrier, cycle-class map, and the fiber slot.
\end{theorem}
\begin{proof}
ATS separates anchor, tensor, and skin. Presentation may vary, but the anchor and tensor data cannot vary while preserving the same endpoint act. A surrogate readout changes tensor, not skin.
\end{proof}
\par\noindent\textbf{Audit note.} This blocks hidden carrier drift.


\begin{theorem}[UEAP no-overreporting]
A Hodge-standing, surrogate, obstruction, or support report cannot be issued as endpoint readout or endpoint defeat without the corresponding bridge or endpoint-use role.
\end{theorem}
\begin{proof}
UEAP preserves target, carrier, dependencies, blockers, and report level. Lower-status material remains lawful but cannot be reported as higher endpoint status by wording alone.
\end{proof}
\par\noindent\textbf{Audit note.} This protects both positive and negative support content.



\section{Counterexample force as endpoint use}

This section supplies the front-loaded dependency in theorem form.  The distinction is not optional: the manuscript does not classify an object, formula, or theorem as endpoint use merely because it is negative.  It classifies the act of using that content to defeat the official endpoint as endpoint use.  The distinction between content and use is the hinge that prevents the proof from becoming a schema of the form $\neg P\Rightarrow P$.

\begin{definition}[Bare empty-fiber content]
For fixed $X,p,\alpha$, the bare empty-fiber content is the raw mathematical formula
\[
\operatorname{BareEmptyFiberContent}(X,p,\alpha)
:=\bigl(\alpha\in\Hdg^p(X)\wedge F_X(\alpha)=\varnothing\bigr).
\]
It is content at the bare proposition layer.  It is not by itself theoremhood, endpoint use, official counterexample force, reportability, downstream reuse, or branch exclusion.
\end{definition}

\begin{definition}[Ordinary empty-fiber theoremhood]
\[
\operatorname{OrdEmptyFiberThm}_R:=\operatorname{Proof}_R\bigl(\exists X,p,\alpha(\alpha\in\Hdg^p(X)\wedge F_X(\alpha)=\varnothing)\bigr).
\]
This is ordinary negative theoremhood.  It is not the rejected discriminator merely by being negative.
\end{definition}

\begin{definition}[Official empty-fiber counterexample force]
$\operatorname{OfficialEmptyFiberCounterexampleForce}_R(X,p,\alpha)$ means that empty-fiber content for the same fixed Hodge/cycle carrier is used to defeat or settle the official Hodge endpoint negatively.  The use fixes the official target, the variety $X$, codimension $p$, rational class $\alpha$, the Hodge-standing predicate, $\CH^p(X)_\QQ$, the rational cycle-class map, the fiber $F_X(\alpha)$, the negative branch, reportability as endpoint-defeating content, downstream reuse as counterexample material, exclusion of the positive branch, act-time finality, and lawful-redescription stability.
\end{definition}

\begin{definition}[Local empty-fiber counterexample force]
$\operatorname{LocalEmptyFiberCounterexampleForce}_{\HC}(X,p,\alpha)$ means that the local assumption $[F_X(\alpha)=\varnothing]_i$ has been opened as the live exact-complement countercase against the endpoint target $F_X(\alpha)\neq\varnothing$ for the same fixed Hodge-standing instance.  This is local proof-role force, not global negative endpoint resolution.
\end{definition}

\begin{theorem}[Bare content is not endpoint use]
\[
\operatorname{BareEmptyFiberContent}(X,p,\alpha)\not\Rightarrow \EndpointUse_R(F_X(\alpha)=\varnothing).
\]
\end{theorem}
\begin{proof}
The formula $F_X(\alpha)=\varnothing$ may appear as a hypothesis, as a possible obstruction, as the conclusion of a negative theorem, or as a raw branch in a proof by cases.  None of these uses alone fixes the official endpoint role of the negative branch.  Endpoint use requires the additional act of offering that branch as endpoint-defeating or endpoint-resolving content on the fixed carrier.  Bare content therefore does not supply endpoint use.
\end{proof}
\par\noindent\textbf{Audit note.} This theorem protects the lawful mathematical status of empty-fiber syntax.  The proof does not ban the formula $F_X(\alpha)=\varnothing$.

\begin{theorem}[Ordinary negative theoremhood is kernel-internal but not the discriminator]
If $\operatorname{OrdEmptyFiberThm}_R$ is genuine theoremhood rather than raw trace, then it is a governed derivation and hence kernel-internal.  It is not, merely as theoremhood, $D_{\mathrm{fib}}(X,p,\alpha)$.
\end{theorem}
\begin{proof}
A genuine theorem is a standing-bearing derivation.  It therefore already presupposes reference, standing, admissibility, and irreversibility.  The sign of the theorem does not make it defective.  The discriminator at issue later is not the fact that a theorem is negative; it is the endpoint-governing use of empty-fiber non-occupation to defeat the official Hodge endpoint on the same fixed carrier.
\end{proof}
\par\noindent\textbf{Audit note.} Negative algebraic geometry is preserved.  Only official endpoint-defeating use enters the separator route.

\begin{theorem}[Official counterexample force is endpoint use]\label{thm:official-counterexample-force-endpoint-use}
For every fixed endpoint context $R$ and branch $B$,
\[
\OfficialCounterexampleForce_R(B)\Rightarrow\EndpointUse_R(B).
\]
\end{theorem}
\begin{proof}
Official counterexample force is not the bare existence of an object.  It is the mathematical act of using the object, theorem, or branch to defeat the official endpoint.  Such an act fixes the target it defeats; otherwise it is not a counterexample to that endpoint.  It fixes the carrier; otherwise it is a counterexample to a different problem.  It fixes the branch; otherwise there is no determinate negative endpoint role.  It fixes exclusion of the positive branch; otherwise it does not defeat the endpoint.  It fixes reportability and downstream reuse; otherwise the alleged counterexample is not usable as official endpoint-defeating content.  It fixes act-time finality; otherwise later repair or reclassification would be the same act.  Finally, it must be stable under lawful redescription of the same carrier; otherwise it is a carrier shift.  These are precisely the features of theorem-bearing endpoint use.  Therefore official counterexample force is endpoint use.
\end{proof}
\par\noindent\textbf{Audit note.} The theorem classifies the use, not the object.  An object satisfying a negative formula is not an act; offering it as official endpoint defeat is an act.

\begin{corollary}[No autonomous counterexample endpoint role]\label{cor:no-autonomous-counterexample-role}
There is no branch $B$ such that $\OfficialCounterexampleForce_R(B)$ and $\neg\EndpointUse_R(B)$ both hold:
\[
\neg\exists B\,\bigl(\OfficialCounterexampleForce_R(B)\wedge\neg\EndpointUse_R(B)\bigr).
\]
\end{corollary}
\begin{proof}
If such a branch existed, Theorem~\ref{thm:official-counterexample-force-endpoint-use} would give $\EndpointUse_R(B)$, contradicting $\neg\EndpointUse_R(B)$.
\end{proof}
\par\noindent\textbf{Audit note.} This is the no-fifth-role lock at the general counterexample-force level.

\begin{theorem}[Official empty-fiber counterexample force is endpoint use]
\[
\operatorname{OfficialEmptyFiberCounterexampleForce}_R(X,p,\alpha)
\Rightarrow \EndpointUse_R(F_X(\alpha)=\varnothing).
\]
\end{theorem}
\begin{proof}
The premise is the Hodge specialization of official counterexample force: the empty-fiber branch is used to defeat the official endpoint for the same $X,p,\alpha$ and the same rational cycle-class map.  Applying Theorem~\ref{thm:official-counterexample-force-endpoint-use} gives endpoint use of the empty-fiber branch.
\end{proof}
\par\noindent\textbf{Audit note.} This is the Hodge specialization of the front-loaded dependency.

\begin{theorem}[Local exact-complement counterexample force is local endpoint use]
In the local reductio subproof,
\[
\operatorname{ReductioCountercase}_{\HC}(F_X(\alpha)=\varnothing;F_X(\alpha)\neq\varnothing)
\wedge \operatorname{EndpointCounterforce}_{\HC}(X,p,\alpha)
\Rightarrow \operatorname{LocalEndpointUse}_R(F_X(\alpha)=\varnothing).
\]
\end{theorem}
\begin{proof}
The reductio annotation records that the sole opened assumption is the live exact complement of the target.  Endpoint counterforce records that this assumption is being used to block the endpoint closure on the fixed carrier.  Together they give the local version of official counterexample force: the branch is endpoint-facing, carrier-preserving, target-fixing, and branch-excluding inside the subproof.  It is local endpoint use.  The conclusion is local; it does not promote the assumption to global official negative resolution.
\end{proof}
\par\noindent\textbf{Audit note.} Local endpoint use triggers the local $A^+$ consequence layer without treating a local assumption as a global theorem.

\begin{corollary}[Local endpoint use triggers local A+]
\[
\operatorname{LocalEndpointUse}_R(F_X(\alpha)=\varnothing)\wedge \FixedDomain(R)
\Rightarrow A^+_{R,\mathrm{loc}}(F_X(\alpha)=\varnothing).
\]
\end{corollary}
\begin{proof}
Endpoint use is governed construction at the local proof-act level.  Fixed-domain closure supplies the same kernel-forced consequences locally: no hidden selector, no second same-domain endpoint classifier, report preservation, and no carrier-shift repair.  These consequences do not suspend inside a reductio subproof.
\end{proof}
\par\noindent\textbf{Audit note.} The local contradiction later uses this local $A^+$ layer, not a global negative endpoint outcome.

\section{Structural-to-cycle-fiber equivalence bridge}


\begin{theorem}[Hodge standing fixes the cycle-fiber readout slot]
For a fixed Hodge-standing class $\alpha$, the algebraic-readout slot is $F_X(\alpha)$.
\end{theorem}
\begin{proof}
The official Hodge route sends rational Hodge standing through the rational cycle-class map. The target readout is therefore the rational Chow fiber over $\alpha$.
\end{proof}
\par\noindent\textbf{Audit note.} Slot identification is not slot occupation.


\begin{theorem}[Algebraic readout equivalence]
$\ReadAlg(X,p,\alpha)\leftrightarrow F_X(\alpha)\neq\varnothing$.
\end{theorem}
\begin{proof}
Both sides say that some rational Chow class maps to $\alpha$ under $\clQ$.
\end{proof}
\par\noindent\textbf{Audit note.} This is the equivalence bridge.


\begin{theorem}[Empty-fiber equivalence]
$\neg\ReadAlg(X,p,\alpha)\leftrightarrow F_X(\alpha)=\varnothing$.
\end{theorem}
\begin{proof}
Negating the nonempty-fiber condition gives empty fiber. This is the exact local complement used in reductio.
\end{proof}
\par\noindent\textbf{Audit note.} The negative branch is exact complement, not a surrogate.


\begin{theorem}[Projection-role content]
Under the official Hodge route, the exact empty-fiber countercase has cycle-fiber non-occupation as its endpoint-role content.
\end{theorem}
\begin{proof}
The official route fixes the positive endpoint image as fiber occupation. Therefore the live complement is non-occupation of that image. The statement remains non-explosive until endpoint counterforce is added.
\end{proof}
\par\noindent\textbf{Audit note.} This is the bridge from raw negation to endpoint role content.


\section{Single cycle-fiber endpoint interface}


\begin{theorem}[Single cycle-fiber interface]
On the fixed Hodge endpoint image, algebraic readout has exactly one occupation interface: $F_X(\alpha)\neq\varnothing$.
\end{theorem}
\begin{proof}
Algebraic readout is defined by nonempty rational cycle-class fiber. Surrogates may support bridges, but they do not occupy the readout slot without providing an element of that fiber.
\end{proof}
\par\noindent\textbf{Audit note.} This is the single-interface theorem.


\begin{corollary}[Empty fiber is non-occupation, not coequal occupation]
$F_X(\alpha)=\varnothing$ implies non-occupation of the cycle-fiber endpoint image.
\end{corollary}
\begin{proof}
If the fiber is empty, the endpoint image $F_X(\alpha)\neq\varnothing$ is not occupied. This is not a second occupant of the same image.
\end{proof}
\par\noindent\textbf{Audit note.} The broad outcome carrier $\{HC,\neg HC\}$ remains legitimate but is not the endpoint-image carrier.


\begin{theorem}[Empty-fiber non-explosiveness]
$F_X(\alpha)=\varnothing$ does not imply contradiction.
\end{theorem}
\begin{proof}
Empty-fiber content may be raw, descriptive, obstruction-level, support-level, hypothetical, or ordinary theoremhood. The contradiction arises only after endpoint counterforce and separator occupation are in force.
\end{proof}
\par\noindent\textbf{Audit note.} This protects negative algebraic-geometric content.



\section{Native empty-fiber counterforce exhaustion}

This is the Hodge-local technical core.  The theorem to prove is not that empty fiber is impossible as raw content.  It is that empty-fiber content used with endpoint-defeating force on the fixed Hodge/cycle carrier has no autonomous ordinary-negative endpoint role outside the endpoint-role trichotomy.  The possible roles are exhausted into bridge-neutralized status, support-only status, bookkeeping, domain shift, or the independent cycle-fiber discriminator.

\begin{definition}[Native empty-fiber counterforce]
$\operatorname{NativeEmptyFiberCounterforce}(X,p,\alpha)$ means that the empty-fiber content $F_X(\alpha)=\varnothing$ is used as endpoint-defeating force while preserving the same smooth projective variety $X$, codimension $p$, rational class $\alpha$, Hodge-standing predicate, rational Chow group, rational cycle-class map, and cycle-fiber readout slot.  It is native because it uses the official Hodge/cycle carrier; it is counterforce because it is offered to defeat the positive endpoint.
\end{definition}

\begin{definition}[Bridge-neutralized route]
$\operatorname{BridgeNeutralized}_{\mathrm{fib}}(X,p,\alpha)$ means that the apparent empty-fiber counterforce is neutralized by a lawful bridge: a witness $Z\in\CH^p(X)_\QQ$ with $\clQ(Z)=\alpha$, or a certificate of nonempty fiber.  Bridge-neutralization is compatible with positive endpoint readout and incompatible with live endpoint-defeating empty-fiber force as the same local countercase.
\end{definition}

\begin{definition}[Support-only and bookkeeping routes]
$\operatorname{SupportOnly}_{\mathrm{fib}}(X,p,\alpha)$ means the empty-fiber material is hypothetical, obstruction-level, local, or lower-status support and is not used to defeat the official endpoint.  $\operatorname{Bookkeeping}_{\mathrm{fib}}(X,p,\alpha)$ means the empty-fiber label changes no endpoint standing, reportability, downstream reuse, branch exclusion, or admissible continuation.
\end{definition}

\begin{definition}[Domain-shift route]
$\operatorname{DomainShift}_{\mathrm{fib}}(X,p,\alpha)$ means that the negative branch changes at least one admissibility-relevant component of the fixed Hodge endpoint: the variety $X$, codimension $p$, class $\alpha$, coefficient field, Chow group, equivalence relation, cycle-class map, cohomology theory, readout carrier, or surrogate category.  A domain-shift route may be mathematically useful, but it is not the same fixed endpoint.
\end{definition}

\begin{definition}[Fiber discriminator]
$D_{\mathrm{fib}}(X,p,\alpha)$ is the independent same-domain endpoint-status discriminator induced when empty-fiber non-occupation is used to defeat cycle-fiber endpoint occupation on the same fixed carrier while not being bridge-neutralized, support-only, bookkeeping, or domain-shifting.
\end{definition}


\subsection{The no-sixth-role closure for ordinary native empty-fiber falsification}

The previous definitions identify the four harmless or exiting roles and the discriminator role. The hostile objection that remains after those definitions is sharper: perhaps there is an additional role called ordinary native falsification. This subsection proves that there is no such sixth role. The proof is not a dismissal of native algebraic geometry. It is a role-exhaustion theorem for a specific use: same-carrier empty-fiber content with endpoint-defeating force.

\begin{definition}[Endpoint-impact bundle]
For the fixed Hodge endpoint instance $(R,X,p,\alpha)$, let
\[
\mathcal E_{\HC}(X,p,\alpha)
\]
denote the endpoint-impact bundle consisting of the following features: target fixation, carrier fixation, branch fixation, endpoint standing, endpoint reportability, downstream reuse, exclusion of the positive fiber-occupation branch, act-time finality, and lawful-redescription stability. A use of empty-fiber content is endpoint-impacting when changing or removing that use changes at least one component of this bundle.
\end{definition}

\begin{definition}[Autonomous ordinary-negative endpoint role]
An autonomous ordinary-negative endpoint role for $(X,p,\alpha)$ is a use of $F_X(\alpha)=\varnothing$ satisfying all of the following:
\begin{enumerate}
\item it is native to the declared Hodge/cycle vocabulary;
\item it preserves the same $X,p,\alpha$, Chow carrier, rational cycle-class map, and fiber slot;
\item it has endpoint-defeating force against $F_X(\alpha)\neq\varnothing$;
\item it is not bridge-neutralized, support-only, bookkeeping, or domain-shifting;
\item it is asserted not to be $D_{\mathrm{fib}}(X,p,\alpha)$.
\end{enumerate}
This definition isolates the hostile-referee sixth role. It is not assumed to be impossible by definition; it is excluded below.
\end{definition}

\begin{lemma}[Endpoint impact versus bookkeeping]
If a use of $F_X(\alpha)=\varnothing$ changes no component of $\mathcal E_{\HC}(X,p,\alpha)$, then it is bookkeeping for endpoint purposes.
\end{lemma}
\begin{proof}
The endpoint-impact bundle lists exactly the features by which a fixed endpoint act differs from inert description: target, carrier, branch, standing, reportability, reuse, exclusion, act-time status, and lawful-redescription behavior. If a use of empty-fiber content changes none of these, then it does no endpoint work. A role that does no endpoint work may still be useful notation, a reminder, a syntactic label, or local organization, but it is bookkeeping with respect to the endpoint. Therefore any non-bookkeeping empty-fiber role must change at least one component of the bundle.
\end{proof}
\par\noindent\textbf{Audit note.} This lemma prevents a hidden sixth role from doing endpoint work while being described as inert ordinary language.

\begin{lemma}[Endpoint impact plus carrier change is domain shift]
If endpoint-impacting empty-fiber use changes $X$, $p$, $\alpha$, the coefficient field, the Chow group, the rational equivalence relation, the cycle-class map, the fiber slot, or the declared Hodge/cycle readout carrier, then it is $\operatorname{DomainShift}_{\mathrm{fib}}(X,p,\alpha)$.
\end{lemma}
\begin{proof}
The fixed Hodge endpoint is individuated by these data. A change in any one of them changes the object whose fiber is being evaluated or the map by which readout is typed. The resulting use may be a lawful comparison, a bridge program, or a new endpoint. It is not the same endpoint act. Thus endpoint-impacting use that changes the invariant bundle is domain shift.
\end{proof}
\par\noindent\textbf{Audit note.} Motives, absolute Hodge classes, and alternative equivalence relations are protected as separate carriers, not silently identified with the Chow-fiber endpoint.

\begin{lemma}[Endpoint impact without carrier change is endpoint-status work]
Assume a use of $F_X(\alpha)=\varnothing$ is endpoint-impacting and preserves the fixed Hodge/cycle carrier. Then it performs endpoint-status work unless it is support-only or bridge-neutralized.
\end{lemma}
\begin{proof}
Because the carrier is preserved, the use is not merely a neighboring problem. Because it is endpoint-impacting, it changes some member of the endpoint-impact bundle. If the change does not affect branch standing, reportability, reuse, exclusion, or admissible continuation, then by the previous bookkeeping lemma it would be endpoint-inert, contrary to endpoint impact. Hence it affects the status of the endpoint act. If the use remains lower-level obstruction or local evidence, it is support-only and does not have endpoint force. If a lawful bridge supplies nonempty fiber or neutralizes the asserted non-occupation, it is bridge-neutralized. Outside those two cases, the use is endpoint-status work on the same carrier.
\end{proof}
\par\noindent\textbf{Audit note.} This lemma is the role-theoretic bridge from ``ordinary falsification used at endpoint level'' to endpoint-status work.

\begin{lemma}[Endpoint-status work by empty fiber is not positive occupation]
Same-carrier endpoint-status work by $F_X(\alpha)=\varnothing$ is not governance-equivalent to positive cycle-fiber occupation unless bridge-neutralized.
\end{lemma}
\begin{proof}
Positive cycle-fiber occupation is $F_X(\alpha)\neq\varnothing$. Empty-fiber content is $F_X(\alpha)=\varnothing$. On the fixed fiber these are exact complements. If empty-fiber work becomes governance-equivalent to positive occupation, then a bridge has neutralized or overturned the empty-fiber role by supplying a nonempty-fiber certificate. That is the bridge-neutralized branch. If no such bridge is supplied, empty-fiber endpoint-status work excludes rather than occupies the positive endpoint image and is not governance-equivalent to it.
\end{proof}
\par\noindent\textbf{Audit note.} This theorem does not assert nonempty fiber. It states the role difference between occupation and non-occupation on the same fixed fiber.

\begin{lemma}[Endpoint-status work by empty fiber is not lawful coequal occupation]
Same-carrier endpoint-status work by $F_X(\alpha)=\varnothing$ is not lawful coequal occupation of the cycle-fiber endpoint image.
\end{lemma}
\begin{proof}
The cycle-fiber endpoint image has one occupation interface: nonempty fiber. A broad problem-outcome carrier may record two outcomes, Hodge and non-Hodge. That broad outcome carrier is not the endpoint image carrier whose occupant is an element or nonempty certificate of $F_X(\alpha)$. Empty fiber supplies non-occupation of that image. Treating non-occupation as coequal occupation changes the interface from cycle-fiber occupation to a symmetric outcome ledger. That is an interface shift, not a same-image occupation.
\end{proof}
\par\noindent\textbf{Audit note.} This is the Hodge counterpart of the SAT/RH single-interface discipline.

\begin{lemma}[Endpoint-status work not equivalent, not coequal, not shifting is independent]
If empty-fiber endpoint-status work preserves the fixed carrier, is not bridge-neutralized, is not support-only, is not bookkeeping, is not domain-shifting, is not governance-equivalent to positive occupation, and is not lawful coequal occupation of the same endpoint image, then it is $D_{\mathrm{fib}}(X,p,\alpha)$.
\end{lemma}
\begin{proof}
The endpoint-discriminator trichotomy for a fixed endpoint carrier has four possibilities for a status-determining role: bookkeeping, governance-equivalence to the admitted interface, domain shift, or independent same-domain endpoint-status work. The hypotheses exclude bookkeeping and domain shift. The previous two lemmas exclude governance-equivalence to positive occupation and lawful coequal image occupation. The role remains endpoint-status work on the same carrier. Therefore the only remaining trichotomy branch is independent same-domain endpoint-status discrimination. By definition this is $D_{\mathrm{fib}}(X,p,\alpha)$.
\end{proof}
\par\noindent\textbf{Audit note.} This is where ``ordinary native falsification'' is forced to declare its role. If it is endpoint-status work and not one of the harmless/exiting roles, it is the discriminator.

\begin{theorem}[No autonomous ordinary-negative endpoint role]\label{thm:no-autonomous-ordinary-negative-endpoint-role}
For the fixed Hodge/cycle endpoint, there is no autonomous ordinary-negative endpoint role for empty-fiber counterforce. Equivalently, any native same-carrier use of $F_X(\alpha)=\varnothing$ with endpoint-defeating force is bridge-neutralized, support-only, bookkeeping, domain-shifting, or $D_{\mathrm{fib}}(X,p,\alpha)$.
\end{theorem}
\begin{proof}
Let a native same-carrier empty-fiber use have endpoint-defeating force. If a lawful bridge supplies or restores nonempty-fiber occupation, the use is bridge-neutralized. If the use does not actually carry endpoint defeat but only evidential, hypothetical, obstruction, or local analytic content, it is support-only. If it changes no component of the endpoint-impact bundle, it is bookkeeping. If it changes any component of the fixed carrier or readout invariant bundle, it is domain shift. Suppose none of these alternatives holds. Then the use is endpoint-impacting, fixed-carrier, endpoint-status work. It is not governance-equivalent to positive occupation unless bridge-neutralized, already excluded. It is not lawful coequal occupation of the same cycle-fiber endpoint image, because empty fiber is non-occupation of that image. The endpoint-discriminator trichotomy therefore leaves independent same-domain endpoint-status work, namely $D_{\mathrm{fib}}(X,p,\alpha)$. Thus the alleged autonomous ordinary-negative role collapses into one of the five listed roles after all.
\end{proof}
\par\noindent\textbf{Audit note.} This is the main technical closure. A critic must now produce a same-carrier endpoint-defeating empty-fiber role that escapes all five branches; merely naming it ``ordinary falsification'' does not add a sixth role.

\begin{corollary}[Ordinary native falsification is status-disambiguated]
The phrase ``ordinary native empty-fiber falsification'' is not a primitive sixth endpoint role. If it has no endpoint-defeating force, it is raw content, theoremhood, support, or bookkeeping. If it has endpoint-defeating force on the fixed Hodge/cycle carrier, Theorem~\ref{thm:no-autonomous-ordinary-negative-endpoint-role} places it in the five-role exhaustion.
\end{corollary}
\begin{proof}
If the use lacks endpoint-defeating force, it does not settle or oppose the endpoint and therefore falls below endpoint status. If it has endpoint-defeating force, it is a native same-carrier empty-fiber endpoint use and the no-autonomous-role theorem applies. Hence the phrase does not name a sixth endpoint role.
\end{proof}
\par\noindent\textbf{Audit note.} This closes the hostile objection that empty fiber is simply native falsification outside AASC governance.

\begin{lemma}[Bridge-neutralized force is incompatible with the live empty-fiber countercase]\label{lem:bridge-neutralized-incompatible}
Inside the live exact-complement subproof,
\[
F_X(\alpha)=\varnothing\wedge \operatorname{BridgeNeutralized}_{\mathrm{fib}}(X,p,\alpha)\Rightarrow \bot.
\]
\end{lemma}
\begin{proof}
Bridge-neutralization supplies the lawful positive bridge: a nonempty fiber witness or certificate.  The live countercase asserts that the same fixed fiber is empty.  Because the carrier, class, cycle-class map, and fiber slot are fixed, these cannot both be the same local endpoint branch.  If the bridge is admitted, the live empty-fiber counterforce has been neutralized and is no longer the opened exact-complement countercase.
\end{proof}
\par\noindent\textbf{Audit note.} The contradiction is local to the fixed fiber.  It does not say that bridge programs are illicit.

\begin{lemma}[Support-only lacks endpoint counterforce]\label{lem:support-lacks-counterforce}
\[
\operatorname{SupportOnly}_{\mathrm{fib}}(X,p,\alpha)
\Rightarrow \neg\operatorname{EndpointCounterforce}_{\HC}(X,p,\alpha).
\]
\end{lemma}
\begin{proof}
Support-only material may motivate, obstruct, suggest, or locally analyze the endpoint.  By definition it is not used to defeat the official endpoint.  Endpoint counterforce is the role of blocking or displacing the endpoint closure.  The two roles therefore conflict.
\end{proof}
\par\noindent\textbf{Audit note.} Empty-fiber support content is lawful, but it does not close the endpoint negatively.

\begin{lemma}[Bookkeeping lacks endpoint counterforce]\label{lem:bookkeeping-lacks-counterforce}
\[
\operatorname{Bookkeeping}_{\mathrm{fib}}(X,p,\alpha)
\Rightarrow \neg\operatorname{EndpointCounterforce}_{\HC}(X,p,\alpha).
\]
\end{lemma}
\begin{proof}
Bookkeeping affects no endpoint standing, reportability, downstream reuse, branch exclusion, or admissible continuation.  Endpoint counterforce affects the endpoint by opposing the positive branch.  A label that affects none of those features cannot at the same time be endpoint counterforce.
\end{proof}
\par\noindent\textbf{Audit note.} This prevents hidden fifth-role labels from doing endpoint work without admitting it.

\begin{lemma}[Domain shift leaves the fixed endpoint]\label{lem:domain-shift-not-fixed}
\[
\operatorname{DomainShift}_{\mathrm{fib}}(X,p,\alpha)
\Rightarrow \neg\FixedDomain_{\HC}(R,X,p,\alpha).
\]
\end{lemma}
\begin{proof}
The fixed Hodge endpoint is individuated by $X,p,\alpha$, the rational Hodge-standing predicate, $\CH^p(X)_\QQ$, the rational cycle-class map, and the fiber $F_X(\alpha)$.  Changing any of these changes the admissibility-relevant invariant bundle.  Such a route may define a new problem or a surrogate bridge program, but it is not the same endpoint act.
\end{proof}
\par\noindent\textbf{Audit note.} This protects motives, absolute Hodge structures, and other frameworks as bridge programs while denying them silent endpoint identity.

\begin{lemma}[Endpoint-governing empty-fiber work is discriminator work]\label{lem:endpoint-governing-empty-fiber-discriminator}
Suppose empty-fiber content is endpoint-governing on the fixed Hodge/cycle carrier and is not bridge-neutralized, not support-only, not bookkeeping, and not a domain shift.  Then it induces $D_{\mathrm{fib}}(X,p,\alpha)$.
\end{lemma}
\begin{proof}
Endpoint-governing empty-fiber work determines that the positive cycle-fiber endpoint image is not occupied.  Because the carrier is fixed, it is not a different problem.  Because it is not support-only or bookkeeping, it changes endpoint standing, reportability, downstream reuse, or branch exclusion.  Because it is not bridge-neutralized, it does not collapse into positive fiber occupation.  Because it is not domain shift, it remains on the same Hodge/cycle endpoint image.  On that image the unique positive occupation interface is $F_X(\alpha)\neq\varnothing$.  Empty-fiber endpoint force is therefore a same-domain status determination that is not governance-equivalent to positive occupation.  By the endpoint-discriminator trichotomy, the remaining role is independent same-domain endpoint-status discrimination, namely $D_{\mathrm{fib}}(X,p,\alpha)$.
\end{proof}
\par\noindent\textbf{Audit note.} This is the place where ordinary native falsification is collapsed into the endpoint-status trichotomy.  A critic must produce a sixth role to avoid the conclusion.

\begin{theorem}[Native empty-fiber counterforce exhaustion]\label{thm:native-empty-fiber-exhaustion}
For the fixed Hodge/cycle endpoint,
\[
\begin{aligned}
\operatorname{NativeEmptyFiberCounterforce}(X,p,\alpha)\Rightarrow{}&
\operatorname{BridgeNeutralized}_{\mathrm{fib}}
\vee \operatorname{SupportOnly}_{\mathrm{fib}}
\vee \operatorname{Bookkeeping}_{\mathrm{fib}}\\
&{}\vee \operatorname{DomainShift}_{\mathrm{fib}}
\vee D_{\mathrm{fib}}(X,p,\alpha).
\end{aligned}
\]
\end{theorem}
\begin{proof}
Take native empty-fiber counterforce.  If a lawful bridge neutralizes it, the first branch holds.  If it has no endpoint-defeating role and merely supplies obstruction, hypothesis, or local material, it is support-only.  If it changes no endpoint-relevant feature, it is bookkeeping.  If it avoids the fixed fiber by changing the variety, class, coefficient field, equivalence relation, Chow carrier, cycle-class map, or readout carrier, it is domain shift.  If none of these four applies, then the empty-fiber content remains same-domain, endpoint-governing, non-bridge-neutralized, non-support, and non-bookkeeping.  Lemma~\ref{lem:endpoint-governing-empty-fiber-discriminator} gives $D_{\mathrm{fib}}(X,p,\alpha)$.
\end{proof}
\par\noindent\textbf{Audit note.} This is the no-autonomous-native-falsification theorem for Hodge.

\begin{theorem}[Native empty-fiber counterforce collapse]\label{thm:native-empty-fiber-collapse}
In the live exact-complement fixed-domain subproof,
\[
F_X(\alpha)=\varnothing\wedge \operatorname{EndpointCounterforce}_{\HC}(X,p,\alpha)
\wedge \FixedDomain_{\HC}(R,X,p,\alpha)
\Rightarrow D_{\mathrm{fib}}(X,p,\alpha).
\]
\end{theorem}
\begin{proof}
Endpoint counterforce on the fixed Hodge/cycle carrier gives native empty-fiber counterforce.  Apply Theorem~\ref{thm:native-empty-fiber-exhaustion}.  The bridge-neutralized branch is incompatible with the live empty-fiber assumption by Lemma~\ref{lem:bridge-neutralized-incompatible}.  The support-only branch is incompatible with endpoint counterforce by Lemma~\ref{lem:support-lacks-counterforce}.  The bookkeeping branch is incompatible with endpoint counterforce by Lemma~\ref{lem:bookkeeping-lacks-counterforce}.  The domain-shift branch is incompatible with fixed-domain evaluation by Lemma~\ref{lem:domain-shift-not-fixed}.  The only remaining branch is $D_{\mathrm{fib}}(X,p,\alpha)$.
\end{proof}
\par\noindent\textbf{Audit note.} This is the technical closure of the ordinary-negative fifth-role objection.

\begin{corollary}[No ordinary-negative fifth role]
There is no same-carrier endpoint-defeating empty-fiber role that is not bridge-neutralized, support-only, bookkeeping, domain shift, or $D_{\mathrm{fib}}(X,p,\alpha)$.
\end{corollary}
\begin{proof}
Such a role would be native empty-fiber counterforce but would avoid every branch in Theorem~\ref{thm:native-empty-fiber-exhaustion}.  That is impossible.
\end{proof}
\par\noindent\textbf{Audit note.} The phrase ``ordinary native falsification'' names a use.  If the use defeats the endpoint, it is in the exhaustion; if it does not, it is support or bookkeeping.

\subsection*{Main-body closure of the ordinary-native-falsification objection}

The prior corollary is the compact closure. This subsection expands it because this is the precise hostile-referee pressure point. The objector's proposed sixth role is not a motive, a selector, a repair channel, or a known-case transfer. It is the claim that the native formula
\[
F_X(\alpha)=\varnothing
\]
should be allowed to stand as ordinary falsification of the positive endpoint without becoming endpoint-use discipline or independent discriminator work. The proof below separates content from use and then exhausts the use.

\begin{definition}[Native falsification package]
A native falsification package for $(X,p,\alpha)$ is a proposed endpoint role $\mathcal O_{\mathrm{nat}}(X,p,\alpha)$ carrying the same $X$, the same $p$, the same rational class $\alpha$, the same Chow group $\CH^p(X)_\QQ$, the same rational cycle-class map $\clQ$, the same fiber $F_X(\alpha)$, the negative content $F_X(\alpha)=\varnothing$, and claimed endpoint-defeating force against $F_X(\alpha)\neq\varnothing$.
\end{definition}

\begin{lemma}[Native falsification below endpoint use]
A native falsification package without claimed endpoint-defeating force is below endpoint use. It is raw content, ordinary theoremhood, obstruction information, support, or bookkeeping.
\end{lemma}
\begin{proof}
Without endpoint-defeating force, the package does not exclude the positive endpoint branch, does not settle the local endpoint, and does not license downstream reuse as a counterexample. It may remain mathematically meaningful: it may be a formula, a theorem, an obstruction statement, a local hypothesis, or a support datum. But it does not yet have the role of defeating the official Hodge endpoint. Therefore it remains below endpoint use.
\end{proof}

\begin{lemma}[Deployment fixes official counterexample force]
If a native falsification package is deployed to defeat $F_X(\alpha)\neq\varnothing$ on the fixed Hodge/cycle carrier, then it is official empty-fiber counterexample force.
\end{lemma}
\begin{proof}
Deployment to defeat the endpoint means that the package is not merely recorded but used against the official positive branch. Because the package preserves the same $X,p,\alpha$, the same Chow group, the same cycle-class map, and the same fiber, the counterexample force is same-carrier force. Because the negative content is $F_X(\alpha)=\varnothing$, the branch is the exact endpoint complement. Thus the deployment is official empty-fiber counterexample force.
\end{proof}

\begin{lemma}[Deployed native falsification is endpoint use]
A deployed native falsification package is endpoint use.
\end{lemma}
\begin{proof}
By the previous lemma, a deployed native falsification package is official empty-fiber counterexample force. Official counterexample force is endpoint use by Theorem~\ref{thm:official-counterexample-force-endpoint-use}. Therefore the deployed native falsification package is endpoint use.
\end{proof}

\begin{lemma}[Endpoint-use native falsification enters the endpoint-status trichotomy]
A deployed native falsification package that is endpoint use is bookkeeping, governance-equivalent to the positive endpoint interface, independent same-domain endpoint-status work, or a domain shift.
\end{lemma}
\begin{proof}
Endpoint use fixes the carrier and endpoint role. A status factor attached to that use either changes no endpoint-governance verdict, in which case it is bookkeeping; reproduces the already admitted positive interface, in which case it is governance-equivalent; changes endpoint status on the same carrier without being equivalent, in which case it is independent same-domain endpoint-status work; or changes the admissibility-relevant carrier, in which case it is domain shift. These cases exhaust whether the alleged role changes endpoint work and whether it remains on the fixed carrier.
\end{proof}

\begin{lemma}[Deployed empty-fiber falsification is not governance-equivalent to positive occupation]
A deployed native empty-fiber falsification package cannot be governance-equivalent to the positive cycle-fiber endpoint interface while preserving endpoint-defeating negative force.
\end{lemma}
\begin{proof}
The positive endpoint interface is $F_X(\alpha)\neq\varnothing$. The deployed native falsification package asserts and uses $F_X(\alpha)=\varnothing$ to defeat that interface. If it were governance-equivalent to positive occupation, it would not exclude the positive endpoint branch. If it excludes the positive endpoint branch, it performs a different endpoint-status role. Hence, while retaining endpoint-defeating force, it is not governance-equivalent to positive fiber occupation.
\end{proof}

\begin{theorem}[Ordinary native falsification collapse, expanded form]\label{thm:ordinary-native-collapse-expanded}
A native falsification package for $(X,p,\alpha)$ has exactly the following same-mode possibilities: it is raw, support, or theorem-level content below endpoint use; it is bridge-neutralized; it is bookkeeping; it is a domain shift; or it is the independent cycle-fiber endpoint discriminator $D_{\mathrm{fib}}(X,p,\alpha)$. There is no autonomous ordinary-negative endpoint role outside this list.
\end{theorem}
\begin{proof}
If the package is not deployed to defeat the endpoint, it is below endpoint use by the first lemma. If a lawful bridge neutralizes it, it is bridge-neutralized. If it changes no endpoint-relevant feature, it is bookkeeping. If it changes the carrier, class, Chow group, cycle-class map, coefficient field, equivalence relation, or readout slot, it is domain shift. It remains to consider the package deployed to defeat the endpoint while preserving the fixed carrier and while not being bridge-neutralized or bookkeeping. By the deployment lemma and the endpoint-use lemma, it is endpoint use. By the endpoint-status trichotomy, it must be bookkeeping, governance-equivalent, independent same-domain work, or domain shift. Bookkeeping and domain shift have already been removed. Governance-equivalence to the positive endpoint is impossible by the previous lemma. Therefore the remaining case is independent same-domain endpoint-status work. In the fixed Hodge/cycle endpoint image, such work induced by empty-fiber non-occupation is precisely $D_{\mathrm{fib}}(X,p,\alpha)$.
\end{proof}

\begin{corollary}[No sixth role in the live subproof, expanded form]
In the exact-complement reductio subproof, the opened assumption $[F_X(\alpha)=\varnothing]_i$ cannot be deployed as ordinary native endpoint falsification without yielding $D_{\mathrm{fib}}(X,p,\alpha)$.
\end{corollary}
\begin{proof}
The opened assumption is used as the live endpoint countercase against $F_X(\alpha)\neq\varnothing$. Thus it is deployed to defeat the endpoint, not merely recorded as raw syntax. The fixed-domain hypotheses remove domain shift. Endpoint counterforce removes support-only and bookkeeping readings. The live empty-fiber assumption removes bridge-neutralized positive occupation. By the expanded collapse theorem, the remaining role is $D_{\mathrm{fib}}(X,p,\alpha)$.
\end{proof}

\paragraph{Technical closure achieved.} After this subsection, the objection ``empty fiber is just ordinary native falsification'' is not an unaddressed case. It is a challenge to a named theorem: deployment fixes official counterexample force; official counterexample force is endpoint use; endpoint-use native falsification enters the trichotomy; and a deployed empty-fiber branch is not governance-equivalent to positive fiber occupation.


\section{Cycle-fiber separator and discriminator closure}

\begin{definition}[Local empty-fiber separator occupation]
$\EmptySep^{\mathrm{loc}}(X,p,\alpha)$ means that cycle-fiber non-occupation is used inside the exact-complement reductio subproof as the endpoint-defeating negative branch on the fixed Hodge/cycle carrier.
\end{definition}

\begin{definition}[Global empty-fiber separator occupation]
$\EmptySep^{\mathrm{occ}}(X,p,\alpha)$ means that cycle-fiber non-occupation is used as the official negative endpoint resolution on the fixed Hodge/cycle carrier.
\end{definition}

\begin{theorem}[Endpoint counterforce to local separator]\label{thm:endpoint-counterforce-to-local-separator}
\[
\operatorname{ProjectionRoleContent}_{\mathrm{cyc}}(F_X(\alpha)=\varnothing,X,p,\alpha)
\wedge \operatorname{EndpointCounterforce}_{\HC}(X,p,\alpha)
\Rightarrow \EmptySep^{\mathrm{loc}}(X,p,\alpha).
\]
\end{theorem}
\begin{proof}
Projection-role content binds the local negative branch to non-occupation of the cycle-fiber endpoint image.  Endpoint counterforce says this non-occupation is being used to defeat the positive endpoint inside the subproof.  That is precisely local empty-fiber separator occupation.
\end{proof}
\par\noindent\textbf{Audit note.} Raw non-occupation alone is not separator occupation; endpoint counterforce is required.

\begin{theorem}[Local separator induces the fiber discriminator]\label{thm:local-separator-induces-discriminator}
\[
\EmptySep^{\mathrm{loc}}(X,p,\alpha)
\Rightarrow D_{\mathrm{fib}}(X,p,\alpha).
\]
\end{theorem}
\begin{proof}
Local separator occupation is endpoint-defeating empty-fiber use on the fixed carrier.  Therefore the hypotheses of Theorem~\ref{thm:native-empty-fiber-collapse} are met inside the local subproof, and $D_{\mathrm{fib}}(X,p,\alpha)$ follows.
\end{proof}
\par\noindent\textbf{Audit note.} The discriminator is not a descriptive label; it is endpoint-status work.

\begin{theorem}[Endpoint use triggers no-independent fiber closure]\label{thm:endpoint-use-no-independent-fiber}
\[
\operatorname{LocalEndpointUse}_R(F_X(\alpha)=\varnothing)
\wedge \Aplus_{R,\mathrm{loc}}
\Rightarrow N_{\mathrm{fib}}(X,p,\alpha).
\]
\end{theorem}
\begin{proof}
Local endpoint use instantiates the fixed-domain kernel, and the local $A^+$ package includes no independent same-domain endpoint-status discriminator.  The Hodge specialization of that closure is $N_{\mathrm{fib}}(X,p,\alpha)$.
\end{proof}
\par\noindent\textbf{Audit note.} The no-independent closure binds local countercase use as well as global official resolution.

\begin{theorem}[No-independent closure excludes the fiber discriminator]\label{thm:no-independent-excludes-fiber}
\[
N_{\mathrm{fib}}(X,p,\alpha)\Rightarrow \neg D_{\mathrm{fib}}(X,p,\alpha).
\]
\end{theorem}
\begin{proof}
$N_{\mathrm{fib}}$ is defined as the absence of independent same-domain cycle-fiber endpoint-status discrimination.  $D_{\mathrm{fib}}$ is exactly such a discriminator.  Therefore $N_{\mathrm{fib}}$ excludes it.
\end{proof}
\par\noindent\textbf{Audit note.} This is the Hodge specialization of the $A^+$ no-independent-discriminator closure.

\begin{corollary}[Local empty-fiber separator impossible]\label{cor:local-empty-fiber-separator-impossible}
\[
\operatorname{LocalEndpointUse}_R(F_X(\alpha)=\varnothing)\wedge \Aplus_{R,\mathrm{loc}}
\Rightarrow \neg\EmptySep^{\mathrm{loc}}(X,p,\alpha).
\]
\end{corollary}
\begin{proof}
Assume local separator occupation.  Theorem~\ref{thm:local-separator-induces-discriminator} gives $D_{\mathrm{fib}}(X,p,\alpha)$.  Theorem~\ref{thm:endpoint-use-no-independent-fiber} gives $N_{\mathrm{fib}}(X,p,\alpha)$, and Theorem~\ref{thm:no-independent-excludes-fiber} gives $\neg D_{\mathrm{fib}}(X,p,\alpha)$.  Contradiction.
\end{proof}
\par\noindent\textbf{Audit note.} Empty-fiber non-occupation remains lawful as content; only local separator occupation is impossible.


\section{Local reductio and universal Hodge endpoint}

The final endpoint step is local.  It does not assume a global negative endpoint outcome and does not infer endpoint standing from raw truth.  It opens the exact complement of the target for a fixed Hodge-standing instance, routes that local countercase through endpoint-counterforce discipline, obtains the local separator contradiction, and discharges the original assumption.

\begin{definition}[Exact-complement reductio annotation]
$\operatorname{ReductioCountercase}_{\HC}(F_X(\alpha)=\varnothing;F_X(\alpha)\neq\varnothing)$ is the proof-act annotation generated when the sole local assumption $[F_X(\alpha)=\varnothing]_i$ is opened as the live exact complement of the target $F_X(\alpha)\neq\varnothing$ for the fixed Hodge-standing instance.  The annotation is not a second object-level hypothesis.
\end{definition}

\begin{theorem}[Annotation is not a premise]\label{thm:annotation-not-premise}
The exact-complement reductio annotation is a proof-act role, not an additional object-level assumption.
\end{theorem}
\begin{proof}
The only object-level formula opened in the subproof is $F_X(\alpha)=\varnothing$.  The annotation records the role of that assumption: it is the live exact complement of the target.  The later contradiction therefore discharges the original assumption, not a strengthened conjunction.
\end{proof}
\par\noindent\textbf{Audit note.} This blocks the objection that the proof proves only $\neg(F_X(\alpha)=\varnothing\wedge\mathrm{annotation})$.

\begin{theorem}[Exact-complement assumption has local countercase standing]\label{thm:exact-complement-local-countercase}
\[
\operatorname{ReductioCountercase}_{\HC}(F_X(\alpha)=\varnothing;F_X(\alpha)\neq\varnothing)
\Rightarrow \operatorname{LocalEmptyFiberCountercase}(X,p,\alpha).
\]
\end{theorem}
\begin{proof}
In the fixed endpoint proof, the complement of nonempty fiber is empty fiber.  Opening that complement as the live branch whose survival would block the target gives it local countercase standing.  This is temporary derivational standing inside the subproof, not theoremhood of the negative branch.
\end{proof}
\par\noindent\textbf{Audit note.} Inert assumptions remain inert; only the exact-complement endpoint countercase receives this local role.

\begin{theorem}[Local countercase has endpoint counterforce]\label{thm:local-countercase-endpoint-counterforce}
\[
\operatorname{LocalEmptyFiberCountercase}(X,p,\alpha)
\Rightarrow \operatorname{EndpointCounterforce}_{\HC}(X,p,\alpha).
\]
\end{theorem}
\begin{proof}
The local countercase is introduced precisely to defeat the target $F_X(\alpha)\neq\varnothing$ if it survives.  That endpoint-facing role is endpoint counterforce.  No global official negative endpoint outcome is inferred.
\end{proof}
\par\noindent\textbf{Audit note.} This is local proof discipline, not a raw-truth-to-standing principle.

\begin{theorem}[Local empty-fiber countercase contradiction]\label{thm:local-empty-fiber-contradiction}
Under fixed Hodge endpoint evaluation, the local assumption $[F_X(\alpha)=\varnothing]_i$ leads to contradiction.
\end{theorem}
\begin{proof}
Open $[F_X(\alpha)=\varnothing]_i$ as the exact complement of $F_X(\alpha)\neq\varnothing$.  By Theorem~\ref{thm:exact-complement-local-countercase}, it has local empty-fiber countercase standing; by Theorem~\ref{thm:local-countercase-endpoint-counterforce}, it has endpoint counterforce.  The official Hodge route gives projection-role content for cycle-fiber non-occupation.  Theorem~\ref{thm:endpoint-counterforce-to-local-separator} gives local separator occupation.  Theorem~\ref{thm:local-separator-induces-discriminator} gives $D_{\mathrm{fib}}(X,p,\alpha)$.  Local endpoint use and the local $A^+$ package give $N_{\mathrm{fib}}(X,p,\alpha)$ by Theorem~\ref{thm:endpoint-use-no-independent-fiber}; Theorem~\ref{thm:no-independent-excludes-fiber} gives $\neg D_{\mathrm{fib}}(X,p,\alpha)$.  Thus the subproof contains contradiction.
\end{proof}
\par\noindent\textbf{Audit note.} The contradiction is local and discharges the exact-complement assumption; it does not promote the assumption into a global endpoint outcome.

\begin{theorem}[Annotation discharge]\label{thm:annotation-discharge}
If the annotated exact-complement subproof from $[F_X(\alpha)=\varnothing]_i$ yields contradiction, then the discharged result is $F_X(\alpha)\neq\varnothing$.
\end{theorem}
\begin{proof}
By Theorem~\ref{thm:annotation-not-premise}, the annotation is not an additional object-level assumption.  The sole discharged object-level assumption is $F_X(\alpha)=\varnothing$.  Classical reductio therefore yields its negation, i.e. $F_X(\alpha)\neq\varnothing$.
\end{proof}
\par\noindent\textbf{Audit note.} This is the Hodge analogue of the exact-complement annotation-discharge rule.

\begin{theorem}[Arbitrary fixed-instance closeout]\label{thm:fixed-instance-closeout}
For arbitrary fixed $X,p,\alpha$,
\[
\alpha\in\Hdg^p(X)\Rightarrow F_X(\alpha)\neq\varnothing.
\]
\end{theorem}
\begin{proof}
Assume $\alpha\in\Hdg^p(X)$ and work in the official Hodge endpoint context for $(X,p,\alpha)$.  Open the exact complement $[F_X(\alpha)=\varnothing]_i$.  Theorem~\ref{thm:local-empty-fiber-contradiction} gives contradiction in the annotated subproof.  Theorem~\ref{thm:annotation-discharge} discharges the original assumption and yields $F_X(\alpha)\neq\varnothing$.
\end{proof}
\par\noindent\textbf{Audit note.} The proof is by local endpoint-countercase exclusion, not by analytic cycle construction.

\begin{theorem}[Universal Hodge endpoint]\label{thm:universal-hodge-endpoint}
\[
\forall X,p,\alpha\,\bigl(\alpha\in\Hdg^p(X)\Rightarrow F_X(\alpha)\neq\varnothing\bigr).
\]
\end{theorem}
\begin{proof}
The fixed-instance closeout was proved for arbitrary smooth projective complex $X$, codimension $p$, and rational class $\alpha$ in Hodge standing.  Universal generalization gives the displayed theorem, which is the Hodge Conjecture in rational cycle-fiber form.
\end{proof}
\par\noindent\textbf{Audit note.} This is the official endpoint.  No context clause asserted it upstream.

\section{Ordinary practice, route tests, and incompleteness firewall}


\subsection{Bridge test: Lefschetz (1,1)}


\par\noindent\textbf{Lawful role.} local bridge completion in codimension one.\par\noindent\textbf{Failure if overreported.} universal transfer beyond scope.\par\noindent\textbf{AASC disposition.} The route is preserved as support or bridge completion when it supplies the cycle-fiber. It is blocked only when used to occupy or defeat the fixed endpoint without the declared bridge.


\subsection{Bridge test: Hard Lefschetz}


\par\noindent\textbf{Lawful role.} cohomological tensor structure.\par\noindent\textbf{Failure if overreported.} cohomological structure reported as Chow fiber.\par\noindent\textbf{AASC disposition.} The route is preserved as support or bridge completion when it supplies the cycle-fiber. It is blocked only when used to occupy or defeat the fixed endpoint without the declared bridge.


\subsection{Bridge test: Hodge loci}


\par\noindent\textbf{Lawful role.} standing persistence in families.\par\noindent\textbf{Failure if overreported.} standing persistence reported as cycle witness.\par\noindent\textbf{AASC disposition.} The route is preserved as support or bridge completion when it supplies the cycle-fiber. It is blocked only when used to occupy or defeat the fixed endpoint without the declared bridge.


\subsection{Bridge test: Intermediate Jacobians}


\par\noindent\textbf{Lawful role.} obstruction/refinement carrier.\par\noindent\textbf{Failure if overreported.} obstruction data reported as positive readout.\par\noindent\textbf{AASC disposition.} The route is preserved as support or bridge completion when it supplies the cycle-fiber. It is blocked only when used to occupy or defeat the fixed endpoint without the declared bridge.


\subsection{Bridge test: Motives}


\par\noindent\textbf{Lawful role.} surrogate/enriched carrier.\par\noindent\textbf{Failure if overreported.} carrier transfer without Chow bridge.\par\noindent\textbf{AASC disposition.} The route is preserved as support or bridge completion when it supplies the cycle-fiber. It is blocked only when used to occupy or defeat the fixed endpoint without the declared bridge.


\subsection{Bridge test: Absolute Hodge classes}


\par\noindent\textbf{Lawful role.} stronger invariance.\par\noindent\textbf{Failure if overreported.} stronger standing reported as algebraic readout.\par\noindent\textbf{AASC disposition.} The route is preserved as support or bridge completion when it supplies the cycle-fiber. It is blocked only when used to occupy or defeat the fixed endpoint without the declared bridge.


\subsection{Bridge test: Known special cases}


\par\noindent\textbf{Lawful role.} local bridge completion.\par\noindent\textbf{Failure if overreported.} universal transfer without target gating.\par\noindent\textbf{AASC disposition.} The route is preserved as support or bridge completion when it supplies the cycle-fiber. It is blocked only when used to occupy or defeat the fixed endpoint without the declared bridge.


\subsection{Bridge test: Standard conjectures}


\par\noindent\textbf{Lawful role.} bridge program.\par\noindent\textbf{Failure if overreported.} programmatic expectation reported as completed bridge.\par\noindent\textbf{AASC disposition.} The route is preserved as support or bridge completion when it supplies the cycle-fiber. It is blocked only when used to occupy or defeat the fixed endpoint without the declared bridge.


\subsection{Bridge test: Deformation and spread}


\par\noindent\textbf{Lawful role.} family route.\par\noindent\textbf{Failure if overreported.} local data treated as global fiber without compatibility.\par\noindent\textbf{AASC disposition.} The route is preserved as support or bridge completion when it supplies the cycle-fiber. It is blocked only when used to occupy or defeat the fixed endpoint without the declared bridge.


\begin{theorem}[Same-domain incompleteness firewall]
Independence, undecidability, stronger-theory, or model-theoretic distinctions affect the endpoint only if they change endpoint standing, reportability, downstream reuse, branch exclusion, or admissible continuation.
\end{theorem}
\begin{proof}
If no endpoint feature changes, the distinction is endpoint-inert. If an endpoint feature changes, it is gate work and falls under A+ endpoint-discriminator discipline.
\end{proof}
\par\noindent\textbf{Audit note.} This does not deny incompleteness phenomena; it denies a third endpoint role.


\section{Exhaustion, weakening resistance, and denial burden}



\begin{longtable}{p{0.22\textwidth}p{0.38\textwidth}p{0.28\textwidth}}
\toprule
Locus & Hodge form & Disposition \\
\midrule
\endhead
Hodge-standing promotion & $\alpha\in\Hdg^p(X)$ reported as $F_X(\alpha)\neq\varnothing$ & UEAP/no-generator \\
Fiber declaration & Nonempty fiber asserted without bridge & report preservation \\
Empty-fiber counterforce & Empty fiber used as endpoint defeat & native-counterforce collapse \\
Motive surrogate & Motivic object replaces Chow fiber & no carrier/silent redescription \\
Absolute-Hodge surrogate & Stronger invariance replaces cycle witness & surrogate separation \\
Known-case transfer & Local theorem used universally & no-carrier \\
Local-to-global assembly & Local cycles treated as global fiber & non-factorizability \\
Selector/canonical representative & Preferred cycle chosen before fiber fixed & AMetric boundary \\
Repair & Later cycle repairs earlier report & irreversibility \\
Hidden fifth negative role & Endpoint-defeating empty fiber not endpoint use & no autonomous counterexample endpoint role \\
\bottomrule
\end{longtable}

\begin{theorem}[Weakening resistance]
Any strict same-carrier weakening that permits endpoint-defeating empty-fiber counterforce either lowers it to support/bookkeeping, changes carrier, or reintroduces $D_{\mathrm{fib}}$.
\end{theorem}
\begin{proof}
A weakening that changes no endpoint verdict is bookkeeping. If it has endpoint force on the same carrier, it is endpoint-status work. If it avoids this by changing carrier or readout, it is domain shift.
\end{proof}
\par\noindent\textbf{Audit note.} This prevents a weakened packet from reinstating the negative branch.


\subsection{Right-mode objection table}



\begin{longtable}{p{0.25\textwidth}p{0.63\textwidth}}
\toprule
Target & Successful objection must produce\\
\midrule
\endhead
Context non-smuggling & A clause in $C_{HC}$ that asserts nonempty fiber, empty fiber, $A^+$, or $D_{\mathrm{fib}}$.\\
Model adequacy & A mismatch between $\ReadAlg(X,p,\alpha)$ and $F_X(\alpha)\neq\varnothing$, or a different official Hodge carrier.\\
Counterexample force & A same-carrier official counterexample act that defeats the endpoint while not fixing endpoint-use features.\\
Cycle-fiber projection & A Hodge endpoint route in which algebraic readout is not cycle-fiber occupation.\\
Native empty-fiber exhaustion & A fifth role for endpoint-defeating empty-fiber counterforce.\\
A+ closure & An independent same-domain cycle-fiber endpoint discriminator compatible with no-independent-discriminator closure.\\
Local reductio & A proof that the exact-complement annotation is an added object-level premise or discharges the wrong assumption.\\
Universalization & A failure of arbitrary fixed-instance closure to yield the universal Hodge statement.\\
\bottomrule
\end{longtable}

\section{Appendix A: theorem ladder and dependency audit}


\begin{longtable}{p{0.28\textwidth}p{0.58\textwidth}}\toprule Ladder item & Function \\\midrule\endhead

Mathematical-status lock & Proof-order component in the endpoint closeout. \\

Front-loaded dependency lock & Proof-order component in the endpoint closeout. \\

Context non-smuggling & Proof-order component in the endpoint closeout. \\

Model adequacy & Proof-order component in the endpoint closeout. \\

Kernel forcing & Proof-order component in the endpoint closeout. \\

Endpoint use forces A+ & Proof-order component in the endpoint closeout. \\

Counterexample force is endpoint use & Proof-order component in the endpoint closeout. \\

No autonomous counterexample role & Proof-order component in the endpoint closeout. \\

Cycle-fiber equivalence bridge & Proof-order component in the endpoint closeout. \\

Single cycle-fiber interface & Proof-order component in the endpoint closeout. \\

Native empty-fiber exhaustion & Proof-order component in the endpoint closeout. \\

Separator induction & Proof-order component in the endpoint closeout. \\

A+ exclusion & Proof-order component in the endpoint closeout. \\

Local reductio discharge & Proof-order component in the endpoint closeout. \\

Universal endpoint & Proof-order component in the endpoint closeout. \\

\bottomrule\end{longtable}


\section{Appendix B: anti-circularity and hostile-misreading audit}


\par\noindent\textbf{Context smuggles nonempty fiber.} Context non-smuggling theorem says carrier and slots only.


\par\noindent\textbf{Empty fiber is illegal.} Non-explosiveness preserves empty-fiber content.


\par\noindent\textbf{Ordinary theoremhood is banned.} Ordinary negative theoremhood is protected.


\par\noindent\textbf{Counterexample object treated as act.} Object/use distinction is explicit.


\par\noindent\textbf{Arbitrary negation becomes separator.} Official route, projection content, and endpoint counterforce are required.


\par\noindent\textbf{Surrogates are dismissed.} They remain bridge programs or support.


\par\noindent\textbf{A+ does algebraic geometry.} A+ excludes endpoint discriminator work; it does not construct cycles.


\par\noindent\textbf{Endpoint changed.} The final theorem is universal nonempty rational cycle-class fiber.


\section{Appendix C: route-locus exhaustion capsules}


\subsection{Status promotion}


This capsule records a route-locus rather than an example. Any unbridged passage from Hodge standing to endpoint readout must alter status, codomain, existence, carrier, time, locality, selector authority, or endpoint classifier. The corresponding AASC closure blocks that alteration unless a bridge is supplied or an explicit domain reset is declared.


\subsection{Codomain substitution}


This capsule records a route-locus rather than an example. Any unbridged passage from Hodge standing to endpoint readout must alter status, codomain, existence, carrier, time, locality, selector authority, or endpoint classifier. The corresponding AASC closure blocks that alteration unless a bridge is supplied or an explicit domain reset is declared.


\subsection{Existence creation}


This capsule records a route-locus rather than an example. Any unbridged passage from Hodge standing to endpoint readout must alter status, codomain, existence, carrier, time, locality, selector authority, or endpoint classifier. The corresponding AASC closure blocks that alteration unless a bridge is supplied or an explicit domain reset is declared.


\subsection{Carrier transfer}


This capsule records a route-locus rather than an example. Any unbridged passage from Hodge standing to endpoint readout must alter status, codomain, existence, carrier, time, locality, selector authority, or endpoint classifier. The corresponding AASC closure blocks that alteration unless a bridge is supplied or an explicit domain reset is declared.


\subsection{Temporal repair}


This capsule records a route-locus rather than an example. Any unbridged passage from Hodge standing to endpoint readout must alter status, codomain, existence, carrier, time, locality, selector authority, or endpoint classifier. The corresponding AASC closure blocks that alteration unless a bridge is supplied or an explicit domain reset is declared.


\subsection{Local assembly}


This capsule records a route-locus rather than an example. Any unbridged passage from Hodge standing to endpoint readout must alter status, codomain, existence, carrier, time, locality, selector authority, or endpoint classifier. The corresponding AASC closure blocks that alteration unless a bridge is supplied or an explicit domain reset is declared.


\subsection{Selector authority}


This capsule records a route-locus rather than an example. Any unbridged passage from Hodge standing to endpoint readout must alter status, codomain, existence, carrier, time, locality, selector authority, or endpoint classifier. The corresponding AASC closure blocks that alteration unless a bridge is supplied or an explicit domain reset is declared.


\subsection{Second classifier}


This capsule records a route-locus rather than an example. Any unbridged passage from Hodge standing to endpoint readout must alter status, codomain, existence, carrier, time, locality, selector authority, or endpoint classifier. The corresponding AASC closure blocks that alteration unless a bridge is supplied or an explicit domain reset is declared.


\subsection{Hidden fifth role}


This capsule records a route-locus rather than an example. Any unbridged passage from Hodge standing to endpoint readout must alter status, codomain, existence, carrier, time, locality, selector authority, or endpoint classifier. The corresponding AASC closure blocks that alteration unless a bridge is supplied or an explicit domain reset is declared.


\section{Appendix D: notation ledger}



\begin{longtable}{p{0.25\textwidth}p{0.65\textwidth}}
\toprule Symbol & Meaning\\ \midrule\endhead
$X$ & smooth projective complex variety\\
$p$ & codimension\\
$\alpha$ & rational cohomology class\\
$\Hdg^p(X)$ & rational Hodge classes $H^{2p}(X,\QQ)\cap H^{p,p}(X)$\\
$\CH^p(X)_\QQ$ & rational Chow group\\
$\clQ$ & rational cycle-class map\\
$F_X(\alpha)$ & cycle-class fiber\\
$\ReadAlg$ & algebraic-readout predicate\\
$D_{\mathrm{fib}}$ & independent cycle-fiber endpoint-status discriminator\\
$N_{\mathrm{fib}}$ & no-independent cycle-fiber closure\\
$\Bridge$ & cycle-realization bridge object\\
\bottomrule
\end{longtable}

\section{Appendix E: corpus support ledger}



\begin{longtable}{p{0.15\textwidth}p{0.22\textwidth}p{0.37\textwidth}p{0.18\textwidth}}
\toprule Local role & Source ID & Paper & Object\\ \midrule\endhead
Kernel & ERK20260527-New\_erk\_files\_May\_C3E3B7-LEM-001 & A Classification Theorem for Admissible Standard-Model Structures Unique & Envelope necessity \\ 
Kernel & ERK20260527-New\_erk\_files\_May\_C3E3B7-THM-004 & A Classification Theorem for Admissible Standard-Model Structures Unique & Kernel necessity as a derived structural constraint \\ 
Kernel & ERK20260527-New\_erk\_files\_May\_C3E3B7-COR-002 & A Classification Theorem for Admissible Standard-Model Structures Unique & Kernel non-optionality \\ 
UEAP & ERK20260527-New\_erk\_files\_May\_AF0C6B-DEF-001 & Claim Standing and Legitimacy A Universal Audit Protocol for Laundering  & ixes an audit object or status rule \\ 
UEAP & ERK20260527-New\_erk\_files\_May\_AF0C6B-PRI-001 & Claim Standing and Legitimacy A Universal Audit Protocol for Laundering  & Counterexample burden \\ 
UEAP & ERK20260527-New\_erk\_files\_May\_AF0C6B-REM-002 & Claim Standing and Legitimacy A Universal Audit Protocol for Laundering  & Why this is the central spine \\ 
Claim Standing & ERK20260527-New\_erk\_files\_May\_2490B6-PRO-012 & Admissibility and the Common Structural Origin of Gauge, Constraint, and & Template character of the fixed-domain closure argument \\ 
Claim Standing & ERK20260527-New\_erk\_files\_May\_2490B6-THM-019 & Admissibility and the Common Structural Origin of Gauge, Constraint, and & Collapse classification under AER-3 \\ 
Claim Standing & ERK20260527-New\_erk\_files\_May\_2490B6-PRO-016 & Admissibility and the Common Structural Origin of Gauge, Constraint, and & Conditional summary for unitary multiplicity under AER \\ 
ATS & ERK20260527-New\_erk\_files\_May\_4C9E73-REM-006 & A Fixed-Domain Exhaustion Theorem Standing, Admissibility, and the Deter & Undefinedness is not a third truth value \\ 
ATS & ERK20260527-New\_erk\_files\_May\_4C9E73-REM-009 & A Fixed-Domain Exhaustion Theorem Standing, Admissibility, and the Deter & The axioms below are the explicit premises extracted from th \\ 
ATS & ERK20260527-New\_erk\_files\_May\_4C9E73-REM-010 & A Fixed-Domain Exhaustion Theorem Standing, Admissibility, and the Deter & Detectability is derived, not assumed \\ 
AMetric & ERK20260527-New\_erk\_files\_May\_C3E3B7-DEF-017 & A Classification Theorem for Admissible Standard-Model Structures Unique & AMetric Boundary for admissible bookkeeping \\ 
AMetric & ERK20260527-New\_erk\_files\_May\_C3E3B7-THM-011 & A Classification Theorem for Admissible Standard-Model Structures Unique & AMetric Boundary forcing \\ 
AMetric & ERK20260527-New\_erk\_files\_May\_C3E3B7-DEF-020 & A Classification Theorem for Admissible Standard-Model Structures Unique & Admissibly distinct construction trajectory \\ 
No Post-Selection & ERK20260527-New\_erk\_files\_May\_311F0C-THM-009 & Boundary-Trace Fixation of Continuation Loci under AASC A Role-Anchor Th & No post-selection repair \\ 
No Post-Selection & ERK20260527-New\_erk\_files\_May\_045695-COR-006 & Boundary-Trace Fixation under Admissibility Forced Environmental Couplin & No post-selection environmental repair \\ 
No Post-Selection & ERK20260527-New\_erk\_files\_May\_D2FE05-COR-002 & Constraint Architectures for Coherent Evaluability Quotient Interfaces,  & No post-selection boundary repair \\ 
Silent Redescription & ERK20260527-New\_erk\_files\_May\_FBEDF2-REM-005 & Admissible Construction and the Non-Instantiability of Gödel Architectur & Seminar and blackboard practice \\ 
Silent Redescription & ERK20260527-New\_erk\_files\_May\_35B594-COR-007 & Non-Degenerate Construction and the Kernel of Admissibility An Exhaustio & No silent redescription \\ 
Silent Redescription & ERK20260527-New\_erk\_files\_May\_9ECDFF-LEM-011 & Rigidity of the Black-Scholes Valuation Operator Representative-Invarian & No silent redescription \\ 
Role-Occupancy & ERK20260527-New\_erk\_files\_May\_DCCC53-REM-001 & Role-Occupancy Closure under Admissibility Why role-cardinality matching & The answer is not chosen \\ 
Role-Occupancy & ERK20260527-New\_erk\_files\_May\_DCCC53-DEF-001 & Role-Occupancy Closure under Admissibility Why role-cardinality matching & Fixed certificate \\ 
Role-Occupancy & ERK20260527-New\_erk\_files\_May\_DCCC53-DEF-002 & Role-Occupancy Closure under Admissibility Why role-cardinality matching & Candidate response class \\ 
No generators & ERK20260527-New\_erk\_files\_May\_FBEDF2-THM-049 & Admissible Construction and the Non-Instantiability of Gödel Architectur & No generators \\ 
No generators & ERK20260527-New\_erk\_files\_May\_3D011A-THM-043 & Forty-Two Exhaustion, Closure, and Downstream Consequence & No generators \\ 
No generators & ERK20260527-New\_erk\_files\_May\_35B594-COR-011 & Non-Degenerate Construction and the Kernel of Admissibility An Exhaustio & No generators \\ 
No carriers & ERK20260527-New\_erk\_files\_May\_41605D-THM-029 & The AMetric Boundary as a Unique Regime-Invariant Constraint on Admissib & Admissibility constraint on No carriers (numeric or structur \\ 
No carriers & ERK20260527-New\_erk\_files\_May\_8B37FD-COR-004 & The Bivalence Theorem for Non-Degenerate Construction Standing Conservat & No Carriers \\ 
No carriers & ERK20260527-New\_erk\_files\_May\_EA512E-THM-003 & Unus Solus Possibilis Est Corpus Closure and Uniqueness of the Admissibl & No carriers (numeric or structural \\ 

\end{longtable}

\section{Appendix F: Lean support boundary}



The first Lean audit target should formalize the local endpoint-role bridge, not all Hodge theory from foundations.  The audit surface should include the following theorem names:
\begin{verbatim}
officialCounterexampleForce_endpointUse
noAutonomousCounterexampleEndpointRole
algReadout_iff_cycleFiberNonempty
emptyFiber_iff_notAlgReadout
projectionRoleContent_of_emptyFiberCountercase
emptyFiberCounterforce_exhaustion
emptyFiberCounterforce_collapse
emptyFiberSeparator_induces_fiberDiscriminator
endpointUse_forces_noIndependentFiberDiscriminator
localEmptyFiberSeparator_impossible
cycleFiberNonempty_of_hodgeStanding
hodgeConjecture_endpoint
\end{verbatim}
The Lean layer should state clearly that it audits the AASC/Hodge endpoint bridge and closeout surface, conditional on declared Hodge/cycle model adequacy.

\section{Appendix G: full hostile-referee audit expansion}


This appendix repeats the load-bearing theorem set in multiple audit modes. The repetition is intentional. A Clay-facing endpoint manuscript must make every branch of the proof visible in theorem form, audit form, and denial-burden form.


\subsection{Audit packet 1: Context non-smuggling (dependency replay)}


\par\noindent\textbf{Claim under audit.} Hodge context fixes carrier and slots without branch occupation.


\par\noindent\textbf{Dependency path.} This packet belongs to the dependency replay. It depends on the official Hodge route, the fixed cycle-fiber endpoint image, the object/use distinction, and the kernel/A+ endpoint-use machinery. It is never used before the carrier and branch slots are fixed.


\par\noindent\textbf{Hostile-referee risk.} A hostile reader may attempt to treat the relevant content as raw mathematical syntax, ordinary negative theoremhood, surrogate support, or a classical counterexample object rather than endpoint-use material.


\par\noindent\textbf{Closure and denial burden.} The closure response is role-sensitive. If the content is raw, support, bookkeeping, or outside the fixed carrier, it has no endpoint-defeating force. If it keeps endpoint-defeating force on the same carrier, it enters endpoint-use discipline and route exhaustion. A successful objection must preserve the same $X,p,\alpha$, the same Hodge-standing predicate, the same Chow carrier $\CH^p(X)_\QQ$, and the same rational cycle-class map $\clQ$.


\par\noindent\textbf{Use in the proof spine.} The packet is reused only after the local exact-complement assumption has been opened. It marks a role transition, not a hidden assertion of nonempty fiber.


\subsection{Audit packet 2: Model adequacy (dependency replay)}


\par\noindent\textbf{Claim under audit.} Readout and counterreadout are identified with nonempty and empty cycle fiber.


\par\noindent\textbf{Dependency path.} This packet belongs to the dependency replay. It depends on the official Hodge route, the fixed cycle-fiber endpoint image, the object/use distinction, and the kernel/A+ endpoint-use machinery. It is never used before the carrier and branch slots are fixed.


\par\noindent\textbf{Hostile-referee risk.} A hostile reader may attempt to treat the relevant content as raw mathematical syntax, ordinary negative theoremhood, surrogate support, or a classical counterexample object rather than endpoint-use material.


\par\noindent\textbf{Closure and denial burden.} The closure response is role-sensitive. If the content is raw, support, bookkeeping, or outside the fixed carrier, it has no endpoint-defeating force. If it keeps endpoint-defeating force on the same carrier, it enters endpoint-use discipline and route exhaustion. A successful objection must preserve the same $X,p,\alpha$, the same Hodge-standing predicate, the same Chow carrier $\CH^p(X)_\QQ$, and the same rational cycle-class map $\clQ$.


\par\noindent\textbf{Use in the proof spine.} The packet is reused only after the local exact-complement assumption has been opened. It marks a role transition, not a hidden assertion of nonempty fiber.


\subsection{Audit packet 3: Counterexample force as endpoint use (dependency replay)}


\par\noindent\textbf{Claim under audit.} Official endpoint defeat is endpoint use.


\par\noindent\textbf{Dependency path.} This packet belongs to the dependency replay. It depends on the official Hodge route, the fixed cycle-fiber endpoint image, the object/use distinction, and the kernel/A+ endpoint-use machinery. It is never used before the carrier and branch slots are fixed.


\par\noindent\textbf{Hostile-referee risk.} A hostile reader may attempt to treat the relevant content as raw mathematical syntax, ordinary negative theoremhood, surrogate support, or a classical counterexample object rather than endpoint-use material.


\par\noindent\textbf{Closure and denial burden.} The closure response is role-sensitive. If the content is raw, support, bookkeeping, or outside the fixed carrier, it has no endpoint-defeating force. If it keeps endpoint-defeating force on the same carrier, it enters endpoint-use discipline and route exhaustion. A successful objection must preserve the same $X,p,\alpha$, the same Hodge-standing predicate, the same Chow carrier $\CH^p(X)_\QQ$, and the same rational cycle-class map $\clQ$.


\par\noindent\textbf{Use in the proof spine.} The packet is reused only after the local exact-complement assumption has been opened. It marks a role transition, not a hidden assertion of nonempty fiber.


\subsection{Audit packet 4: No autonomous counterexample role (dependency replay)}


\par\noindent\textbf{Claim under audit.} There is no endpoint-defeating branch outside endpoint use.


\par\noindent\textbf{Dependency path.} This packet belongs to the dependency replay. It depends on the official Hodge route, the fixed cycle-fiber endpoint image, the object/use distinction, and the kernel/A+ endpoint-use machinery. It is never used before the carrier and branch slots are fixed.


\par\noindent\textbf{Hostile-referee risk.} A hostile reader may attempt to treat the relevant content as raw mathematical syntax, ordinary negative theoremhood, surrogate support, or a classical counterexample object rather than endpoint-use material.


\par\noindent\textbf{Closure and denial burden.} The closure response is role-sensitive. If the content is raw, support, bookkeeping, or outside the fixed carrier, it has no endpoint-defeating force. If it keeps endpoint-defeating force on the same carrier, it enters endpoint-use discipline and route exhaustion. A successful objection must preserve the same $X,p,\alpha$, the same Hodge-standing predicate, the same Chow carrier $\CH^p(X)_\QQ$, and the same rational cycle-class map $\clQ$.


\par\noindent\textbf{Use in the proof spine.} The packet is reused only after the local exact-complement assumption has been opened. It marks a role transition, not a hidden assertion of nonempty fiber.


\subsection{Audit packet 5: ATS locking (dependency replay)}


\par\noindent\textbf{Claim under audit.} The same endpoint act preserves anchor and tensor data.


\par\noindent\textbf{Dependency path.} This packet belongs to the dependency replay. It depends on the official Hodge route, the fixed cycle-fiber endpoint image, the object/use distinction, and the kernel/A+ endpoint-use machinery. It is never used before the carrier and branch slots are fixed.


\par\noindent\textbf{Hostile-referee risk.} A hostile reader may attempt to treat the relevant content as raw mathematical syntax, ordinary negative theoremhood, surrogate support, or a classical counterexample object rather than endpoint-use material.


\par\noindent\textbf{Closure and denial burden.} The closure response is role-sensitive. If the content is raw, support, bookkeeping, or outside the fixed carrier, it has no endpoint-defeating force. If it keeps endpoint-defeating force on the same carrier, it enters endpoint-use discipline and route exhaustion. A successful objection must preserve the same $X,p,\alpha$, the same Hodge-standing predicate, the same Chow carrier $\CH^p(X)_\QQ$, and the same rational cycle-class map $\clQ$.


\par\noindent\textbf{Use in the proof spine.} The packet is reused only after the local exact-complement assumption has been opened. It marks a role transition, not a hidden assertion of nonempty fiber.


\subsection{Audit packet 6: UEAP no-overreporting (dependency replay)}


\par\noindent\textbf{Claim under audit.} Support-level content cannot be reported as endpoint readout or endpoint defeat.


\par\noindent\textbf{Dependency path.} This packet belongs to the dependency replay. It depends on the official Hodge route, the fixed cycle-fiber endpoint image, the object/use distinction, and the kernel/A+ endpoint-use machinery. It is never used before the carrier and branch slots are fixed.


\par\noindent\textbf{Hostile-referee risk.} A hostile reader may attempt to treat the relevant content as raw mathematical syntax, ordinary negative theoremhood, surrogate support, or a classical counterexample object rather than endpoint-use material.


\par\noindent\textbf{Closure and denial burden.} The closure response is role-sensitive. If the content is raw, support, bookkeeping, or outside the fixed carrier, it has no endpoint-defeating force. If it keeps endpoint-defeating force on the same carrier, it enters endpoint-use discipline and route exhaustion. A successful objection must preserve the same $X,p,\alpha$, the same Hodge-standing predicate, the same Chow carrier $\CH^p(X)_\QQ$, and the same rational cycle-class map $\clQ$.


\par\noindent\textbf{Use in the proof spine.} The packet is reused only after the local exact-complement assumption has been opened. It marks a role transition, not a hidden assertion of nonempty fiber.


\subsection{Audit packet 7: Single cycle-fiber interface (dependency replay)}


\par\noindent\textbf{Claim under audit.} Cycle-fiber occupation is the only positive endpoint-image interface.


\par\noindent\textbf{Dependency path.} This packet belongs to the dependency replay. It depends on the official Hodge route, the fixed cycle-fiber endpoint image, the object/use distinction, and the kernel/A+ endpoint-use machinery. It is never used before the carrier and branch slots are fixed.


\par\noindent\textbf{Hostile-referee risk.} A hostile reader may attempt to treat the relevant content as raw mathematical syntax, ordinary negative theoremhood, surrogate support, or a classical counterexample object rather than endpoint-use material.


\par\noindent\textbf{Closure and denial burden.} The closure response is role-sensitive. If the content is raw, support, bookkeeping, or outside the fixed carrier, it has no endpoint-defeating force. If it keeps endpoint-defeating force on the same carrier, it enters endpoint-use discipline and route exhaustion. A successful objection must preserve the same $X,p,\alpha$, the same Hodge-standing predicate, the same Chow carrier $\CH^p(X)_\QQ$, and the same rational cycle-class map $\clQ$.


\par\noindent\textbf{Use in the proof spine.} The packet is reused only after the local exact-complement assumption has been opened. It marks a role transition, not a hidden assertion of nonempty fiber.


\subsection{Audit packet 8: Native counterforce exhaustion (dependency replay)}


\par\noindent\textbf{Claim under audit.} Native empty-fiber force has no fifth role.


\par\noindent\textbf{Dependency path.} This packet belongs to the dependency replay. It depends on the official Hodge route, the fixed cycle-fiber endpoint image, the object/use distinction, and the kernel/A+ endpoint-use machinery. It is never used before the carrier and branch slots are fixed.


\par\noindent\textbf{Hostile-referee risk.} A hostile reader may attempt to treat the relevant content as raw mathematical syntax, ordinary negative theoremhood, surrogate support, or a classical counterexample object rather than endpoint-use material.


\par\noindent\textbf{Closure and denial burden.} The closure response is role-sensitive. If the content is raw, support, bookkeeping, or outside the fixed carrier, it has no endpoint-defeating force. If it keeps endpoint-defeating force on the same carrier, it enters endpoint-use discipline and route exhaustion. A successful objection must preserve the same $X,p,\alpha$, the same Hodge-standing predicate, the same Chow carrier $\CH^p(X)_\QQ$, and the same rational cycle-class map $\clQ$.


\par\noindent\textbf{Use in the proof spine.} The packet is reused only after the local exact-complement assumption has been opened. It marks a role transition, not a hidden assertion of nonempty fiber.


\subsection{Audit packet 9: Separator induction (dependency replay)}


\par\noindent\textbf{Claim under audit.} Local separator occupation induces the fiber discriminator.


\par\noindent\textbf{Dependency path.} This packet belongs to the dependency replay. It depends on the official Hodge route, the fixed cycle-fiber endpoint image, the object/use distinction, and the kernel/A+ endpoint-use machinery. It is never used before the carrier and branch slots are fixed.


\par\noindent\textbf{Hostile-referee risk.} A hostile reader may attempt to treat the relevant content as raw mathematical syntax, ordinary negative theoremhood, surrogate support, or a classical counterexample object rather than endpoint-use material.


\par\noindent\textbf{Closure and denial burden.} The closure response is role-sensitive. If the content is raw, support, bookkeeping, or outside the fixed carrier, it has no endpoint-defeating force. If it keeps endpoint-defeating force on the same carrier, it enters endpoint-use discipline and route exhaustion. A successful objection must preserve the same $X,p,\alpha$, the same Hodge-standing predicate, the same Chow carrier $\CH^p(X)_\QQ$, and the same rational cycle-class map $\clQ$.


\par\noindent\textbf{Use in the proof spine.} The packet is reused only after the local exact-complement assumption has been opened. It marks a role transition, not a hidden assertion of nonempty fiber.


\subsection{Audit packet 10: A+ no-independent closure (dependency replay)}


\par\noindent\textbf{Claim under audit.} Endpoint use excludes independent same-domain discriminators.


\par\noindent\textbf{Dependency path.} This packet belongs to the dependency replay. It depends on the official Hodge route, the fixed cycle-fiber endpoint image, the object/use distinction, and the kernel/A+ endpoint-use machinery. It is never used before the carrier and branch slots are fixed.


\par\noindent\textbf{Hostile-referee risk.} A hostile reader may attempt to treat the relevant content as raw mathematical syntax, ordinary negative theoremhood, surrogate support, or a classical counterexample object rather than endpoint-use material.


\par\noindent\textbf{Closure and denial burden.} The closure response is role-sensitive. If the content is raw, support, bookkeeping, or outside the fixed carrier, it has no endpoint-defeating force. If it keeps endpoint-defeating force on the same carrier, it enters endpoint-use discipline and route exhaustion. A successful objection must preserve the same $X,p,\alpha$, the same Hodge-standing predicate, the same Chow carrier $\CH^p(X)_\QQ$, and the same rational cycle-class map $\clQ$.


\par\noindent\textbf{Use in the proof spine.} The packet is reused only after the local exact-complement assumption has been opened. It marks a role transition, not a hidden assertion of nonempty fiber.


\subsection{Audit packet 11: Annotation discharge (dependency replay)}


\par\noindent\textbf{Claim under audit.} The exact-complement annotation is not an added premise.


\par\noindent\textbf{Dependency path.} This packet belongs to the dependency replay. It depends on the official Hodge route, the fixed cycle-fiber endpoint image, the object/use distinction, and the kernel/A+ endpoint-use machinery. It is never used before the carrier and branch slots are fixed.


\par\noindent\textbf{Hostile-referee risk.} A hostile reader may attempt to treat the relevant content as raw mathematical syntax, ordinary negative theoremhood, surrogate support, or a classical counterexample object rather than endpoint-use material.


\par\noindent\textbf{Closure and denial burden.} The closure response is role-sensitive. If the content is raw, support, bookkeeping, or outside the fixed carrier, it has no endpoint-defeating force. If it keeps endpoint-defeating force on the same carrier, it enters endpoint-use discipline and route exhaustion. A successful objection must preserve the same $X,p,\alpha$, the same Hodge-standing predicate, the same Chow carrier $\CH^p(X)_\QQ$, and the same rational cycle-class map $\clQ$.


\par\noindent\textbf{Use in the proof spine.} The packet is reused only after the local exact-complement assumption has been opened. It marks a role transition, not a hidden assertion of nonempty fiber.


\subsection{Audit packet 12: Universalization (dependency replay)}


\par\noindent\textbf{Claim under audit.} Arbitrary instance closure yields the global Hodge endpoint.


\par\noindent\textbf{Dependency path.} This packet belongs to the dependency replay. It depends on the official Hodge route, the fixed cycle-fiber endpoint image, the object/use distinction, and the kernel/A+ endpoint-use machinery. It is never used before the carrier and branch slots are fixed.


\par\noindent\textbf{Hostile-referee risk.} A hostile reader may attempt to treat the relevant content as raw mathematical syntax, ordinary negative theoremhood, surrogate support, or a classical counterexample object rather than endpoint-use material.


\par\noindent\textbf{Closure and denial burden.} The closure response is role-sensitive. If the content is raw, support, bookkeeping, or outside the fixed carrier, it has no endpoint-defeating force. If it keeps endpoint-defeating force on the same carrier, it enters endpoint-use discipline and route exhaustion. A successful objection must preserve the same $X,p,\alpha$, the same Hodge-standing predicate, the same Chow carrier $\CH^p(X)_\QQ$, and the same rational cycle-class map $\clQ$.


\par\noindent\textbf{Use in the proof spine.} The packet is reused only after the local exact-complement assumption has been opened. It marks a role transition, not a hidden assertion of nonempty fiber.


\subsection{Audit packet 13: Surrogate separation (dependency replay)}


\par\noindent\textbf{Claim under audit.} Motives and absolute Hodge data are bridge programs unless they supply Chow fiber.


\par\noindent\textbf{Dependency path.} This packet belongs to the dependency replay. It depends on the official Hodge route, the fixed cycle-fiber endpoint image, the object/use distinction, and the kernel/A+ endpoint-use machinery. It is never used before the carrier and branch slots are fixed.


\par\noindent\textbf{Hostile-referee risk.} A hostile reader may attempt to treat the relevant content as raw mathematical syntax, ordinary negative theoremhood, surrogate support, or a classical counterexample object rather than endpoint-use material.


\par\noindent\textbf{Closure and denial burden.} The closure response is role-sensitive. If the content is raw, support, bookkeeping, or outside the fixed carrier, it has no endpoint-defeating force. If it keeps endpoint-defeating force on the same carrier, it enters endpoint-use discipline and route exhaustion. A successful objection must preserve the same $X,p,\alpha$, the same Hodge-standing predicate, the same Chow carrier $\CH^p(X)_\QQ$, and the same rational cycle-class map $\clQ$.


\par\noindent\textbf{Use in the proof spine.} The packet is reused only after the local exact-complement assumption has been opened. It marks a role transition, not a hidden assertion of nonempty fiber.


\subsection{Audit packet 14: Known-case transfer (dependency replay)}


\par\noindent\textbf{Claim under audit.} Known bridge completions do not transfer universally without target gating.


\par\noindent\textbf{Dependency path.} This packet belongs to the dependency replay. It depends on the official Hodge route, the fixed cycle-fiber endpoint image, the object/use distinction, and the kernel/A+ endpoint-use machinery. It is never used before the carrier and branch slots are fixed.


\par\noindent\textbf{Hostile-referee risk.} A hostile reader may attempt to treat the relevant content as raw mathematical syntax, ordinary negative theoremhood, surrogate support, or a classical counterexample object rather than endpoint-use material.


\par\noindent\textbf{Closure and denial burden.} The closure response is role-sensitive. If the content is raw, support, bookkeeping, or outside the fixed carrier, it has no endpoint-defeating force. If it keeps endpoint-defeating force on the same carrier, it enters endpoint-use discipline and route exhaustion. A successful objection must preserve the same $X,p,\alpha$, the same Hodge-standing predicate, the same Chow carrier $\CH^p(X)_\QQ$, and the same rational cycle-class map $\clQ$.


\par\noindent\textbf{Use in the proof spine.} The packet is reused only after the local exact-complement assumption has been opened. It marks a role transition, not a hidden assertion of nonempty fiber.


\subsection{Audit packet 15: Local-to-global nonfactorization (dependency replay)}


\par\noindent\textbf{Claim under audit.} Local cycle data require a global compatibility certificate.


\par\noindent\textbf{Dependency path.} This packet belongs to the dependency replay. It depends on the official Hodge route, the fixed cycle-fiber endpoint image, the object/use distinction, and the kernel/A+ endpoint-use machinery. It is never used before the carrier and branch slots are fixed.


\par\noindent\textbf{Hostile-referee risk.} A hostile reader may attempt to treat the relevant content as raw mathematical syntax, ordinary negative theoremhood, surrogate support, or a classical counterexample object rather than endpoint-use material.


\par\noindent\textbf{Closure and denial burden.} The closure response is role-sensitive. If the content is raw, support, bookkeeping, or outside the fixed carrier, it has no endpoint-defeating force. If it keeps endpoint-defeating force on the same carrier, it enters endpoint-use discipline and route exhaustion. A successful objection must preserve the same $X,p,\alpha$, the same Hodge-standing predicate, the same Chow carrier $\CH^p(X)_\QQ$, and the same rational cycle-class map $\clQ$.


\par\noindent\textbf{Use in the proof spine.} The packet is reused only after the local exact-complement assumption has been opened. It marks a role transition, not a hidden assertion of nonempty fiber.


\subsection{Audit packet 16: No repair (dependency replay)}


\par\noindent\textbf{Claim under audit.} Later cycle discovery creates a new bridge-complete act, not repair of an old one.


\par\noindent\textbf{Dependency path.} This packet belongs to the dependency replay. It depends on the official Hodge route, the fixed cycle-fiber endpoint image, the object/use distinction, and the kernel/A+ endpoint-use machinery. It is never used before the carrier and branch slots are fixed.


\par\noindent\textbf{Hostile-referee risk.} A hostile reader may attempt to treat the relevant content as raw mathematical syntax, ordinary negative theoremhood, surrogate support, or a classical counterexample object rather than endpoint-use material.


\par\noindent\textbf{Closure and denial burden.} The closure response is role-sensitive. If the content is raw, support, bookkeeping, or outside the fixed carrier, it has no endpoint-defeating force. If it keeps endpoint-defeating force on the same carrier, it enters endpoint-use discipline and route exhaustion. A successful objection must preserve the same $X,p,\alpha$, the same Hodge-standing predicate, the same Chow carrier $\CH^p(X)_\QQ$, and the same rational cycle-class map $\clQ$.


\par\noindent\textbf{Use in the proof spine.} The packet is reused only after the local exact-complement assumption has been opened. It marks a role transition, not a hidden assertion of nonempty fiber.


\subsection{Audit packet 17: No selector (dependency replay)}


\par\noindent\textbf{Claim under audit.} Canonical representatives cannot decide readout before the fiber is fixed.


\par\noindent\textbf{Dependency path.} This packet belongs to the dependency replay. It depends on the official Hodge route, the fixed cycle-fiber endpoint image, the object/use distinction, and the kernel/A+ endpoint-use machinery. It is never used before the carrier and branch slots are fixed.


\par\noindent\textbf{Hostile-referee risk.} A hostile reader may attempt to treat the relevant content as raw mathematical syntax, ordinary negative theoremhood, surrogate support, or a classical counterexample object rather than endpoint-use material.


\par\noindent\textbf{Closure and denial burden.} The closure response is role-sensitive. If the content is raw, support, bookkeeping, or outside the fixed carrier, it has no endpoint-defeating force. If it keeps endpoint-defeating force on the same carrier, it enters endpoint-use discipline and route exhaustion. A successful objection must preserve the same $X,p,\alpha$, the same Hodge-standing predicate, the same Chow carrier $\CH^p(X)_\QQ$, and the same rational cycle-class map $\clQ$.


\par\noindent\textbf{Use in the proof spine.} The packet is reused only after the local exact-complement assumption has been opened. It marks a role transition, not a hidden assertion of nonempty fiber.


\subsection{Audit packet 18: Incompleteness firewall (dependency replay)}


\par\noindent\textbf{Claim under audit.} Independence matters only if it re-enters endpoint standing or trajectory.


\par\noindent\textbf{Dependency path.} This packet belongs to the dependency replay. It depends on the official Hodge route, the fixed cycle-fiber endpoint image, the object/use distinction, and the kernel/A+ endpoint-use machinery. It is never used before the carrier and branch slots are fixed.


\par\noindent\textbf{Hostile-referee risk.} A hostile reader may attempt to treat the relevant content as raw mathematical syntax, ordinary negative theoremhood, surrogate support, or a classical counterexample object rather than endpoint-use material.


\par\noindent\textbf{Closure and denial burden.} The closure response is role-sensitive. If the content is raw, support, bookkeeping, or outside the fixed carrier, it has no endpoint-defeating force. If it keeps endpoint-defeating force on the same carrier, it enters endpoint-use discipline and route exhaustion. A successful objection must preserve the same $X,p,\alpha$, the same Hodge-standing predicate, the same Chow carrier $\CH^p(X)_\QQ$, and the same rational cycle-class map $\clQ$.


\par\noindent\textbf{Use in the proof spine.} The packet is reused only after the local exact-complement assumption has been opened. It marks a role transition, not a hidden assertion of nonempty fiber.


\subsection{Audit packet 19: Ordinary theoremhood protection (dependency replay)}


\par\noindent\textbf{Claim under audit.} Negative theoremhood is proof legitimacy, not forbidden endpoint governance.


\par\noindent\textbf{Dependency path.} This packet belongs to the dependency replay. It depends on the official Hodge route, the fixed cycle-fiber endpoint image, the object/use distinction, and the kernel/A+ endpoint-use machinery. It is never used before the carrier and branch slots are fixed.


\par\noindent\textbf{Hostile-referee risk.} A hostile reader may attempt to treat the relevant content as raw mathematical syntax, ordinary negative theoremhood, surrogate support, or a classical counterexample object rather than endpoint-use material.


\par\noindent\textbf{Closure and denial burden.} The closure response is role-sensitive. If the content is raw, support, bookkeeping, or outside the fixed carrier, it has no endpoint-defeating force. If it keeps endpoint-defeating force on the same carrier, it enters endpoint-use discipline and route exhaustion. A successful objection must preserve the same $X,p,\alpha$, the same Hodge-standing predicate, the same Chow carrier $\CH^p(X)_\QQ$, and the same rational cycle-class map $\clQ$.


\par\noindent\textbf{Use in the proof spine.} The packet is reused only after the local exact-complement assumption has been opened. It marks a role transition, not a hidden assertion of nonempty fiber.


\subsection{Audit packet 20: Non-explosiveness (dependency replay)}


\par\noindent\textbf{Claim under audit.} Empty fiber is lawful until endpoint force is added.


\par\noindent\textbf{Dependency path.} This packet belongs to the dependency replay. It depends on the official Hodge route, the fixed cycle-fiber endpoint image, the object/use distinction, and the kernel/A+ endpoint-use machinery. It is never used before the carrier and branch slots are fixed.


\par\noindent\textbf{Hostile-referee risk.} A hostile reader may attempt to treat the relevant content as raw mathematical syntax, ordinary negative theoremhood, surrogate support, or a classical counterexample object rather than endpoint-use material.


\par\noindent\textbf{Closure and denial burden.} The closure response is role-sensitive. If the content is raw, support, bookkeeping, or outside the fixed carrier, it has no endpoint-defeating force. If it keeps endpoint-defeating force on the same carrier, it enters endpoint-use discipline and route exhaustion. A successful objection must preserve the same $X,p,\alpha$, the same Hodge-standing predicate, the same Chow carrier $\CH^p(X)_\QQ$, and the same rational cycle-class map $\clQ$.


\par\noindent\textbf{Use in the proof spine.} The packet is reused only after the local exact-complement assumption has been opened. It marks a role transition, not a hidden assertion of nonempty fiber.


\subsection{Audit packet 21: Context non-smuggling (hostile-reader replay)}


\par\noindent\textbf{Claim under audit.} Hodge context fixes carrier and slots without branch occupation.


\par\noindent\textbf{Dependency path.} This packet belongs to the hostile-reader replay. It depends on the official Hodge route, the fixed cycle-fiber endpoint image, the object/use distinction, and the kernel/A+ endpoint-use machinery. It is never used before the carrier and branch slots are fixed.


\par\noindent\textbf{Hostile-referee risk.} A hostile reader may attempt to treat the relevant content as raw mathematical syntax, ordinary negative theoremhood, surrogate support, or a classical counterexample object rather than endpoint-use material.


\par\noindent\textbf{Closure and denial burden.} The closure response is role-sensitive. If the content is raw, support, bookkeeping, or outside the fixed carrier, it has no endpoint-defeating force. If it keeps endpoint-defeating force on the same carrier, it enters endpoint-use discipline and route exhaustion. A successful objection must preserve the same $X,p,\alpha$, the same Hodge-standing predicate, the same Chow carrier $\CH^p(X)_\QQ$, and the same rational cycle-class map $\clQ$.


\par\noindent\textbf{Use in the proof spine.} The packet is reused only after the local exact-complement assumption has been opened. It marks a role transition, not a hidden assertion of nonempty fiber.


\subsection{Audit packet 22: Model adequacy (hostile-reader replay)}


\par\noindent\textbf{Claim under audit.} Readout and counterreadout are identified with nonempty and empty cycle fiber.


\par\noindent\textbf{Dependency path.} This packet belongs to the hostile-reader replay. It depends on the official Hodge route, the fixed cycle-fiber endpoint image, the object/use distinction, and the kernel/A+ endpoint-use machinery. It is never used before the carrier and branch slots are fixed.


\par\noindent\textbf{Hostile-referee risk.} A hostile reader may attempt to treat the relevant content as raw mathematical syntax, ordinary negative theoremhood, surrogate support, or a classical counterexample object rather than endpoint-use material.


\par\noindent\textbf{Closure and denial burden.} The closure response is role-sensitive. If the content is raw, support, bookkeeping, or outside the fixed carrier, it has no endpoint-defeating force. If it keeps endpoint-defeating force on the same carrier, it enters endpoint-use discipline and route exhaustion. A successful objection must preserve the same $X,p,\alpha$, the same Hodge-standing predicate, the same Chow carrier $\CH^p(X)_\QQ$, and the same rational cycle-class map $\clQ$.


\par\noindent\textbf{Use in the proof spine.} The packet is reused only after the local exact-complement assumption has been opened. It marks a role transition, not a hidden assertion of nonempty fiber.


\subsection{Audit packet 23: Counterexample force as endpoint use (hostile-reader replay)}


\par\noindent\textbf{Claim under audit.} Official endpoint defeat is endpoint use.


\par\noindent\textbf{Dependency path.} This packet belongs to the hostile-reader replay. It depends on the official Hodge route, the fixed cycle-fiber endpoint image, the object/use distinction, and the kernel/A+ endpoint-use machinery. It is never used before the carrier and branch slots are fixed.


\par\noindent\textbf{Hostile-referee risk.} A hostile reader may attempt to treat the relevant content as raw mathematical syntax, ordinary negative theoremhood, surrogate support, or a classical counterexample object rather than endpoint-use material.


\par\noindent\textbf{Closure and denial burden.} The closure response is role-sensitive. If the content is raw, support, bookkeeping, or outside the fixed carrier, it has no endpoint-defeating force. If it keeps endpoint-defeating force on the same carrier, it enters endpoint-use discipline and route exhaustion. A successful objection must preserve the same $X,p,\alpha$, the same Hodge-standing predicate, the same Chow carrier $\CH^p(X)_\QQ$, and the same rational cycle-class map $\clQ$.


\par\noindent\textbf{Use in the proof spine.} The packet is reused only after the local exact-complement assumption has been opened. It marks a role transition, not a hidden assertion of nonempty fiber.


\subsection{Audit packet 24: No autonomous counterexample role (hostile-reader replay)}


\par\noindent\textbf{Claim under audit.} There is no endpoint-defeating branch outside endpoint use.


\par\noindent\textbf{Dependency path.} This packet belongs to the hostile-reader replay. It depends on the official Hodge route, the fixed cycle-fiber endpoint image, the object/use distinction, and the kernel/A+ endpoint-use machinery. It is never used before the carrier and branch slots are fixed.


\par\noindent\textbf{Hostile-referee risk.} A hostile reader may attempt to treat the relevant content as raw mathematical syntax, ordinary negative theoremhood, surrogate support, or a classical counterexample object rather than endpoint-use material.


\par\noindent\textbf{Closure and denial burden.} The closure response is role-sensitive. If the content is raw, support, bookkeeping, or outside the fixed carrier, it has no endpoint-defeating force. If it keeps endpoint-defeating force on the same carrier, it enters endpoint-use discipline and route exhaustion. A successful objection must preserve the same $X,p,\alpha$, the same Hodge-standing predicate, the same Chow carrier $\CH^p(X)_\QQ$, and the same rational cycle-class map $\clQ$.


\par\noindent\textbf{Use in the proof spine.} The packet is reused only after the local exact-complement assumption has been opened. It marks a role transition, not a hidden assertion of nonempty fiber.


\subsection{Audit packet 25: ATS locking (hostile-reader replay)}


\par\noindent\textbf{Claim under audit.} The same endpoint act preserves anchor and tensor data.


\par\noindent\textbf{Dependency path.} This packet belongs to the hostile-reader replay. It depends on the official Hodge route, the fixed cycle-fiber endpoint image, the object/use distinction, and the kernel/A+ endpoint-use machinery. It is never used before the carrier and branch slots are fixed.


\par\noindent\textbf{Hostile-referee risk.} A hostile reader may attempt to treat the relevant content as raw mathematical syntax, ordinary negative theoremhood, surrogate support, or a classical counterexample object rather than endpoint-use material.


\par\noindent\textbf{Closure and denial burden.} The closure response is role-sensitive. If the content is raw, support, bookkeeping, or outside the fixed carrier, it has no endpoint-defeating force. If it keeps endpoint-defeating force on the same carrier, it enters endpoint-use discipline and route exhaustion. A successful objection must preserve the same $X,p,\alpha$, the same Hodge-standing predicate, the same Chow carrier $\CH^p(X)_\QQ$, and the same rational cycle-class map $\clQ$.


\par\noindent\textbf{Use in the proof spine.} The packet is reused only after the local exact-complement assumption has been opened. It marks a role transition, not a hidden assertion of nonempty fiber.


\subsection{Audit packet 26: UEAP no-overreporting (hostile-reader replay)}


\par\noindent\textbf{Claim under audit.} Support-level content cannot be reported as endpoint readout or endpoint defeat.


\par\noindent\textbf{Dependency path.} This packet belongs to the hostile-reader replay. It depends on the official Hodge route, the fixed cycle-fiber endpoint image, the object/use distinction, and the kernel/A+ endpoint-use machinery. It is never used before the carrier and branch slots are fixed.


\par\noindent\textbf{Hostile-referee risk.} A hostile reader may attempt to treat the relevant content as raw mathematical syntax, ordinary negative theoremhood, surrogate support, or a classical counterexample object rather than endpoint-use material.


\par\noindent\textbf{Closure and denial burden.} The closure response is role-sensitive. If the content is raw, support, bookkeeping, or outside the fixed carrier, it has no endpoint-defeating force. If it keeps endpoint-defeating force on the same carrier, it enters endpoint-use discipline and route exhaustion. A successful objection must preserve the same $X,p,\alpha$, the same Hodge-standing predicate, the same Chow carrier $\CH^p(X)_\QQ$, and the same rational cycle-class map $\clQ$.


\par\noindent\textbf{Use in the proof spine.} The packet is reused only after the local exact-complement assumption has been opened. It marks a role transition, not a hidden assertion of nonempty fiber.


\subsection{Audit packet 27: Single cycle-fiber interface (hostile-reader replay)}


\par\noindent\textbf{Claim under audit.} Cycle-fiber occupation is the only positive endpoint-image interface.


\par\noindent\textbf{Dependency path.} This packet belongs to the hostile-reader replay. It depends on the official Hodge route, the fixed cycle-fiber endpoint image, the object/use distinction, and the kernel/A+ endpoint-use machinery. It is never used before the carrier and branch slots are fixed.


\par\noindent\textbf{Hostile-referee risk.} A hostile reader may attempt to treat the relevant content as raw mathematical syntax, ordinary negative theoremhood, surrogate support, or a classical counterexample object rather than endpoint-use material.


\par\noindent\textbf{Closure and denial burden.} The closure response is role-sensitive. If the content is raw, support, bookkeeping, or outside the fixed carrier, it has no endpoint-defeating force. If it keeps endpoint-defeating force on the same carrier, it enters endpoint-use discipline and route exhaustion. A successful objection must preserve the same $X,p,\alpha$, the same Hodge-standing predicate, the same Chow carrier $\CH^p(X)_\QQ$, and the same rational cycle-class map $\clQ$.


\par\noindent\textbf{Use in the proof spine.} The packet is reused only after the local exact-complement assumption has been opened. It marks a role transition, not a hidden assertion of nonempty fiber.


\subsection{Audit packet 28: Native counterforce exhaustion (hostile-reader replay)}


\par\noindent\textbf{Claim under audit.} Native empty-fiber force has no fifth role.


\par\noindent\textbf{Dependency path.} This packet belongs to the hostile-reader replay. It depends on the official Hodge route, the fixed cycle-fiber endpoint image, the object/use distinction, and the kernel/A+ endpoint-use machinery. It is never used before the carrier and branch slots are fixed.


\par\noindent\textbf{Hostile-referee risk.} A hostile reader may attempt to treat the relevant content as raw mathematical syntax, ordinary negative theoremhood, surrogate support, or a classical counterexample object rather than endpoint-use material.


\par\noindent\textbf{Closure and denial burden.} The closure response is role-sensitive. If the content is raw, support, bookkeeping, or outside the fixed carrier, it has no endpoint-defeating force. If it keeps endpoint-defeating force on the same carrier, it enters endpoint-use discipline and route exhaustion. A successful objection must preserve the same $X,p,\alpha$, the same Hodge-standing predicate, the same Chow carrier $\CH^p(X)_\QQ$, and the same rational cycle-class map $\clQ$.


\par\noindent\textbf{Use in the proof spine.} The packet is reused only after the local exact-complement assumption has been opened. It marks a role transition, not a hidden assertion of nonempty fiber.


\subsection{Audit packet 29: Separator induction (hostile-reader replay)}


\par\noindent\textbf{Claim under audit.} Local separator occupation induces the fiber discriminator.


\par\noindent\textbf{Dependency path.} This packet belongs to the hostile-reader replay. It depends on the official Hodge route, the fixed cycle-fiber endpoint image, the object/use distinction, and the kernel/A+ endpoint-use machinery. It is never used before the carrier and branch slots are fixed.


\par\noindent\textbf{Hostile-referee risk.} A hostile reader may attempt to treat the relevant content as raw mathematical syntax, ordinary negative theoremhood, surrogate support, or a classical counterexample object rather than endpoint-use material.


\par\noindent\textbf{Closure and denial burden.} The closure response is role-sensitive. If the content is raw, support, bookkeeping, or outside the fixed carrier, it has no endpoint-defeating force. If it keeps endpoint-defeating force on the same carrier, it enters endpoint-use discipline and route exhaustion. A successful objection must preserve the same $X,p,\alpha$, the same Hodge-standing predicate, the same Chow carrier $\CH^p(X)_\QQ$, and the same rational cycle-class map $\clQ$.


\par\noindent\textbf{Use in the proof spine.} The packet is reused only after the local exact-complement assumption has been opened. It marks a role transition, not a hidden assertion of nonempty fiber.


\subsection{Audit packet 30: A+ no-independent closure (hostile-reader replay)}


\par\noindent\textbf{Claim under audit.} Endpoint use excludes independent same-domain discriminators.


\par\noindent\textbf{Dependency path.} This packet belongs to the hostile-reader replay. It depends on the official Hodge route, the fixed cycle-fiber endpoint image, the object/use distinction, and the kernel/A+ endpoint-use machinery. It is never used before the carrier and branch slots are fixed.


\par\noindent\textbf{Hostile-referee risk.} A hostile reader may attempt to treat the relevant content as raw mathematical syntax, ordinary negative theoremhood, surrogate support, or a classical counterexample object rather than endpoint-use material.


\par\noindent\textbf{Closure and denial burden.} The closure response is role-sensitive. If the content is raw, support, bookkeeping, or outside the fixed carrier, it has no endpoint-defeating force. If it keeps endpoint-defeating force on the same carrier, it enters endpoint-use discipline and route exhaustion. A successful objection must preserve the same $X,p,\alpha$, the same Hodge-standing predicate, the same Chow carrier $\CH^p(X)_\QQ$, and the same rational cycle-class map $\clQ$.


\par\noindent\textbf{Use in the proof spine.} The packet is reused only after the local exact-complement assumption has been opened. It marks a role transition, not a hidden assertion of nonempty fiber.


\subsection{Audit packet 31: Annotation discharge (hostile-reader replay)}


\par\noindent\textbf{Claim under audit.} The exact-complement annotation is not an added premise.


\par\noindent\textbf{Dependency path.} This packet belongs to the hostile-reader replay. It depends on the official Hodge route, the fixed cycle-fiber endpoint image, the object/use distinction, and the kernel/A+ endpoint-use machinery. It is never used before the carrier and branch slots are fixed.


\par\noindent\textbf{Hostile-referee risk.} A hostile reader may attempt to treat the relevant content as raw mathematical syntax, ordinary negative theoremhood, surrogate support, or a classical counterexample object rather than endpoint-use material.


\par\noindent\textbf{Closure and denial burden.} The closure response is role-sensitive. If the content is raw, support, bookkeeping, or outside the fixed carrier, it has no endpoint-defeating force. If it keeps endpoint-defeating force on the same carrier, it enters endpoint-use discipline and route exhaustion. A successful objection must preserve the same $X,p,\alpha$, the same Hodge-standing predicate, the same Chow carrier $\CH^p(X)_\QQ$, and the same rational cycle-class map $\clQ$.


\par\noindent\textbf{Use in the proof spine.} The packet is reused only after the local exact-complement assumption has been opened. It marks a role transition, not a hidden assertion of nonempty fiber.


\subsection{Audit packet 32: Universalization (hostile-reader replay)}


\par\noindent\textbf{Claim under audit.} Arbitrary instance closure yields the global Hodge endpoint.


\par\noindent\textbf{Dependency path.} This packet belongs to the hostile-reader replay. It depends on the official Hodge route, the fixed cycle-fiber endpoint image, the object/use distinction, and the kernel/A+ endpoint-use machinery. It is never used before the carrier and branch slots are fixed.


\par\noindent\textbf{Hostile-referee risk.} A hostile reader may attempt to treat the relevant content as raw mathematical syntax, ordinary negative theoremhood, surrogate support, or a classical counterexample object rather than endpoint-use material.


\par\noindent\textbf{Closure and denial burden.} The closure response is role-sensitive. If the content is raw, support, bookkeeping, or outside the fixed carrier, it has no endpoint-defeating force. If it keeps endpoint-defeating force on the same carrier, it enters endpoint-use discipline and route exhaustion. A successful objection must preserve the same $X,p,\alpha$, the same Hodge-standing predicate, the same Chow carrier $\CH^p(X)_\QQ$, and the same rational cycle-class map $\clQ$.


\par\noindent\textbf{Use in the proof spine.} The packet is reused only after the local exact-complement assumption has been opened. It marks a role transition, not a hidden assertion of nonempty fiber.


\subsection{Audit packet 33: Surrogate separation (hostile-reader replay)}


\par\noindent\textbf{Claim under audit.} Motives and absolute Hodge data are bridge programs unless they supply Chow fiber.


\par\noindent\textbf{Dependency path.} This packet belongs to the hostile-reader replay. It depends on the official Hodge route, the fixed cycle-fiber endpoint image, the object/use distinction, and the kernel/A+ endpoint-use machinery. It is never used before the carrier and branch slots are fixed.


\par\noindent\textbf{Hostile-referee risk.} A hostile reader may attempt to treat the relevant content as raw mathematical syntax, ordinary negative theoremhood, surrogate support, or a classical counterexample object rather than endpoint-use material.


\par\noindent\textbf{Closure and denial burden.} The closure response is role-sensitive. If the content is raw, support, bookkeeping, or outside the fixed carrier, it has no endpoint-defeating force. If it keeps endpoint-defeating force on the same carrier, it enters endpoint-use discipline and route exhaustion. A successful objection must preserve the same $X,p,\alpha$, the same Hodge-standing predicate, the same Chow carrier $\CH^p(X)_\QQ$, and the same rational cycle-class map $\clQ$.


\par\noindent\textbf{Use in the proof spine.} The packet is reused only after the local exact-complement assumption has been opened. It marks a role transition, not a hidden assertion of nonempty fiber.


\subsection{Audit packet 34: Known-case transfer (hostile-reader replay)}


\par\noindent\textbf{Claim under audit.} Known bridge completions do not transfer universally without target gating.


\par\noindent\textbf{Dependency path.} This packet belongs to the hostile-reader replay. It depends on the official Hodge route, the fixed cycle-fiber endpoint image, the object/use distinction, and the kernel/A+ endpoint-use machinery. It is never used before the carrier and branch slots are fixed.


\par\noindent\textbf{Hostile-referee risk.} A hostile reader may attempt to treat the relevant content as raw mathematical syntax, ordinary negative theoremhood, surrogate support, or a classical counterexample object rather than endpoint-use material.


\par\noindent\textbf{Closure and denial burden.} The closure response is role-sensitive. If the content is raw, support, bookkeeping, or outside the fixed carrier, it has no endpoint-defeating force. If it keeps endpoint-defeating force on the same carrier, it enters endpoint-use discipline and route exhaustion. A successful objection must preserve the same $X,p,\alpha$, the same Hodge-standing predicate, the same Chow carrier $\CH^p(X)_\QQ$, and the same rational cycle-class map $\clQ$.


\par\noindent\textbf{Use in the proof spine.} The packet is reused only after the local exact-complement assumption has been opened. It marks a role transition, not a hidden assertion of nonempty fiber.


\subsection{Audit packet 35: Local-to-global nonfactorization (hostile-reader replay)}


\par\noindent\textbf{Claim under audit.} Local cycle data require a global compatibility certificate.


\par\noindent\textbf{Dependency path.} This packet belongs to the hostile-reader replay. It depends on the official Hodge route, the fixed cycle-fiber endpoint image, the object/use distinction, and the kernel/A+ endpoint-use machinery. It is never used before the carrier and branch slots are fixed.


\par\noindent\textbf{Hostile-referee risk.} A hostile reader may attempt to treat the relevant content as raw mathematical syntax, ordinary negative theoremhood, surrogate support, or a classical counterexample object rather than endpoint-use material.


\par\noindent\textbf{Closure and denial burden.} The closure response is role-sensitive. If the content is raw, support, bookkeeping, or outside the fixed carrier, it has no endpoint-defeating force. If it keeps endpoint-defeating force on the same carrier, it enters endpoint-use discipline and route exhaustion. A successful objection must preserve the same $X,p,\alpha$, the same Hodge-standing predicate, the same Chow carrier $\CH^p(X)_\QQ$, and the same rational cycle-class map $\clQ$.


\par\noindent\textbf{Use in the proof spine.} The packet is reused only after the local exact-complement assumption has been opened. It marks a role transition, not a hidden assertion of nonempty fiber.


\subsection{Audit packet 36: No repair (hostile-reader replay)}


\par\noindent\textbf{Claim under audit.} Later cycle discovery creates a new bridge-complete act, not repair of an old one.


\par\noindent\textbf{Dependency path.} This packet belongs to the hostile-reader replay. It depends on the official Hodge route, the fixed cycle-fiber endpoint image, the object/use distinction, and the kernel/A+ endpoint-use machinery. It is never used before the carrier and branch slots are fixed.


\par\noindent\textbf{Hostile-referee risk.} A hostile reader may attempt to treat the relevant content as raw mathematical syntax, ordinary negative theoremhood, surrogate support, or a classical counterexample object rather than endpoint-use material.


\par\noindent\textbf{Closure and denial burden.} The closure response is role-sensitive. If the content is raw, support, bookkeeping, or outside the fixed carrier, it has no endpoint-defeating force. If it keeps endpoint-defeating force on the same carrier, it enters endpoint-use discipline and route exhaustion. A successful objection must preserve the same $X,p,\alpha$, the same Hodge-standing predicate, the same Chow carrier $\CH^p(X)_\QQ$, and the same rational cycle-class map $\clQ$.


\par\noindent\textbf{Use in the proof spine.} The packet is reused only after the local exact-complement assumption has been opened. It marks a role transition, not a hidden assertion of nonempty fiber.


\subsection{Audit packet 37: No selector (hostile-reader replay)}


\par\noindent\textbf{Claim under audit.} Canonical representatives cannot decide readout before the fiber is fixed.


\par\noindent\textbf{Dependency path.} This packet belongs to the hostile-reader replay. It depends on the official Hodge route, the fixed cycle-fiber endpoint image, the object/use distinction, and the kernel/A+ endpoint-use machinery. It is never used before the carrier and branch slots are fixed.


\par\noindent\textbf{Hostile-referee risk.} A hostile reader may attempt to treat the relevant content as raw mathematical syntax, ordinary negative theoremhood, surrogate support, or a classical counterexample object rather than endpoint-use material.


\par\noindent\textbf{Closure and denial burden.} The closure response is role-sensitive. If the content is raw, support, bookkeeping, or outside the fixed carrier, it has no endpoint-defeating force. If it keeps endpoint-defeating force on the same carrier, it enters endpoint-use discipline and route exhaustion. A successful objection must preserve the same $X,p,\alpha$, the same Hodge-standing predicate, the same Chow carrier $\CH^p(X)_\QQ$, and the same rational cycle-class map $\clQ$.


\par\noindent\textbf{Use in the proof spine.} The packet is reused only after the local exact-complement assumption has been opened. It marks a role transition, not a hidden assertion of nonempty fiber.


\subsection{Audit packet 38: Incompleteness firewall (hostile-reader replay)}


\par\noindent\textbf{Claim under audit.} Independence matters only if it re-enters endpoint standing or trajectory.


\par\noindent\textbf{Dependency path.} This packet belongs to the hostile-reader replay. It depends on the official Hodge route, the fixed cycle-fiber endpoint image, the object/use distinction, and the kernel/A+ endpoint-use machinery. It is never used before the carrier and branch slots are fixed.


\par\noindent\textbf{Hostile-referee risk.} A hostile reader may attempt to treat the relevant content as raw mathematical syntax, ordinary negative theoremhood, surrogate support, or a classical counterexample object rather than endpoint-use material.


\par\noindent\textbf{Closure and denial burden.} The closure response is role-sensitive. If the content is raw, support, bookkeeping, or outside the fixed carrier, it has no endpoint-defeating force. If it keeps endpoint-defeating force on the same carrier, it enters endpoint-use discipline and route exhaustion. A successful objection must preserve the same $X,p,\alpha$, the same Hodge-standing predicate, the same Chow carrier $\CH^p(X)_\QQ$, and the same rational cycle-class map $\clQ$.


\par\noindent\textbf{Use in the proof spine.} The packet is reused only after the local exact-complement assumption has been opened. It marks a role transition, not a hidden assertion of nonempty fiber.


\subsection{Audit packet 39: Ordinary theoremhood protection (hostile-reader replay)}


\par\noindent\textbf{Claim under audit.} Negative theoremhood is proof legitimacy, not forbidden endpoint governance.


\par\noindent\textbf{Dependency path.} This packet belongs to the hostile-reader replay. It depends on the official Hodge route, the fixed cycle-fiber endpoint image, the object/use distinction, and the kernel/A+ endpoint-use machinery. It is never used before the carrier and branch slots are fixed.


\par\noindent\textbf{Hostile-referee risk.} A hostile reader may attempt to treat the relevant content as raw mathematical syntax, ordinary negative theoremhood, surrogate support, or a classical counterexample object rather than endpoint-use material.


\par\noindent\textbf{Closure and denial burden.} The closure response is role-sensitive. If the content is raw, support, bookkeeping, or outside the fixed carrier, it has no endpoint-defeating force. If it keeps endpoint-defeating force on the same carrier, it enters endpoint-use discipline and route exhaustion. A successful objection must preserve the same $X,p,\alpha$, the same Hodge-standing predicate, the same Chow carrier $\CH^p(X)_\QQ$, and the same rational cycle-class map $\clQ$.


\par\noindent\textbf{Use in the proof spine.} The packet is reused only after the local exact-complement assumption has been opened. It marks a role transition, not a hidden assertion of nonempty fiber.


\subsection{Audit packet 40: Non-explosiveness (hostile-reader replay)}


\par\noindent\textbf{Claim under audit.} Empty fiber is lawful until endpoint force is added.


\par\noindent\textbf{Dependency path.} This packet belongs to the hostile-reader replay. It depends on the official Hodge route, the fixed cycle-fiber endpoint image, the object/use distinction, and the kernel/A+ endpoint-use machinery. It is never used before the carrier and branch slots are fixed.


\par\noindent\textbf{Hostile-referee risk.} A hostile reader may attempt to treat the relevant content as raw mathematical syntax, ordinary negative theoremhood, surrogate support, or a classical counterexample object rather than endpoint-use material.


\par\noindent\textbf{Closure and denial burden.} The closure response is role-sensitive. If the content is raw, support, bookkeeping, or outside the fixed carrier, it has no endpoint-defeating force. If it keeps endpoint-defeating force on the same carrier, it enters endpoint-use discipline and route exhaustion. A successful objection must preserve the same $X,p,\alpha$, the same Hodge-standing predicate, the same Chow carrier $\CH^p(X)_\QQ$, and the same rational cycle-class map $\clQ$.


\par\noindent\textbf{Use in the proof spine.} The packet is reused only after the local exact-complement assumption has been opened. It marks a role transition, not a hidden assertion of nonempty fiber.


\subsection{Audit packet 41: Context non-smuggling (same-mode denial replay)}


\par\noindent\textbf{Claim under audit.} Hodge context fixes carrier and slots without branch occupation.


\par\noindent\textbf{Dependency path.} This packet belongs to the same-mode denial replay. It depends on the official Hodge route, the fixed cycle-fiber endpoint image, the object/use distinction, and the kernel/A+ endpoint-use machinery. It is never used before the carrier and branch slots are fixed.


\par\noindent\textbf{Hostile-referee risk.} A hostile reader may attempt to treat the relevant content as raw mathematical syntax, ordinary negative theoremhood, surrogate support, or a classical counterexample object rather than endpoint-use material.


\par\noindent\textbf{Closure and denial burden.} The closure response is role-sensitive. If the content is raw, support, bookkeeping, or outside the fixed carrier, it has no endpoint-defeating force. If it keeps endpoint-defeating force on the same carrier, it enters endpoint-use discipline and route exhaustion. A successful objection must preserve the same $X,p,\alpha$, the same Hodge-standing predicate, the same Chow carrier $\CH^p(X)_\QQ$, and the same rational cycle-class map $\clQ$.


\par\noindent\textbf{Use in the proof spine.} The packet is reused only after the local exact-complement assumption has been opened. It marks a role transition, not a hidden assertion of nonempty fiber.


\subsection{Audit packet 42: Model adequacy (same-mode denial replay)}


\par\noindent\textbf{Claim under audit.} Readout and counterreadout are identified with nonempty and empty cycle fiber.


\par\noindent\textbf{Dependency path.} This packet belongs to the same-mode denial replay. It depends on the official Hodge route, the fixed cycle-fiber endpoint image, the object/use distinction, and the kernel/A+ endpoint-use machinery. It is never used before the carrier and branch slots are fixed.


\par\noindent\textbf{Hostile-referee risk.} A hostile reader may attempt to treat the relevant content as raw mathematical syntax, ordinary negative theoremhood, surrogate support, or a classical counterexample object rather than endpoint-use material.


\par\noindent\textbf{Closure and denial burden.} The closure response is role-sensitive. If the content is raw, support, bookkeeping, or outside the fixed carrier, it has no endpoint-defeating force. If it keeps endpoint-defeating force on the same carrier, it enters endpoint-use discipline and route exhaustion. A successful objection must preserve the same $X,p,\alpha$, the same Hodge-standing predicate, the same Chow carrier $\CH^p(X)_\QQ$, and the same rational cycle-class map $\clQ$.


\par\noindent\textbf{Use in the proof spine.} The packet is reused only after the local exact-complement assumption has been opened. It marks a role transition, not a hidden assertion of nonempty fiber.


\subsection{Audit packet 43: Counterexample force as endpoint use (same-mode denial replay)}


\par\noindent\textbf{Claim under audit.} Official endpoint defeat is endpoint use.


\par\noindent\textbf{Dependency path.} This packet belongs to the same-mode denial replay. It depends on the official Hodge route, the fixed cycle-fiber endpoint image, the object/use distinction, and the kernel/A+ endpoint-use machinery. It is never used before the carrier and branch slots are fixed.


\par\noindent\textbf{Hostile-referee risk.} A hostile reader may attempt to treat the relevant content as raw mathematical syntax, ordinary negative theoremhood, surrogate support, or a classical counterexample object rather than endpoint-use material.


\par\noindent\textbf{Closure and denial burden.} The closure response is role-sensitive. If the content is raw, support, bookkeeping, or outside the fixed carrier, it has no endpoint-defeating force. If it keeps endpoint-defeating force on the same carrier, it enters endpoint-use discipline and route exhaustion. A successful objection must preserve the same $X,p,\alpha$, the same Hodge-standing predicate, the same Chow carrier $\CH^p(X)_\QQ$, and the same rational cycle-class map $\clQ$.


\par\noindent\textbf{Use in the proof spine.} The packet is reused only after the local exact-complement assumption has been opened. It marks a role transition, not a hidden assertion of nonempty fiber.


\subsection{Audit packet 44: No autonomous counterexample role (same-mode denial replay)}


\par\noindent\textbf{Claim under audit.} There is no endpoint-defeating branch outside endpoint use.


\par\noindent\textbf{Dependency path.} This packet belongs to the same-mode denial replay. It depends on the official Hodge route, the fixed cycle-fiber endpoint image, the object/use distinction, and the kernel/A+ endpoint-use machinery. It is never used before the carrier and branch slots are fixed.


\par\noindent\textbf{Hostile-referee risk.} A hostile reader may attempt to treat the relevant content as raw mathematical syntax, ordinary negative theoremhood, surrogate support, or a classical counterexample object rather than endpoint-use material.


\par\noindent\textbf{Closure and denial burden.} The closure response is role-sensitive. If the content is raw, support, bookkeeping, or outside the fixed carrier, it has no endpoint-defeating force. If it keeps endpoint-defeating force on the same carrier, it enters endpoint-use discipline and route exhaustion. A successful objection must preserve the same $X,p,\alpha$, the same Hodge-standing predicate, the same Chow carrier $\CH^p(X)_\QQ$, and the same rational cycle-class map $\clQ$.


\par\noindent\textbf{Use in the proof spine.} The packet is reused only after the local exact-complement assumption has been opened. It marks a role transition, not a hidden assertion of nonempty fiber.


\subsection{Audit packet 45: ATS locking (same-mode denial replay)}


\par\noindent\textbf{Claim under audit.} The same endpoint act preserves anchor and tensor data.


\par\noindent\textbf{Dependency path.} This packet belongs to the same-mode denial replay. It depends on the official Hodge route, the fixed cycle-fiber endpoint image, the object/use distinction, and the kernel/A+ endpoint-use machinery. It is never used before the carrier and branch slots are fixed.


\par\noindent\textbf{Hostile-referee risk.} A hostile reader may attempt to treat the relevant content as raw mathematical syntax, ordinary negative theoremhood, surrogate support, or a classical counterexample object rather than endpoint-use material.


\par\noindent\textbf{Closure and denial burden.} The closure response is role-sensitive. If the content is raw, support, bookkeeping, or outside the fixed carrier, it has no endpoint-defeating force. If it keeps endpoint-defeating force on the same carrier, it enters endpoint-use discipline and route exhaustion. A successful objection must preserve the same $X,p,\alpha$, the same Hodge-standing predicate, the same Chow carrier $\CH^p(X)_\QQ$, and the same rational cycle-class map $\clQ$.


\par\noindent\textbf{Use in the proof spine.} The packet is reused only after the local exact-complement assumption has been opened. It marks a role transition, not a hidden assertion of nonempty fiber.


\subsection{Audit packet 46: UEAP no-overreporting (same-mode denial replay)}


\par\noindent\textbf{Claim under audit.} Support-level content cannot be reported as endpoint readout or endpoint defeat.


\par\noindent\textbf{Dependency path.} This packet belongs to the same-mode denial replay. It depends on the official Hodge route, the fixed cycle-fiber endpoint image, the object/use distinction, and the kernel/A+ endpoint-use machinery. It is never used before the carrier and branch slots are fixed.


\par\noindent\textbf{Hostile-referee risk.} A hostile reader may attempt to treat the relevant content as raw mathematical syntax, ordinary negative theoremhood, surrogate support, or a classical counterexample object rather than endpoint-use material.


\par\noindent\textbf{Closure and denial burden.} The closure response is role-sensitive. If the content is raw, support, bookkeeping, or outside the fixed carrier, it has no endpoint-defeating force. If it keeps endpoint-defeating force on the same carrier, it enters endpoint-use discipline and route exhaustion. A successful objection must preserve the same $X,p,\alpha$, the same Hodge-standing predicate, the same Chow carrier $\CH^p(X)_\QQ$, and the same rational cycle-class map $\clQ$.


\par\noindent\textbf{Use in the proof spine.} The packet is reused only after the local exact-complement assumption has been opened. It marks a role transition, not a hidden assertion of nonempty fiber.


\subsection{Audit packet 47: Single cycle-fiber interface (same-mode denial replay)}


\par\noindent\textbf{Claim under audit.} Cycle-fiber occupation is the only positive endpoint-image interface.


\par\noindent\textbf{Dependency path.} This packet belongs to the same-mode denial replay. It depends on the official Hodge route, the fixed cycle-fiber endpoint image, the object/use distinction, and the kernel/A+ endpoint-use machinery. It is never used before the carrier and branch slots are fixed.


\par\noindent\textbf{Hostile-referee risk.} A hostile reader may attempt to treat the relevant content as raw mathematical syntax, ordinary negative theoremhood, surrogate support, or a classical counterexample object rather than endpoint-use material.


\par\noindent\textbf{Closure and denial burden.} The closure response is role-sensitive. If the content is raw, support, bookkeeping, or outside the fixed carrier, it has no endpoint-defeating force. If it keeps endpoint-defeating force on the same carrier, it enters endpoint-use discipline and route exhaustion. A successful objection must preserve the same $X,p,\alpha$, the same Hodge-standing predicate, the same Chow carrier $\CH^p(X)_\QQ$, and the same rational cycle-class map $\clQ$.


\par\noindent\textbf{Use in the proof spine.} The packet is reused only after the local exact-complement assumption has been opened. It marks a role transition, not a hidden assertion of nonempty fiber.


\subsection{Audit packet 48: Native counterforce exhaustion (same-mode denial replay)}


\par\noindent\textbf{Claim under audit.} Native empty-fiber force has no fifth role.


\par\noindent\textbf{Dependency path.} This packet belongs to the same-mode denial replay. It depends on the official Hodge route, the fixed cycle-fiber endpoint image, the object/use distinction, and the kernel/A+ endpoint-use machinery. It is never used before the carrier and branch slots are fixed.


\par\noindent\textbf{Hostile-referee risk.} A hostile reader may attempt to treat the relevant content as raw mathematical syntax, ordinary negative theoremhood, surrogate support, or a classical counterexample object rather than endpoint-use material.


\par\noindent\textbf{Closure and denial burden.} The closure response is role-sensitive. If the content is raw, support, bookkeeping, or outside the fixed carrier, it has no endpoint-defeating force. If it keeps endpoint-defeating force on the same carrier, it enters endpoint-use discipline and route exhaustion. A successful objection must preserve the same $X,p,\alpha$, the same Hodge-standing predicate, the same Chow carrier $\CH^p(X)_\QQ$, and the same rational cycle-class map $\clQ$.


\par\noindent\textbf{Use in the proof spine.} The packet is reused only after the local exact-complement assumption has been opened. It marks a role transition, not a hidden assertion of nonempty fiber.


\subsection{Audit packet 49: Separator induction (same-mode denial replay)}


\par\noindent\textbf{Claim under audit.} Local separator occupation induces the fiber discriminator.


\par\noindent\textbf{Dependency path.} This packet belongs to the same-mode denial replay. It depends on the official Hodge route, the fixed cycle-fiber endpoint image, the object/use distinction, and the kernel/A+ endpoint-use machinery. It is never used before the carrier and branch slots are fixed.


\par\noindent\textbf{Hostile-referee risk.} A hostile reader may attempt to treat the relevant content as raw mathematical syntax, ordinary negative theoremhood, surrogate support, or a classical counterexample object rather than endpoint-use material.


\par\noindent\textbf{Closure and denial burden.} The closure response is role-sensitive. If the content is raw, support, bookkeeping, or outside the fixed carrier, it has no endpoint-defeating force. If it keeps endpoint-defeating force on the same carrier, it enters endpoint-use discipline and route exhaustion. A successful objection must preserve the same $X,p,\alpha$, the same Hodge-standing predicate, the same Chow carrier $\CH^p(X)_\QQ$, and the same rational cycle-class map $\clQ$.


\par\noindent\textbf{Use in the proof spine.} The packet is reused only after the local exact-complement assumption has been opened. It marks a role transition, not a hidden assertion of nonempty fiber.


\subsection{Audit packet 50: A+ no-independent closure (same-mode denial replay)}


\par\noindent\textbf{Claim under audit.} Endpoint use excludes independent same-domain discriminators.


\par\noindent\textbf{Dependency path.} This packet belongs to the same-mode denial replay. It depends on the official Hodge route, the fixed cycle-fiber endpoint image, the object/use distinction, and the kernel/A+ endpoint-use machinery. It is never used before the carrier and branch slots are fixed.


\par\noindent\textbf{Hostile-referee risk.} A hostile reader may attempt to treat the relevant content as raw mathematical syntax, ordinary negative theoremhood, surrogate support, or a classical counterexample object rather than endpoint-use material.


\par\noindent\textbf{Closure and denial burden.} The closure response is role-sensitive. If the content is raw, support, bookkeeping, or outside the fixed carrier, it has no endpoint-defeating force. If it keeps endpoint-defeating force on the same carrier, it enters endpoint-use discipline and route exhaustion. A successful objection must preserve the same $X,p,\alpha$, the same Hodge-standing predicate, the same Chow carrier $\CH^p(X)_\QQ$, and the same rational cycle-class map $\clQ$.


\par\noindent\textbf{Use in the proof spine.} The packet is reused only after the local exact-complement assumption has been opened. It marks a role transition, not a hidden assertion of nonempty fiber.


\subsection{Audit packet 51: Annotation discharge (same-mode denial replay)}


\par\noindent\textbf{Claim under audit.} The exact-complement annotation is not an added premise.


\par\noindent\textbf{Dependency path.} This packet belongs to the same-mode denial replay. It depends on the official Hodge route, the fixed cycle-fiber endpoint image, the object/use distinction, and the kernel/A+ endpoint-use machinery. It is never used before the carrier and branch slots are fixed.


\par\noindent\textbf{Hostile-referee risk.} A hostile reader may attempt to treat the relevant content as raw mathematical syntax, ordinary negative theoremhood, surrogate support, or a classical counterexample object rather than endpoint-use material.


\par\noindent\textbf{Closure and denial burden.} The closure response is role-sensitive. If the content is raw, support, bookkeeping, or outside the fixed carrier, it has no endpoint-defeating force. If it keeps endpoint-defeating force on the same carrier, it enters endpoint-use discipline and route exhaustion. A successful objection must preserve the same $X,p,\alpha$, the same Hodge-standing predicate, the same Chow carrier $\CH^p(X)_\QQ$, and the same rational cycle-class map $\clQ$.


\par\noindent\textbf{Use in the proof spine.} The packet is reused only after the local exact-complement assumption has been opened. It marks a role transition, not a hidden assertion of nonempty fiber.


\subsection{Audit packet 52: Universalization (same-mode denial replay)}


\par\noindent\textbf{Claim under audit.} Arbitrary instance closure yields the global Hodge endpoint.


\par\noindent\textbf{Dependency path.} This packet belongs to the same-mode denial replay. It depends on the official Hodge route, the fixed cycle-fiber endpoint image, the object/use distinction, and the kernel/A+ endpoint-use machinery. It is never used before the carrier and branch slots are fixed.


\par\noindent\textbf{Hostile-referee risk.} A hostile reader may attempt to treat the relevant content as raw mathematical syntax, ordinary negative theoremhood, surrogate support, or a classical counterexample object rather than endpoint-use material.


\par\noindent\textbf{Closure and denial burden.} The closure response is role-sensitive. If the content is raw, support, bookkeeping, or outside the fixed carrier, it has no endpoint-defeating force. If it keeps endpoint-defeating force on the same carrier, it enters endpoint-use discipline and route exhaustion. A successful objection must preserve the same $X,p,\alpha$, the same Hodge-standing predicate, the same Chow carrier $\CH^p(X)_\QQ$, and the same rational cycle-class map $\clQ$.


\par\noindent\textbf{Use in the proof spine.} The packet is reused only after the local exact-complement assumption has been opened. It marks a role transition, not a hidden assertion of nonempty fiber.


\subsection{Audit packet 53: Surrogate separation (same-mode denial replay)}


\par\noindent\textbf{Claim under audit.} Motives and absolute Hodge data are bridge programs unless they supply Chow fiber.


\par\noindent\textbf{Dependency path.} This packet belongs to the same-mode denial replay. It depends on the official Hodge route, the fixed cycle-fiber endpoint image, the object/use distinction, and the kernel/A+ endpoint-use machinery. It is never used before the carrier and branch slots are fixed.


\par\noindent\textbf{Hostile-referee risk.} A hostile reader may attempt to treat the relevant content as raw mathematical syntax, ordinary negative theoremhood, surrogate support, or a classical counterexample object rather than endpoint-use material.


\par\noindent\textbf{Closure and denial burden.} The closure response is role-sensitive. If the content is raw, support, bookkeeping, or outside the fixed carrier, it has no endpoint-defeating force. If it keeps endpoint-defeating force on the same carrier, it enters endpoint-use discipline and route exhaustion. A successful objection must preserve the same $X,p,\alpha$, the same Hodge-standing predicate, the same Chow carrier $\CH^p(X)_\QQ$, and the same rational cycle-class map $\clQ$.


\par\noindent\textbf{Use in the proof spine.} The packet is reused only after the local exact-complement assumption has been opened. It marks a role transition, not a hidden assertion of nonempty fiber.


\subsection{Audit packet 54: Known-case transfer (same-mode denial replay)}


\par\noindent\textbf{Claim under audit.} Known bridge completions do not transfer universally without target gating.


\par\noindent\textbf{Dependency path.} This packet belongs to the same-mode denial replay. It depends on the official Hodge route, the fixed cycle-fiber endpoint image, the object/use distinction, and the kernel/A+ endpoint-use machinery. It is never used before the carrier and branch slots are fixed.


\par\noindent\textbf{Hostile-referee risk.} A hostile reader may attempt to treat the relevant content as raw mathematical syntax, ordinary negative theoremhood, surrogate support, or a classical counterexample object rather than endpoint-use material.


\par\noindent\textbf{Closure and denial burden.} The closure response is role-sensitive. If the content is raw, support, bookkeeping, or outside the fixed carrier, it has no endpoint-defeating force. If it keeps endpoint-defeating force on the same carrier, it enters endpoint-use discipline and route exhaustion. A successful objection must preserve the same $X,p,\alpha$, the same Hodge-standing predicate, the same Chow carrier $\CH^p(X)_\QQ$, and the same rational cycle-class map $\clQ$.


\par\noindent\textbf{Use in the proof spine.} The packet is reused only after the local exact-complement assumption has been opened. It marks a role transition, not a hidden assertion of nonempty fiber.


\subsection{Audit packet 55: Local-to-global nonfactorization (same-mode denial replay)}


\par\noindent\textbf{Claim under audit.} Local cycle data require a global compatibility certificate.


\par\noindent\textbf{Dependency path.} This packet belongs to the same-mode denial replay. It depends on the official Hodge route, the fixed cycle-fiber endpoint image, the object/use distinction, and the kernel/A+ endpoint-use machinery. It is never used before the carrier and branch slots are fixed.


\par\noindent\textbf{Hostile-referee risk.} A hostile reader may attempt to treat the relevant content as raw mathematical syntax, ordinary negative theoremhood, surrogate support, or a classical counterexample object rather than endpoint-use material.


\par\noindent\textbf{Closure and denial burden.} The closure response is role-sensitive. If the content is raw, support, bookkeeping, or outside the fixed carrier, it has no endpoint-defeating force. If it keeps endpoint-defeating force on the same carrier, it enters endpoint-use discipline and route exhaustion. A successful objection must preserve the same $X,p,\alpha$, the same Hodge-standing predicate, the same Chow carrier $\CH^p(X)_\QQ$, and the same rational cycle-class map $\clQ$.


\par\noindent\textbf{Use in the proof spine.} The packet is reused only after the local exact-complement assumption has been opened. It marks a role transition, not a hidden assertion of nonempty fiber.


\subsection{Audit packet 56: No repair (same-mode denial replay)}


\par\noindent\textbf{Claim under audit.} Later cycle discovery creates a new bridge-complete act, not repair of an old one.


\par\noindent\textbf{Dependency path.} This packet belongs to the same-mode denial replay. It depends on the official Hodge route, the fixed cycle-fiber endpoint image, the object/use distinction, and the kernel/A+ endpoint-use machinery. It is never used before the carrier and branch slots are fixed.


\par\noindent\textbf{Hostile-referee risk.} A hostile reader may attempt to treat the relevant content as raw mathematical syntax, ordinary negative theoremhood, surrogate support, or a classical counterexample object rather than endpoint-use material.


\par\noindent\textbf{Closure and denial burden.} The closure response is role-sensitive. If the content is raw, support, bookkeeping, or outside the fixed carrier, it has no endpoint-defeating force. If it keeps endpoint-defeating force on the same carrier, it enters endpoint-use discipline and route exhaustion. A successful objection must preserve the same $X,p,\alpha$, the same Hodge-standing predicate, the same Chow carrier $\CH^p(X)_\QQ$, and the same rational cycle-class map $\clQ$.


\par\noindent\textbf{Use in the proof spine.} The packet is reused only after the local exact-complement assumption has been opened. It marks a role transition, not a hidden assertion of nonempty fiber.


\subsection{Audit packet 57: No selector (same-mode denial replay)}


\par\noindent\textbf{Claim under audit.} Canonical representatives cannot decide readout before the fiber is fixed.


\par\noindent\textbf{Dependency path.} This packet belongs to the same-mode denial replay. It depends on the official Hodge route, the fixed cycle-fiber endpoint image, the object/use distinction, and the kernel/A+ endpoint-use machinery. It is never used before the carrier and branch slots are fixed.


\par\noindent\textbf{Hostile-referee risk.} A hostile reader may attempt to treat the relevant content as raw mathematical syntax, ordinary negative theoremhood, surrogate support, or a classical counterexample object rather than endpoint-use material.


\par\noindent\textbf{Closure and denial burden.} The closure response is role-sensitive. If the content is raw, support, bookkeeping, or outside the fixed carrier, it has no endpoint-defeating force. If it keeps endpoint-defeating force on the same carrier, it enters endpoint-use discipline and route exhaustion. A successful objection must preserve the same $X,p,\alpha$, the same Hodge-standing predicate, the same Chow carrier $\CH^p(X)_\QQ$, and the same rational cycle-class map $\clQ$.


\par\noindent\textbf{Use in the proof spine.} The packet is reused only after the local exact-complement assumption has been opened. It marks a role transition, not a hidden assertion of nonempty fiber.


\subsection{Audit packet 58: Incompleteness firewall (same-mode denial replay)}


\par\noindent\textbf{Claim under audit.} Independence matters only if it re-enters endpoint standing or trajectory.


\par\noindent\textbf{Dependency path.} This packet belongs to the same-mode denial replay. It depends on the official Hodge route, the fixed cycle-fiber endpoint image, the object/use distinction, and the kernel/A+ endpoint-use machinery. It is never used before the carrier and branch slots are fixed.


\par\noindent\textbf{Hostile-referee risk.} A hostile reader may attempt to treat the relevant content as raw mathematical syntax, ordinary negative theoremhood, surrogate support, or a classical counterexample object rather than endpoint-use material.


\par\noindent\textbf{Closure and denial burden.} The closure response is role-sensitive. If the content is raw, support, bookkeeping, or outside the fixed carrier, it has no endpoint-defeating force. If it keeps endpoint-defeating force on the same carrier, it enters endpoint-use discipline and route exhaustion. A successful objection must preserve the same $X,p,\alpha$, the same Hodge-standing predicate, the same Chow carrier $\CH^p(X)_\QQ$, and the same rational cycle-class map $\clQ$.


\par\noindent\textbf{Use in the proof spine.} The packet is reused only after the local exact-complement assumption has been opened. It marks a role transition, not a hidden assertion of nonempty fiber.


\subsection{Audit packet 59: Ordinary theoremhood protection (same-mode denial replay)}


\par\noindent\textbf{Claim under audit.} Negative theoremhood is proof legitimacy, not forbidden endpoint governance.


\par\noindent\textbf{Dependency path.} This packet belongs to the same-mode denial replay. It depends on the official Hodge route, the fixed cycle-fiber endpoint image, the object/use distinction, and the kernel/A+ endpoint-use machinery. It is never used before the carrier and branch slots are fixed.


\par\noindent\textbf{Hostile-referee risk.} A hostile reader may attempt to treat the relevant content as raw mathematical syntax, ordinary negative theoremhood, surrogate support, or a classical counterexample object rather than endpoint-use material.


\par\noindent\textbf{Closure and denial burden.} The closure response is role-sensitive. If the content is raw, support, bookkeeping, or outside the fixed carrier, it has no endpoint-defeating force. If it keeps endpoint-defeating force on the same carrier, it enters endpoint-use discipline and route exhaustion. A successful objection must preserve the same $X,p,\alpha$, the same Hodge-standing predicate, the same Chow carrier $\CH^p(X)_\QQ$, and the same rational cycle-class map $\clQ$.


\par\noindent\textbf{Use in the proof spine.} The packet is reused only after the local exact-complement assumption has been opened. It marks a role transition, not a hidden assertion of nonempty fiber.


\subsection{Audit packet 60: Non-explosiveness (same-mode denial replay)}


\par\noindent\textbf{Claim under audit.} Empty fiber is lawful until endpoint force is added.


\par\noindent\textbf{Dependency path.} This packet belongs to the same-mode denial replay. It depends on the official Hodge route, the fixed cycle-fiber endpoint image, the object/use distinction, and the kernel/A+ endpoint-use machinery. It is never used before the carrier and branch slots are fixed.


\par\noindent\textbf{Hostile-referee risk.} A hostile reader may attempt to treat the relevant content as raw mathematical syntax, ordinary negative theoremhood, surrogate support, or a classical counterexample object rather than endpoint-use material.


\par\noindent\textbf{Closure and denial burden.} The closure response is role-sensitive. If the content is raw, support, bookkeeping, or outside the fixed carrier, it has no endpoint-defeating force. If it keeps endpoint-defeating force on the same carrier, it enters endpoint-use discipline and route exhaustion. A successful objection must preserve the same $X,p,\alpha$, the same Hodge-standing predicate, the same Chow carrier $\CH^p(X)_\QQ$, and the same rational cycle-class map $\clQ$.


\par\noindent\textbf{Use in the proof spine.} The packet is reused only after the local exact-complement assumption has been opened. It marks a role transition, not a hidden assertion of nonempty fiber.


\subsection{Audit packet 61: Context non-smuggling (endpoint-use replay)}


\par\noindent\textbf{Claim under audit.} Hodge context fixes carrier and slots without branch occupation.


\par\noindent\textbf{Dependency path.} This packet belongs to the endpoint-use replay. It depends on the official Hodge route, the fixed cycle-fiber endpoint image, the object/use distinction, and the kernel/A+ endpoint-use machinery. It is never used before the carrier and branch slots are fixed.


\par\noindent\textbf{Hostile-referee risk.} A hostile reader may attempt to treat the relevant content as raw mathematical syntax, ordinary negative theoremhood, surrogate support, or a classical counterexample object rather than endpoint-use material.


\par\noindent\textbf{Closure and denial burden.} The closure response is role-sensitive. If the content is raw, support, bookkeeping, or outside the fixed carrier, it has no endpoint-defeating force. If it keeps endpoint-defeating force on the same carrier, it enters endpoint-use discipline and route exhaustion. A successful objection must preserve the same $X,p,\alpha$, the same Hodge-standing predicate, the same Chow carrier $\CH^p(X)_\QQ$, and the same rational cycle-class map $\clQ$.


\par\noindent\textbf{Use in the proof spine.} The packet is reused only after the local exact-complement assumption has been opened. It marks a role transition, not a hidden assertion of nonempty fiber.


\subsection{Audit packet 62: Model adequacy (endpoint-use replay)}


\par\noindent\textbf{Claim under audit.} Readout and counterreadout are identified with nonempty and empty cycle fiber.


\par\noindent\textbf{Dependency path.} This packet belongs to the endpoint-use replay. It depends on the official Hodge route, the fixed cycle-fiber endpoint image, the object/use distinction, and the kernel/A+ endpoint-use machinery. It is never used before the carrier and branch slots are fixed.


\par\noindent\textbf{Hostile-referee risk.} A hostile reader may attempt to treat the relevant content as raw mathematical syntax, ordinary negative theoremhood, surrogate support, or a classical counterexample object rather than endpoint-use material.


\par\noindent\textbf{Closure and denial burden.} The closure response is role-sensitive. If the content is raw, support, bookkeeping, or outside the fixed carrier, it has no endpoint-defeating force. If it keeps endpoint-defeating force on the same carrier, it enters endpoint-use discipline and route exhaustion. A successful objection must preserve the same $X,p,\alpha$, the same Hodge-standing predicate, the same Chow carrier $\CH^p(X)_\QQ$, and the same rational cycle-class map $\clQ$.


\par\noindent\textbf{Use in the proof spine.} The packet is reused only after the local exact-complement assumption has been opened. It marks a role transition, not a hidden assertion of nonempty fiber.


\subsection{Audit packet 63: Counterexample force as endpoint use (endpoint-use replay)}


\par\noindent\textbf{Claim under audit.} Official endpoint defeat is endpoint use.


\par\noindent\textbf{Dependency path.} This packet belongs to the endpoint-use replay. It depends on the official Hodge route, the fixed cycle-fiber endpoint image, the object/use distinction, and the kernel/A+ endpoint-use machinery. It is never used before the carrier and branch slots are fixed.


\par\noindent\textbf{Hostile-referee risk.} A hostile reader may attempt to treat the relevant content as raw mathematical syntax, ordinary negative theoremhood, surrogate support, or a classical counterexample object rather than endpoint-use material.


\par\noindent\textbf{Closure and denial burden.} The closure response is role-sensitive. If the content is raw, support, bookkeeping, or outside the fixed carrier, it has no endpoint-defeating force. If it keeps endpoint-defeating force on the same carrier, it enters endpoint-use discipline and route exhaustion. A successful objection must preserve the same $X,p,\alpha$, the same Hodge-standing predicate, the same Chow carrier $\CH^p(X)_\QQ$, and the same rational cycle-class map $\clQ$.


\par\noindent\textbf{Use in the proof spine.} The packet is reused only after the local exact-complement assumption has been opened. It marks a role transition, not a hidden assertion of nonempty fiber.


\subsection{Audit packet 64: No autonomous counterexample role (endpoint-use replay)}


\par\noindent\textbf{Claim under audit.} There is no endpoint-defeating branch outside endpoint use.


\par\noindent\textbf{Dependency path.} This packet belongs to the endpoint-use replay. It depends on the official Hodge route, the fixed cycle-fiber endpoint image, the object/use distinction, and the kernel/A+ endpoint-use machinery. It is never used before the carrier and branch slots are fixed.


\par\noindent\textbf{Hostile-referee risk.} A hostile reader may attempt to treat the relevant content as raw mathematical syntax, ordinary negative theoremhood, surrogate support, or a classical counterexample object rather than endpoint-use material.


\par\noindent\textbf{Closure and denial burden.} The closure response is role-sensitive. If the content is raw, support, bookkeeping, or outside the fixed carrier, it has no endpoint-defeating force. If it keeps endpoint-defeating force on the same carrier, it enters endpoint-use discipline and route exhaustion. A successful objection must preserve the same $X,p,\alpha$, the same Hodge-standing predicate, the same Chow carrier $\CH^p(X)_\QQ$, and the same rational cycle-class map $\clQ$.


\par\noindent\textbf{Use in the proof spine.} The packet is reused only after the local exact-complement assumption has been opened. It marks a role transition, not a hidden assertion of nonempty fiber.


\subsection{Audit packet 65: ATS locking (endpoint-use replay)}


\par\noindent\textbf{Claim under audit.} The same endpoint act preserves anchor and tensor data.


\par\noindent\textbf{Dependency path.} This packet belongs to the endpoint-use replay. It depends on the official Hodge route, the fixed cycle-fiber endpoint image, the object/use distinction, and the kernel/A+ endpoint-use machinery. It is never used before the carrier and branch slots are fixed.


\par\noindent\textbf{Hostile-referee risk.} A hostile reader may attempt to treat the relevant content as raw mathematical syntax, ordinary negative theoremhood, surrogate support, or a classical counterexample object rather than endpoint-use material.


\par\noindent\textbf{Closure and denial burden.} The closure response is role-sensitive. If the content is raw, support, bookkeeping, or outside the fixed carrier, it has no endpoint-defeating force. If it keeps endpoint-defeating force on the same carrier, it enters endpoint-use discipline and route exhaustion. A successful objection must preserve the same $X,p,\alpha$, the same Hodge-standing predicate, the same Chow carrier $\CH^p(X)_\QQ$, and the same rational cycle-class map $\clQ$.


\par\noindent\textbf{Use in the proof spine.} The packet is reused only after the local exact-complement assumption has been opened. It marks a role transition, not a hidden assertion of nonempty fiber.


\subsection{Audit packet 66: UEAP no-overreporting (endpoint-use replay)}


\par\noindent\textbf{Claim under audit.} Support-level content cannot be reported as endpoint readout or endpoint defeat.


\par\noindent\textbf{Dependency path.} This packet belongs to the endpoint-use replay. It depends on the official Hodge route, the fixed cycle-fiber endpoint image, the object/use distinction, and the kernel/A+ endpoint-use machinery. It is never used before the carrier and branch slots are fixed.


\par\noindent\textbf{Hostile-referee risk.} A hostile reader may attempt to treat the relevant content as raw mathematical syntax, ordinary negative theoremhood, surrogate support, or a classical counterexample object rather than endpoint-use material.


\par\noindent\textbf{Closure and denial burden.} The closure response is role-sensitive. If the content is raw, support, bookkeeping, or outside the fixed carrier, it has no endpoint-defeating force. If it keeps endpoint-defeating force on the same carrier, it enters endpoint-use discipline and route exhaustion. A successful objection must preserve the same $X,p,\alpha$, the same Hodge-standing predicate, the same Chow carrier $\CH^p(X)_\QQ$, and the same rational cycle-class map $\clQ$.


\par\noindent\textbf{Use in the proof spine.} The packet is reused only after the local exact-complement assumption has been opened. It marks a role transition, not a hidden assertion of nonempty fiber.


\subsection{Audit packet 67: Single cycle-fiber interface (endpoint-use replay)}


\par\noindent\textbf{Claim under audit.} Cycle-fiber occupation is the only positive endpoint-image interface.


\par\noindent\textbf{Dependency path.} This packet belongs to the endpoint-use replay. It depends on the official Hodge route, the fixed cycle-fiber endpoint image, the object/use distinction, and the kernel/A+ endpoint-use machinery. It is never used before the carrier and branch slots are fixed.


\par\noindent\textbf{Hostile-referee risk.} A hostile reader may attempt to treat the relevant content as raw mathematical syntax, ordinary negative theoremhood, surrogate support, or a classical counterexample object rather than endpoint-use material.


\par\noindent\textbf{Closure and denial burden.} The closure response is role-sensitive. If the content is raw, support, bookkeeping, or outside the fixed carrier, it has no endpoint-defeating force. If it keeps endpoint-defeating force on the same carrier, it enters endpoint-use discipline and route exhaustion. A successful objection must preserve the same $X,p,\alpha$, the same Hodge-standing predicate, the same Chow carrier $\CH^p(X)_\QQ$, and the same rational cycle-class map $\clQ$.


\par\noindent\textbf{Use in the proof spine.} The packet is reused only after the local exact-complement assumption has been opened. It marks a role transition, not a hidden assertion of nonempty fiber.


\subsection{Audit packet 68: Native counterforce exhaustion (endpoint-use replay)}


\par\noindent\textbf{Claim under audit.} Native empty-fiber force has no fifth role.


\par\noindent\textbf{Dependency path.} This packet belongs to the endpoint-use replay. It depends on the official Hodge route, the fixed cycle-fiber endpoint image, the object/use distinction, and the kernel/A+ endpoint-use machinery. It is never used before the carrier and branch slots are fixed.


\par\noindent\textbf{Hostile-referee risk.} A hostile reader may attempt to treat the relevant content as raw mathematical syntax, ordinary negative theoremhood, surrogate support, or a classical counterexample object rather than endpoint-use material.


\par\noindent\textbf{Closure and denial burden.} The closure response is role-sensitive. If the content is raw, support, bookkeeping, or outside the fixed carrier, it has no endpoint-defeating force. If it keeps endpoint-defeating force on the same carrier, it enters endpoint-use discipline and route exhaustion. A successful objection must preserve the same $X,p,\alpha$, the same Hodge-standing predicate, the same Chow carrier $\CH^p(X)_\QQ$, and the same rational cycle-class map $\clQ$.


\par\noindent\textbf{Use in the proof spine.} The packet is reused only after the local exact-complement assumption has been opened. It marks a role transition, not a hidden assertion of nonempty fiber.


\subsection{Audit packet 69: Separator induction (endpoint-use replay)}


\par\noindent\textbf{Claim under audit.} Local separator occupation induces the fiber discriminator.


\par\noindent\textbf{Dependency path.} This packet belongs to the endpoint-use replay. It depends on the official Hodge route, the fixed cycle-fiber endpoint image, the object/use distinction, and the kernel/A+ endpoint-use machinery. It is never used before the carrier and branch slots are fixed.


\par\noindent\textbf{Hostile-referee risk.} A hostile reader may attempt to treat the relevant content as raw mathematical syntax, ordinary negative theoremhood, surrogate support, or a classical counterexample object rather than endpoint-use material.


\par\noindent\textbf{Closure and denial burden.} The closure response is role-sensitive. If the content is raw, support, bookkeeping, or outside the fixed carrier, it has no endpoint-defeating force. If it keeps endpoint-defeating force on the same carrier, it enters endpoint-use discipline and route exhaustion. A successful objection must preserve the same $X,p,\alpha$, the same Hodge-standing predicate, the same Chow carrier $\CH^p(X)_\QQ$, and the same rational cycle-class map $\clQ$.


\par\noindent\textbf{Use in the proof spine.} The packet is reused only after the local exact-complement assumption has been opened. It marks a role transition, not a hidden assertion of nonempty fiber.


\subsection{Audit packet 70: A+ no-independent closure (endpoint-use replay)}


\par\noindent\textbf{Claim under audit.} Endpoint use excludes independent same-domain discriminators.


\par\noindent\textbf{Dependency path.} This packet belongs to the endpoint-use replay. It depends on the official Hodge route, the fixed cycle-fiber endpoint image, the object/use distinction, and the kernel/A+ endpoint-use machinery. It is never used before the carrier and branch slots are fixed.


\par\noindent\textbf{Hostile-referee risk.} A hostile reader may attempt to treat the relevant content as raw mathematical syntax, ordinary negative theoremhood, surrogate support, or a classical counterexample object rather than endpoint-use material.


\par\noindent\textbf{Closure and denial burden.} The closure response is role-sensitive. If the content is raw, support, bookkeeping, or outside the fixed carrier, it has no endpoint-defeating force. If it keeps endpoint-defeating force on the same carrier, it enters endpoint-use discipline and route exhaustion. A successful objection must preserve the same $X,p,\alpha$, the same Hodge-standing predicate, the same Chow carrier $\CH^p(X)_\QQ$, and the same rational cycle-class map $\clQ$.


\par\noindent\textbf{Use in the proof spine.} The packet is reused only after the local exact-complement assumption has been opened. It marks a role transition, not a hidden assertion of nonempty fiber.


\subsection{Audit packet 71: Annotation discharge (endpoint-use replay)}


\par\noindent\textbf{Claim under audit.} The exact-complement annotation is not an added premise.


\par\noindent\textbf{Dependency path.} This packet belongs to the endpoint-use replay. It depends on the official Hodge route, the fixed cycle-fiber endpoint image, the object/use distinction, and the kernel/A+ endpoint-use machinery. It is never used before the carrier and branch slots are fixed.


\par\noindent\textbf{Hostile-referee risk.} A hostile reader may attempt to treat the relevant content as raw mathematical syntax, ordinary negative theoremhood, surrogate support, or a classical counterexample object rather than endpoint-use material.


\par\noindent\textbf{Closure and denial burden.} The closure response is role-sensitive. If the content is raw, support, bookkeeping, or outside the fixed carrier, it has no endpoint-defeating force. If it keeps endpoint-defeating force on the same carrier, it enters endpoint-use discipline and route exhaustion. A successful objection must preserve the same $X,p,\alpha$, the same Hodge-standing predicate, the same Chow carrier $\CH^p(X)_\QQ$, and the same rational cycle-class map $\clQ$.


\par\noindent\textbf{Use in the proof spine.} The packet is reused only after the local exact-complement assumption has been opened. It marks a role transition, not a hidden assertion of nonempty fiber.


\subsection{Audit packet 72: Universalization (endpoint-use replay)}


\par\noindent\textbf{Claim under audit.} Arbitrary instance closure yields the global Hodge endpoint.


\par\noindent\textbf{Dependency path.} This packet belongs to the endpoint-use replay. It depends on the official Hodge route, the fixed cycle-fiber endpoint image, the object/use distinction, and the kernel/A+ endpoint-use machinery. It is never used before the carrier and branch slots are fixed.


\par\noindent\textbf{Hostile-referee risk.} A hostile reader may attempt to treat the relevant content as raw mathematical syntax, ordinary negative theoremhood, surrogate support, or a classical counterexample object rather than endpoint-use material.


\par\noindent\textbf{Closure and denial burden.} The closure response is role-sensitive. If the content is raw, support, bookkeeping, or outside the fixed carrier, it has no endpoint-defeating force. If it keeps endpoint-defeating force on the same carrier, it enters endpoint-use discipline and route exhaustion. A successful objection must preserve the same $X,p,\alpha$, the same Hodge-standing predicate, the same Chow carrier $\CH^p(X)_\QQ$, and the same rational cycle-class map $\clQ$.


\par\noindent\textbf{Use in the proof spine.} The packet is reused only after the local exact-complement assumption has been opened. It marks a role transition, not a hidden assertion of nonempty fiber.


\subsection{Audit packet 73: Surrogate separation (endpoint-use replay)}


\par\noindent\textbf{Claim under audit.} Motives and absolute Hodge data are bridge programs unless they supply Chow fiber.


\par\noindent\textbf{Dependency path.} This packet belongs to the endpoint-use replay. It depends on the official Hodge route, the fixed cycle-fiber endpoint image, the object/use distinction, and the kernel/A+ endpoint-use machinery. It is never used before the carrier and branch slots are fixed.


\par\noindent\textbf{Hostile-referee risk.} A hostile reader may attempt to treat the relevant content as raw mathematical syntax, ordinary negative theoremhood, surrogate support, or a classical counterexample object rather than endpoint-use material.


\par\noindent\textbf{Closure and denial burden.} The closure response is role-sensitive. If the content is raw, support, bookkeeping, or outside the fixed carrier, it has no endpoint-defeating force. If it keeps endpoint-defeating force on the same carrier, it enters endpoint-use discipline and route exhaustion. A successful objection must preserve the same $X,p,\alpha$, the same Hodge-standing predicate, the same Chow carrier $\CH^p(X)_\QQ$, and the same rational cycle-class map $\clQ$.


\par\noindent\textbf{Use in the proof spine.} The packet is reused only after the local exact-complement assumption has been opened. It marks a role transition, not a hidden assertion of nonempty fiber.


\subsection{Audit packet 74: Known-case transfer (endpoint-use replay)}


\par\noindent\textbf{Claim under audit.} Known bridge completions do not transfer universally without target gating.


\par\noindent\textbf{Dependency path.} This packet belongs to the endpoint-use replay. It depends on the official Hodge route, the fixed cycle-fiber endpoint image, the object/use distinction, and the kernel/A+ endpoint-use machinery. It is never used before the carrier and branch slots are fixed.


\par\noindent\textbf{Hostile-referee risk.} A hostile reader may attempt to treat the relevant content as raw mathematical syntax, ordinary negative theoremhood, surrogate support, or a classical counterexample object rather than endpoint-use material.


\par\noindent\textbf{Closure and denial burden.} The closure response is role-sensitive. If the content is raw, support, bookkeeping, or outside the fixed carrier, it has no endpoint-defeating force. If it keeps endpoint-defeating force on the same carrier, it enters endpoint-use discipline and route exhaustion. A successful objection must preserve the same $X,p,\alpha$, the same Hodge-standing predicate, the same Chow carrier $\CH^p(X)_\QQ$, and the same rational cycle-class map $\clQ$.


\par\noindent\textbf{Use in the proof spine.} The packet is reused only after the local exact-complement assumption has been opened. It marks a role transition, not a hidden assertion of nonempty fiber.


\subsection{Audit packet 75: Local-to-global nonfactorization (endpoint-use replay)}


\par\noindent\textbf{Claim under audit.} Local cycle data require a global compatibility certificate.


\par\noindent\textbf{Dependency path.} This packet belongs to the endpoint-use replay. It depends on the official Hodge route, the fixed cycle-fiber endpoint image, the object/use distinction, and the kernel/A+ endpoint-use machinery. It is never used before the carrier and branch slots are fixed.


\par\noindent\textbf{Hostile-referee risk.} A hostile reader may attempt to treat the relevant content as raw mathematical syntax, ordinary negative theoremhood, surrogate support, or a classical counterexample object rather than endpoint-use material.


\par\noindent\textbf{Closure and denial burden.} The closure response is role-sensitive. If the content is raw, support, bookkeeping, or outside the fixed carrier, it has no endpoint-defeating force. If it keeps endpoint-defeating force on the same carrier, it enters endpoint-use discipline and route exhaustion. A successful objection must preserve the same $X,p,\alpha$, the same Hodge-standing predicate, the same Chow carrier $\CH^p(X)_\QQ$, and the same rational cycle-class map $\clQ$.


\par\noindent\textbf{Use in the proof spine.} The packet is reused only after the local exact-complement assumption has been opened. It marks a role transition, not a hidden assertion of nonempty fiber.


\subsection{Audit packet 76: No repair (endpoint-use replay)}


\par\noindent\textbf{Claim under audit.} Later cycle discovery creates a new bridge-complete act, not repair of an old one.


\par\noindent\textbf{Dependency path.} This packet belongs to the endpoint-use replay. It depends on the official Hodge route, the fixed cycle-fiber endpoint image, the object/use distinction, and the kernel/A+ endpoint-use machinery. It is never used before the carrier and branch slots are fixed.


\par\noindent\textbf{Hostile-referee risk.} A hostile reader may attempt to treat the relevant content as raw mathematical syntax, ordinary negative theoremhood, surrogate support, or a classical counterexample object rather than endpoint-use material.


\par\noindent\textbf{Closure and denial burden.} The closure response is role-sensitive. If the content is raw, support, bookkeeping, or outside the fixed carrier, it has no endpoint-defeating force. If it keeps endpoint-defeating force on the same carrier, it enters endpoint-use discipline and route exhaustion. A successful objection must preserve the same $X,p,\alpha$, the same Hodge-standing predicate, the same Chow carrier $\CH^p(X)_\QQ$, and the same rational cycle-class map $\clQ$.


\par\noindent\textbf{Use in the proof spine.} The packet is reused only after the local exact-complement assumption has been opened. It marks a role transition, not a hidden assertion of nonempty fiber.


\subsection{Audit packet 77: No selector (endpoint-use replay)}


\par\noindent\textbf{Claim under audit.} Canonical representatives cannot decide readout before the fiber is fixed.


\par\noindent\textbf{Dependency path.} This packet belongs to the endpoint-use replay. It depends on the official Hodge route, the fixed cycle-fiber endpoint image, the object/use distinction, and the kernel/A+ endpoint-use machinery. It is never used before the carrier and branch slots are fixed.


\par\noindent\textbf{Hostile-referee risk.} A hostile reader may attempt to treat the relevant content as raw mathematical syntax, ordinary negative theoremhood, surrogate support, or a classical counterexample object rather than endpoint-use material.


\par\noindent\textbf{Closure and denial burden.} The closure response is role-sensitive. If the content is raw, support, bookkeeping, or outside the fixed carrier, it has no endpoint-defeating force. If it keeps endpoint-defeating force on the same carrier, it enters endpoint-use discipline and route exhaustion. A successful objection must preserve the same $X,p,\alpha$, the same Hodge-standing predicate, the same Chow carrier $\CH^p(X)_\QQ$, and the same rational cycle-class map $\clQ$.


\par\noindent\textbf{Use in the proof spine.} The packet is reused only after the local exact-complement assumption has been opened. It marks a role transition, not a hidden assertion of nonempty fiber.


\subsection{Audit packet 78: Incompleteness firewall (endpoint-use replay)}


\par\noindent\textbf{Claim under audit.} Independence matters only if it re-enters endpoint standing or trajectory.


\par\noindent\textbf{Dependency path.} This packet belongs to the endpoint-use replay. It depends on the official Hodge route, the fixed cycle-fiber endpoint image, the object/use distinction, and the kernel/A+ endpoint-use machinery. It is never used before the carrier and branch slots are fixed.


\par\noindent\textbf{Hostile-referee risk.} A hostile reader may attempt to treat the relevant content as raw mathematical syntax, ordinary negative theoremhood, surrogate support, or a classical counterexample object rather than endpoint-use material.


\par\noindent\textbf{Closure and denial burden.} The closure response is role-sensitive. If the content is raw, support, bookkeeping, or outside the fixed carrier, it has no endpoint-defeating force. If it keeps endpoint-defeating force on the same carrier, it enters endpoint-use discipline and route exhaustion. A successful objection must preserve the same $X,p,\alpha$, the same Hodge-standing predicate, the same Chow carrier $\CH^p(X)_\QQ$, and the same rational cycle-class map $\clQ$.


\par\noindent\textbf{Use in the proof spine.} The packet is reused only after the local exact-complement assumption has been opened. It marks a role transition, not a hidden assertion of nonempty fiber.


\subsection{Audit packet 79: Ordinary theoremhood protection (endpoint-use replay)}


\par\noindent\textbf{Claim under audit.} Negative theoremhood is proof legitimacy, not forbidden endpoint governance.


\par\noindent\textbf{Dependency path.} This packet belongs to the endpoint-use replay. It depends on the official Hodge route, the fixed cycle-fiber endpoint image, the object/use distinction, and the kernel/A+ endpoint-use machinery. It is never used before the carrier and branch slots are fixed.


\par\noindent\textbf{Hostile-referee risk.} A hostile reader may attempt to treat the relevant content as raw mathematical syntax, ordinary negative theoremhood, surrogate support, or a classical counterexample object rather than endpoint-use material.


\par\noindent\textbf{Closure and denial burden.} The closure response is role-sensitive. If the content is raw, support, bookkeeping, or outside the fixed carrier, it has no endpoint-defeating force. If it keeps endpoint-defeating force on the same carrier, it enters endpoint-use discipline and route exhaustion. A successful objection must preserve the same $X,p,\alpha$, the same Hodge-standing predicate, the same Chow carrier $\CH^p(X)_\QQ$, and the same rational cycle-class map $\clQ$.


\par\noindent\textbf{Use in the proof spine.} The packet is reused only after the local exact-complement assumption has been opened. It marks a role transition, not a hidden assertion of nonempty fiber.


\subsection{Audit packet 80: Non-explosiveness (endpoint-use replay)}


\par\noindent\textbf{Claim under audit.} Empty fiber is lawful until endpoint force is added.


\par\noindent\textbf{Dependency path.} This packet belongs to the endpoint-use replay. It depends on the official Hodge route, the fixed cycle-fiber endpoint image, the object/use distinction, and the kernel/A+ endpoint-use machinery. It is never used before the carrier and branch slots are fixed.


\par\noindent\textbf{Hostile-referee risk.} A hostile reader may attempt to treat the relevant content as raw mathematical syntax, ordinary negative theoremhood, surrogate support, or a classical counterexample object rather than endpoint-use material.


\par\noindent\textbf{Closure and denial burden.} The closure response is role-sensitive. If the content is raw, support, bookkeeping, or outside the fixed carrier, it has no endpoint-defeating force. If it keeps endpoint-defeating force on the same carrier, it enters endpoint-use discipline and route exhaustion. A successful objection must preserve the same $X,p,\alpha$, the same Hodge-standing predicate, the same Chow carrier $\CH^p(X)_\QQ$, and the same rational cycle-class map $\clQ$.


\par\noindent\textbf{Use in the proof spine.} The packet is reused only after the local exact-complement assumption has been opened. It marks a role transition, not a hidden assertion of nonempty fiber.


\subsection{Audit packet 81: Context non-smuggling (Lean-audit replay)}


\par\noindent\textbf{Claim under audit.} Hodge context fixes carrier and slots without branch occupation.


\par\noindent\textbf{Dependency path.} This packet belongs to the Lean-audit replay. It depends on the official Hodge route, the fixed cycle-fiber endpoint image, the object/use distinction, and the kernel/A+ endpoint-use machinery. It is never used before the carrier and branch slots are fixed.


\par\noindent\textbf{Hostile-referee risk.} A hostile reader may attempt to treat the relevant content as raw mathematical syntax, ordinary negative theoremhood, surrogate support, or a classical counterexample object rather than endpoint-use material.


\par\noindent\textbf{Closure and denial burden.} The closure response is role-sensitive. If the content is raw, support, bookkeeping, or outside the fixed carrier, it has no endpoint-defeating force. If it keeps endpoint-defeating force on the same carrier, it enters endpoint-use discipline and route exhaustion. A successful objection must preserve the same $X,p,\alpha$, the same Hodge-standing predicate, the same Chow carrier $\CH^p(X)_\QQ$, and the same rational cycle-class map $\clQ$.


\par\noindent\textbf{Use in the proof spine.} The packet is reused only after the local exact-complement assumption has been opened. It marks a role transition, not a hidden assertion of nonempty fiber.


\subsection{Audit packet 82: Model adequacy (Lean-audit replay)}


\par\noindent\textbf{Claim under audit.} Readout and counterreadout are identified with nonempty and empty cycle fiber.


\par\noindent\textbf{Dependency path.} This packet belongs to the Lean-audit replay. It depends on the official Hodge route, the fixed cycle-fiber endpoint image, the object/use distinction, and the kernel/A+ endpoint-use machinery. It is never used before the carrier and branch slots are fixed.


\par\noindent\textbf{Hostile-referee risk.} A hostile reader may attempt to treat the relevant content as raw mathematical syntax, ordinary negative theoremhood, surrogate support, or a classical counterexample object rather than endpoint-use material.


\par\noindent\textbf{Closure and denial burden.} The closure response is role-sensitive. If the content is raw, support, bookkeeping, or outside the fixed carrier, it has no endpoint-defeating force. If it keeps endpoint-defeating force on the same carrier, it enters endpoint-use discipline and route exhaustion. A successful objection must preserve the same $X,p,\alpha$, the same Hodge-standing predicate, the same Chow carrier $\CH^p(X)_\QQ$, and the same rational cycle-class map $\clQ$.


\par\noindent\textbf{Use in the proof spine.} The packet is reused only after the local exact-complement assumption has been opened. It marks a role transition, not a hidden assertion of nonempty fiber.


\subsection{Audit packet 83: Counterexample force as endpoint use (Lean-audit replay)}


\par\noindent\textbf{Claim under audit.} Official endpoint defeat is endpoint use.


\par\noindent\textbf{Dependency path.} This packet belongs to the Lean-audit replay. It depends on the official Hodge route, the fixed cycle-fiber endpoint image, the object/use distinction, and the kernel/A+ endpoint-use machinery. It is never used before the carrier and branch slots are fixed.


\par\noindent\textbf{Hostile-referee risk.} A hostile reader may attempt to treat the relevant content as raw mathematical syntax, ordinary negative theoremhood, surrogate support, or a classical counterexample object rather than endpoint-use material.


\par\noindent\textbf{Closure and denial burden.} The closure response is role-sensitive. If the content is raw, support, bookkeeping, or outside the fixed carrier, it has no endpoint-defeating force. If it keeps endpoint-defeating force on the same carrier, it enters endpoint-use discipline and route exhaustion. A successful objection must preserve the same $X,p,\alpha$, the same Hodge-standing predicate, the same Chow carrier $\CH^p(X)_\QQ$, and the same rational cycle-class map $\clQ$.


\par\noindent\textbf{Use in the proof spine.} The packet is reused only after the local exact-complement assumption has been opened. It marks a role transition, not a hidden assertion of nonempty fiber.


\subsection{Audit packet 84: No autonomous counterexample role (Lean-audit replay)}


\par\noindent\textbf{Claim under audit.} There is no endpoint-defeating branch outside endpoint use.


\par\noindent\textbf{Dependency path.} This packet belongs to the Lean-audit replay. It depends on the official Hodge route, the fixed cycle-fiber endpoint image, the object/use distinction, and the kernel/A+ endpoint-use machinery. It is never used before the carrier and branch slots are fixed.


\par\noindent\textbf{Hostile-referee risk.} A hostile reader may attempt to treat the relevant content as raw mathematical syntax, ordinary negative theoremhood, surrogate support, or a classical counterexample object rather than endpoint-use material.


\par\noindent\textbf{Closure and denial burden.} The closure response is role-sensitive. If the content is raw, support, bookkeeping, or outside the fixed carrier, it has no endpoint-defeating force. If it keeps endpoint-defeating force on the same carrier, it enters endpoint-use discipline and route exhaustion. A successful objection must preserve the same $X,p,\alpha$, the same Hodge-standing predicate, the same Chow carrier $\CH^p(X)_\QQ$, and the same rational cycle-class map $\clQ$.


\par\noindent\textbf{Use in the proof spine.} The packet is reused only after the local exact-complement assumption has been opened. It marks a role transition, not a hidden assertion of nonempty fiber.


\subsection{Audit packet 85: ATS locking (Lean-audit replay)}


\par\noindent\textbf{Claim under audit.} The same endpoint act preserves anchor and tensor data.


\par\noindent\textbf{Dependency path.} This packet belongs to the Lean-audit replay. It depends on the official Hodge route, the fixed cycle-fiber endpoint image, the object/use distinction, and the kernel/A+ endpoint-use machinery. It is never used before the carrier and branch slots are fixed.


\par\noindent\textbf{Hostile-referee risk.} A hostile reader may attempt to treat the relevant content as raw mathematical syntax, ordinary negative theoremhood, surrogate support, or a classical counterexample object rather than endpoint-use material.


\par\noindent\textbf{Closure and denial burden.} The closure response is role-sensitive. If the content is raw, support, bookkeeping, or outside the fixed carrier, it has no endpoint-defeating force. If it keeps endpoint-defeating force on the same carrier, it enters endpoint-use discipline and route exhaustion. A successful objection must preserve the same $X,p,\alpha$, the same Hodge-standing predicate, the same Chow carrier $\CH^p(X)_\QQ$, and the same rational cycle-class map $\clQ$.


\par\noindent\textbf{Use in the proof spine.} The packet is reused only after the local exact-complement assumption has been opened. It marks a role transition, not a hidden assertion of nonempty fiber.


\subsection{Audit packet 86: UEAP no-overreporting (Lean-audit replay)}


\par\noindent\textbf{Claim under audit.} Support-level content cannot be reported as endpoint readout or endpoint defeat.


\par\noindent\textbf{Dependency path.} This packet belongs to the Lean-audit replay. It depends on the official Hodge route, the fixed cycle-fiber endpoint image, the object/use distinction, and the kernel/A+ endpoint-use machinery. It is never used before the carrier and branch slots are fixed.


\par\noindent\textbf{Hostile-referee risk.} A hostile reader may attempt to treat the relevant content as raw mathematical syntax, ordinary negative theoremhood, surrogate support, or a classical counterexample object rather than endpoint-use material.


\par\noindent\textbf{Closure and denial burden.} The closure response is role-sensitive. If the content is raw, support, bookkeeping, or outside the fixed carrier, it has no endpoint-defeating force. If it keeps endpoint-defeating force on the same carrier, it enters endpoint-use discipline and route exhaustion. A successful objection must preserve the same $X,p,\alpha$, the same Hodge-standing predicate, the same Chow carrier $\CH^p(X)_\QQ$, and the same rational cycle-class map $\clQ$.


\par\noindent\textbf{Use in the proof spine.} The packet is reused only after the local exact-complement assumption has been opened. It marks a role transition, not a hidden assertion of nonempty fiber.


\subsection{Audit packet 87: Single cycle-fiber interface (Lean-audit replay)}


\par\noindent\textbf{Claim under audit.} Cycle-fiber occupation is the only positive endpoint-image interface.


\par\noindent\textbf{Dependency path.} This packet belongs to the Lean-audit replay. It depends on the official Hodge route, the fixed cycle-fiber endpoint image, the object/use distinction, and the kernel/A+ endpoint-use machinery. It is never used before the carrier and branch slots are fixed.


\par\noindent\textbf{Hostile-referee risk.} A hostile reader may attempt to treat the relevant content as raw mathematical syntax, ordinary negative theoremhood, surrogate support, or a classical counterexample object rather than endpoint-use material.


\par\noindent\textbf{Closure and denial burden.} The closure response is role-sensitive. If the content is raw, support, bookkeeping, or outside the fixed carrier, it has no endpoint-defeating force. If it keeps endpoint-defeating force on the same carrier, it enters endpoint-use discipline and route exhaustion. A successful objection must preserve the same $X,p,\alpha$, the same Hodge-standing predicate, the same Chow carrier $\CH^p(X)_\QQ$, and the same rational cycle-class map $\clQ$.


\par\noindent\textbf{Use in the proof spine.} The packet is reused only after the local exact-complement assumption has been opened. It marks a role transition, not a hidden assertion of nonempty fiber.


\subsection{Audit packet 88: Native counterforce exhaustion (Lean-audit replay)}


\par\noindent\textbf{Claim under audit.} Native empty-fiber force has no fifth role.


\par\noindent\textbf{Dependency path.} This packet belongs to the Lean-audit replay. It depends on the official Hodge route, the fixed cycle-fiber endpoint image, the object/use distinction, and the kernel/A+ endpoint-use machinery. It is never used before the carrier and branch slots are fixed.


\par\noindent\textbf{Hostile-referee risk.} A hostile reader may attempt to treat the relevant content as raw mathematical syntax, ordinary negative theoremhood, surrogate support, or a classical counterexample object rather than endpoint-use material.


\par\noindent\textbf{Closure and denial burden.} The closure response is role-sensitive. If the content is raw, support, bookkeeping, or outside the fixed carrier, it has no endpoint-defeating force. If it keeps endpoint-defeating force on the same carrier, it enters endpoint-use discipline and route exhaustion. A successful objection must preserve the same $X,p,\alpha$, the same Hodge-standing predicate, the same Chow carrier $\CH^p(X)_\QQ$, and the same rational cycle-class map $\clQ$.


\par\noindent\textbf{Use in the proof spine.} The packet is reused only after the local exact-complement assumption has been opened. It marks a role transition, not a hidden assertion of nonempty fiber.


\subsection{Audit packet 89: Separator induction (Lean-audit replay)}


\par\noindent\textbf{Claim under audit.} Local separator occupation induces the fiber discriminator.


\par\noindent\textbf{Dependency path.} This packet belongs to the Lean-audit replay. It depends on the official Hodge route, the fixed cycle-fiber endpoint image, the object/use distinction, and the kernel/A+ endpoint-use machinery. It is never used before the carrier and branch slots are fixed.


\par\noindent\textbf{Hostile-referee risk.} A hostile reader may attempt to treat the relevant content as raw mathematical syntax, ordinary negative theoremhood, surrogate support, or a classical counterexample object rather than endpoint-use material.


\par\noindent\textbf{Closure and denial burden.} The closure response is role-sensitive. If the content is raw, support, bookkeeping, or outside the fixed carrier, it has no endpoint-defeating force. If it keeps endpoint-defeating force on the same carrier, it enters endpoint-use discipline and route exhaustion. A successful objection must preserve the same $X,p,\alpha$, the same Hodge-standing predicate, the same Chow carrier $\CH^p(X)_\QQ$, and the same rational cycle-class map $\clQ$.


\par\noindent\textbf{Use in the proof spine.} The packet is reused only after the local exact-complement assumption has been opened. It marks a role transition, not a hidden assertion of nonempty fiber.


\subsection{Audit packet 90: A+ no-independent closure (Lean-audit replay)}


\par\noindent\textbf{Claim under audit.} Endpoint use excludes independent same-domain discriminators.


\par\noindent\textbf{Dependency path.} This packet belongs to the Lean-audit replay. It depends on the official Hodge route, the fixed cycle-fiber endpoint image, the object/use distinction, and the kernel/A+ endpoint-use machinery. It is never used before the carrier and branch slots are fixed.


\par\noindent\textbf{Hostile-referee risk.} A hostile reader may attempt to treat the relevant content as raw mathematical syntax, ordinary negative theoremhood, surrogate support, or a classical counterexample object rather than endpoint-use material.


\par\noindent\textbf{Closure and denial burden.} The closure response is role-sensitive. If the content is raw, support, bookkeeping, or outside the fixed carrier, it has no endpoint-defeating force. If it keeps endpoint-defeating force on the same carrier, it enters endpoint-use discipline and route exhaustion. A successful objection must preserve the same $X,p,\alpha$, the same Hodge-standing predicate, the same Chow carrier $\CH^p(X)_\QQ$, and the same rational cycle-class map $\clQ$.


\par\noindent\textbf{Use in the proof spine.} The packet is reused only after the local exact-complement assumption has been opened. It marks a role transition, not a hidden assertion of nonempty fiber.


\subsection{Audit packet 91: Annotation discharge (Lean-audit replay)}


\par\noindent\textbf{Claim under audit.} The exact-complement annotation is not an added premise.


\par\noindent\textbf{Dependency path.} This packet belongs to the Lean-audit replay. It depends on the official Hodge route, the fixed cycle-fiber endpoint image, the object/use distinction, and the kernel/A+ endpoint-use machinery. It is never used before the carrier and branch slots are fixed.


\par\noindent\textbf{Hostile-referee risk.} A hostile reader may attempt to treat the relevant content as raw mathematical syntax, ordinary negative theoremhood, surrogate support, or a classical counterexample object rather than endpoint-use material.


\par\noindent\textbf{Closure and denial burden.} The closure response is role-sensitive. If the content is raw, support, bookkeeping, or outside the fixed carrier, it has no endpoint-defeating force. If it keeps endpoint-defeating force on the same carrier, it enters endpoint-use discipline and route exhaustion. A successful objection must preserve the same $X,p,\alpha$, the same Hodge-standing predicate, the same Chow carrier $\CH^p(X)_\QQ$, and the same rational cycle-class map $\clQ$.


\par\noindent\textbf{Use in the proof spine.} The packet is reused only after the local exact-complement assumption has been opened. It marks a role transition, not a hidden assertion of nonempty fiber.


\subsection{Audit packet 92: Universalization (Lean-audit replay)}


\par\noindent\textbf{Claim under audit.} Arbitrary instance closure yields the global Hodge endpoint.


\par\noindent\textbf{Dependency path.} This packet belongs to the Lean-audit replay. It depends on the official Hodge route, the fixed cycle-fiber endpoint image, the object/use distinction, and the kernel/A+ endpoint-use machinery. It is never used before the carrier and branch slots are fixed.


\par\noindent\textbf{Hostile-referee risk.} A hostile reader may attempt to treat the relevant content as raw mathematical syntax, ordinary negative theoremhood, surrogate support, or a classical counterexample object rather than endpoint-use material.


\par\noindent\textbf{Closure and denial burden.} The closure response is role-sensitive. If the content is raw, support, bookkeeping, or outside the fixed carrier, it has no endpoint-defeating force. If it keeps endpoint-defeating force on the same carrier, it enters endpoint-use discipline and route exhaustion. A successful objection must preserve the same $X,p,\alpha$, the same Hodge-standing predicate, the same Chow carrier $\CH^p(X)_\QQ$, and the same rational cycle-class map $\clQ$.


\par\noindent\textbf{Use in the proof spine.} The packet is reused only after the local exact-complement assumption has been opened. It marks a role transition, not a hidden assertion of nonempty fiber.


\subsection{Audit packet 93: Surrogate separation (Lean-audit replay)}


\par\noindent\textbf{Claim under audit.} Motives and absolute Hodge data are bridge programs unless they supply Chow fiber.


\par\noindent\textbf{Dependency path.} This packet belongs to the Lean-audit replay. It depends on the official Hodge route, the fixed cycle-fiber endpoint image, the object/use distinction, and the kernel/A+ endpoint-use machinery. It is never used before the carrier and branch slots are fixed.


\par\noindent\textbf{Hostile-referee risk.} A hostile reader may attempt to treat the relevant content as raw mathematical syntax, ordinary negative theoremhood, surrogate support, or a classical counterexample object rather than endpoint-use material.


\par\noindent\textbf{Closure and denial burden.} The closure response is role-sensitive. If the content is raw, support, bookkeeping, or outside the fixed carrier, it has no endpoint-defeating force. If it keeps endpoint-defeating force on the same carrier, it enters endpoint-use discipline and route exhaustion. A successful objection must preserve the same $X,p,\alpha$, the same Hodge-standing predicate, the same Chow carrier $\CH^p(X)_\QQ$, and the same rational cycle-class map $\clQ$.


\par\noindent\textbf{Use in the proof spine.} The packet is reused only after the local exact-complement assumption has been opened. It marks a role transition, not a hidden assertion of nonempty fiber.


\subsection{Audit packet 94: Known-case transfer (Lean-audit replay)}


\par\noindent\textbf{Claim under audit.} Known bridge completions do not transfer universally without target gating.


\par\noindent\textbf{Dependency path.} This packet belongs to the Lean-audit replay. It depends on the official Hodge route, the fixed cycle-fiber endpoint image, the object/use distinction, and the kernel/A+ endpoint-use machinery. It is never used before the carrier and branch slots are fixed.


\par\noindent\textbf{Hostile-referee risk.} A hostile reader may attempt to treat the relevant content as raw mathematical syntax, ordinary negative theoremhood, surrogate support, or a classical counterexample object rather than endpoint-use material.


\par\noindent\textbf{Closure and denial burden.} The closure response is role-sensitive. If the content is raw, support, bookkeeping, or outside the fixed carrier, it has no endpoint-defeating force. If it keeps endpoint-defeating force on the same carrier, it enters endpoint-use discipline and route exhaustion. A successful objection must preserve the same $X,p,\alpha$, the same Hodge-standing predicate, the same Chow carrier $\CH^p(X)_\QQ$, and the same rational cycle-class map $\clQ$.


\par\noindent\textbf{Use in the proof spine.} The packet is reused only after the local exact-complement assumption has been opened. It marks a role transition, not a hidden assertion of nonempty fiber.


\subsection{Audit packet 95: Local-to-global nonfactorization (Lean-audit replay)}


\par\noindent\textbf{Claim under audit.} Local cycle data require a global compatibility certificate.


\par\noindent\textbf{Dependency path.} This packet belongs to the Lean-audit replay. It depends on the official Hodge route, the fixed cycle-fiber endpoint image, the object/use distinction, and the kernel/A+ endpoint-use machinery. It is never used before the carrier and branch slots are fixed.


\par\noindent\textbf{Hostile-referee risk.} A hostile reader may attempt to treat the relevant content as raw mathematical syntax, ordinary negative theoremhood, surrogate support, or a classical counterexample object rather than endpoint-use material.


\par\noindent\textbf{Closure and denial burden.} The closure response is role-sensitive. If the content is raw, support, bookkeeping, or outside the fixed carrier, it has no endpoint-defeating force. If it keeps endpoint-defeating force on the same carrier, it enters endpoint-use discipline and route exhaustion. A successful objection must preserve the same $X,p,\alpha$, the same Hodge-standing predicate, the same Chow carrier $\CH^p(X)_\QQ$, and the same rational cycle-class map $\clQ$.


\par\noindent\textbf{Use in the proof spine.} The packet is reused only after the local exact-complement assumption has been opened. It marks a role transition, not a hidden assertion of nonempty fiber.


\subsection{Audit packet 96: No repair (Lean-audit replay)}


\par\noindent\textbf{Claim under audit.} Later cycle discovery creates a new bridge-complete act, not repair of an old one.


\par\noindent\textbf{Dependency path.} This packet belongs to the Lean-audit replay. It depends on the official Hodge route, the fixed cycle-fiber endpoint image, the object/use distinction, and the kernel/A+ endpoint-use machinery. It is never used before the carrier and branch slots are fixed.


\par\noindent\textbf{Hostile-referee risk.} A hostile reader may attempt to treat the relevant content as raw mathematical syntax, ordinary negative theoremhood, surrogate support, or a classical counterexample object rather than endpoint-use material.


\par\noindent\textbf{Closure and denial burden.} The closure response is role-sensitive. If the content is raw, support, bookkeeping, or outside the fixed carrier, it has no endpoint-defeating force. If it keeps endpoint-defeating force on the same carrier, it enters endpoint-use discipline and route exhaustion. A successful objection must preserve the same $X,p,\alpha$, the same Hodge-standing predicate, the same Chow carrier $\CH^p(X)_\QQ$, and the same rational cycle-class map $\clQ$.


\par\noindent\textbf{Use in the proof spine.} The packet is reused only after the local exact-complement assumption has been opened. It marks a role transition, not a hidden assertion of nonempty fiber.


\subsection{Audit packet 97: No selector (Lean-audit replay)}


\par\noindent\textbf{Claim under audit.} Canonical representatives cannot decide readout before the fiber is fixed.


\par\noindent\textbf{Dependency path.} This packet belongs to the Lean-audit replay. It depends on the official Hodge route, the fixed cycle-fiber endpoint image, the object/use distinction, and the kernel/A+ endpoint-use machinery. It is never used before the carrier and branch slots are fixed.


\par\noindent\textbf{Hostile-referee risk.} A hostile reader may attempt to treat the relevant content as raw mathematical syntax, ordinary negative theoremhood, surrogate support, or a classical counterexample object rather than endpoint-use material.


\par\noindent\textbf{Closure and denial burden.} The closure response is role-sensitive. If the content is raw, support, bookkeeping, or outside the fixed carrier, it has no endpoint-defeating force. If it keeps endpoint-defeating force on the same carrier, it enters endpoint-use discipline and route exhaustion. A successful objection must preserve the same $X,p,\alpha$, the same Hodge-standing predicate, the same Chow carrier $\CH^p(X)_\QQ$, and the same rational cycle-class map $\clQ$.


\par\noindent\textbf{Use in the proof spine.} The packet is reused only after the local exact-complement assumption has been opened. It marks a role transition, not a hidden assertion of nonempty fiber.


\subsection{Audit packet 98: Incompleteness firewall (Lean-audit replay)}


\par\noindent\textbf{Claim under audit.} Independence matters only if it re-enters endpoint standing or trajectory.


\par\noindent\textbf{Dependency path.} This packet belongs to the Lean-audit replay. It depends on the official Hodge route, the fixed cycle-fiber endpoint image, the object/use distinction, and the kernel/A+ endpoint-use machinery. It is never used before the carrier and branch slots are fixed.


\par\noindent\textbf{Hostile-referee risk.} A hostile reader may attempt to treat the relevant content as raw mathematical syntax, ordinary negative theoremhood, surrogate support, or a classical counterexample object rather than endpoint-use material.


\par\noindent\textbf{Closure and denial burden.} The closure response is role-sensitive. If the content is raw, support, bookkeeping, or outside the fixed carrier, it has no endpoint-defeating force. If it keeps endpoint-defeating force on the same carrier, it enters endpoint-use discipline and route exhaustion. A successful objection must preserve the same $X,p,\alpha$, the same Hodge-standing predicate, the same Chow carrier $\CH^p(X)_\QQ$, and the same rational cycle-class map $\clQ$.


\par\noindent\textbf{Use in the proof spine.} The packet is reused only after the local exact-complement assumption has been opened. It marks a role transition, not a hidden assertion of nonempty fiber.


\subsection{Audit packet 99: Ordinary theoremhood protection (Lean-audit replay)}


\par\noindent\textbf{Claim under audit.} Negative theoremhood is proof legitimacy, not forbidden endpoint governance.


\par\noindent\textbf{Dependency path.} This packet belongs to the Lean-audit replay. It depends on the official Hodge route, the fixed cycle-fiber endpoint image, the object/use distinction, and the kernel/A+ endpoint-use machinery. It is never used before the carrier and branch slots are fixed.


\par\noindent\textbf{Hostile-referee risk.} A hostile reader may attempt to treat the relevant content as raw mathematical syntax, ordinary negative theoremhood, surrogate support, or a classical counterexample object rather than endpoint-use material.


\par\noindent\textbf{Closure and denial burden.} The closure response is role-sensitive. If the content is raw, support, bookkeeping, or outside the fixed carrier, it has no endpoint-defeating force. If it keeps endpoint-defeating force on the same carrier, it enters endpoint-use discipline and route exhaustion. A successful objection must preserve the same $X,p,\alpha$, the same Hodge-standing predicate, the same Chow carrier $\CH^p(X)_\QQ$, and the same rational cycle-class map $\clQ$.


\par\noindent\textbf{Use in the proof spine.} The packet is reused only after the local exact-complement assumption has been opened. It marks a role transition, not a hidden assertion of nonempty fiber.


\subsection{Audit packet 100: Non-explosiveness (Lean-audit replay)}


\par\noindent\textbf{Claim under audit.} Empty fiber is lawful until endpoint force is added.


\par\noindent\textbf{Dependency path.} This packet belongs to the Lean-audit replay. It depends on the official Hodge route, the fixed cycle-fiber endpoint image, the object/use distinction, and the kernel/A+ endpoint-use machinery. It is never used before the carrier and branch slots are fixed.


\par\noindent\textbf{Hostile-referee risk.} A hostile reader may attempt to treat the relevant content as raw mathematical syntax, ordinary negative theoremhood, surrogate support, or a classical counterexample object rather than endpoint-use material.


\par\noindent\textbf{Closure and denial burden.} The closure response is role-sensitive. If the content is raw, support, bookkeeping, or outside the fixed carrier, it has no endpoint-defeating force. If it keeps endpoint-defeating force on the same carrier, it enters endpoint-use discipline and route exhaustion. A successful objection must preserve the same $X,p,\alpha$, the same Hodge-standing predicate, the same Chow carrier $\CH^p(X)_\QQ$, and the same rational cycle-class map $\clQ$.


\par\noindent\textbf{Use in the proof spine.} The packet is reused only after the local exact-complement assumption has been opened. It marks a role transition, not a hidden assertion of nonempty fiber.


\subsection{Audit packet 101: Context non-smuggling (route-exhaustion replay)}


\par\noindent\textbf{Claim under audit.} Hodge context fixes carrier and slots without branch occupation.


\par\noindent\textbf{Dependency path.} This packet belongs to the route-exhaustion replay. It depends on the official Hodge route, the fixed cycle-fiber endpoint image, the object/use distinction, and the kernel/A+ endpoint-use machinery. It is never used before the carrier and branch slots are fixed.


\par\noindent\textbf{Hostile-referee risk.} A hostile reader may attempt to treat the relevant content as raw mathematical syntax, ordinary negative theoremhood, surrogate support, or a classical counterexample object rather than endpoint-use material.


\par\noindent\textbf{Closure and denial burden.} The closure response is role-sensitive. If the content is raw, support, bookkeeping, or outside the fixed carrier, it has no endpoint-defeating force. If it keeps endpoint-defeating force on the same carrier, it enters endpoint-use discipline and route exhaustion. A successful objection must preserve the same $X,p,\alpha$, the same Hodge-standing predicate, the same Chow carrier $\CH^p(X)_\QQ$, and the same rational cycle-class map $\clQ$.


\par\noindent\textbf{Use in the proof spine.} The packet is reused only after the local exact-complement assumption has been opened. It marks a role transition, not a hidden assertion of nonempty fiber.


\subsection{Audit packet 102: Model adequacy (route-exhaustion replay)}


\par\noindent\textbf{Claim under audit.} Readout and counterreadout are identified with nonempty and empty cycle fiber.


\par\noindent\textbf{Dependency path.} This packet belongs to the route-exhaustion replay. It depends on the official Hodge route, the fixed cycle-fiber endpoint image, the object/use distinction, and the kernel/A+ endpoint-use machinery. It is never used before the carrier and branch slots are fixed.


\par\noindent\textbf{Hostile-referee risk.} A hostile reader may attempt to treat the relevant content as raw mathematical syntax, ordinary negative theoremhood, surrogate support, or a classical counterexample object rather than endpoint-use material.


\par\noindent\textbf{Closure and denial burden.} The closure response is role-sensitive. If the content is raw, support, bookkeeping, or outside the fixed carrier, it has no endpoint-defeating force. If it keeps endpoint-defeating force on the same carrier, it enters endpoint-use discipline and route exhaustion. A successful objection must preserve the same $X,p,\alpha$, the same Hodge-standing predicate, the same Chow carrier $\CH^p(X)_\QQ$, and the same rational cycle-class map $\clQ$.


\par\noindent\textbf{Use in the proof spine.} The packet is reused only after the local exact-complement assumption has been opened. It marks a role transition, not a hidden assertion of nonempty fiber.


\subsection{Audit packet 103: Counterexample force as endpoint use (route-exhaustion replay)}


\par\noindent\textbf{Claim under audit.} Official endpoint defeat is endpoint use.


\par\noindent\textbf{Dependency path.} This packet belongs to the route-exhaustion replay. It depends on the official Hodge route, the fixed cycle-fiber endpoint image, the object/use distinction, and the kernel/A+ endpoint-use machinery. It is never used before the carrier and branch slots are fixed.


\par\noindent\textbf{Hostile-referee risk.} A hostile reader may attempt to treat the relevant content as raw mathematical syntax, ordinary negative theoremhood, surrogate support, or a classical counterexample object rather than endpoint-use material.


\par\noindent\textbf{Closure and denial burden.} The closure response is role-sensitive. If the content is raw, support, bookkeeping, or outside the fixed carrier, it has no endpoint-defeating force. If it keeps endpoint-defeating force on the same carrier, it enters endpoint-use discipline and route exhaustion. A successful objection must preserve the same $X,p,\alpha$, the same Hodge-standing predicate, the same Chow carrier $\CH^p(X)_\QQ$, and the same rational cycle-class map $\clQ$.


\par\noindent\textbf{Use in the proof spine.} The packet is reused only after the local exact-complement assumption has been opened. It marks a role transition, not a hidden assertion of nonempty fiber.


\subsection{Audit packet 104: No autonomous counterexample role (route-exhaustion replay)}


\par\noindent\textbf{Claim under audit.} There is no endpoint-defeating branch outside endpoint use.


\par\noindent\textbf{Dependency path.} This packet belongs to the route-exhaustion replay. It depends on the official Hodge route, the fixed cycle-fiber endpoint image, the object/use distinction, and the kernel/A+ endpoint-use machinery. It is never used before the carrier and branch slots are fixed.


\par\noindent\textbf{Hostile-referee risk.} A hostile reader may attempt to treat the relevant content as raw mathematical syntax, ordinary negative theoremhood, surrogate support, or a classical counterexample object rather than endpoint-use material.


\par\noindent\textbf{Closure and denial burden.} The closure response is role-sensitive. If the content is raw, support, bookkeeping, or outside the fixed carrier, it has no endpoint-defeating force. If it keeps endpoint-defeating force on the same carrier, it enters endpoint-use discipline and route exhaustion. A successful objection must preserve the same $X,p,\alpha$, the same Hodge-standing predicate, the same Chow carrier $\CH^p(X)_\QQ$, and the same rational cycle-class map $\clQ$.


\par\noindent\textbf{Use in the proof spine.} The packet is reused only after the local exact-complement assumption has been opened. It marks a role transition, not a hidden assertion of nonempty fiber.


\subsection{Audit packet 105: ATS locking (route-exhaustion replay)}


\par\noindent\textbf{Claim under audit.} The same endpoint act preserves anchor and tensor data.


\par\noindent\textbf{Dependency path.} This packet belongs to the route-exhaustion replay. It depends on the official Hodge route, the fixed cycle-fiber endpoint image, the object/use distinction, and the kernel/A+ endpoint-use machinery. It is never used before the carrier and branch slots are fixed.


\par\noindent\textbf{Hostile-referee risk.} A hostile reader may attempt to treat the relevant content as raw mathematical syntax, ordinary negative theoremhood, surrogate support, or a classical counterexample object rather than endpoint-use material.


\par\noindent\textbf{Closure and denial burden.} The closure response is role-sensitive. If the content is raw, support, bookkeeping, or outside the fixed carrier, it has no endpoint-defeating force. If it keeps endpoint-defeating force on the same carrier, it enters endpoint-use discipline and route exhaustion. A successful objection must preserve the same $X,p,\alpha$, the same Hodge-standing predicate, the same Chow carrier $\CH^p(X)_\QQ$, and the same rational cycle-class map $\clQ$.


\par\noindent\textbf{Use in the proof spine.} The packet is reused only after the local exact-complement assumption has been opened. It marks a role transition, not a hidden assertion of nonempty fiber.


\subsection{Audit packet 106: UEAP no-overreporting (route-exhaustion replay)}


\par\noindent\textbf{Claim under audit.} Support-level content cannot be reported as endpoint readout or endpoint defeat.


\par\noindent\textbf{Dependency path.} This packet belongs to the route-exhaustion replay. It depends on the official Hodge route, the fixed cycle-fiber endpoint image, the object/use distinction, and the kernel/A+ endpoint-use machinery. It is never used before the carrier and branch slots are fixed.


\par\noindent\textbf{Hostile-referee risk.} A hostile reader may attempt to treat the relevant content as raw mathematical syntax, ordinary negative theoremhood, surrogate support, or a classical counterexample object rather than endpoint-use material.


\par\noindent\textbf{Closure and denial burden.} The closure response is role-sensitive. If the content is raw, support, bookkeeping, or outside the fixed carrier, it has no endpoint-defeating force. If it keeps endpoint-defeating force on the same carrier, it enters endpoint-use discipline and route exhaustion. A successful objection must preserve the same $X,p,\alpha$, the same Hodge-standing predicate, the same Chow carrier $\CH^p(X)_\QQ$, and the same rational cycle-class map $\clQ$.


\par\noindent\textbf{Use in the proof spine.} The packet is reused only after the local exact-complement assumption has been opened. It marks a role transition, not a hidden assertion of nonempty fiber.


\subsection{Audit packet 107: Single cycle-fiber interface (route-exhaustion replay)}


\par\noindent\textbf{Claim under audit.} Cycle-fiber occupation is the only positive endpoint-image interface.


\par\noindent\textbf{Dependency path.} This packet belongs to the route-exhaustion replay. It depends on the official Hodge route, the fixed cycle-fiber endpoint image, the object/use distinction, and the kernel/A+ endpoint-use machinery. It is never used before the carrier and branch slots are fixed.


\par\noindent\textbf{Hostile-referee risk.} A hostile reader may attempt to treat the relevant content as raw mathematical syntax, ordinary negative theoremhood, surrogate support, or a classical counterexample object rather than endpoint-use material.


\par\noindent\textbf{Closure and denial burden.} The closure response is role-sensitive. If the content is raw, support, bookkeeping, or outside the fixed carrier, it has no endpoint-defeating force. If it keeps endpoint-defeating force on the same carrier, it enters endpoint-use discipline and route exhaustion. A successful objection must preserve the same $X,p,\alpha$, the same Hodge-standing predicate, the same Chow carrier $\CH^p(X)_\QQ$, and the same rational cycle-class map $\clQ$.


\par\noindent\textbf{Use in the proof spine.} The packet is reused only after the local exact-complement assumption has been opened. It marks a role transition, not a hidden assertion of nonempty fiber.


\subsection{Audit packet 108: Native counterforce exhaustion (route-exhaustion replay)}


\par\noindent\textbf{Claim under audit.} Native empty-fiber force has no fifth role.


\par\noindent\textbf{Dependency path.} This packet belongs to the route-exhaustion replay. It depends on the official Hodge route, the fixed cycle-fiber endpoint image, the object/use distinction, and the kernel/A+ endpoint-use machinery. It is never used before the carrier and branch slots are fixed.


\par\noindent\textbf{Hostile-referee risk.} A hostile reader may attempt to treat the relevant content as raw mathematical syntax, ordinary negative theoremhood, surrogate support, or a classical counterexample object rather than endpoint-use material.


\par\noindent\textbf{Closure and denial burden.} The closure response is role-sensitive. If the content is raw, support, bookkeeping, or outside the fixed carrier, it has no endpoint-defeating force. If it keeps endpoint-defeating force on the same carrier, it enters endpoint-use discipline and route exhaustion. A successful objection must preserve the same $X,p,\alpha$, the same Hodge-standing predicate, the same Chow carrier $\CH^p(X)_\QQ$, and the same rational cycle-class map $\clQ$.


\par\noindent\textbf{Use in the proof spine.} The packet is reused only after the local exact-complement assumption has been opened. It marks a role transition, not a hidden assertion of nonempty fiber.


\subsection{Audit packet 109: Separator induction (route-exhaustion replay)}


\par\noindent\textbf{Claim under audit.} Local separator occupation induces the fiber discriminator.


\par\noindent\textbf{Dependency path.} This packet belongs to the route-exhaustion replay. It depends on the official Hodge route, the fixed cycle-fiber endpoint image, the object/use distinction, and the kernel/A+ endpoint-use machinery. It is never used before the carrier and branch slots are fixed.


\par\noindent\textbf{Hostile-referee risk.} A hostile reader may attempt to treat the relevant content as raw mathematical syntax, ordinary negative theoremhood, surrogate support, or a classical counterexample object rather than endpoint-use material.


\par\noindent\textbf{Closure and denial burden.} The closure response is role-sensitive. If the content is raw, support, bookkeeping, or outside the fixed carrier, it has no endpoint-defeating force. If it keeps endpoint-defeating force on the same carrier, it enters endpoint-use discipline and route exhaustion. A successful objection must preserve the same $X,p,\alpha$, the same Hodge-standing predicate, the same Chow carrier $\CH^p(X)_\QQ$, and the same rational cycle-class map $\clQ$.


\par\noindent\textbf{Use in the proof spine.} The packet is reused only after the local exact-complement assumption has been opened. It marks a role transition, not a hidden assertion of nonempty fiber.


\subsection{Audit packet 110: A+ no-independent closure (route-exhaustion replay)}


\par\noindent\textbf{Claim under audit.} Endpoint use excludes independent same-domain discriminators.


\par\noindent\textbf{Dependency path.} This packet belongs to the route-exhaustion replay. It depends on the official Hodge route, the fixed cycle-fiber endpoint image, the object/use distinction, and the kernel/A+ endpoint-use machinery. It is never used before the carrier and branch slots are fixed.


\par\noindent\textbf{Hostile-referee risk.} A hostile reader may attempt to treat the relevant content as raw mathematical syntax, ordinary negative theoremhood, surrogate support, or a classical counterexample object rather than endpoint-use material.


\par\noindent\textbf{Closure and denial burden.} The closure response is role-sensitive. If the content is raw, support, bookkeeping, or outside the fixed carrier, it has no endpoint-defeating force. If it keeps endpoint-defeating force on the same carrier, it enters endpoint-use discipline and route exhaustion. A successful objection must preserve the same $X,p,\alpha$, the same Hodge-standing predicate, the same Chow carrier $\CH^p(X)_\QQ$, and the same rational cycle-class map $\clQ$.


\par\noindent\textbf{Use in the proof spine.} The packet is reused only after the local exact-complement assumption has been opened. It marks a role transition, not a hidden assertion of nonempty fiber.


\subsection{Audit packet 111: Annotation discharge (route-exhaustion replay)}


\par\noindent\textbf{Claim under audit.} The exact-complement annotation is not an added premise.


\par\noindent\textbf{Dependency path.} This packet belongs to the route-exhaustion replay. It depends on the official Hodge route, the fixed cycle-fiber endpoint image, the object/use distinction, and the kernel/A+ endpoint-use machinery. It is never used before the carrier and branch slots are fixed.


\par\noindent\textbf{Hostile-referee risk.} A hostile reader may attempt to treat the relevant content as raw mathematical syntax, ordinary negative theoremhood, surrogate support, or a classical counterexample object rather than endpoint-use material.


\par\noindent\textbf{Closure and denial burden.} The closure response is role-sensitive. If the content is raw, support, bookkeeping, or outside the fixed carrier, it has no endpoint-defeating force. If it keeps endpoint-defeating force on the same carrier, it enters endpoint-use discipline and route exhaustion. A successful objection must preserve the same $X,p,\alpha$, the same Hodge-standing predicate, the same Chow carrier $\CH^p(X)_\QQ$, and the same rational cycle-class map $\clQ$.


\par\noindent\textbf{Use in the proof spine.} The packet is reused only after the local exact-complement assumption has been opened. It marks a role transition, not a hidden assertion of nonempty fiber.


\subsection{Audit packet 112: Universalization (route-exhaustion replay)}


\par\noindent\textbf{Claim under audit.} Arbitrary instance closure yields the global Hodge endpoint.


\par\noindent\textbf{Dependency path.} This packet belongs to the route-exhaustion replay. It depends on the official Hodge route, the fixed cycle-fiber endpoint image, the object/use distinction, and the kernel/A+ endpoint-use machinery. It is never used before the carrier and branch slots are fixed.


\par\noindent\textbf{Hostile-referee risk.} A hostile reader may attempt to treat the relevant content as raw mathematical syntax, ordinary negative theoremhood, surrogate support, or a classical counterexample object rather than endpoint-use material.


\par\noindent\textbf{Closure and denial burden.} The closure response is role-sensitive. If the content is raw, support, bookkeeping, or outside the fixed carrier, it has no endpoint-defeating force. If it keeps endpoint-defeating force on the same carrier, it enters endpoint-use discipline and route exhaustion. A successful objection must preserve the same $X,p,\alpha$, the same Hodge-standing predicate, the same Chow carrier $\CH^p(X)_\QQ$, and the same rational cycle-class map $\clQ$.


\par\noindent\textbf{Use in the proof spine.} The packet is reused only after the local exact-complement assumption has been opened. It marks a role transition, not a hidden assertion of nonempty fiber.


\subsection{Audit packet 113: Surrogate separation (route-exhaustion replay)}


\par\noindent\textbf{Claim under audit.} Motives and absolute Hodge data are bridge programs unless they supply Chow fiber.


\par\noindent\textbf{Dependency path.} This packet belongs to the route-exhaustion replay. It depends on the official Hodge route, the fixed cycle-fiber endpoint image, the object/use distinction, and the kernel/A+ endpoint-use machinery. It is never used before the carrier and branch slots are fixed.


\par\noindent\textbf{Hostile-referee risk.} A hostile reader may attempt to treat the relevant content as raw mathematical syntax, ordinary negative theoremhood, surrogate support, or a classical counterexample object rather than endpoint-use material.


\par\noindent\textbf{Closure and denial burden.} The closure response is role-sensitive. If the content is raw, support, bookkeeping, or outside the fixed carrier, it has no endpoint-defeating force. If it keeps endpoint-defeating force on the same carrier, it enters endpoint-use discipline and route exhaustion. A successful objection must preserve the same $X,p,\alpha$, the same Hodge-standing predicate, the same Chow carrier $\CH^p(X)_\QQ$, and the same rational cycle-class map $\clQ$.


\par\noindent\textbf{Use in the proof spine.} The packet is reused only after the local exact-complement assumption has been opened. It marks a role transition, not a hidden assertion of nonempty fiber.


\subsection{Audit packet 114: Known-case transfer (route-exhaustion replay)}


\par\noindent\textbf{Claim under audit.} Known bridge completions do not transfer universally without target gating.


\par\noindent\textbf{Dependency path.} This packet belongs to the route-exhaustion replay. It depends on the official Hodge route, the fixed cycle-fiber endpoint image, the object/use distinction, and the kernel/A+ endpoint-use machinery. It is never used before the carrier and branch slots are fixed.


\par\noindent\textbf{Hostile-referee risk.} A hostile reader may attempt to treat the relevant content as raw mathematical syntax, ordinary negative theoremhood, surrogate support, or a classical counterexample object rather than endpoint-use material.


\par\noindent\textbf{Closure and denial burden.} The closure response is role-sensitive. If the content is raw, support, bookkeeping, or outside the fixed carrier, it has no endpoint-defeating force. If it keeps endpoint-defeating force on the same carrier, it enters endpoint-use discipline and route exhaustion. A successful objection must preserve the same $X,p,\alpha$, the same Hodge-standing predicate, the same Chow carrier $\CH^p(X)_\QQ$, and the same rational cycle-class map $\clQ$.


\par\noindent\textbf{Use in the proof spine.} The packet is reused only after the local exact-complement assumption has been opened. It marks a role transition, not a hidden assertion of nonempty fiber.


\subsection{Audit packet 115: Local-to-global nonfactorization (route-exhaustion replay)}


\par\noindent\textbf{Claim under audit.} Local cycle data require a global compatibility certificate.


\par\noindent\textbf{Dependency path.} This packet belongs to the route-exhaustion replay. It depends on the official Hodge route, the fixed cycle-fiber endpoint image, the object/use distinction, and the kernel/A+ endpoint-use machinery. It is never used before the carrier and branch slots are fixed.


\par\noindent\textbf{Hostile-referee risk.} A hostile reader may attempt to treat the relevant content as raw mathematical syntax, ordinary negative theoremhood, surrogate support, or a classical counterexample object rather than endpoint-use material.


\par\noindent\textbf{Closure and denial burden.} The closure response is role-sensitive. If the content is raw, support, bookkeeping, or outside the fixed carrier, it has no endpoint-defeating force. If it keeps endpoint-defeating force on the same carrier, it enters endpoint-use discipline and route exhaustion. A successful objection must preserve the same $X,p,\alpha$, the same Hodge-standing predicate, the same Chow carrier $\CH^p(X)_\QQ$, and the same rational cycle-class map $\clQ$.


\par\noindent\textbf{Use in the proof spine.} The packet is reused only after the local exact-complement assumption has been opened. It marks a role transition, not a hidden assertion of nonempty fiber.


\subsection{Audit packet 116: No repair (route-exhaustion replay)}


\par\noindent\textbf{Claim under audit.} Later cycle discovery creates a new bridge-complete act, not repair of an old one.


\par\noindent\textbf{Dependency path.} This packet belongs to the route-exhaustion replay. It depends on the official Hodge route, the fixed cycle-fiber endpoint image, the object/use distinction, and the kernel/A+ endpoint-use machinery. It is never used before the carrier and branch slots are fixed.


\par\noindent\textbf{Hostile-referee risk.} A hostile reader may attempt to treat the relevant content as raw mathematical syntax, ordinary negative theoremhood, surrogate support, or a classical counterexample object rather than endpoint-use material.


\par\noindent\textbf{Closure and denial burden.} The closure response is role-sensitive. If the content is raw, support, bookkeeping, or outside the fixed carrier, it has no endpoint-defeating force. If it keeps endpoint-defeating force on the same carrier, it enters endpoint-use discipline and route exhaustion. A successful objection must preserve the same $X,p,\alpha$, the same Hodge-standing predicate, the same Chow carrier $\CH^p(X)_\QQ$, and the same rational cycle-class map $\clQ$.


\par\noindent\textbf{Use in the proof spine.} The packet is reused only after the local exact-complement assumption has been opened. It marks a role transition, not a hidden assertion of nonempty fiber.


\subsection{Audit packet 117: No selector (route-exhaustion replay)}


\par\noindent\textbf{Claim under audit.} Canonical representatives cannot decide readout before the fiber is fixed.


\par\noindent\textbf{Dependency path.} This packet belongs to the route-exhaustion replay. It depends on the official Hodge route, the fixed cycle-fiber endpoint image, the object/use distinction, and the kernel/A+ endpoint-use machinery. It is never used before the carrier and branch slots are fixed.


\par\noindent\textbf{Hostile-referee risk.} A hostile reader may attempt to treat the relevant content as raw mathematical syntax, ordinary negative theoremhood, surrogate support, or a classical counterexample object rather than endpoint-use material.


\par\noindent\textbf{Closure and denial burden.} The closure response is role-sensitive. If the content is raw, support, bookkeeping, or outside the fixed carrier, it has no endpoint-defeating force. If it keeps endpoint-defeating force on the same carrier, it enters endpoint-use discipline and route exhaustion. A successful objection must preserve the same $X,p,\alpha$, the same Hodge-standing predicate, the same Chow carrier $\CH^p(X)_\QQ$, and the same rational cycle-class map $\clQ$.


\par\noindent\textbf{Use in the proof spine.} The packet is reused only after the local exact-complement assumption has been opened. It marks a role transition, not a hidden assertion of nonempty fiber.


\subsection{Audit packet 118: Incompleteness firewall (route-exhaustion replay)}


\par\noindent\textbf{Claim under audit.} Independence matters only if it re-enters endpoint standing or trajectory.


\par\noindent\textbf{Dependency path.} This packet belongs to the route-exhaustion replay. It depends on the official Hodge route, the fixed cycle-fiber endpoint image, the object/use distinction, and the kernel/A+ endpoint-use machinery. It is never used before the carrier and branch slots are fixed.


\par\noindent\textbf{Hostile-referee risk.} A hostile reader may attempt to treat the relevant content as raw mathematical syntax, ordinary negative theoremhood, surrogate support, or a classical counterexample object rather than endpoint-use material.


\par\noindent\textbf{Closure and denial burden.} The closure response is role-sensitive. If the content is raw, support, bookkeeping, or outside the fixed carrier, it has no endpoint-defeating force. If it keeps endpoint-defeating force on the same carrier, it enters endpoint-use discipline and route exhaustion. A successful objection must preserve the same $X,p,\alpha$, the same Hodge-standing predicate, the same Chow carrier $\CH^p(X)_\QQ$, and the same rational cycle-class map $\clQ$.


\par\noindent\textbf{Use in the proof spine.} The packet is reused only after the local exact-complement assumption has been opened. It marks a role transition, not a hidden assertion of nonempty fiber.


\subsection{Audit packet 119: Ordinary theoremhood protection (route-exhaustion replay)}


\par\noindent\textbf{Claim under audit.} Negative theoremhood is proof legitimacy, not forbidden endpoint governance.


\par\noindent\textbf{Dependency path.} This packet belongs to the route-exhaustion replay. It depends on the official Hodge route, the fixed cycle-fiber endpoint image, the object/use distinction, and the kernel/A+ endpoint-use machinery. It is never used before the carrier and branch slots are fixed.


\par\noindent\textbf{Hostile-referee risk.} A hostile reader may attempt to treat the relevant content as raw mathematical syntax, ordinary negative theoremhood, surrogate support, or a classical counterexample object rather than endpoint-use material.


\par\noindent\textbf{Closure and denial burden.} The closure response is role-sensitive. If the content is raw, support, bookkeeping, or outside the fixed carrier, it has no endpoint-defeating force. If it keeps endpoint-defeating force on the same carrier, it enters endpoint-use discipline and route exhaustion. A successful objection must preserve the same $X,p,\alpha$, the same Hodge-standing predicate, the same Chow carrier $\CH^p(X)_\QQ$, and the same rational cycle-class map $\clQ$.


\par\noindent\textbf{Use in the proof spine.} The packet is reused only after the local exact-complement assumption has been opened. It marks a role transition, not a hidden assertion of nonempty fiber.


\subsection{Audit packet 120: Non-explosiveness (route-exhaustion replay)}


\par\noindent\textbf{Claim under audit.} Empty fiber is lawful until endpoint force is added.


\par\noindent\textbf{Dependency path.} This packet belongs to the route-exhaustion replay. It depends on the official Hodge route, the fixed cycle-fiber endpoint image, the object/use distinction, and the kernel/A+ endpoint-use machinery. It is never used before the carrier and branch slots are fixed.


\par\noindent\textbf{Hostile-referee risk.} A hostile reader may attempt to treat the relevant content as raw mathematical syntax, ordinary negative theoremhood, surrogate support, or a classical counterexample object rather than endpoint-use material.


\par\noindent\textbf{Closure and denial burden.} The closure response is role-sensitive. If the content is raw, support, bookkeeping, or outside the fixed carrier, it has no endpoint-defeating force. If it keeps endpoint-defeating force on the same carrier, it enters endpoint-use discipline and route exhaustion. A successful objection must preserve the same $X,p,\alpha$, the same Hodge-standing predicate, the same Chow carrier $\CH^p(X)_\QQ$, and the same rational cycle-class map $\clQ$.


\par\noindent\textbf{Use in the proof spine.} The packet is reused only after the local exact-complement assumption has been opened. It marks a role transition, not a hidden assertion of nonempty fiber.


\subsection{Audit packet 121: Context non-smuggling (anti-circularity replay)}


\par\noindent\textbf{Claim under audit.} Hodge context fixes carrier and slots without branch occupation.


\par\noindent\textbf{Dependency path.} This packet belongs to the anti-circularity replay. It depends on the official Hodge route, the fixed cycle-fiber endpoint image, the object/use distinction, and the kernel/A+ endpoint-use machinery. It is never used before the carrier and branch slots are fixed.


\par\noindent\textbf{Hostile-referee risk.} A hostile reader may attempt to treat the relevant content as raw mathematical syntax, ordinary negative theoremhood, surrogate support, or a classical counterexample object rather than endpoint-use material.


\par\noindent\textbf{Closure and denial burden.} The closure response is role-sensitive. If the content is raw, support, bookkeeping, or outside the fixed carrier, it has no endpoint-defeating force. If it keeps endpoint-defeating force on the same carrier, it enters endpoint-use discipline and route exhaustion. A successful objection must preserve the same $X,p,\alpha$, the same Hodge-standing predicate, the same Chow carrier $\CH^p(X)_\QQ$, and the same rational cycle-class map $\clQ$.


\par\noindent\textbf{Use in the proof spine.} The packet is reused only after the local exact-complement assumption has been opened. It marks a role transition, not a hidden assertion of nonempty fiber.


\subsection{Audit packet 122: Model adequacy (anti-circularity replay)}


\par\noindent\textbf{Claim under audit.} Readout and counterreadout are identified with nonempty and empty cycle fiber.


\par\noindent\textbf{Dependency path.} This packet belongs to the anti-circularity replay. It depends on the official Hodge route, the fixed cycle-fiber endpoint image, the object/use distinction, and the kernel/A+ endpoint-use machinery. It is never used before the carrier and branch slots are fixed.


\par\noindent\textbf{Hostile-referee risk.} A hostile reader may attempt to treat the relevant content as raw mathematical syntax, ordinary negative theoremhood, surrogate support, or a classical counterexample object rather than endpoint-use material.


\par\noindent\textbf{Closure and denial burden.} The closure response is role-sensitive. If the content is raw, support, bookkeeping, or outside the fixed carrier, it has no endpoint-defeating force. If it keeps endpoint-defeating force on the same carrier, it enters endpoint-use discipline and route exhaustion. A successful objection must preserve the same $X,p,\alpha$, the same Hodge-standing predicate, the same Chow carrier $\CH^p(X)_\QQ$, and the same rational cycle-class map $\clQ$.


\par\noindent\textbf{Use in the proof spine.} The packet is reused only after the local exact-complement assumption has been opened. It marks a role transition, not a hidden assertion of nonempty fiber.


\subsection{Audit packet 123: Counterexample force as endpoint use (anti-circularity replay)}


\par\noindent\textbf{Claim under audit.} Official endpoint defeat is endpoint use.


\par\noindent\textbf{Dependency path.} This packet belongs to the anti-circularity replay. It depends on the official Hodge route, the fixed cycle-fiber endpoint image, the object/use distinction, and the kernel/A+ endpoint-use machinery. It is never used before the carrier and branch slots are fixed.


\par\noindent\textbf{Hostile-referee risk.} A hostile reader may attempt to treat the relevant content as raw mathematical syntax, ordinary negative theoremhood, surrogate support, or a classical counterexample object rather than endpoint-use material.


\par\noindent\textbf{Closure and denial burden.} The closure response is role-sensitive. If the content is raw, support, bookkeeping, or outside the fixed carrier, it has no endpoint-defeating force. If it keeps endpoint-defeating force on the same carrier, it enters endpoint-use discipline and route exhaustion. A successful objection must preserve the same $X,p,\alpha$, the same Hodge-standing predicate, the same Chow carrier $\CH^p(X)_\QQ$, and the same rational cycle-class map $\clQ$.


\par\noindent\textbf{Use in the proof spine.} The packet is reused only after the local exact-complement assumption has been opened. It marks a role transition, not a hidden assertion of nonempty fiber.


\subsection{Audit packet 124: No autonomous counterexample role (anti-circularity replay)}


\par\noindent\textbf{Claim under audit.} There is no endpoint-defeating branch outside endpoint use.


\par\noindent\textbf{Dependency path.} This packet belongs to the anti-circularity replay. It depends on the official Hodge route, the fixed cycle-fiber endpoint image, the object/use distinction, and the kernel/A+ endpoint-use machinery. It is never used before the carrier and branch slots are fixed.


\par\noindent\textbf{Hostile-referee risk.} A hostile reader may attempt to treat the relevant content as raw mathematical syntax, ordinary negative theoremhood, surrogate support, or a classical counterexample object rather than endpoint-use material.


\par\noindent\textbf{Closure and denial burden.} The closure response is role-sensitive. If the content is raw, support, bookkeeping, or outside the fixed carrier, it has no endpoint-defeating force. If it keeps endpoint-defeating force on the same carrier, it enters endpoint-use discipline and route exhaustion. A successful objection must preserve the same $X,p,\alpha$, the same Hodge-standing predicate, the same Chow carrier $\CH^p(X)_\QQ$, and the same rational cycle-class map $\clQ$.


\par\noindent\textbf{Use in the proof spine.} The packet is reused only after the local exact-complement assumption has been opened. It marks a role transition, not a hidden assertion of nonempty fiber.


\subsection{Audit packet 125: ATS locking (anti-circularity replay)}


\par\noindent\textbf{Claim under audit.} The same endpoint act preserves anchor and tensor data.


\par\noindent\textbf{Dependency path.} This packet belongs to the anti-circularity replay. It depends on the official Hodge route, the fixed cycle-fiber endpoint image, the object/use distinction, and the kernel/A+ endpoint-use machinery. It is never used before the carrier and branch slots are fixed.


\par\noindent\textbf{Hostile-referee risk.} A hostile reader may attempt to treat the relevant content as raw mathematical syntax, ordinary negative theoremhood, surrogate support, or a classical counterexample object rather than endpoint-use material.


\par\noindent\textbf{Closure and denial burden.} The closure response is role-sensitive. If the content is raw, support, bookkeeping, or outside the fixed carrier, it has no endpoint-defeating force. If it keeps endpoint-defeating force on the same carrier, it enters endpoint-use discipline and route exhaustion. A successful objection must preserve the same $X,p,\alpha$, the same Hodge-standing predicate, the same Chow carrier $\CH^p(X)_\QQ$, and the same rational cycle-class map $\clQ$.


\par\noindent\textbf{Use in the proof spine.} The packet is reused only after the local exact-complement assumption has been opened. It marks a role transition, not a hidden assertion of nonempty fiber.


\subsection{Audit packet 126: UEAP no-overreporting (anti-circularity replay)}


\par\noindent\textbf{Claim under audit.} Support-level content cannot be reported as endpoint readout or endpoint defeat.


\par\noindent\textbf{Dependency path.} This packet belongs to the anti-circularity replay. It depends on the official Hodge route, the fixed cycle-fiber endpoint image, the object/use distinction, and the kernel/A+ endpoint-use machinery. It is never used before the carrier and branch slots are fixed.


\par\noindent\textbf{Hostile-referee risk.} A hostile reader may attempt to treat the relevant content as raw mathematical syntax, ordinary negative theoremhood, surrogate support, or a classical counterexample object rather than endpoint-use material.


\par\noindent\textbf{Closure and denial burden.} The closure response is role-sensitive. If the content is raw, support, bookkeeping, or outside the fixed carrier, it has no endpoint-defeating force. If it keeps endpoint-defeating force on the same carrier, it enters endpoint-use discipline and route exhaustion. A successful objection must preserve the same $X,p,\alpha$, the same Hodge-standing predicate, the same Chow carrier $\CH^p(X)_\QQ$, and the same rational cycle-class map $\clQ$.


\par\noindent\textbf{Use in the proof spine.} The packet is reused only after the local exact-complement assumption has been opened. It marks a role transition, not a hidden assertion of nonempty fiber.


\subsection{Audit packet 127: Single cycle-fiber interface (anti-circularity replay)}


\par\noindent\textbf{Claim under audit.} Cycle-fiber occupation is the only positive endpoint-image interface.


\par\noindent\textbf{Dependency path.} This packet belongs to the anti-circularity replay. It depends on the official Hodge route, the fixed cycle-fiber endpoint image, the object/use distinction, and the kernel/A+ endpoint-use machinery. It is never used before the carrier and branch slots are fixed.


\par\noindent\textbf{Hostile-referee risk.} A hostile reader may attempt to treat the relevant content as raw mathematical syntax, ordinary negative theoremhood, surrogate support, or a classical counterexample object rather than endpoint-use material.


\par\noindent\textbf{Closure and denial burden.} The closure response is role-sensitive. If the content is raw, support, bookkeeping, or outside the fixed carrier, it has no endpoint-defeating force. If it keeps endpoint-defeating force on the same carrier, it enters endpoint-use discipline and route exhaustion. A successful objection must preserve the same $X,p,\alpha$, the same Hodge-standing predicate, the same Chow carrier $\CH^p(X)_\QQ$, and the same rational cycle-class map $\clQ$.


\par\noindent\textbf{Use in the proof spine.} The packet is reused only after the local exact-complement assumption has been opened. It marks a role transition, not a hidden assertion of nonempty fiber.


\subsection{Audit packet 128: Native counterforce exhaustion (anti-circularity replay)}


\par\noindent\textbf{Claim under audit.} Native empty-fiber force has no fifth role.


\par\noindent\textbf{Dependency path.} This packet belongs to the anti-circularity replay. It depends on the official Hodge route, the fixed cycle-fiber endpoint image, the object/use distinction, and the kernel/A+ endpoint-use machinery. It is never used before the carrier and branch slots are fixed.


\par\noindent\textbf{Hostile-referee risk.} A hostile reader may attempt to treat the relevant content as raw mathematical syntax, ordinary negative theoremhood, surrogate support, or a classical counterexample object rather than endpoint-use material.


\par\noindent\textbf{Closure and denial burden.} The closure response is role-sensitive. If the content is raw, support, bookkeeping, or outside the fixed carrier, it has no endpoint-defeating force. If it keeps endpoint-defeating force on the same carrier, it enters endpoint-use discipline and route exhaustion. A successful objection must preserve the same $X,p,\alpha$, the same Hodge-standing predicate, the same Chow carrier $\CH^p(X)_\QQ$, and the same rational cycle-class map $\clQ$.


\par\noindent\textbf{Use in the proof spine.} The packet is reused only after the local exact-complement assumption has been opened. It marks a role transition, not a hidden assertion of nonempty fiber.


\subsection{Audit packet 129: Separator induction (anti-circularity replay)}


\par\noindent\textbf{Claim under audit.} Local separator occupation induces the fiber discriminator.


\par\noindent\textbf{Dependency path.} This packet belongs to the anti-circularity replay. It depends on the official Hodge route, the fixed cycle-fiber endpoint image, the object/use distinction, and the kernel/A+ endpoint-use machinery. It is never used before the carrier and branch slots are fixed.


\par\noindent\textbf{Hostile-referee risk.} A hostile reader may attempt to treat the relevant content as raw mathematical syntax, ordinary negative theoremhood, surrogate support, or a classical counterexample object rather than endpoint-use material.


\par\noindent\textbf{Closure and denial burden.} The closure response is role-sensitive. If the content is raw, support, bookkeeping, or outside the fixed carrier, it has no endpoint-defeating force. If it keeps endpoint-defeating force on the same carrier, it enters endpoint-use discipline and route exhaustion. A successful objection must preserve the same $X,p,\alpha$, the same Hodge-standing predicate, the same Chow carrier $\CH^p(X)_\QQ$, and the same rational cycle-class map $\clQ$.


\par\noindent\textbf{Use in the proof spine.} The packet is reused only after the local exact-complement assumption has been opened. It marks a role transition, not a hidden assertion of nonempty fiber.


\subsection{Audit packet 130: A+ no-independent closure (anti-circularity replay)}


\par\noindent\textbf{Claim under audit.} Endpoint use excludes independent same-domain discriminators.


\par\noindent\textbf{Dependency path.} This packet belongs to the anti-circularity replay. It depends on the official Hodge route, the fixed cycle-fiber endpoint image, the object/use distinction, and the kernel/A+ endpoint-use machinery. It is never used before the carrier and branch slots are fixed.


\par\noindent\textbf{Hostile-referee risk.} A hostile reader may attempt to treat the relevant content as raw mathematical syntax, ordinary negative theoremhood, surrogate support, or a classical counterexample object rather than endpoint-use material.


\par\noindent\textbf{Closure and denial burden.} The closure response is role-sensitive. If the content is raw, support, bookkeeping, or outside the fixed carrier, it has no endpoint-defeating force. If it keeps endpoint-defeating force on the same carrier, it enters endpoint-use discipline and route exhaustion. A successful objection must preserve the same $X,p,\alpha$, the same Hodge-standing predicate, the same Chow carrier $\CH^p(X)_\QQ$, and the same rational cycle-class map $\clQ$.


\par\noindent\textbf{Use in the proof spine.} The packet is reused only after the local exact-complement assumption has been opened. It marks a role transition, not a hidden assertion of nonempty fiber.


\subsection{Audit packet 131: Annotation discharge (anti-circularity replay)}


\par\noindent\textbf{Claim under audit.} The exact-complement annotation is not an added premise.


\par\noindent\textbf{Dependency path.} This packet belongs to the anti-circularity replay. It depends on the official Hodge route, the fixed cycle-fiber endpoint image, the object/use distinction, and the kernel/A+ endpoint-use machinery. It is never used before the carrier and branch slots are fixed.


\par\noindent\textbf{Hostile-referee risk.} A hostile reader may attempt to treat the relevant content as raw mathematical syntax, ordinary negative theoremhood, surrogate support, or a classical counterexample object rather than endpoint-use material.


\par\noindent\textbf{Closure and denial burden.} The closure response is role-sensitive. If the content is raw, support, bookkeeping, or outside the fixed carrier, it has no endpoint-defeating force. If it keeps endpoint-defeating force on the same carrier, it enters endpoint-use discipline and route exhaustion. A successful objection must preserve the same $X,p,\alpha$, the same Hodge-standing predicate, the same Chow carrier $\CH^p(X)_\QQ$, and the same rational cycle-class map $\clQ$.


\par\noindent\textbf{Use in the proof spine.} The packet is reused only after the local exact-complement assumption has been opened. It marks a role transition, not a hidden assertion of nonempty fiber.


\subsection{Audit packet 132: Universalization (anti-circularity replay)}


\par\noindent\textbf{Claim under audit.} Arbitrary instance closure yields the global Hodge endpoint.


\par\noindent\textbf{Dependency path.} This packet belongs to the anti-circularity replay. It depends on the official Hodge route, the fixed cycle-fiber endpoint image, the object/use distinction, and the kernel/A+ endpoint-use machinery. It is never used before the carrier and branch slots are fixed.


\par\noindent\textbf{Hostile-referee risk.} A hostile reader may attempt to treat the relevant content as raw mathematical syntax, ordinary negative theoremhood, surrogate support, or a classical counterexample object rather than endpoint-use material.


\par\noindent\textbf{Closure and denial burden.} The closure response is role-sensitive. If the content is raw, support, bookkeeping, or outside the fixed carrier, it has no endpoint-defeating force. If it keeps endpoint-defeating force on the same carrier, it enters endpoint-use discipline and route exhaustion. A successful objection must preserve the same $X,p,\alpha$, the same Hodge-standing predicate, the same Chow carrier $\CH^p(X)_\QQ$, and the same rational cycle-class map $\clQ$.


\par\noindent\textbf{Use in the proof spine.} The packet is reused only after the local exact-complement assumption has been opened. It marks a role transition, not a hidden assertion of nonempty fiber.


\subsection{Audit packet 133: Surrogate separation (anti-circularity replay)}


\par\noindent\textbf{Claim under audit.} Motives and absolute Hodge data are bridge programs unless they supply Chow fiber.


\par\noindent\textbf{Dependency path.} This packet belongs to the anti-circularity replay. It depends on the official Hodge route, the fixed cycle-fiber endpoint image, the object/use distinction, and the kernel/A+ endpoint-use machinery. It is never used before the carrier and branch slots are fixed.


\par\noindent\textbf{Hostile-referee risk.} A hostile reader may attempt to treat the relevant content as raw mathematical syntax, ordinary negative theoremhood, surrogate support, or a classical counterexample object rather than endpoint-use material.


\par\noindent\textbf{Closure and denial burden.} The closure response is role-sensitive. If the content is raw, support, bookkeeping, or outside the fixed carrier, it has no endpoint-defeating force. If it keeps endpoint-defeating force on the same carrier, it enters endpoint-use discipline and route exhaustion. A successful objection must preserve the same $X,p,\alpha$, the same Hodge-standing predicate, the same Chow carrier $\CH^p(X)_\QQ$, and the same rational cycle-class map $\clQ$.


\par\noindent\textbf{Use in the proof spine.} The packet is reused only after the local exact-complement assumption has been opened. It marks a role transition, not a hidden assertion of nonempty fiber.


\subsection{Audit packet 134: Known-case transfer (anti-circularity replay)}


\par\noindent\textbf{Claim under audit.} Known bridge completions do not transfer universally without target gating.


\par\noindent\textbf{Dependency path.} This packet belongs to the anti-circularity replay. It depends on the official Hodge route, the fixed cycle-fiber endpoint image, the object/use distinction, and the kernel/A+ endpoint-use machinery. It is never used before the carrier and branch slots are fixed.


\par\noindent\textbf{Hostile-referee risk.} A hostile reader may attempt to treat the relevant content as raw mathematical syntax, ordinary negative theoremhood, surrogate support, or a classical counterexample object rather than endpoint-use material.


\par\noindent\textbf{Closure and denial burden.} The closure response is role-sensitive. If the content is raw, support, bookkeeping, or outside the fixed carrier, it has no endpoint-defeating force. If it keeps endpoint-defeating force on the same carrier, it enters endpoint-use discipline and route exhaustion. A successful objection must preserve the same $X,p,\alpha$, the same Hodge-standing predicate, the same Chow carrier $\CH^p(X)_\QQ$, and the same rational cycle-class map $\clQ$.


\par\noindent\textbf{Use in the proof spine.} The packet is reused only after the local exact-complement assumption has been opened. It marks a role transition, not a hidden assertion of nonempty fiber.


\subsection{Audit packet 135: Local-to-global nonfactorization (anti-circularity replay)}


\par\noindent\textbf{Claim under audit.} Local cycle data require a global compatibility certificate.


\par\noindent\textbf{Dependency path.} This packet belongs to the anti-circularity replay. It depends on the official Hodge route, the fixed cycle-fiber endpoint image, the object/use distinction, and the kernel/A+ endpoint-use machinery. It is never used before the carrier and branch slots are fixed.


\par\noindent\textbf{Hostile-referee risk.} A hostile reader may attempt to treat the relevant content as raw mathematical syntax, ordinary negative theoremhood, surrogate support, or a classical counterexample object rather than endpoint-use material.


\par\noindent\textbf{Closure and denial burden.} The closure response is role-sensitive. If the content is raw, support, bookkeeping, or outside the fixed carrier, it has no endpoint-defeating force. If it keeps endpoint-defeating force on the same carrier, it enters endpoint-use discipline and route exhaustion. A successful objection must preserve the same $X,p,\alpha$, the same Hodge-standing predicate, the same Chow carrier $\CH^p(X)_\QQ$, and the same rational cycle-class map $\clQ$.


\par\noindent\textbf{Use in the proof spine.} The packet is reused only after the local exact-complement assumption has been opened. It marks a role transition, not a hidden assertion of nonempty fiber.


\subsection{Audit packet 136: No repair (anti-circularity replay)}


\par\noindent\textbf{Claim under audit.} Later cycle discovery creates a new bridge-complete act, not repair of an old one.


\par\noindent\textbf{Dependency path.} This packet belongs to the anti-circularity replay. It depends on the official Hodge route, the fixed cycle-fiber endpoint image, the object/use distinction, and the kernel/A+ endpoint-use machinery. It is never used before the carrier and branch slots are fixed.


\par\noindent\textbf{Hostile-referee risk.} A hostile reader may attempt to treat the relevant content as raw mathematical syntax, ordinary negative theoremhood, surrogate support, or a classical counterexample object rather than endpoint-use material.


\par\noindent\textbf{Closure and denial burden.} The closure response is role-sensitive. If the content is raw, support, bookkeeping, or outside the fixed carrier, it has no endpoint-defeating force. If it keeps endpoint-defeating force on the same carrier, it enters endpoint-use discipline and route exhaustion. A successful objection must preserve the same $X,p,\alpha$, the same Hodge-standing predicate, the same Chow carrier $\CH^p(X)_\QQ$, and the same rational cycle-class map $\clQ$.


\par\noindent\textbf{Use in the proof spine.} The packet is reused only after the local exact-complement assumption has been opened. It marks a role transition, not a hidden assertion of nonempty fiber.


\subsection{Audit packet 137: No selector (anti-circularity replay)}


\par\noindent\textbf{Claim under audit.} Canonical representatives cannot decide readout before the fiber is fixed.


\par\noindent\textbf{Dependency path.} This packet belongs to the anti-circularity replay. It depends on the official Hodge route, the fixed cycle-fiber endpoint image, the object/use distinction, and the kernel/A+ endpoint-use machinery. It is never used before the carrier and branch slots are fixed.


\par\noindent\textbf{Hostile-referee risk.} A hostile reader may attempt to treat the relevant content as raw mathematical syntax, ordinary negative theoremhood, surrogate support, or a classical counterexample object rather than endpoint-use material.


\par\noindent\textbf{Closure and denial burden.} The closure response is role-sensitive. If the content is raw, support, bookkeeping, or outside the fixed carrier, it has no endpoint-defeating force. If it keeps endpoint-defeating force on the same carrier, it enters endpoint-use discipline and route exhaustion. A successful objection must preserve the same $X,p,\alpha$, the same Hodge-standing predicate, the same Chow carrier $\CH^p(X)_\QQ$, and the same rational cycle-class map $\clQ$.


\par\noindent\textbf{Use in the proof spine.} The packet is reused only after the local exact-complement assumption has been opened. It marks a role transition, not a hidden assertion of nonempty fiber.


\subsection{Audit packet 138: Incompleteness firewall (anti-circularity replay)}


\par\noindent\textbf{Claim under audit.} Independence matters only if it re-enters endpoint standing or trajectory.


\par\noindent\textbf{Dependency path.} This packet belongs to the anti-circularity replay. It depends on the official Hodge route, the fixed cycle-fiber endpoint image, the object/use distinction, and the kernel/A+ endpoint-use machinery. It is never used before the carrier and branch slots are fixed.


\par\noindent\textbf{Hostile-referee risk.} A hostile reader may attempt to treat the relevant content as raw mathematical syntax, ordinary negative theoremhood, surrogate support, or a classical counterexample object rather than endpoint-use material.


\par\noindent\textbf{Closure and denial burden.} The closure response is role-sensitive. If the content is raw, support, bookkeeping, or outside the fixed carrier, it has no endpoint-defeating force. If it keeps endpoint-defeating force on the same carrier, it enters endpoint-use discipline and route exhaustion. A successful objection must preserve the same $X,p,\alpha$, the same Hodge-standing predicate, the same Chow carrier $\CH^p(X)_\QQ$, and the same rational cycle-class map $\clQ$.


\par\noindent\textbf{Use in the proof spine.} The packet is reused only after the local exact-complement assumption has been opened. It marks a role transition, not a hidden assertion of nonempty fiber.


\subsection{Audit packet 139: Ordinary theoremhood protection (anti-circularity replay)}


\par\noindent\textbf{Claim under audit.} Negative theoremhood is proof legitimacy, not forbidden endpoint governance.


\par\noindent\textbf{Dependency path.} This packet belongs to the anti-circularity replay. It depends on the official Hodge route, the fixed cycle-fiber endpoint image, the object/use distinction, and the kernel/A+ endpoint-use machinery. It is never used before the carrier and branch slots are fixed.


\par\noindent\textbf{Hostile-referee risk.} A hostile reader may attempt to treat the relevant content as raw mathematical syntax, ordinary negative theoremhood, surrogate support, or a classical counterexample object rather than endpoint-use material.


\par\noindent\textbf{Closure and denial burden.} The closure response is role-sensitive. If the content is raw, support, bookkeeping, or outside the fixed carrier, it has no endpoint-defeating force. If it keeps endpoint-defeating force on the same carrier, it enters endpoint-use discipline and route exhaustion. A successful objection must preserve the same $X,p,\alpha$, the same Hodge-standing predicate, the same Chow carrier $\CH^p(X)_\QQ$, and the same rational cycle-class map $\clQ$.


\par\noindent\textbf{Use in the proof spine.} The packet is reused only after the local exact-complement assumption has been opened. It marks a role transition, not a hidden assertion of nonempty fiber.


\subsection{Audit packet 140: Non-explosiveness (anti-circularity replay)}


\par\noindent\textbf{Claim under audit.} Empty fiber is lawful until endpoint force is added.


\par\noindent\textbf{Dependency path.} This packet belongs to the anti-circularity replay. It depends on the official Hodge route, the fixed cycle-fiber endpoint image, the object/use distinction, and the kernel/A+ endpoint-use machinery. It is never used before the carrier and branch slots are fixed.


\par\noindent\textbf{Hostile-referee risk.} A hostile reader may attempt to treat the relevant content as raw mathematical syntax, ordinary negative theoremhood, surrogate support, or a classical counterexample object rather than endpoint-use material.


\par\noindent\textbf{Closure and denial burden.} The closure response is role-sensitive. If the content is raw, support, bookkeeping, or outside the fixed carrier, it has no endpoint-defeating force. If it keeps endpoint-defeating force on the same carrier, it enters endpoint-use discipline and route exhaustion. A successful objection must preserve the same $X,p,\alpha$, the same Hodge-standing predicate, the same Chow carrier $\CH^p(X)_\QQ$, and the same rational cycle-class map $\clQ$.


\par\noindent\textbf{Use in the proof spine.} The packet is reused only after the local exact-complement assumption has been opened. It marks a role transition, not a hidden assertion of nonempty fiber.


\subsection{Audit packet 141: Context non-smuggling (counterexample-force replay)}


\par\noindent\textbf{Claim under audit.} Hodge context fixes carrier and slots without branch occupation.


\par\noindent\textbf{Dependency path.} This packet belongs to the counterexample-force replay. It depends on the official Hodge route, the fixed cycle-fiber endpoint image, the object/use distinction, and the kernel/A+ endpoint-use machinery. It is never used before the carrier and branch slots are fixed.


\par\noindent\textbf{Hostile-referee risk.} A hostile reader may attempt to treat the relevant content as raw mathematical syntax, ordinary negative theoremhood, surrogate support, or a classical counterexample object rather than endpoint-use material.


\par\noindent\textbf{Closure and denial burden.} The closure response is role-sensitive. If the content is raw, support, bookkeeping, or outside the fixed carrier, it has no endpoint-defeating force. If it keeps endpoint-defeating force on the same carrier, it enters endpoint-use discipline and route exhaustion. A successful objection must preserve the same $X,p,\alpha$, the same Hodge-standing predicate, the same Chow carrier $\CH^p(X)_\QQ$, and the same rational cycle-class map $\clQ$.


\par\noindent\textbf{Use in the proof spine.} The packet is reused only after the local exact-complement assumption has been opened. It marks a role transition, not a hidden assertion of nonempty fiber.


\subsection{Audit packet 142: Model adequacy (counterexample-force replay)}


\par\noindent\textbf{Claim under audit.} Readout and counterreadout are identified with nonempty and empty cycle fiber.


\par\noindent\textbf{Dependency path.} This packet belongs to the counterexample-force replay. It depends on the official Hodge route, the fixed cycle-fiber endpoint image, the object/use distinction, and the kernel/A+ endpoint-use machinery. It is never used before the carrier and branch slots are fixed.


\par\noindent\textbf{Hostile-referee risk.} A hostile reader may attempt to treat the relevant content as raw mathematical syntax, ordinary negative theoremhood, surrogate support, or a classical counterexample object rather than endpoint-use material.


\par\noindent\textbf{Closure and denial burden.} The closure response is role-sensitive. If the content is raw, support, bookkeeping, or outside the fixed carrier, it has no endpoint-defeating force. If it keeps endpoint-defeating force on the same carrier, it enters endpoint-use discipline and route exhaustion. A successful objection must preserve the same $X,p,\alpha$, the same Hodge-standing predicate, the same Chow carrier $\CH^p(X)_\QQ$, and the same rational cycle-class map $\clQ$.


\par\noindent\textbf{Use in the proof spine.} The packet is reused only after the local exact-complement assumption has been opened. It marks a role transition, not a hidden assertion of nonempty fiber.


\subsection{Audit packet 143: Counterexample force as endpoint use (counterexample-force replay)}


\par\noindent\textbf{Claim under audit.} Official endpoint defeat is endpoint use.


\par\noindent\textbf{Dependency path.} This packet belongs to the counterexample-force replay. It depends on the official Hodge route, the fixed cycle-fiber endpoint image, the object/use distinction, and the kernel/A+ endpoint-use machinery. It is never used before the carrier and branch slots are fixed.


\par\noindent\textbf{Hostile-referee risk.} A hostile reader may attempt to treat the relevant content as raw mathematical syntax, ordinary negative theoremhood, surrogate support, or a classical counterexample object rather than endpoint-use material.


\par\noindent\textbf{Closure and denial burden.} The closure response is role-sensitive. If the content is raw, support, bookkeeping, or outside the fixed carrier, it has no endpoint-defeating force. If it keeps endpoint-defeating force on the same carrier, it enters endpoint-use discipline and route exhaustion. A successful objection must preserve the same $X,p,\alpha$, the same Hodge-standing predicate, the same Chow carrier $\CH^p(X)_\QQ$, and the same rational cycle-class map $\clQ$.


\par\noindent\textbf{Use in the proof spine.} The packet is reused only after the local exact-complement assumption has been opened. It marks a role transition, not a hidden assertion of nonempty fiber.


\subsection{Audit packet 144: No autonomous counterexample role (counterexample-force replay)}


\par\noindent\textbf{Claim under audit.} There is no endpoint-defeating branch outside endpoint use.


\par\noindent\textbf{Dependency path.} This packet belongs to the counterexample-force replay. It depends on the official Hodge route, the fixed cycle-fiber endpoint image, the object/use distinction, and the kernel/A+ endpoint-use machinery. It is never used before the carrier and branch slots are fixed.


\par\noindent\textbf{Hostile-referee risk.} A hostile reader may attempt to treat the relevant content as raw mathematical syntax, ordinary negative theoremhood, surrogate support, or a classical counterexample object rather than endpoint-use material.


\par\noindent\textbf{Closure and denial burden.} The closure response is role-sensitive. If the content is raw, support, bookkeeping, or outside the fixed carrier, it has no endpoint-defeating force. If it keeps endpoint-defeating force on the same carrier, it enters endpoint-use discipline and route exhaustion. A successful objection must preserve the same $X,p,\alpha$, the same Hodge-standing predicate, the same Chow carrier $\CH^p(X)_\QQ$, and the same rational cycle-class map $\clQ$.


\par\noindent\textbf{Use in the proof spine.} The packet is reused only after the local exact-complement assumption has been opened. It marks a role transition, not a hidden assertion of nonempty fiber.


\subsection{Audit packet 145: ATS locking (counterexample-force replay)}


\par\noindent\textbf{Claim under audit.} The same endpoint act preserves anchor and tensor data.


\par\noindent\textbf{Dependency path.} This packet belongs to the counterexample-force replay. It depends on the official Hodge route, the fixed cycle-fiber endpoint image, the object/use distinction, and the kernel/A+ endpoint-use machinery. It is never used before the carrier and branch slots are fixed.


\par\noindent\textbf{Hostile-referee risk.} A hostile reader may attempt to treat the relevant content as raw mathematical syntax, ordinary negative theoremhood, surrogate support, or a classical counterexample object rather than endpoint-use material.


\par\noindent\textbf{Closure and denial burden.} The closure response is role-sensitive. If the content is raw, support, bookkeeping, or outside the fixed carrier, it has no endpoint-defeating force. If it keeps endpoint-defeating force on the same carrier, it enters endpoint-use discipline and route exhaustion. A successful objection must preserve the same $X,p,\alpha$, the same Hodge-standing predicate, the same Chow carrier $\CH^p(X)_\QQ$, and the same rational cycle-class map $\clQ$.


\par\noindent\textbf{Use in the proof spine.} The packet is reused only after the local exact-complement assumption has been opened. It marks a role transition, not a hidden assertion of nonempty fiber.


\subsection{Audit packet 146: UEAP no-overreporting (counterexample-force replay)}


\par\noindent\textbf{Claim under audit.} Support-level content cannot be reported as endpoint readout or endpoint defeat.


\par\noindent\textbf{Dependency path.} This packet belongs to the counterexample-force replay. It depends on the official Hodge route, the fixed cycle-fiber endpoint image, the object/use distinction, and the kernel/A+ endpoint-use machinery. It is never used before the carrier and branch slots are fixed.


\par\noindent\textbf{Hostile-referee risk.} A hostile reader may attempt to treat the relevant content as raw mathematical syntax, ordinary negative theoremhood, surrogate support, or a classical counterexample object rather than endpoint-use material.


\par\noindent\textbf{Closure and denial burden.} The closure response is role-sensitive. If the content is raw, support, bookkeeping, or outside the fixed carrier, it has no endpoint-defeating force. If it keeps endpoint-defeating force on the same carrier, it enters endpoint-use discipline and route exhaustion. A successful objection must preserve the same $X,p,\alpha$, the same Hodge-standing predicate, the same Chow carrier $\CH^p(X)_\QQ$, and the same rational cycle-class map $\clQ$.


\par\noindent\textbf{Use in the proof spine.} The packet is reused only after the local exact-complement assumption has been opened. It marks a role transition, not a hidden assertion of nonempty fiber.


\subsection{Audit packet 147: Single cycle-fiber interface (counterexample-force replay)}


\par\noindent\textbf{Claim under audit.} Cycle-fiber occupation is the only positive endpoint-image interface.


\par\noindent\textbf{Dependency path.} This packet belongs to the counterexample-force replay. It depends on the official Hodge route, the fixed cycle-fiber endpoint image, the object/use distinction, and the kernel/A+ endpoint-use machinery. It is never used before the carrier and branch slots are fixed.


\par\noindent\textbf{Hostile-referee risk.} A hostile reader may attempt to treat the relevant content as raw mathematical syntax, ordinary negative theoremhood, surrogate support, or a classical counterexample object rather than endpoint-use material.


\par\noindent\textbf{Closure and denial burden.} The closure response is role-sensitive. If the content is raw, support, bookkeeping, or outside the fixed carrier, it has no endpoint-defeating force. If it keeps endpoint-defeating force on the same carrier, it enters endpoint-use discipline and route exhaustion. A successful objection must preserve the same $X,p,\alpha$, the same Hodge-standing predicate, the same Chow carrier $\CH^p(X)_\QQ$, and the same rational cycle-class map $\clQ$.


\par\noindent\textbf{Use in the proof spine.} The packet is reused only after the local exact-complement assumption has been opened. It marks a role transition, not a hidden assertion of nonempty fiber.


\subsection{Audit packet 148: Native counterforce exhaustion (counterexample-force replay)}


\par\noindent\textbf{Claim under audit.} Native empty-fiber force has no fifth role.


\par\noindent\textbf{Dependency path.} This packet belongs to the counterexample-force replay. It depends on the official Hodge route, the fixed cycle-fiber endpoint image, the object/use distinction, and the kernel/A+ endpoint-use machinery. It is never used before the carrier and branch slots are fixed.


\par\noindent\textbf{Hostile-referee risk.} A hostile reader may attempt to treat the relevant content as raw mathematical syntax, ordinary negative theoremhood, surrogate support, or a classical counterexample object rather than endpoint-use material.


\par\noindent\textbf{Closure and denial burden.} The closure response is role-sensitive. If the content is raw, support, bookkeeping, or outside the fixed carrier, it has no endpoint-defeating force. If it keeps endpoint-defeating force on the same carrier, it enters endpoint-use discipline and route exhaustion. A successful objection must preserve the same $X,p,\alpha$, the same Hodge-standing predicate, the same Chow carrier $\CH^p(X)_\QQ$, and the same rational cycle-class map $\clQ$.


\par\noindent\textbf{Use in the proof spine.} The packet is reused only after the local exact-complement assumption has been opened. It marks a role transition, not a hidden assertion of nonempty fiber.


\subsection{Audit packet 149: Separator induction (counterexample-force replay)}


\par\noindent\textbf{Claim under audit.} Local separator occupation induces the fiber discriminator.


\par\noindent\textbf{Dependency path.} This packet belongs to the counterexample-force replay. It depends on the official Hodge route, the fixed cycle-fiber endpoint image, the object/use distinction, and the kernel/A+ endpoint-use machinery. It is never used before the carrier and branch slots are fixed.


\par\noindent\textbf{Hostile-referee risk.} A hostile reader may attempt to treat the relevant content as raw mathematical syntax, ordinary negative theoremhood, surrogate support, or a classical counterexample object rather than endpoint-use material.


\par\noindent\textbf{Closure and denial burden.} The closure response is role-sensitive. If the content is raw, support, bookkeeping, or outside the fixed carrier, it has no endpoint-defeating force. If it keeps endpoint-defeating force on the same carrier, it enters endpoint-use discipline and route exhaustion. A successful objection must preserve the same $X,p,\alpha$, the same Hodge-standing predicate, the same Chow carrier $\CH^p(X)_\QQ$, and the same rational cycle-class map $\clQ$.


\par\noindent\textbf{Use in the proof spine.} The packet is reused only after the local exact-complement assumption has been opened. It marks a role transition, not a hidden assertion of nonempty fiber.


\subsection{Audit packet 150: A+ no-independent closure (counterexample-force replay)}


\par\noindent\textbf{Claim under audit.} Endpoint use excludes independent same-domain discriminators.


\par\noindent\textbf{Dependency path.} This packet belongs to the counterexample-force replay. It depends on the official Hodge route, the fixed cycle-fiber endpoint image, the object/use distinction, and the kernel/A+ endpoint-use machinery. It is never used before the carrier and branch slots are fixed.


\par\noindent\textbf{Hostile-referee risk.} A hostile reader may attempt to treat the relevant content as raw mathematical syntax, ordinary negative theoremhood, surrogate support, or a classical counterexample object rather than endpoint-use material.


\par\noindent\textbf{Closure and denial burden.} The closure response is role-sensitive. If the content is raw, support, bookkeeping, or outside the fixed carrier, it has no endpoint-defeating force. If it keeps endpoint-defeating force on the same carrier, it enters endpoint-use discipline and route exhaustion. A successful objection must preserve the same $X,p,\alpha$, the same Hodge-standing predicate, the same Chow carrier $\CH^p(X)_\QQ$, and the same rational cycle-class map $\clQ$.


\par\noindent\textbf{Use in the proof spine.} The packet is reused only after the local exact-complement assumption has been opened. It marks a role transition, not a hidden assertion of nonempty fiber.


\subsection{Audit packet 151: Annotation discharge (counterexample-force replay)}


\par\noindent\textbf{Claim under audit.} The exact-complement annotation is not an added premise.


\par\noindent\textbf{Dependency path.} This packet belongs to the counterexample-force replay. It depends on the official Hodge route, the fixed cycle-fiber endpoint image, the object/use distinction, and the kernel/A+ endpoint-use machinery. It is never used before the carrier and branch slots are fixed.


\par\noindent\textbf{Hostile-referee risk.} A hostile reader may attempt to treat the relevant content as raw mathematical syntax, ordinary negative theoremhood, surrogate support, or a classical counterexample object rather than endpoint-use material.


\par\noindent\textbf{Closure and denial burden.} The closure response is role-sensitive. If the content is raw, support, bookkeeping, or outside the fixed carrier, it has no endpoint-defeating force. If it keeps endpoint-defeating force on the same carrier, it enters endpoint-use discipline and route exhaustion. A successful objection must preserve the same $X,p,\alpha$, the same Hodge-standing predicate, the same Chow carrier $\CH^p(X)_\QQ$, and the same rational cycle-class map $\clQ$.


\par\noindent\textbf{Use in the proof spine.} The packet is reused only after the local exact-complement assumption has been opened. It marks a role transition, not a hidden assertion of nonempty fiber.


\subsection{Audit packet 152: Universalization (counterexample-force replay)}


\par\noindent\textbf{Claim under audit.} Arbitrary instance closure yields the global Hodge endpoint.


\par\noindent\textbf{Dependency path.} This packet belongs to the counterexample-force replay. It depends on the official Hodge route, the fixed cycle-fiber endpoint image, the object/use distinction, and the kernel/A+ endpoint-use machinery. It is never used before the carrier and branch slots are fixed.


\par\noindent\textbf{Hostile-referee risk.} A hostile reader may attempt to treat the relevant content as raw mathematical syntax, ordinary negative theoremhood, surrogate support, or a classical counterexample object rather than endpoint-use material.


\par\noindent\textbf{Closure and denial burden.} The closure response is role-sensitive. If the content is raw, support, bookkeeping, or outside the fixed carrier, it has no endpoint-defeating force. If it keeps endpoint-defeating force on the same carrier, it enters endpoint-use discipline and route exhaustion. A successful objection must preserve the same $X,p,\alpha$, the same Hodge-standing predicate, the same Chow carrier $\CH^p(X)_\QQ$, and the same rational cycle-class map $\clQ$.


\par\noindent\textbf{Use in the proof spine.} The packet is reused only after the local exact-complement assumption has been opened. It marks a role transition, not a hidden assertion of nonempty fiber.


\subsection{Audit packet 153: Surrogate separation (counterexample-force replay)}


\par\noindent\textbf{Claim under audit.} Motives and absolute Hodge data are bridge programs unless they supply Chow fiber.


\par\noindent\textbf{Dependency path.} This packet belongs to the counterexample-force replay. It depends on the official Hodge route, the fixed cycle-fiber endpoint image, the object/use distinction, and the kernel/A+ endpoint-use machinery. It is never used before the carrier and branch slots are fixed.


\par\noindent\textbf{Hostile-referee risk.} A hostile reader may attempt to treat the relevant content as raw mathematical syntax, ordinary negative theoremhood, surrogate support, or a classical counterexample object rather than endpoint-use material.


\par\noindent\textbf{Closure and denial burden.} The closure response is role-sensitive. If the content is raw, support, bookkeeping, or outside the fixed carrier, it has no endpoint-defeating force. If it keeps endpoint-defeating force on the same carrier, it enters endpoint-use discipline and route exhaustion. A successful objection must preserve the same $X,p,\alpha$, the same Hodge-standing predicate, the same Chow carrier $\CH^p(X)_\QQ$, and the same rational cycle-class map $\clQ$.


\par\noindent\textbf{Use in the proof spine.} The packet is reused only after the local exact-complement assumption has been opened. It marks a role transition, not a hidden assertion of nonempty fiber.


\subsection{Audit packet 154: Known-case transfer (counterexample-force replay)}


\par\noindent\textbf{Claim under audit.} Known bridge completions do not transfer universally without target gating.


\par\noindent\textbf{Dependency path.} This packet belongs to the counterexample-force replay. It depends on the official Hodge route, the fixed cycle-fiber endpoint image, the object/use distinction, and the kernel/A+ endpoint-use machinery. It is never used before the carrier and branch slots are fixed.


\par\noindent\textbf{Hostile-referee risk.} A hostile reader may attempt to treat the relevant content as raw mathematical syntax, ordinary negative theoremhood, surrogate support, or a classical counterexample object rather than endpoint-use material.


\par\noindent\textbf{Closure and denial burden.} The closure response is role-sensitive. If the content is raw, support, bookkeeping, or outside the fixed carrier, it has no endpoint-defeating force. If it keeps endpoint-defeating force on the same carrier, it enters endpoint-use discipline and route exhaustion. A successful objection must preserve the same $X,p,\alpha$, the same Hodge-standing predicate, the same Chow carrier $\CH^p(X)_\QQ$, and the same rational cycle-class map $\clQ$.


\par\noindent\textbf{Use in the proof spine.} The packet is reused only after the local exact-complement assumption has been opened. It marks a role transition, not a hidden assertion of nonempty fiber.


\subsection{Audit packet 155: Local-to-global nonfactorization (counterexample-force replay)}


\par\noindent\textbf{Claim under audit.} Local cycle data require a global compatibility certificate.


\par\noindent\textbf{Dependency path.} This packet belongs to the counterexample-force replay. It depends on the official Hodge route, the fixed cycle-fiber endpoint image, the object/use distinction, and the kernel/A+ endpoint-use machinery. It is never used before the carrier and branch slots are fixed.


\par\noindent\textbf{Hostile-referee risk.} A hostile reader may attempt to treat the relevant content as raw mathematical syntax, ordinary negative theoremhood, surrogate support, or a classical counterexample object rather than endpoint-use material.


\par\noindent\textbf{Closure and denial burden.} The closure response is role-sensitive. If the content is raw, support, bookkeeping, or outside the fixed carrier, it has no endpoint-defeating force. If it keeps endpoint-defeating force on the same carrier, it enters endpoint-use discipline and route exhaustion. A successful objection must preserve the same $X,p,\alpha$, the same Hodge-standing predicate, the same Chow carrier $\CH^p(X)_\QQ$, and the same rational cycle-class map $\clQ$.


\par\noindent\textbf{Use in the proof spine.} The packet is reused only after the local exact-complement assumption has been opened. It marks a role transition, not a hidden assertion of nonempty fiber.


\subsection{Audit packet 156: No repair (counterexample-force replay)}


\par\noindent\textbf{Claim under audit.} Later cycle discovery creates a new bridge-complete act, not repair of an old one.


\par\noindent\textbf{Dependency path.} This packet belongs to the counterexample-force replay. It depends on the official Hodge route, the fixed cycle-fiber endpoint image, the object/use distinction, and the kernel/A+ endpoint-use machinery. It is never used before the carrier and branch slots are fixed.


\par\noindent\textbf{Hostile-referee risk.} A hostile reader may attempt to treat the relevant content as raw mathematical syntax, ordinary negative theoremhood, surrogate support, or a classical counterexample object rather than endpoint-use material.


\par\noindent\textbf{Closure and denial burden.} The closure response is role-sensitive. If the content is raw, support, bookkeeping, or outside the fixed carrier, it has no endpoint-defeating force. If it keeps endpoint-defeating force on the same carrier, it enters endpoint-use discipline and route exhaustion. A successful objection must preserve the same $X,p,\alpha$, the same Hodge-standing predicate, the same Chow carrier $\CH^p(X)_\QQ$, and the same rational cycle-class map $\clQ$.


\par\noindent\textbf{Use in the proof spine.} The packet is reused only after the local exact-complement assumption has been opened. It marks a role transition, not a hidden assertion of nonempty fiber.


\subsection{Audit packet 157: No selector (counterexample-force replay)}


\par\noindent\textbf{Claim under audit.} Canonical representatives cannot decide readout before the fiber is fixed.


\par\noindent\textbf{Dependency path.} This packet belongs to the counterexample-force replay. It depends on the official Hodge route, the fixed cycle-fiber endpoint image, the object/use distinction, and the kernel/A+ endpoint-use machinery. It is never used before the carrier and branch slots are fixed.


\par\noindent\textbf{Hostile-referee risk.} A hostile reader may attempt to treat the relevant content as raw mathematical syntax, ordinary negative theoremhood, surrogate support, or a classical counterexample object rather than endpoint-use material.


\par\noindent\textbf{Closure and denial burden.} The closure response is role-sensitive. If the content is raw, support, bookkeeping, or outside the fixed carrier, it has no endpoint-defeating force. If it keeps endpoint-defeating force on the same carrier, it enters endpoint-use discipline and route exhaustion. A successful objection must preserve the same $X,p,\alpha$, the same Hodge-standing predicate, the same Chow carrier $\CH^p(X)_\QQ$, and the same rational cycle-class map $\clQ$.


\par\noindent\textbf{Use in the proof spine.} The packet is reused only after the local exact-complement assumption has been opened. It marks a role transition, not a hidden assertion of nonempty fiber.


\subsection{Audit packet 158: Incompleteness firewall (counterexample-force replay)}


\par\noindent\textbf{Claim under audit.} Independence matters only if it re-enters endpoint standing or trajectory.


\par\noindent\textbf{Dependency path.} This packet belongs to the counterexample-force replay. It depends on the official Hodge route, the fixed cycle-fiber endpoint image, the object/use distinction, and the kernel/A+ endpoint-use machinery. It is never used before the carrier and branch slots are fixed.


\par\noindent\textbf{Hostile-referee risk.} A hostile reader may attempt to treat the relevant content as raw mathematical syntax, ordinary negative theoremhood, surrogate support, or a classical counterexample object rather than endpoint-use material.


\par\noindent\textbf{Closure and denial burden.} The closure response is role-sensitive. If the content is raw, support, bookkeeping, or outside the fixed carrier, it has no endpoint-defeating force. If it keeps endpoint-defeating force on the same carrier, it enters endpoint-use discipline and route exhaustion. A successful objection must preserve the same $X,p,\alpha$, the same Hodge-standing predicate, the same Chow carrier $\CH^p(X)_\QQ$, and the same rational cycle-class map $\clQ$.


\par\noindent\textbf{Use in the proof spine.} The packet is reused only after the local exact-complement assumption has been opened. It marks a role transition, not a hidden assertion of nonempty fiber.


\subsection{Audit packet 159: Ordinary theoremhood protection (counterexample-force replay)}


\par\noindent\textbf{Claim under audit.} Negative theoremhood is proof legitimacy, not forbidden endpoint governance.


\par\noindent\textbf{Dependency path.} This packet belongs to the counterexample-force replay. It depends on the official Hodge route, the fixed cycle-fiber endpoint image, the object/use distinction, and the kernel/A+ endpoint-use machinery. It is never used before the carrier and branch slots are fixed.


\par\noindent\textbf{Hostile-referee risk.} A hostile reader may attempt to treat the relevant content as raw mathematical syntax, ordinary negative theoremhood, surrogate support, or a classical counterexample object rather than endpoint-use material.


\par\noindent\textbf{Closure and denial burden.} The closure response is role-sensitive. If the content is raw, support, bookkeeping, or outside the fixed carrier, it has no endpoint-defeating force. If it keeps endpoint-defeating force on the same carrier, it enters endpoint-use discipline and route exhaustion. A successful objection must preserve the same $X,p,\alpha$, the same Hodge-standing predicate, the same Chow carrier $\CH^p(X)_\QQ$, and the same rational cycle-class map $\clQ$.


\par\noindent\textbf{Use in the proof spine.} The packet is reused only after the local exact-complement assumption has been opened. It marks a role transition, not a hidden assertion of nonempty fiber.


\subsection{Audit packet 160: Non-explosiveness (counterexample-force replay)}


\par\noindent\textbf{Claim under audit.} Empty fiber is lawful until endpoint force is added.


\par\noindent\textbf{Dependency path.} This packet belongs to the counterexample-force replay. It depends on the official Hodge route, the fixed cycle-fiber endpoint image, the object/use distinction, and the kernel/A+ endpoint-use machinery. It is never used before the carrier and branch slots are fixed.


\par\noindent\textbf{Hostile-referee risk.} A hostile reader may attempt to treat the relevant content as raw mathematical syntax, ordinary negative theoremhood, surrogate support, or a classical counterexample object rather than endpoint-use material.


\par\noindent\textbf{Closure and denial burden.} The closure response is role-sensitive. If the content is raw, support, bookkeeping, or outside the fixed carrier, it has no endpoint-defeating force. If it keeps endpoint-defeating force on the same carrier, it enters endpoint-use discipline and route exhaustion. A successful objection must preserve the same $X,p,\alpha$, the same Hodge-standing predicate, the same Chow carrier $\CH^p(X)_\QQ$, and the same rational cycle-class map $\clQ$.


\par\noindent\textbf{Use in the proof spine.} The packet is reused only after the local exact-complement assumption has been opened. It marks a role transition, not a hidden assertion of nonempty fiber.


\section{Appendix H: expanded faithful counterexample worksheet}


\subsection{Worksheet item 1: Autonomous endpoint counterexample}


A successful objection must Produce endpoint-defeating empty-fiber force that is not endpoint use. It must remain on the fixed Hodge/cycle carrier. If it changes the carrier, it is domain shift. If it changes no endpoint feature, it is support or bookkeeping. If it preserves endpoint-defeating force, it falls under endpoint-use and native-counterforce exhaustion.


\subsection{Worksheet item 2: Fifth native role}


A successful objection must Produce same-carrier empty-fiber endpoint force that is not bridge-neutralized, support-only, bookkeeping, domain shift, or $D_{\mathrm{fib}}$. It must remain on the fixed Hodge/cycle carrier. If it changes the carrier, it is domain shift. If it changes no endpoint feature, it is support or bookkeeping. If it preserves endpoint-defeating force, it falls under endpoint-use and native-counterforce exhaustion.


\subsection{Worksheet item 3: Surrogate occupation}


A successful objection must Produce a non-Chow surrogate that occupies $F_X(\alpha)$ without bridge. It must remain on the fixed Hodge/cycle carrier. If it changes the carrier, it is domain shift. If it changes no endpoint feature, it is support or bookkeeping. If it preserves endpoint-defeating force, it falls under endpoint-use and native-counterforce exhaustion.


\subsection{Worksheet item 4: Selector authority}


A successful objection must Produce canonical representative selection that decides readout before fiber occupation. It must remain on the fixed Hodge/cycle carrier. If it changes the carrier, it is domain shift. If it changes no endpoint feature, it is support or bookkeeping. If it preserves endpoint-defeating force, it falls under endpoint-use and native-counterforce exhaustion.


\subsection{Worksheet item 5: Repair}


A successful objection must Produce later bridge data that repairs an earlier unbridged report as the same act. It must remain on the fixed Hodge/cycle carrier. If it changes the carrier, it is domain shift. If it changes no endpoint feature, it is support or bookkeeping. If it preserves endpoint-defeating force, it falls under endpoint-use and native-counterforce exhaustion.


\subsection{Worksheet item 6: Local factorization}


A successful objection must Produce local cycle data that becomes a global Chow fiber without compatibility certificate. It must remain on the fixed Hodge/cycle carrier. If it changes the carrier, it is domain shift. If it changes no endpoint feature, it is support or bookkeeping. If it preserves endpoint-defeating force, it falls under endpoint-use and native-counterforce exhaustion.


\subsection{Worksheet item 7: A+ failure}


A successful objection must Produce a same-domain independent discriminator compatible with no-independent closure. It must remain on the fixed Hodge/cycle carrier. If it changes the carrier, it is domain shift. If it changes no endpoint feature, it is support or bookkeeping. If it preserves endpoint-defeating force, it falls under endpoint-use and native-counterforce exhaustion.


\subsection{Worksheet item 8: Projection failure}


A successful objection must Show algebraic readout is not nonempty cycle-fiber status. It must remain on the fixed Hodge/cycle carrier. If it changes the carrier, it is domain shift. If it changes no endpoint feature, it is support or bookkeeping. If it preserves endpoint-defeating force, it falls under endpoint-use and native-counterforce exhaustion.


\subsection{Worksheet item 9: Annotation failure}


A successful objection must Show the local annotation adds an object-level premise. It must remain on the fixed Hodge/cycle carrier. If it changes the carrier, it is domain shift. If it changes no endpoint feature, it is support or bookkeeping. If it preserves endpoint-defeating force, it falls under endpoint-use and native-counterforce exhaustion.


\subsection{Worksheet item 10: Universalization failure}


A successful objection must Show arbitrary fixed-instance nonempty fiber does not yield the universal Hodge endpoint. It must remain on the fixed Hodge/cycle carrier. If it changes the carrier, it is domain shift. If it changes no endpoint feature, it is support or bookkeeping. If it preserves endpoint-defeating force, it falls under endpoint-use and native-counterforce exhaustion.



\section{Appendix I: expanded load-bearing stress-test capsules}

This appendix repeats the load-bearing Hodge endpoint moves under multiple hostile-referee reconstructions. It is intentionally redundant. The purpose is to prevent any decisive transition from appearing only once, or only as a slogan. Each capsule keeps the raw Hodge-theoretic content lawful and audits only the endpoint-use role.


\subsection{Stress test: Official endpoint identity (endpoint-force lens)}

\par\noindent\textbf{Load-bearing claim.} This stress test audits the identification of the Clay-facing endpoint with universal nonempty rational cycle-class fibers. The claim is not allowed to function by implication, atmosphere, or familiarity with AASC vocabulary. It must be available as a same-carrier theorem target inside the Hodge endpoint proof and must be reused only after the carrier, branch, and report slots have been instantiated.

\par\noindent\textbf{Hostile reconstruction.} A hostile referee will try to reconstruct this point in the strongest ordinary Hodge-theoretic language: rational Hodge classes are native cohomological objects, Chow groups and cycle-class maps are ordinary algebraic-geometric carriers, and an empty fiber would be ordinary mathematical falsification. The audit therefore never answers by declaring the underlying mathematical content illicit. It answers only by separating raw content, support status, theoremhood, endpoint force, and the fixed cycle-fiber endpoint image.

\par\noindent\textbf{Closure condition.} The claim survives this stress test only if it preserves the same $X$, $p$, and $\alpha$, the same rational Hodge-standing predicate, the same group $\CH^p(X)_\QQ$, the same map $\clQ$, and the same fiber $F_X(\alpha)$. A proposed escape that changes any of these is a domain or carrier shift. A proposed escape that preserves them while retaining endpoint-defeating force enters endpoint-use discipline and must pass through the native empty-fiber exhaustion.

\par\noindent\textbf{Audit result.} Under the fixed endpoint assumptions, the hostile reconstruction leaves only the declared branches: bridge completion, support-only content, bookkeeping, explicit domain shift, or independent same-domain cycle-fiber endpoint discrimination. This packet is therefore not decorative; it is a repeatable closure capsule for the theorem ladder and denial-burden table.


\subsection{Stress test: Context non-smuggling (endpoint-force lens)}

\par\noindent\textbf{Load-bearing claim.} This stress test audits the claim that the Hodge endpoint context fixes the carrier and slots without asserting nonempty fiber. The claim is not allowed to function by implication, atmosphere, or familiarity with AASC vocabulary. It must be available as a same-carrier theorem target inside the Hodge endpoint proof and must be reused only after the carrier, branch, and report slots have been instantiated.

\par\noindent\textbf{Hostile reconstruction.} A hostile referee will try to reconstruct this point in the strongest ordinary Hodge-theoretic language: rational Hodge classes are native cohomological objects, Chow groups and cycle-class maps are ordinary algebraic-geometric carriers, and an empty fiber would be ordinary mathematical falsification. The audit therefore never answers by declaring the underlying mathematical content illicit. It answers only by separating raw content, support status, theoremhood, endpoint force, and the fixed cycle-fiber endpoint image.

\par\noindent\textbf{Closure condition.} The claim survives this stress test only if it preserves the same $X$, $p$, and $\alpha$, the same rational Hodge-standing predicate, the same group $\CH^p(X)_\QQ$, the same map $\clQ$, and the same fiber $F_X(\alpha)$. A proposed escape that changes any of these is a domain or carrier shift. A proposed escape that preserves them while retaining endpoint-defeating force enters endpoint-use discipline and must pass through the native empty-fiber exhaustion.

\par\noindent\textbf{Audit result.} Under the fixed endpoint assumptions, the hostile reconstruction leaves only the declared branches: bridge completion, support-only content, bookkeeping, explicit domain shift, or independent same-domain cycle-fiber endpoint discrimination. This packet is therefore not decorative; it is a repeatable closure capsule for the theorem ladder and denial-burden table.


\subsection{Stress test: Hodge-standing exactness (endpoint-force lens)}

\par\noindent\textbf{Load-bearing claim.} This stress test audits the representation of rational Hodge standing as cohomological type status. The claim is not allowed to function by implication, atmosphere, or familiarity with AASC vocabulary. It must be available as a same-carrier theorem target inside the Hodge endpoint proof and must be reused only after the carrier, branch, and report slots have been instantiated.

\par\noindent\textbf{Hostile reconstruction.} A hostile referee will try to reconstruct this point in the strongest ordinary Hodge-theoretic language: rational Hodge classes are native cohomological objects, Chow groups and cycle-class maps are ordinary algebraic-geometric carriers, and an empty fiber would be ordinary mathematical falsification. The audit therefore never answers by declaring the underlying mathematical content illicit. It answers only by separating raw content, support status, theoremhood, endpoint force, and the fixed cycle-fiber endpoint image.

\par\noindent\textbf{Closure condition.} The claim survives this stress test only if it preserves the same $X$, $p$, and $\alpha$, the same rational Hodge-standing predicate, the same group $\CH^p(X)_\QQ$, the same map $\clQ$, and the same fiber $F_X(\alpha)$. A proposed escape that changes any of these is a domain or carrier shift. A proposed escape that preserves them while retaining endpoint-defeating force enters endpoint-use discipline and must pass through the native empty-fiber exhaustion.

\par\noindent\textbf{Audit result.} Under the fixed endpoint assumptions, the hostile reconstruction leaves only the declared branches: bridge completion, support-only content, bookkeeping, explicit domain shift, or independent same-domain cycle-fiber endpoint discrimination. This packet is therefore not decorative; it is a repeatable closure capsule for the theorem ladder and denial-burden table.


\subsection{Stress test: Cycle-fiber exactness (endpoint-force lens)}

\par\noindent\textbf{Load-bearing claim.} This stress test audits the representation of algebraic readout by the fiber $F_X(\alpha)$. The claim is not allowed to function by implication, atmosphere, or familiarity with AASC vocabulary. It must be available as a same-carrier theorem target inside the Hodge endpoint proof and must be reused only after the carrier, branch, and report slots have been instantiated.

\par\noindent\textbf{Hostile reconstruction.} A hostile referee will try to reconstruct this point in the strongest ordinary Hodge-theoretic language: rational Hodge classes are native cohomological objects, Chow groups and cycle-class maps are ordinary algebraic-geometric carriers, and an empty fiber would be ordinary mathematical falsification. The audit therefore never answers by declaring the underlying mathematical content illicit. It answers only by separating raw content, support status, theoremhood, endpoint force, and the fixed cycle-fiber endpoint image.

\par\noindent\textbf{Closure condition.} The claim survives this stress test only if it preserves the same $X$, $p$, and $\alpha$, the same rational Hodge-standing predicate, the same group $\CH^p(X)_\QQ$, the same map $\clQ$, and the same fiber $F_X(\alpha)$. A proposed escape that changes any of these is a domain or carrier shift. A proposed escape that preserves them while retaining endpoint-defeating force enters endpoint-use discipline and must pass through the native empty-fiber exhaustion.

\par\noindent\textbf{Audit result.} Under the fixed endpoint assumptions, the hostile reconstruction leaves only the declared branches: bridge completion, support-only content, bookkeeping, explicit domain shift, or independent same-domain cycle-fiber endpoint discrimination. This packet is therefore not decorative; it is a repeatable closure capsule for the theorem ladder and denial-burden table.


\subsection{Stress test: Empty-fiber exactness (endpoint-force lens)}

\par\noindent\textbf{Load-bearing claim.} This stress test audits the negative local branch as fiber non-occupation. The claim is not allowed to function by implication, atmosphere, or familiarity with AASC vocabulary. It must be available as a same-carrier theorem target inside the Hodge endpoint proof and must be reused only after the carrier, branch, and report slots have been instantiated.

\par\noindent\textbf{Hostile reconstruction.} A hostile referee will try to reconstruct this point in the strongest ordinary Hodge-theoretic language: rational Hodge classes are native cohomological objects, Chow groups and cycle-class maps are ordinary algebraic-geometric carriers, and an empty fiber would be ordinary mathematical falsification. The audit therefore never answers by declaring the underlying mathematical content illicit. It answers only by separating raw content, support status, theoremhood, endpoint force, and the fixed cycle-fiber endpoint image.

\par\noindent\textbf{Closure condition.} The claim survives this stress test only if it preserves the same $X$, $p$, and $\alpha$, the same rational Hodge-standing predicate, the same group $\CH^p(X)_\QQ$, the same map $\clQ$, and the same fiber $F_X(\alpha)$. A proposed escape that changes any of these is a domain or carrier shift. A proposed escape that preserves them while retaining endpoint-defeating force enters endpoint-use discipline and must pass through the native empty-fiber exhaustion.

\par\noindent\textbf{Audit result.} Under the fixed endpoint assumptions, the hostile reconstruction leaves only the declared branches: bridge completion, support-only content, bookkeeping, explicit domain shift, or independent same-domain cycle-fiber endpoint discrimination. This packet is therefore not decorative; it is a repeatable closure capsule for the theorem ladder and denial-burden table.


\subsection{Stress test: Counterexample object/use distinction (endpoint-force lens)}

\par\noindent\textbf{Load-bearing claim.} This stress test audits the separation between an object satisfying the negative branch and the act of offering it as endpoint defeat. The claim is not allowed to function by implication, atmosphere, or familiarity with AASC vocabulary. It must be available as a same-carrier theorem target inside the Hodge endpoint proof and must be reused only after the carrier, branch, and report slots have been instantiated.

\par\noindent\textbf{Hostile reconstruction.} A hostile referee will try to reconstruct this point in the strongest ordinary Hodge-theoretic language: rational Hodge classes are native cohomological objects, Chow groups and cycle-class maps are ordinary algebraic-geometric carriers, and an empty fiber would be ordinary mathematical falsification. The audit therefore never answers by declaring the underlying mathematical content illicit. It answers only by separating raw content, support status, theoremhood, endpoint force, and the fixed cycle-fiber endpoint image.

\par\noindent\textbf{Closure condition.} The claim survives this stress test only if it preserves the same $X$, $p$, and $\alpha$, the same rational Hodge-standing predicate, the same group $\CH^p(X)_\QQ$, the same map $\clQ$, and the same fiber $F_X(\alpha)$. A proposed escape that changes any of these is a domain or carrier shift. A proposed escape that preserves them while retaining endpoint-defeating force enters endpoint-use discipline and must pass through the native empty-fiber exhaustion.

\par\noindent\textbf{Audit result.} Under the fixed endpoint assumptions, the hostile reconstruction leaves only the declared branches: bridge completion, support-only content, bookkeeping, explicit domain shift, or independent same-domain cycle-fiber endpoint discrimination. This packet is therefore not decorative; it is a repeatable closure capsule for the theorem ladder and denial-burden table.


\subsection{Stress test: Official counterexample force (endpoint-force lens)}

\par\noindent\textbf{Load-bearing claim.} This stress test audits the nonclassical dependency identifying endpoint-defeating counterexample force with theorem-bearing endpoint use. The claim is not allowed to function by implication, atmosphere, or familiarity with AASC vocabulary. It must be available as a same-carrier theorem target inside the Hodge endpoint proof and must be reused only after the carrier, branch, and report slots have been instantiated.

\par\noindent\textbf{Hostile reconstruction.} A hostile referee will try to reconstruct this point in the strongest ordinary Hodge-theoretic language: rational Hodge classes are native cohomological objects, Chow groups and cycle-class maps are ordinary algebraic-geometric carriers, and an empty fiber would be ordinary mathematical falsification. The audit therefore never answers by declaring the underlying mathematical content illicit. It answers only by separating raw content, support status, theoremhood, endpoint force, and the fixed cycle-fiber endpoint image.

\par\noindent\textbf{Closure condition.} The claim survives this stress test only if it preserves the same $X$, $p$, and $\alpha$, the same rational Hodge-standing predicate, the same group $\CH^p(X)_\QQ$, the same map $\clQ$, and the same fiber $F_X(\alpha)$. A proposed escape that changes any of these is a domain or carrier shift. A proposed escape that preserves them while retaining endpoint-defeating force enters endpoint-use discipline and must pass through the native empty-fiber exhaustion.

\par\noindent\textbf{Audit result.} Under the fixed endpoint assumptions, the hostile reconstruction leaves only the declared branches: bridge completion, support-only content, bookkeeping, explicit domain shift, or independent same-domain cycle-fiber endpoint discrimination. This packet is therefore not decorative; it is a repeatable closure capsule for the theorem ladder and denial-burden table.


\subsection{Stress test: No autonomous counterexample endpoint role (endpoint-force lens)}

\par\noindent\textbf{Load-bearing claim.} This stress test audits the denial of a fifth role outside endpoint use for official endpoint defeat. The claim is not allowed to function by implication, atmosphere, or familiarity with AASC vocabulary. It must be available as a same-carrier theorem target inside the Hodge endpoint proof and must be reused only after the carrier, branch, and report slots have been instantiated.

\par\noindent\textbf{Hostile reconstruction.} A hostile referee will try to reconstruct this point in the strongest ordinary Hodge-theoretic language: rational Hodge classes are native cohomological objects, Chow groups and cycle-class maps are ordinary algebraic-geometric carriers, and an empty fiber would be ordinary mathematical falsification. The audit therefore never answers by declaring the underlying mathematical content illicit. It answers only by separating raw content, support status, theoremhood, endpoint force, and the fixed cycle-fiber endpoint image.

\par\noindent\textbf{Closure condition.} The claim survives this stress test only if it preserves the same $X$, $p$, and $\alpha$, the same rational Hodge-standing predicate, the same group $\CH^p(X)_\QQ$, the same map $\clQ$, and the same fiber $F_X(\alpha)$. A proposed escape that changes any of these is a domain or carrier shift. A proposed escape that preserves them while retaining endpoint-defeating force enters endpoint-use discipline and must pass through the native empty-fiber exhaustion.

\par\noindent\textbf{Audit result.} Under the fixed endpoint assumptions, the hostile reconstruction leaves only the declared branches: bridge completion, support-only content, bookkeeping, explicit domain shift, or independent same-domain cycle-fiber endpoint discrimination. This packet is therefore not decorative; it is a repeatable closure capsule for the theorem ladder and denial-burden table.


\subsection{Stress test: Kernel forcing (endpoint-force lens)}

\par\noindent\textbf{Load-bearing claim.} This stress test audits the derivation of reference, standing, admissibility, and irreversibility from target adequacy. The claim is not allowed to function by implication, atmosphere, or familiarity with AASC vocabulary. It must be available as a same-carrier theorem target inside the Hodge endpoint proof and must be reused only after the carrier, branch, and report slots have been instantiated.

\par\noindent\textbf{Hostile reconstruction.} A hostile referee will try to reconstruct this point in the strongest ordinary Hodge-theoretic language: rational Hodge classes are native cohomological objects, Chow groups and cycle-class maps are ordinary algebraic-geometric carriers, and an empty fiber would be ordinary mathematical falsification. The audit therefore never answers by declaring the underlying mathematical content illicit. It answers only by separating raw content, support status, theoremhood, endpoint force, and the fixed cycle-fiber endpoint image.

\par\noindent\textbf{Closure condition.} The claim survives this stress test only if it preserves the same $X$, $p$, and $\alpha$, the same rational Hodge-standing predicate, the same group $\CH^p(X)_\QQ$, the same map $\clQ$, and the same fiber $F_X(\alpha)$. A proposed escape that changes any of these is a domain or carrier shift. A proposed escape that preserves them while retaining endpoint-defeating force enters endpoint-use discipline and must pass through the native empty-fiber exhaustion.

\par\noindent\textbf{Audit result.} Under the fixed endpoint assumptions, the hostile reconstruction leaves only the declared branches: bridge completion, support-only content, bookkeeping, explicit domain shift, or independent same-domain cycle-fiber endpoint discrimination. This packet is therefore not decorative; it is a repeatable closure capsule for the theorem ladder and denial-burden table.


\subsection{Stress test: A+ consequence layer (endpoint-force lens)}

\par\noindent\textbf{Load-bearing claim.} This stress test audits the fixed-domain consequence package binding endpoint use. The claim is not allowed to function by implication, atmosphere, or familiarity with AASC vocabulary. It must be available as a same-carrier theorem target inside the Hodge endpoint proof and must be reused only after the carrier, branch, and report slots have been instantiated.

\par\noindent\textbf{Hostile reconstruction.} A hostile referee will try to reconstruct this point in the strongest ordinary Hodge-theoretic language: rational Hodge classes are native cohomological objects, Chow groups and cycle-class maps are ordinary algebraic-geometric carriers, and an empty fiber would be ordinary mathematical falsification. The audit therefore never answers by declaring the underlying mathematical content illicit. It answers only by separating raw content, support status, theoremhood, endpoint force, and the fixed cycle-fiber endpoint image.

\par\noindent\textbf{Closure condition.} The claim survives this stress test only if it preserves the same $X$, $p$, and $\alpha$, the same rational Hodge-standing predicate, the same group $\CH^p(X)_\QQ$, the same map $\clQ$, and the same fiber $F_X(\alpha)$. A proposed escape that changes any of these is a domain or carrier shift. A proposed escape that preserves them while retaining endpoint-defeating force enters endpoint-use discipline and must pass through the native empty-fiber exhaustion.

\par\noindent\textbf{Audit result.} Under the fixed endpoint assumptions, the hostile reconstruction leaves only the declared branches: bridge completion, support-only content, bookkeeping, explicit domain shift, or independent same-domain cycle-fiber endpoint discrimination. This packet is therefore not decorative; it is a repeatable closure capsule for the theorem ladder and denial-burden table.


\subsection{Stress test: ATS locking (endpoint-force lens)}

\par\noindent\textbf{Load-bearing claim.} This stress test audits same target, carrier, branch, readout trajectory, and no second classifier. The claim is not allowed to function by implication, atmosphere, or familiarity with AASC vocabulary. It must be available as a same-carrier theorem target inside the Hodge endpoint proof and must be reused only after the carrier, branch, and report slots have been instantiated.

\par\noindent\textbf{Hostile reconstruction.} A hostile referee will try to reconstruct this point in the strongest ordinary Hodge-theoretic language: rational Hodge classes are native cohomological objects, Chow groups and cycle-class maps are ordinary algebraic-geometric carriers, and an empty fiber would be ordinary mathematical falsification. The audit therefore never answers by declaring the underlying mathematical content illicit. It answers only by separating raw content, support status, theoremhood, endpoint force, and the fixed cycle-fiber endpoint image.

\par\noindent\textbf{Closure condition.} The claim survives this stress test only if it preserves the same $X$, $p$, and $\alpha$, the same rational Hodge-standing predicate, the same group $\CH^p(X)_\QQ$, the same map $\clQ$, and the same fiber $F_X(\alpha)$. A proposed escape that changes any of these is a domain or carrier shift. A proposed escape that preserves them while retaining endpoint-defeating force enters endpoint-use discipline and must pass through the native empty-fiber exhaustion.

\par\noindent\textbf{Audit result.} Under the fixed endpoint assumptions, the hostile reconstruction leaves only the declared branches: bridge completion, support-only content, bookkeeping, explicit domain shift, or independent same-domain cycle-fiber endpoint discrimination. This packet is therefore not decorative; it is a repeatable closure capsule for the theorem ladder and denial-burden table.


\subsection{Stress test: UEAP report discipline (endpoint-force lens)}

\par\noindent\textbf{Load-bearing claim.} This stress test audits no report exceeding carrier, bridge, and support status. The claim is not allowed to function by implication, atmosphere, or familiarity with AASC vocabulary. It must be available as a same-carrier theorem target inside the Hodge endpoint proof and must be reused only after the carrier, branch, and report slots have been instantiated.

\par\noindent\textbf{Hostile reconstruction.} A hostile referee will try to reconstruct this point in the strongest ordinary Hodge-theoretic language: rational Hodge classes are native cohomological objects, Chow groups and cycle-class maps are ordinary algebraic-geometric carriers, and an empty fiber would be ordinary mathematical falsification. The audit therefore never answers by declaring the underlying mathematical content illicit. It answers only by separating raw content, support status, theoremhood, endpoint force, and the fixed cycle-fiber endpoint image.

\par\noindent\textbf{Closure condition.} The claim survives this stress test only if it preserves the same $X$, $p$, and $\alpha$, the same rational Hodge-standing predicate, the same group $\CH^p(X)_\QQ$, the same map $\clQ$, and the same fiber $F_X(\alpha)$. A proposed escape that changes any of these is a domain or carrier shift. A proposed escape that preserves them while retaining endpoint-defeating force enters endpoint-use discipline and must pass through the native empty-fiber exhaustion.

\par\noindent\textbf{Audit result.} Under the fixed endpoint assumptions, the hostile reconstruction leaves only the declared branches: bridge completion, support-only content, bookkeeping, explicit domain shift, or independent same-domain cycle-fiber endpoint discrimination. This packet is therefore not decorative; it is a repeatable closure capsule for the theorem ladder and denial-burden table.


\subsection{Stress test: AMetric no-selector boundary (endpoint-force lens)}

\par\noindent\textbf{Load-bearing claim.} This stress test audits no preferred representative or scalar ranking as pre-bridge authority. The claim is not allowed to function by implication, atmosphere, or familiarity with AASC vocabulary. It must be available as a same-carrier theorem target inside the Hodge endpoint proof and must be reused only after the carrier, branch, and report slots have been instantiated.

\par\noindent\textbf{Hostile reconstruction.} A hostile referee will try to reconstruct this point in the strongest ordinary Hodge-theoretic language: rational Hodge classes are native cohomological objects, Chow groups and cycle-class maps are ordinary algebraic-geometric carriers, and an empty fiber would be ordinary mathematical falsification. The audit therefore never answers by declaring the underlying mathematical content illicit. It answers only by separating raw content, support status, theoremhood, endpoint force, and the fixed cycle-fiber endpoint image.

\par\noindent\textbf{Closure condition.} The claim survives this stress test only if it preserves the same $X$, $p$, and $\alpha$, the same rational Hodge-standing predicate, the same group $\CH^p(X)_\QQ$, the same map $\clQ$, and the same fiber $F_X(\alpha)$. A proposed escape that changes any of these is a domain or carrier shift. A proposed escape that preserves them while retaining endpoint-defeating force enters endpoint-use discipline and must pass through the native empty-fiber exhaustion.

\par\noindent\textbf{Audit result.} Under the fixed endpoint assumptions, the hostile reconstruction leaves only the declared branches: bridge completion, support-only content, bookkeeping, explicit domain shift, or independent same-domain cycle-fiber endpoint discrimination. This packet is therefore not decorative; it is a repeatable closure capsule for the theorem ladder and denial-burden table.


\subsection{Stress test: Surrogate carrier separation (endpoint-force lens)}

\par\noindent\textbf{Load-bearing claim.} This stress test audits motivic, variational, or obstruction objects as bridge support rather than readout. The claim is not allowed to function by implication, atmosphere, or familiarity with AASC vocabulary. It must be available as a same-carrier theorem target inside the Hodge endpoint proof and must be reused only after the carrier, branch, and report slots have been instantiated.

\par\noindent\textbf{Hostile reconstruction.} A hostile referee will try to reconstruct this point in the strongest ordinary Hodge-theoretic language: rational Hodge classes are native cohomological objects, Chow groups and cycle-class maps are ordinary algebraic-geometric carriers, and an empty fiber would be ordinary mathematical falsification. The audit therefore never answers by declaring the underlying mathematical content illicit. It answers only by separating raw content, support status, theoremhood, endpoint force, and the fixed cycle-fiber endpoint image.

\par\noindent\textbf{Closure condition.} The claim survives this stress test only if it preserves the same $X$, $p$, and $\alpha$, the same rational Hodge-standing predicate, the same group $\CH^p(X)_\QQ$, the same map $\clQ$, and the same fiber $F_X(\alpha)$. A proposed escape that changes any of these is a domain or carrier shift. A proposed escape that preserves them while retaining endpoint-defeating force enters endpoint-use discipline and must pass through the native empty-fiber exhaustion.

\par\noindent\textbf{Audit result.} Under the fixed endpoint assumptions, the hostile reconstruction leaves only the declared branches: bridge completion, support-only content, bookkeeping, explicit domain shift, or independent same-domain cycle-fiber endpoint discrimination. This packet is therefore not decorative; it is a repeatable closure capsule for the theorem ladder and denial-burden table.


\subsection{Stress test: No post-selection repair (endpoint-force lens)}

\par\noindent\textbf{Load-bearing claim.} This stress test audits later cycle discovery as fresh bridge completion rather than same-act repair. The claim is not allowed to function by implication, atmosphere, or familiarity with AASC vocabulary. It must be available as a same-carrier theorem target inside the Hodge endpoint proof and must be reused only after the carrier, branch, and report slots have been instantiated.

\par\noindent\textbf{Hostile reconstruction.} A hostile referee will try to reconstruct this point in the strongest ordinary Hodge-theoretic language: rational Hodge classes are native cohomological objects, Chow groups and cycle-class maps are ordinary algebraic-geometric carriers, and an empty fiber would be ordinary mathematical falsification. The audit therefore never answers by declaring the underlying mathematical content illicit. It answers only by separating raw content, support status, theoremhood, endpoint force, and the fixed cycle-fiber endpoint image.

\par\noindent\textbf{Closure condition.} The claim survives this stress test only if it preserves the same $X$, $p$, and $\alpha$, the same rational Hodge-standing predicate, the same group $\CH^p(X)_\QQ$, the same map $\clQ$, and the same fiber $F_X(\alpha)$. A proposed escape that changes any of these is a domain or carrier shift. A proposed escape that preserves them while retaining endpoint-defeating force enters endpoint-use discipline and must pass through the native empty-fiber exhaustion.

\par\noindent\textbf{Audit result.} Under the fixed endpoint assumptions, the hostile reconstruction leaves only the declared branches: bridge completion, support-only content, bookkeeping, explicit domain shift, or independent same-domain cycle-fiber endpoint discrimination. This packet is therefore not decorative; it is a repeatable closure capsule for the theorem ladder and denial-burden table.


\subsection{Stress test: No carrier transfer (endpoint-force lens)}

\par\noindent\textbf{Load-bearing claim.} This stress test audits known cases and neighboring domains not carrying endpoint standing without target inclusion. The claim is not allowed to function by implication, atmosphere, or familiarity with AASC vocabulary. It must be available as a same-carrier theorem target inside the Hodge endpoint proof and must be reused only after the carrier, branch, and report slots have been instantiated.

\par\noindent\textbf{Hostile reconstruction.} A hostile referee will try to reconstruct this point in the strongest ordinary Hodge-theoretic language: rational Hodge classes are native cohomological objects, Chow groups and cycle-class maps are ordinary algebraic-geometric carriers, and an empty fiber would be ordinary mathematical falsification. The audit therefore never answers by declaring the underlying mathematical content illicit. It answers only by separating raw content, support status, theoremhood, endpoint force, and the fixed cycle-fiber endpoint image.

\par\noindent\textbf{Closure condition.} The claim survives this stress test only if it preserves the same $X$, $p$, and $\alpha$, the same rational Hodge-standing predicate, the same group $\CH^p(X)_\QQ$, the same map $\clQ$, and the same fiber $F_X(\alpha)$. A proposed escape that changes any of these is a domain or carrier shift. A proposed escape that preserves them while retaining endpoint-defeating force enters endpoint-use discipline and must pass through the native empty-fiber exhaustion.

\par\noindent\textbf{Audit result.} Under the fixed endpoint assumptions, the hostile reconstruction leaves only the declared branches: bridge completion, support-only content, bookkeeping, explicit domain shift, or independent same-domain cycle-fiber endpoint discrimination. This packet is therefore not decorative; it is a repeatable closure capsule for the theorem ladder and denial-burden table.


\subsection{Stress test: No generators (endpoint-force lens)}

\par\noindent\textbf{Load-bearing claim.} This stress test audits Hodge standing not generating Chow-fiber nonemptiness. The claim is not allowed to function by implication, atmosphere, or familiarity with AASC vocabulary. It must be available as a same-carrier theorem target inside the Hodge endpoint proof and must be reused only after the carrier, branch, and report slots have been instantiated.

\par\noindent\textbf{Hostile reconstruction.} A hostile referee will try to reconstruct this point in the strongest ordinary Hodge-theoretic language: rational Hodge classes are native cohomological objects, Chow groups and cycle-class maps are ordinary algebraic-geometric carriers, and an empty fiber would be ordinary mathematical falsification. The audit therefore never answers by declaring the underlying mathematical content illicit. It answers only by separating raw content, support status, theoremhood, endpoint force, and the fixed cycle-fiber endpoint image.

\par\noindent\textbf{Closure condition.} The claim survives this stress test only if it preserves the same $X$, $p$, and $\alpha$, the same rational Hodge-standing predicate, the same group $\CH^p(X)_\QQ$, the same map $\clQ$, and the same fiber $F_X(\alpha)$. A proposed escape that changes any of these is a domain or carrier shift. A proposed escape that preserves them while retaining endpoint-defeating force enters endpoint-use discipline and must pass through the native empty-fiber exhaustion.

\par\noindent\textbf{Audit result.} Under the fixed endpoint assumptions, the hostile reconstruction leaves only the declared branches: bridge completion, support-only content, bookkeeping, explicit domain shift, or independent same-domain cycle-fiber endpoint discrimination. This packet is therefore not decorative; it is a repeatable closure capsule for the theorem ladder and denial-burden table.


\subsection{Stress test: No silent redescription (endpoint-force lens)}

\par\noindent\textbf{Load-bearing claim.} This stress test audits no replacement of Chow-fiber readout by another carrier under same label. The claim is not allowed to function by implication, atmosphere, or familiarity with AASC vocabulary. It must be available as a same-carrier theorem target inside the Hodge endpoint proof and must be reused only after the carrier, branch, and report slots have been instantiated.

\par\noindent\textbf{Hostile reconstruction.} A hostile referee will try to reconstruct this point in the strongest ordinary Hodge-theoretic language: rational Hodge classes are native cohomological objects, Chow groups and cycle-class maps are ordinary algebraic-geometric carriers, and an empty fiber would be ordinary mathematical falsification. The audit therefore never answers by declaring the underlying mathematical content illicit. It answers only by separating raw content, support status, theoremhood, endpoint force, and the fixed cycle-fiber endpoint image.

\par\noindent\textbf{Closure condition.} The claim survives this stress test only if it preserves the same $X$, $p$, and $\alpha$, the same rational Hodge-standing predicate, the same group $\CH^p(X)_\QQ$, the same map $\clQ$, and the same fiber $F_X(\alpha)$. A proposed escape that changes any of these is a domain or carrier shift. A proposed escape that preserves them while retaining endpoint-defeating force enters endpoint-use discipline and must pass through the native empty-fiber exhaustion.

\par\noindent\textbf{Audit result.} Under the fixed endpoint assumptions, the hostile reconstruction leaves only the declared branches: bridge completion, support-only content, bookkeeping, explicit domain shift, or independent same-domain cycle-fiber endpoint discrimination. This packet is therefore not decorative; it is a repeatable closure capsule for the theorem ladder and denial-burden table.


\subsection{Stress test: Non-factorizability (endpoint-force lens)}

\par\noindent\textbf{Load-bearing claim.} This stress test audits local bridge fragments not composing into global readout without compatibility. The claim is not allowed to function by implication, atmosphere, or familiarity with AASC vocabulary. It must be available as a same-carrier theorem target inside the Hodge endpoint proof and must be reused only after the carrier, branch, and report slots have been instantiated.

\par\noindent\textbf{Hostile reconstruction.} A hostile referee will try to reconstruct this point in the strongest ordinary Hodge-theoretic language: rational Hodge classes are native cohomological objects, Chow groups and cycle-class maps are ordinary algebraic-geometric carriers, and an empty fiber would be ordinary mathematical falsification. The audit therefore never answers by declaring the underlying mathematical content illicit. It answers only by separating raw content, support status, theoremhood, endpoint force, and the fixed cycle-fiber endpoint image.

\par\noindent\textbf{Closure condition.} The claim survives this stress test only if it preserves the same $X$, $p$, and $\alpha$, the same rational Hodge-standing predicate, the same group $\CH^p(X)_\QQ$, the same map $\clQ$, and the same fiber $F_X(\alpha)$. A proposed escape that changes any of these is a domain or carrier shift. A proposed escape that preserves them while retaining endpoint-defeating force enters endpoint-use discipline and must pass through the native empty-fiber exhaustion.

\par\noindent\textbf{Audit result.} Under the fixed endpoint assumptions, the hostile reconstruction leaves only the declared branches: bridge completion, support-only content, bookkeeping, explicit domain shift, or independent same-domain cycle-fiber endpoint discrimination. This packet is therefore not decorative; it is a repeatable closure capsule for the theorem ladder and denial-burden table.


\subsection{Stress test: Role-occupancy closure (endpoint-force lens)}

\par\noindent\textbf{Load-bearing claim.} This stress test audits nonempty fiber as role occupation rather than unique representative selection. The claim is not allowed to function by implication, atmosphere, or familiarity with AASC vocabulary. It must be available as a same-carrier theorem target inside the Hodge endpoint proof and must be reused only after the carrier, branch, and report slots have been instantiated.

\par\noindent\textbf{Hostile reconstruction.} A hostile referee will try to reconstruct this point in the strongest ordinary Hodge-theoretic language: rational Hodge classes are native cohomological objects, Chow groups and cycle-class maps are ordinary algebraic-geometric carriers, and an empty fiber would be ordinary mathematical falsification. The audit therefore never answers by declaring the underlying mathematical content illicit. It answers only by separating raw content, support status, theoremhood, endpoint force, and the fixed cycle-fiber endpoint image.

\par\noindent\textbf{Closure condition.} The claim survives this stress test only if it preserves the same $X$, $p$, and $\alpha$, the same rational Hodge-standing predicate, the same group $\CH^p(X)_\QQ$, the same map $\clQ$, and the same fiber $F_X(\alpha)$. A proposed escape that changes any of these is a domain or carrier shift. A proposed escape that preserves them while retaining endpoint-defeating force enters endpoint-use discipline and must pass through the native empty-fiber exhaustion.

\par\noindent\textbf{Audit result.} Under the fixed endpoint assumptions, the hostile reconstruction leaves only the declared branches: bridge completion, support-only content, bookkeeping, explicit domain shift, or independent same-domain cycle-fiber endpoint discrimination. This packet is therefore not decorative; it is a repeatable closure capsule for the theorem ladder and denial-burden table.


\subsection{Stress test: Cycle-fiber single interface (endpoint-force lens)}

\par\noindent\textbf{Load-bearing claim.} This stress test audits the single endpoint-image interface for algebraic readout. The claim is not allowed to function by implication, atmosphere, or familiarity with AASC vocabulary. It must be available as a same-carrier theorem target inside the Hodge endpoint proof and must be reused only after the carrier, branch, and report slots have been instantiated.

\par\noindent\textbf{Hostile reconstruction.} A hostile referee will try to reconstruct this point in the strongest ordinary Hodge-theoretic language: rational Hodge classes are native cohomological objects, Chow groups and cycle-class maps are ordinary algebraic-geometric carriers, and an empty fiber would be ordinary mathematical falsification. The audit therefore never answers by declaring the underlying mathematical content illicit. It answers only by separating raw content, support status, theoremhood, endpoint force, and the fixed cycle-fiber endpoint image.

\par\noindent\textbf{Closure condition.} The claim survives this stress test only if it preserves the same $X$, $p$, and $\alpha$, the same rational Hodge-standing predicate, the same group $\CH^p(X)_\QQ$, the same map $\clQ$, and the same fiber $F_X(\alpha)$. A proposed escape that changes any of these is a domain or carrier shift. A proposed escape that preserves them while retaining endpoint-defeating force enters endpoint-use discipline and must pass through the native empty-fiber exhaustion.

\par\noindent\textbf{Audit result.} Under the fixed endpoint assumptions, the hostile reconstruction leaves only the declared branches: bridge completion, support-only content, bookkeeping, explicit domain shift, or independent same-domain cycle-fiber endpoint discrimination. This packet is therefore not decorative; it is a repeatable closure capsule for the theorem ladder and denial-burden table.


\subsection{Stress test: Coequal carrier distinction (endpoint-force lens)}

\par\noindent\textbf{Load-bearing claim.} This stress test audits separating broad outcome carrier from cycle-fiber image carrier. The claim is not allowed to function by implication, atmosphere, or familiarity with AASC vocabulary. It must be available as a same-carrier theorem target inside the Hodge endpoint proof and must be reused only after the carrier, branch, and report slots have been instantiated.

\par\noindent\textbf{Hostile reconstruction.} A hostile referee will try to reconstruct this point in the strongest ordinary Hodge-theoretic language: rational Hodge classes are native cohomological objects, Chow groups and cycle-class maps are ordinary algebraic-geometric carriers, and an empty fiber would be ordinary mathematical falsification. The audit therefore never answers by declaring the underlying mathematical content illicit. It answers only by separating raw content, support status, theoremhood, endpoint force, and the fixed cycle-fiber endpoint image.

\par\noindent\textbf{Closure condition.} The claim survives this stress test only if it preserves the same $X$, $p$, and $\alpha$, the same rational Hodge-standing predicate, the same group $\CH^p(X)_\QQ$, the same map $\clQ$, and the same fiber $F_X(\alpha)$. A proposed escape that changes any of these is a domain or carrier shift. A proposed escape that preserves them while retaining endpoint-defeating force enters endpoint-use discipline and must pass through the native empty-fiber exhaustion.

\par\noindent\textbf{Audit result.} Under the fixed endpoint assumptions, the hostile reconstruction leaves only the declared branches: bridge completion, support-only content, bookkeeping, explicit domain shift, or independent same-domain cycle-fiber endpoint discrimination. This packet is therefore not decorative; it is a repeatable closure capsule for the theorem ladder and denial-burden table.


\subsection{Stress test: Native empty-fiber exhaustion (endpoint-force lens)}

\par\noindent\textbf{Load-bearing claim.} This stress test audits bridge-neutralized, support-only, bookkeeping, domain shift, or discriminator. The claim is not allowed to function by implication, atmosphere, or familiarity with AASC vocabulary. It must be available as a same-carrier theorem target inside the Hodge endpoint proof and must be reused only after the carrier, branch, and report slots have been instantiated.

\par\noindent\textbf{Hostile reconstruction.} A hostile referee will try to reconstruct this point in the strongest ordinary Hodge-theoretic language: rational Hodge classes are native cohomological objects, Chow groups and cycle-class maps are ordinary algebraic-geometric carriers, and an empty fiber would be ordinary mathematical falsification. The audit therefore never answers by declaring the underlying mathematical content illicit. It answers only by separating raw content, support status, theoremhood, endpoint force, and the fixed cycle-fiber endpoint image.

\par\noindent\textbf{Closure condition.} The claim survives this stress test only if it preserves the same $X$, $p$, and $\alpha$, the same rational Hodge-standing predicate, the same group $\CH^p(X)_\QQ$, the same map $\clQ$, and the same fiber $F_X(\alpha)$. A proposed escape that changes any of these is a domain or carrier shift. A proposed escape that preserves them while retaining endpoint-defeating force enters endpoint-use discipline and must pass through the native empty-fiber exhaustion.

\par\noindent\textbf{Audit result.} Under the fixed endpoint assumptions, the hostile reconstruction leaves only the declared branches: bridge completion, support-only content, bookkeeping, explicit domain shift, or independent same-domain cycle-fiber endpoint discrimination. This packet is therefore not decorative; it is a repeatable closure capsule for the theorem ladder and denial-burden table.


\subsection{Stress test: Bridge-neutralized branch (endpoint-force lens)}

\par\noindent\textbf{Load-bearing claim.} This stress test audits why a bridge-neutralized negative branch is not live endpoint counterforce. The claim is not allowed to function by implication, atmosphere, or familiarity with AASC vocabulary. It must be available as a same-carrier theorem target inside the Hodge endpoint proof and must be reused only after the carrier, branch, and report slots have been instantiated.

\par\noindent\textbf{Hostile reconstruction.} A hostile referee will try to reconstruct this point in the strongest ordinary Hodge-theoretic language: rational Hodge classes are native cohomological objects, Chow groups and cycle-class maps are ordinary algebraic-geometric carriers, and an empty fiber would be ordinary mathematical falsification. The audit therefore never answers by declaring the underlying mathematical content illicit. It answers only by separating raw content, support status, theoremhood, endpoint force, and the fixed cycle-fiber endpoint image.

\par\noindent\textbf{Closure condition.} The claim survives this stress test only if it preserves the same $X$, $p$, and $\alpha$, the same rational Hodge-standing predicate, the same group $\CH^p(X)_\QQ$, the same map $\clQ$, and the same fiber $F_X(\alpha)$. A proposed escape that changes any of these is a domain or carrier shift. A proposed escape that preserves them while retaining endpoint-defeating force enters endpoint-use discipline and must pass through the native empty-fiber exhaustion.

\par\noindent\textbf{Audit result.} Under the fixed endpoint assumptions, the hostile reconstruction leaves only the declared branches: bridge completion, support-only content, bookkeeping, explicit domain shift, or independent same-domain cycle-fiber endpoint discrimination. This packet is therefore not decorative; it is a repeatable closure capsule for the theorem ladder and denial-burden table.


\subsection{Stress test: Support-only branch (endpoint-force lens)}

\par\noindent\textbf{Load-bearing claim.} This stress test audits why support status cannot defeat endpoint closure. The claim is not allowed to function by implication, atmosphere, or familiarity with AASC vocabulary. It must be available as a same-carrier theorem target inside the Hodge endpoint proof and must be reused only after the carrier, branch, and report slots have been instantiated.

\par\noindent\textbf{Hostile reconstruction.} A hostile referee will try to reconstruct this point in the strongest ordinary Hodge-theoretic language: rational Hodge classes are native cohomological objects, Chow groups and cycle-class maps are ordinary algebraic-geometric carriers, and an empty fiber would be ordinary mathematical falsification. The audit therefore never answers by declaring the underlying mathematical content illicit. It answers only by separating raw content, support status, theoremhood, endpoint force, and the fixed cycle-fiber endpoint image.

\par\noindent\textbf{Closure condition.} The claim survives this stress test only if it preserves the same $X$, $p$, and $\alpha$, the same rational Hodge-standing predicate, the same group $\CH^p(X)_\QQ$, the same map $\clQ$, and the same fiber $F_X(\alpha)$. A proposed escape that changes any of these is a domain or carrier shift. A proposed escape that preserves them while retaining endpoint-defeating force enters endpoint-use discipline and must pass through the native empty-fiber exhaustion.

\par\noindent\textbf{Audit result.} Under the fixed endpoint assumptions, the hostile reconstruction leaves only the declared branches: bridge completion, support-only content, bookkeeping, explicit domain shift, or independent same-domain cycle-fiber endpoint discrimination. This packet is therefore not decorative; it is a repeatable closure capsule for the theorem ladder and denial-burden table.


\subsection{Stress test: Bookkeeping branch (endpoint-force lens)}

\par\noindent\textbf{Load-bearing claim.} This stress test audits why labels and ledgers do not determine endpoint status. The claim is not allowed to function by implication, atmosphere, or familiarity with AASC vocabulary. It must be available as a same-carrier theorem target inside the Hodge endpoint proof and must be reused only after the carrier, branch, and report slots have been instantiated.

\par\noindent\textbf{Hostile reconstruction.} A hostile referee will try to reconstruct this point in the strongest ordinary Hodge-theoretic language: rational Hodge classes are native cohomological objects, Chow groups and cycle-class maps are ordinary algebraic-geometric carriers, and an empty fiber would be ordinary mathematical falsification. The audit therefore never answers by declaring the underlying mathematical content illicit. It answers only by separating raw content, support status, theoremhood, endpoint force, and the fixed cycle-fiber endpoint image.

\par\noindent\textbf{Closure condition.} The claim survives this stress test only if it preserves the same $X$, $p$, and $\alpha$, the same rational Hodge-standing predicate, the same group $\CH^p(X)_\QQ$, the same map $\clQ$, and the same fiber $F_X(\alpha)$. A proposed escape that changes any of these is a domain or carrier shift. A proposed escape that preserves them while retaining endpoint-defeating force enters endpoint-use discipline and must pass through the native empty-fiber exhaustion.

\par\noindent\textbf{Audit result.} Under the fixed endpoint assumptions, the hostile reconstruction leaves only the declared branches: bridge completion, support-only content, bookkeeping, explicit domain shift, or independent same-domain cycle-fiber endpoint discrimination. This packet is therefore not decorative; it is a repeatable closure capsule for the theorem ladder and denial-burden table.


\subsection{Stress test: Domain-shift branch (endpoint-force lens)}

\par\noindent\textbf{Load-bearing claim.} This stress test audits why changing X,p,alpha, coefficients, equivalence, or carrier exits the fixed endpoint. The claim is not allowed to function by implication, atmosphere, or familiarity with AASC vocabulary. It must be available as a same-carrier theorem target inside the Hodge endpoint proof and must be reused only after the carrier, branch, and report slots have been instantiated.

\par\noindent\textbf{Hostile reconstruction.} A hostile referee will try to reconstruct this point in the strongest ordinary Hodge-theoretic language: rational Hodge classes are native cohomological objects, Chow groups and cycle-class maps are ordinary algebraic-geometric carriers, and an empty fiber would be ordinary mathematical falsification. The audit therefore never answers by declaring the underlying mathematical content illicit. It answers only by separating raw content, support status, theoremhood, endpoint force, and the fixed cycle-fiber endpoint image.

\par\noindent\textbf{Closure condition.} The claim survives this stress test only if it preserves the same $X$, $p$, and $\alpha$, the same rational Hodge-standing predicate, the same group $\CH^p(X)_\QQ$, the same map $\clQ$, and the same fiber $F_X(\alpha)$. A proposed escape that changes any of these is a domain or carrier shift. A proposed escape that preserves them while retaining endpoint-defeating force enters endpoint-use discipline and must pass through the native empty-fiber exhaustion.

\par\noindent\textbf{Audit result.} Under the fixed endpoint assumptions, the hostile reconstruction leaves only the declared branches: bridge completion, support-only content, bookkeeping, explicit domain shift, or independent same-domain cycle-fiber endpoint discrimination. This packet is therefore not decorative; it is a repeatable closure capsule for the theorem ladder and denial-burden table.


\subsection{Stress test: Cycle-fiber discriminator (endpoint-force lens)}

\par\noindent\textbf{Load-bearing claim.} This stress test audits the independent same-domain discriminator induced by endpoint-standing empty-fiber force. The claim is not allowed to function by implication, atmosphere, or familiarity with AASC vocabulary. It must be available as a same-carrier theorem target inside the Hodge endpoint proof and must be reused only after the carrier, branch, and report slots have been instantiated.

\par\noindent\textbf{Hostile reconstruction.} A hostile referee will try to reconstruct this point in the strongest ordinary Hodge-theoretic language: rational Hodge classes are native cohomological objects, Chow groups and cycle-class maps are ordinary algebraic-geometric carriers, and an empty fiber would be ordinary mathematical falsification. The audit therefore never answers by declaring the underlying mathematical content illicit. It answers only by separating raw content, support status, theoremhood, endpoint force, and the fixed cycle-fiber endpoint image.

\par\noindent\textbf{Closure condition.} The claim survives this stress test only if it preserves the same $X$, $p$, and $\alpha$, the same rational Hodge-standing predicate, the same group $\CH^p(X)_\QQ$, the same map $\clQ$, and the same fiber $F_X(\alpha)$. A proposed escape that changes any of these is a domain or carrier shift. A proposed escape that preserves them while retaining endpoint-defeating force enters endpoint-use discipline and must pass through the native empty-fiber exhaustion.

\par\noindent\textbf{Audit result.} Under the fixed endpoint assumptions, the hostile reconstruction leaves only the declared branches: bridge completion, support-only content, bookkeeping, explicit domain shift, or independent same-domain cycle-fiber endpoint discrimination. This packet is therefore not decorative; it is a repeatable closure capsule for the theorem ladder and denial-burden table.


\subsection{Stress test: No-independent closure (endpoint-force lens)}

\par\noindent\textbf{Load-bearing claim.} This stress test audits the A+ consequence excluding the independent discriminator. The claim is not allowed to function by implication, atmosphere, or familiarity with AASC vocabulary. It must be available as a same-carrier theorem target inside the Hodge endpoint proof and must be reused only after the carrier, branch, and report slots have been instantiated.

\par\noindent\textbf{Hostile reconstruction.} A hostile referee will try to reconstruct this point in the strongest ordinary Hodge-theoretic language: rational Hodge classes are native cohomological objects, Chow groups and cycle-class maps are ordinary algebraic-geometric carriers, and an empty fiber would be ordinary mathematical falsification. The audit therefore never answers by declaring the underlying mathematical content illicit. It answers only by separating raw content, support status, theoremhood, endpoint force, and the fixed cycle-fiber endpoint image.

\par\noindent\textbf{Closure condition.} The claim survives this stress test only if it preserves the same $X$, $p$, and $\alpha$, the same rational Hodge-standing predicate, the same group $\CH^p(X)_\QQ$, the same map $\clQ$, and the same fiber $F_X(\alpha)$. A proposed escape that changes any of these is a domain or carrier shift. A proposed escape that preserves them while retaining endpoint-defeating force enters endpoint-use discipline and must pass through the native empty-fiber exhaustion.

\par\noindent\textbf{Audit result.} Under the fixed endpoint assumptions, the hostile reconstruction leaves only the declared branches: bridge completion, support-only content, bookkeeping, explicit domain shift, or independent same-domain cycle-fiber endpoint discrimination. This packet is therefore not decorative; it is a repeatable closure capsule for the theorem ladder and denial-burden table.


\subsection{Stress test: Local separator occupation (endpoint-force lens)}

\par\noindent\textbf{Load-bearing claim.} This stress test audits the local empty-fiber countercase as endpoint-defeating non-occupation. The claim is not allowed to function by implication, atmosphere, or familiarity with AASC vocabulary. It must be available as a same-carrier theorem target inside the Hodge endpoint proof and must be reused only after the carrier, branch, and report slots have been instantiated.

\par\noindent\textbf{Hostile reconstruction.} A hostile referee will try to reconstruct this point in the strongest ordinary Hodge-theoretic language: rational Hodge classes are native cohomological objects, Chow groups and cycle-class maps are ordinary algebraic-geometric carriers, and an empty fiber would be ordinary mathematical falsification. The audit therefore never answers by declaring the underlying mathematical content illicit. It answers only by separating raw content, support status, theoremhood, endpoint force, and the fixed cycle-fiber endpoint image.

\par\noindent\textbf{Closure condition.} The claim survives this stress test only if it preserves the same $X$, $p$, and $\alpha$, the same rational Hodge-standing predicate, the same group $\CH^p(X)_\QQ$, the same map $\clQ$, and the same fiber $F_X(\alpha)$. A proposed escape that changes any of these is a domain or carrier shift. A proposed escape that preserves them while retaining endpoint-defeating force enters endpoint-use discipline and must pass through the native empty-fiber exhaustion.

\par\noindent\textbf{Audit result.} Under the fixed endpoint assumptions, the hostile reconstruction leaves only the declared branches: bridge completion, support-only content, bookkeeping, explicit domain shift, or independent same-domain cycle-fiber endpoint discrimination. This packet is therefore not decorative; it is a repeatable closure capsule for the theorem ladder and denial-burden table.


\subsection{Stress test: Reductio annotation (endpoint-force lens)}

\par\noindent\textbf{Load-bearing claim.} This stress test audits proof-act annotation not object-level premise. The claim is not allowed to function by implication, atmosphere, or familiarity with AASC vocabulary. It must be available as a same-carrier theorem target inside the Hodge endpoint proof and must be reused only after the carrier, branch, and report slots have been instantiated.

\par\noindent\textbf{Hostile reconstruction.} A hostile referee will try to reconstruct this point in the strongest ordinary Hodge-theoretic language: rational Hodge classes are native cohomological objects, Chow groups and cycle-class maps are ordinary algebraic-geometric carriers, and an empty fiber would be ordinary mathematical falsification. The audit therefore never answers by declaring the underlying mathematical content illicit. It answers only by separating raw content, support status, theoremhood, endpoint force, and the fixed cycle-fiber endpoint image.

\par\noindent\textbf{Closure condition.} The claim survives this stress test only if it preserves the same $X$, $p$, and $\alpha$, the same rational Hodge-standing predicate, the same group $\CH^p(X)_\QQ$, the same map $\clQ$, and the same fiber $F_X(\alpha)$. A proposed escape that changes any of these is a domain or carrier shift. A proposed escape that preserves them while retaining endpoint-defeating force enters endpoint-use discipline and must pass through the native empty-fiber exhaustion.

\par\noindent\textbf{Audit result.} Under the fixed endpoint assumptions, the hostile reconstruction leaves only the declared branches: bridge completion, support-only content, bookkeeping, explicit domain shift, or independent same-domain cycle-fiber endpoint discrimination. This packet is therefore not decorative; it is a repeatable closure capsule for the theorem ladder and denial-burden table.


\subsection{Stress test: Annotation discharge (endpoint-force lens)}

\par\noindent\textbf{Load-bearing claim.} This stress test audits the original empty-fiber assumption is discharged. The claim is not allowed to function by implication, atmosphere, or familiarity with AASC vocabulary. It must be available as a same-carrier theorem target inside the Hodge endpoint proof and must be reused only after the carrier, branch, and report slots have been instantiated.

\par\noindent\textbf{Hostile reconstruction.} A hostile referee will try to reconstruct this point in the strongest ordinary Hodge-theoretic language: rational Hodge classes are native cohomological objects, Chow groups and cycle-class maps are ordinary algebraic-geometric carriers, and an empty fiber would be ordinary mathematical falsification. The audit therefore never answers by declaring the underlying mathematical content illicit. It answers only by separating raw content, support status, theoremhood, endpoint force, and the fixed cycle-fiber endpoint image.

\par\noindent\textbf{Closure condition.} The claim survives this stress test only if it preserves the same $X$, $p$, and $\alpha$, the same rational Hodge-standing predicate, the same group $\CH^p(X)_\QQ$, the same map $\clQ$, and the same fiber $F_X(\alpha)$. A proposed escape that changes any of these is a domain or carrier shift. A proposed escape that preserves them while retaining endpoint-defeating force enters endpoint-use discipline and must pass through the native empty-fiber exhaustion.

\par\noindent\textbf{Audit result.} Under the fixed endpoint assumptions, the hostile reconstruction leaves only the declared branches: bridge completion, support-only content, bookkeeping, explicit domain shift, or independent same-domain cycle-fiber endpoint discrimination. This packet is therefore not decorative; it is a repeatable closure capsule for the theorem ladder and denial-burden table.


\subsection{Stress test: Ordinary theoremhood protection (endpoint-force lens)}

\par\noindent\textbf{Load-bearing claim.} This stress test audits negative theoremhood is allowed until used as official endpoint resolution. The claim is not allowed to function by implication, atmosphere, or familiarity with AASC vocabulary. It must be available as a same-carrier theorem target inside the Hodge endpoint proof and must be reused only after the carrier, branch, and report slots have been instantiated.

\par\noindent\textbf{Hostile reconstruction.} A hostile referee will try to reconstruct this point in the strongest ordinary Hodge-theoretic language: rational Hodge classes are native cohomological objects, Chow groups and cycle-class maps are ordinary algebraic-geometric carriers, and an empty fiber would be ordinary mathematical falsification. The audit therefore never answers by declaring the underlying mathematical content illicit. It answers only by separating raw content, support status, theoremhood, endpoint force, and the fixed cycle-fiber endpoint image.

\par\noindent\textbf{Closure condition.} The claim survives this stress test only if it preserves the same $X$, $p$, and $\alpha$, the same rational Hodge-standing predicate, the same group $\CH^p(X)_\QQ$, the same map $\clQ$, and the same fiber $F_X(\alpha)$. A proposed escape that changes any of these is a domain or carrier shift. A proposed escape that preserves them while retaining endpoint-defeating force enters endpoint-use discipline and must pass through the native empty-fiber exhaustion.

\par\noindent\textbf{Audit result.} Under the fixed endpoint assumptions, the hostile reconstruction leaves only the declared branches: bridge completion, support-only content, bookkeeping, explicit domain shift, or independent same-domain cycle-fiber endpoint discrimination. This packet is therefore not decorative; it is a repeatable closure capsule for the theorem ladder and denial-burden table.


\subsection{Stress test: Incompleteness firewall (endpoint-force lens)}

\par\noindent\textbf{Load-bearing claim.} This stress test audits same-domain independence matters only if endpoint-standing relevant. The claim is not allowed to function by implication, atmosphere, or familiarity with AASC vocabulary. It must be available as a same-carrier theorem target inside the Hodge endpoint proof and must be reused only after the carrier, branch, and report slots have been instantiated.

\par\noindent\textbf{Hostile reconstruction.} A hostile referee will try to reconstruct this point in the strongest ordinary Hodge-theoretic language: rational Hodge classes are native cohomological objects, Chow groups and cycle-class maps are ordinary algebraic-geometric carriers, and an empty fiber would be ordinary mathematical falsification. The audit therefore never answers by declaring the underlying mathematical content illicit. It answers only by separating raw content, support status, theoremhood, endpoint force, and the fixed cycle-fiber endpoint image.

\par\noindent\textbf{Closure condition.} The claim survives this stress test only if it preserves the same $X$, $p$, and $\alpha$, the same rational Hodge-standing predicate, the same group $\CH^p(X)_\QQ$, the same map $\clQ$, and the same fiber $F_X(\alpha)$. A proposed escape that changes any of these is a domain or carrier shift. A proposed escape that preserves them while retaining endpoint-defeating force enters endpoint-use discipline and must pass through the native empty-fiber exhaustion.

\par\noindent\textbf{Audit result.} Under the fixed endpoint assumptions, the hostile reconstruction leaves only the declared branches: bridge completion, support-only content, bookkeeping, explicit domain shift, or independent same-domain cycle-fiber endpoint discrimination. This packet is therefore not decorative; it is a repeatable closure capsule for the theorem ladder and denial-burden table.


\subsection{Stress test: Route-locus exhaustion (endpoint-force lens)}

\par\noindent\textbf{Load-bearing claim.} This stress test audits status, codomain, existence, carrier, temporal, locality, selector, and second-classifier loci. The claim is not allowed to function by implication, atmosphere, or familiarity with AASC vocabulary. It must be available as a same-carrier theorem target inside the Hodge endpoint proof and must be reused only after the carrier, branch, and report slots have been instantiated.

\par\noindent\textbf{Hostile reconstruction.} A hostile referee will try to reconstruct this point in the strongest ordinary Hodge-theoretic language: rational Hodge classes are native cohomological objects, Chow groups and cycle-class maps are ordinary algebraic-geometric carriers, and an empty fiber would be ordinary mathematical falsification. The audit therefore never answers by declaring the underlying mathematical content illicit. It answers only by separating raw content, support status, theoremhood, endpoint force, and the fixed cycle-fiber endpoint image.

\par\noindent\textbf{Closure condition.} The claim survives this stress test only if it preserves the same $X$, $p$, and $\alpha$, the same rational Hodge-standing predicate, the same group $\CH^p(X)_\QQ$, the same map $\clQ$, and the same fiber $F_X(\alpha)$. A proposed escape that changes any of these is a domain or carrier shift. A proposed escape that preserves them while retaining endpoint-defeating force enters endpoint-use discipline and must pass through the native empty-fiber exhaustion.

\par\noindent\textbf{Audit result.} Under the fixed endpoint assumptions, the hostile reconstruction leaves only the declared branches: bridge completion, support-only content, bookkeeping, explicit domain shift, or independent same-domain cycle-fiber endpoint discrimination. This packet is therefore not decorative; it is a repeatable closure capsule for the theorem ladder and denial-burden table.


\subsection{Stress test: Weakening resistance (endpoint-force lens)}

\par\noindent\textbf{Load-bearing claim.} This stress test audits weaker regimes preserve behavior, shift domain, lower status, or reintroduce forbidden governance. The claim is not allowed to function by implication, atmosphere, or familiarity with AASC vocabulary. It must be available as a same-carrier theorem target inside the Hodge endpoint proof and must be reused only after the carrier, branch, and report slots have been instantiated.

\par\noindent\textbf{Hostile reconstruction.} A hostile referee will try to reconstruct this point in the strongest ordinary Hodge-theoretic language: rational Hodge classes are native cohomological objects, Chow groups and cycle-class maps are ordinary algebraic-geometric carriers, and an empty fiber would be ordinary mathematical falsification. The audit therefore never answers by declaring the underlying mathematical content illicit. It answers only by separating raw content, support status, theoremhood, endpoint force, and the fixed cycle-fiber endpoint image.

\par\noindent\textbf{Closure condition.} The claim survives this stress test only if it preserves the same $X$, $p$, and $\alpha$, the same rational Hodge-standing predicate, the same group $\CH^p(X)_\QQ$, the same map $\clQ$, and the same fiber $F_X(\alpha)$. A proposed escape that changes any of these is a domain or carrier shift. A proposed escape that preserves them while retaining endpoint-defeating force enters endpoint-use discipline and must pass through the native empty-fiber exhaustion.

\par\noindent\textbf{Audit result.} Under the fixed endpoint assumptions, the hostile reconstruction leaves only the declared branches: bridge completion, support-only content, bookkeeping, explicit domain shift, or independent same-domain cycle-fiber endpoint discrimination. This packet is therefore not decorative; it is a repeatable closure capsule for the theorem ladder and denial-burden table.


\subsection{Stress test: Official endpoint identity (route-exhaustion lens)}

\par\noindent\textbf{Load-bearing claim.} This stress test audits the identification of the Clay-facing endpoint with universal nonempty rational cycle-class fibers. The claim is not allowed to function by implication, atmosphere, or familiarity with AASC vocabulary. It must be available as a same-carrier theorem target inside the Hodge endpoint proof and must be reused only after the carrier, branch, and report slots have been instantiated.

\par\noindent\textbf{Hostile reconstruction.} A hostile referee will try to reconstruct this point in the strongest ordinary Hodge-theoretic language: rational Hodge classes are native cohomological objects, Chow groups and cycle-class maps are ordinary algebraic-geometric carriers, and an empty fiber would be ordinary mathematical falsification. The audit therefore never answers by declaring the underlying mathematical content illicit. It answers only by separating raw content, support status, theoremhood, endpoint force, and the fixed cycle-fiber endpoint image.

\par\noindent\textbf{Closure condition.} The claim survives this stress test only if it preserves the same $X$, $p$, and $\alpha$, the same rational Hodge-standing predicate, the same group $\CH^p(X)_\QQ$, the same map $\clQ$, and the same fiber $F_X(\alpha)$. A proposed escape that changes any of these is a domain or carrier shift. A proposed escape that preserves them while retaining endpoint-defeating force enters endpoint-use discipline and must pass through the native empty-fiber exhaustion.

\par\noindent\textbf{Audit result.} Under the fixed endpoint assumptions, the hostile reconstruction leaves only the declared branches: bridge completion, support-only content, bookkeeping, explicit domain shift, or independent same-domain cycle-fiber endpoint discrimination. This packet is therefore not decorative; it is a repeatable closure capsule for the theorem ladder and denial-burden table.


\subsection{Stress test: Context non-smuggling (route-exhaustion lens)}

\par\noindent\textbf{Load-bearing claim.} This stress test audits the claim that the Hodge endpoint context fixes the carrier and slots without asserting nonempty fiber. The claim is not allowed to function by implication, atmosphere, or familiarity with AASC vocabulary. It must be available as a same-carrier theorem target inside the Hodge endpoint proof and must be reused only after the carrier, branch, and report slots have been instantiated.

\par\noindent\textbf{Hostile reconstruction.} A hostile referee will try to reconstruct this point in the strongest ordinary Hodge-theoretic language: rational Hodge classes are native cohomological objects, Chow groups and cycle-class maps are ordinary algebraic-geometric carriers, and an empty fiber would be ordinary mathematical falsification. The audit therefore never answers by declaring the underlying mathematical content illicit. It answers only by separating raw content, support status, theoremhood, endpoint force, and the fixed cycle-fiber endpoint image.

\par\noindent\textbf{Closure condition.} The claim survives this stress test only if it preserves the same $X$, $p$, and $\alpha$, the same rational Hodge-standing predicate, the same group $\CH^p(X)_\QQ$, the same map $\clQ$, and the same fiber $F_X(\alpha)$. A proposed escape that changes any of these is a domain or carrier shift. A proposed escape that preserves them while retaining endpoint-defeating force enters endpoint-use discipline and must pass through the native empty-fiber exhaustion.

\par\noindent\textbf{Audit result.} Under the fixed endpoint assumptions, the hostile reconstruction leaves only the declared branches: bridge completion, support-only content, bookkeeping, explicit domain shift, or independent same-domain cycle-fiber endpoint discrimination. This packet is therefore not decorative; it is a repeatable closure capsule for the theorem ladder and denial-burden table.


\subsection{Stress test: Hodge-standing exactness (route-exhaustion lens)}

\par\noindent\textbf{Load-bearing claim.} This stress test audits the representation of rational Hodge standing as cohomological type status. The claim is not allowed to function by implication, atmosphere, or familiarity with AASC vocabulary. It must be available as a same-carrier theorem target inside the Hodge endpoint proof and must be reused only after the carrier, branch, and report slots have been instantiated.

\par\noindent\textbf{Hostile reconstruction.} A hostile referee will try to reconstruct this point in the strongest ordinary Hodge-theoretic language: rational Hodge classes are native cohomological objects, Chow groups and cycle-class maps are ordinary algebraic-geometric carriers, and an empty fiber would be ordinary mathematical falsification. The audit therefore never answers by declaring the underlying mathematical content illicit. It answers only by separating raw content, support status, theoremhood, endpoint force, and the fixed cycle-fiber endpoint image.

\par\noindent\textbf{Closure condition.} The claim survives this stress test only if it preserves the same $X$, $p$, and $\alpha$, the same rational Hodge-standing predicate, the same group $\CH^p(X)_\QQ$, the same map $\clQ$, and the same fiber $F_X(\alpha)$. A proposed escape that changes any of these is a domain or carrier shift. A proposed escape that preserves them while retaining endpoint-defeating force enters endpoint-use discipline and must pass through the native empty-fiber exhaustion.

\par\noindent\textbf{Audit result.} Under the fixed endpoint assumptions, the hostile reconstruction leaves only the declared branches: bridge completion, support-only content, bookkeeping, explicit domain shift, or independent same-domain cycle-fiber endpoint discrimination. This packet is therefore not decorative; it is a repeatable closure capsule for the theorem ladder and denial-burden table.


\subsection{Stress test: Cycle-fiber exactness (route-exhaustion lens)}

\par\noindent\textbf{Load-bearing claim.} This stress test audits the representation of algebraic readout by the fiber $F_X(\alpha)$. The claim is not allowed to function by implication, atmosphere, or familiarity with AASC vocabulary. It must be available as a same-carrier theorem target inside the Hodge endpoint proof and must be reused only after the carrier, branch, and report slots have been instantiated.

\par\noindent\textbf{Hostile reconstruction.} A hostile referee will try to reconstruct this point in the strongest ordinary Hodge-theoretic language: rational Hodge classes are native cohomological objects, Chow groups and cycle-class maps are ordinary algebraic-geometric carriers, and an empty fiber would be ordinary mathematical falsification. The audit therefore never answers by declaring the underlying mathematical content illicit. It answers only by separating raw content, support status, theoremhood, endpoint force, and the fixed cycle-fiber endpoint image.

\par\noindent\textbf{Closure condition.} The claim survives this stress test only if it preserves the same $X$, $p$, and $\alpha$, the same rational Hodge-standing predicate, the same group $\CH^p(X)_\QQ$, the same map $\clQ$, and the same fiber $F_X(\alpha)$. A proposed escape that changes any of these is a domain or carrier shift. A proposed escape that preserves them while retaining endpoint-defeating force enters endpoint-use discipline and must pass through the native empty-fiber exhaustion.

\par\noindent\textbf{Audit result.} Under the fixed endpoint assumptions, the hostile reconstruction leaves only the declared branches: bridge completion, support-only content, bookkeeping, explicit domain shift, or independent same-domain cycle-fiber endpoint discrimination. This packet is therefore not decorative; it is a repeatable closure capsule for the theorem ladder and denial-burden table.


\subsection{Stress test: Empty-fiber exactness (route-exhaustion lens)}

\par\noindent\textbf{Load-bearing claim.} This stress test audits the negative local branch as fiber non-occupation. The claim is not allowed to function by implication, atmosphere, or familiarity with AASC vocabulary. It must be available as a same-carrier theorem target inside the Hodge endpoint proof and must be reused only after the carrier, branch, and report slots have been instantiated.

\par\noindent\textbf{Hostile reconstruction.} A hostile referee will try to reconstruct this point in the strongest ordinary Hodge-theoretic language: rational Hodge classes are native cohomological objects, Chow groups and cycle-class maps are ordinary algebraic-geometric carriers, and an empty fiber would be ordinary mathematical falsification. The audit therefore never answers by declaring the underlying mathematical content illicit. It answers only by separating raw content, support status, theoremhood, endpoint force, and the fixed cycle-fiber endpoint image.

\par\noindent\textbf{Closure condition.} The claim survives this stress test only if it preserves the same $X$, $p$, and $\alpha$, the same rational Hodge-standing predicate, the same group $\CH^p(X)_\QQ$, the same map $\clQ$, and the same fiber $F_X(\alpha)$. A proposed escape that changes any of these is a domain or carrier shift. A proposed escape that preserves them while retaining endpoint-defeating force enters endpoint-use discipline and must pass through the native empty-fiber exhaustion.

\par\noindent\textbf{Audit result.} Under the fixed endpoint assumptions, the hostile reconstruction leaves only the declared branches: bridge completion, support-only content, bookkeeping, explicit domain shift, or independent same-domain cycle-fiber endpoint discrimination. This packet is therefore not decorative; it is a repeatable closure capsule for the theorem ladder and denial-burden table.


\subsection{Stress test: Counterexample object/use distinction (route-exhaustion lens)}

\par\noindent\textbf{Load-bearing claim.} This stress test audits the separation between an object satisfying the negative branch and the act of offering it as endpoint defeat. The claim is not allowed to function by implication, atmosphere, or familiarity with AASC vocabulary. It must be available as a same-carrier theorem target inside the Hodge endpoint proof and must be reused only after the carrier, branch, and report slots have been instantiated.

\par\noindent\textbf{Hostile reconstruction.} A hostile referee will try to reconstruct this point in the strongest ordinary Hodge-theoretic language: rational Hodge classes are native cohomological objects, Chow groups and cycle-class maps are ordinary algebraic-geometric carriers, and an empty fiber would be ordinary mathematical falsification. The audit therefore never answers by declaring the underlying mathematical content illicit. It answers only by separating raw content, support status, theoremhood, endpoint force, and the fixed cycle-fiber endpoint image.

\par\noindent\textbf{Closure condition.} The claim survives this stress test only if it preserves the same $X$, $p$, and $\alpha$, the same rational Hodge-standing predicate, the same group $\CH^p(X)_\QQ$, the same map $\clQ$, and the same fiber $F_X(\alpha)$. A proposed escape that changes any of these is a domain or carrier shift. A proposed escape that preserves them while retaining endpoint-defeating force enters endpoint-use discipline and must pass through the native empty-fiber exhaustion.

\par\noindent\textbf{Audit result.} Under the fixed endpoint assumptions, the hostile reconstruction leaves only the declared branches: bridge completion, support-only content, bookkeeping, explicit domain shift, or independent same-domain cycle-fiber endpoint discrimination. This packet is therefore not decorative; it is a repeatable closure capsule for the theorem ladder and denial-burden table.


\subsection{Stress test: Official counterexample force (route-exhaustion lens)}

\par\noindent\textbf{Load-bearing claim.} This stress test audits the nonclassical dependency identifying endpoint-defeating counterexample force with theorem-bearing endpoint use. The claim is not allowed to function by implication, atmosphere, or familiarity with AASC vocabulary. It must be available as a same-carrier theorem target inside the Hodge endpoint proof and must be reused only after the carrier, branch, and report slots have been instantiated.

\par\noindent\textbf{Hostile reconstruction.} A hostile referee will try to reconstruct this point in the strongest ordinary Hodge-theoretic language: rational Hodge classes are native cohomological objects, Chow groups and cycle-class maps are ordinary algebraic-geometric carriers, and an empty fiber would be ordinary mathematical falsification. The audit therefore never answers by declaring the underlying mathematical content illicit. It answers only by separating raw content, support status, theoremhood, endpoint force, and the fixed cycle-fiber endpoint image.

\par\noindent\textbf{Closure condition.} The claim survives this stress test only if it preserves the same $X$, $p$, and $\alpha$, the same rational Hodge-standing predicate, the same group $\CH^p(X)_\QQ$, the same map $\clQ$, and the same fiber $F_X(\alpha)$. A proposed escape that changes any of these is a domain or carrier shift. A proposed escape that preserves them while retaining endpoint-defeating force enters endpoint-use discipline and must pass through the native empty-fiber exhaustion.

\par\noindent\textbf{Audit result.} Under the fixed endpoint assumptions, the hostile reconstruction leaves only the declared branches: bridge completion, support-only content, bookkeeping, explicit domain shift, or independent same-domain cycle-fiber endpoint discrimination. This packet is therefore not decorative; it is a repeatable closure capsule for the theorem ladder and denial-burden table.


\subsection{Stress test: No autonomous counterexample endpoint role (route-exhaustion lens)}

\par\noindent\textbf{Load-bearing claim.} This stress test audits the denial of a fifth role outside endpoint use for official endpoint defeat. The claim is not allowed to function by implication, atmosphere, or familiarity with AASC vocabulary. It must be available as a same-carrier theorem target inside the Hodge endpoint proof and must be reused only after the carrier, branch, and report slots have been instantiated.

\par\noindent\textbf{Hostile reconstruction.} A hostile referee will try to reconstruct this point in the strongest ordinary Hodge-theoretic language: rational Hodge classes are native cohomological objects, Chow groups and cycle-class maps are ordinary algebraic-geometric carriers, and an empty fiber would be ordinary mathematical falsification. The audit therefore never answers by declaring the underlying mathematical content illicit. It answers only by separating raw content, support status, theoremhood, endpoint force, and the fixed cycle-fiber endpoint image.

\par\noindent\textbf{Closure condition.} The claim survives this stress test only if it preserves the same $X$, $p$, and $\alpha$, the same rational Hodge-standing predicate, the same group $\CH^p(X)_\QQ$, the same map $\clQ$, and the same fiber $F_X(\alpha)$. A proposed escape that changes any of these is a domain or carrier shift. A proposed escape that preserves them while retaining endpoint-defeating force enters endpoint-use discipline and must pass through the native empty-fiber exhaustion.

\par\noindent\textbf{Audit result.} Under the fixed endpoint assumptions, the hostile reconstruction leaves only the declared branches: bridge completion, support-only content, bookkeeping, explicit domain shift, or independent same-domain cycle-fiber endpoint discrimination. This packet is therefore not decorative; it is a repeatable closure capsule for the theorem ladder and denial-burden table.


\subsection{Stress test: Kernel forcing (route-exhaustion lens)}

\par\noindent\textbf{Load-bearing claim.} This stress test audits the derivation of reference, standing, admissibility, and irreversibility from target adequacy. The claim is not allowed to function by implication, atmosphere, or familiarity with AASC vocabulary. It must be available as a same-carrier theorem target inside the Hodge endpoint proof and must be reused only after the carrier, branch, and report slots have been instantiated.

\par\noindent\textbf{Hostile reconstruction.} A hostile referee will try to reconstruct this point in the strongest ordinary Hodge-theoretic language: rational Hodge classes are native cohomological objects, Chow groups and cycle-class maps are ordinary algebraic-geometric carriers, and an empty fiber would be ordinary mathematical falsification. The audit therefore never answers by declaring the underlying mathematical content illicit. It answers only by separating raw content, support status, theoremhood, endpoint force, and the fixed cycle-fiber endpoint image.

\par\noindent\textbf{Closure condition.} The claim survives this stress test only if it preserves the same $X$, $p$, and $\alpha$, the same rational Hodge-standing predicate, the same group $\CH^p(X)_\QQ$, the same map $\clQ$, and the same fiber $F_X(\alpha)$. A proposed escape that changes any of these is a domain or carrier shift. A proposed escape that preserves them while retaining endpoint-defeating force enters endpoint-use discipline and must pass through the native empty-fiber exhaustion.

\par\noindent\textbf{Audit result.} Under the fixed endpoint assumptions, the hostile reconstruction leaves only the declared branches: bridge completion, support-only content, bookkeeping, explicit domain shift, or independent same-domain cycle-fiber endpoint discrimination. This packet is therefore not decorative; it is a repeatable closure capsule for the theorem ladder and denial-burden table.


\subsection{Stress test: A+ consequence layer (route-exhaustion lens)}

\par\noindent\textbf{Load-bearing claim.} This stress test audits the fixed-domain consequence package binding endpoint use. The claim is not allowed to function by implication, atmosphere, or familiarity with AASC vocabulary. It must be available as a same-carrier theorem target inside the Hodge endpoint proof and must be reused only after the carrier, branch, and report slots have been instantiated.

\par\noindent\textbf{Hostile reconstruction.} A hostile referee will try to reconstruct this point in the strongest ordinary Hodge-theoretic language: rational Hodge classes are native cohomological objects, Chow groups and cycle-class maps are ordinary algebraic-geometric carriers, and an empty fiber would be ordinary mathematical falsification. The audit therefore never answers by declaring the underlying mathematical content illicit. It answers only by separating raw content, support status, theoremhood, endpoint force, and the fixed cycle-fiber endpoint image.

\par\noindent\textbf{Closure condition.} The claim survives this stress test only if it preserves the same $X$, $p$, and $\alpha$, the same rational Hodge-standing predicate, the same group $\CH^p(X)_\QQ$, the same map $\clQ$, and the same fiber $F_X(\alpha)$. A proposed escape that changes any of these is a domain or carrier shift. A proposed escape that preserves them while retaining endpoint-defeating force enters endpoint-use discipline and must pass through the native empty-fiber exhaustion.

\par\noindent\textbf{Audit result.} Under the fixed endpoint assumptions, the hostile reconstruction leaves only the declared branches: bridge completion, support-only content, bookkeeping, explicit domain shift, or independent same-domain cycle-fiber endpoint discrimination. This packet is therefore not decorative; it is a repeatable closure capsule for the theorem ladder and denial-burden table.


\subsection{Stress test: ATS locking (route-exhaustion lens)}

\par\noindent\textbf{Load-bearing claim.} This stress test audits same target, carrier, branch, readout trajectory, and no second classifier. The claim is not allowed to function by implication, atmosphere, or familiarity with AASC vocabulary. It must be available as a same-carrier theorem target inside the Hodge endpoint proof and must be reused only after the carrier, branch, and report slots have been instantiated.

\par\noindent\textbf{Hostile reconstruction.} A hostile referee will try to reconstruct this point in the strongest ordinary Hodge-theoretic language: rational Hodge classes are native cohomological objects, Chow groups and cycle-class maps are ordinary algebraic-geometric carriers, and an empty fiber would be ordinary mathematical falsification. The audit therefore never answers by declaring the underlying mathematical content illicit. It answers only by separating raw content, support status, theoremhood, endpoint force, and the fixed cycle-fiber endpoint image.

\par\noindent\textbf{Closure condition.} The claim survives this stress test only if it preserves the same $X$, $p$, and $\alpha$, the same rational Hodge-standing predicate, the same group $\CH^p(X)_\QQ$, the same map $\clQ$, and the same fiber $F_X(\alpha)$. A proposed escape that changes any of these is a domain or carrier shift. A proposed escape that preserves them while retaining endpoint-defeating force enters endpoint-use discipline and must pass through the native empty-fiber exhaustion.

\par\noindent\textbf{Audit result.} Under the fixed endpoint assumptions, the hostile reconstruction leaves only the declared branches: bridge completion, support-only content, bookkeeping, explicit domain shift, or independent same-domain cycle-fiber endpoint discrimination. This packet is therefore not decorative; it is a repeatable closure capsule for the theorem ladder and denial-burden table.


\subsection{Stress test: UEAP report discipline (route-exhaustion lens)}

\par\noindent\textbf{Load-bearing claim.} This stress test audits no report exceeding carrier, bridge, and support status. The claim is not allowed to function by implication, atmosphere, or familiarity with AASC vocabulary. It must be available as a same-carrier theorem target inside the Hodge endpoint proof and must be reused only after the carrier, branch, and report slots have been instantiated.

\par\noindent\textbf{Hostile reconstruction.} A hostile referee will try to reconstruct this point in the strongest ordinary Hodge-theoretic language: rational Hodge classes are native cohomological objects, Chow groups and cycle-class maps are ordinary algebraic-geometric carriers, and an empty fiber would be ordinary mathematical falsification. The audit therefore never answers by declaring the underlying mathematical content illicit. It answers only by separating raw content, support status, theoremhood, endpoint force, and the fixed cycle-fiber endpoint image.

\par\noindent\textbf{Closure condition.} The claim survives this stress test only if it preserves the same $X$, $p$, and $\alpha$, the same rational Hodge-standing predicate, the same group $\CH^p(X)_\QQ$, the same map $\clQ$, and the same fiber $F_X(\alpha)$. A proposed escape that changes any of these is a domain or carrier shift. A proposed escape that preserves them while retaining endpoint-defeating force enters endpoint-use discipline and must pass through the native empty-fiber exhaustion.

\par\noindent\textbf{Audit result.} Under the fixed endpoint assumptions, the hostile reconstruction leaves only the declared branches: bridge completion, support-only content, bookkeeping, explicit domain shift, or independent same-domain cycle-fiber endpoint discrimination. This packet is therefore not decorative; it is a repeatable closure capsule for the theorem ladder and denial-burden table.


\subsection{Stress test: AMetric no-selector boundary (route-exhaustion lens)}

\par\noindent\textbf{Load-bearing claim.} This stress test audits no preferred representative or scalar ranking as pre-bridge authority. The claim is not allowed to function by implication, atmosphere, or familiarity with AASC vocabulary. It must be available as a same-carrier theorem target inside the Hodge endpoint proof and must be reused only after the carrier, branch, and report slots have been instantiated.

\par\noindent\textbf{Hostile reconstruction.} A hostile referee will try to reconstruct this point in the strongest ordinary Hodge-theoretic language: rational Hodge classes are native cohomological objects, Chow groups and cycle-class maps are ordinary algebraic-geometric carriers, and an empty fiber would be ordinary mathematical falsification. The audit therefore never answers by declaring the underlying mathematical content illicit. It answers only by separating raw content, support status, theoremhood, endpoint force, and the fixed cycle-fiber endpoint image.

\par\noindent\textbf{Closure condition.} The claim survives this stress test only if it preserves the same $X$, $p$, and $\alpha$, the same rational Hodge-standing predicate, the same group $\CH^p(X)_\QQ$, the same map $\clQ$, and the same fiber $F_X(\alpha)$. A proposed escape that changes any of these is a domain or carrier shift. A proposed escape that preserves them while retaining endpoint-defeating force enters endpoint-use discipline and must pass through the native empty-fiber exhaustion.

\par\noindent\textbf{Audit result.} Under the fixed endpoint assumptions, the hostile reconstruction leaves only the declared branches: bridge completion, support-only content, bookkeeping, explicit domain shift, or independent same-domain cycle-fiber endpoint discrimination. This packet is therefore not decorative; it is a repeatable closure capsule for the theorem ladder and denial-burden table.


\subsection{Stress test: Surrogate carrier separation (route-exhaustion lens)}

\par\noindent\textbf{Load-bearing claim.} This stress test audits motivic, variational, or obstruction objects as bridge support rather than readout. The claim is not allowed to function by implication, atmosphere, or familiarity with AASC vocabulary. It must be available as a same-carrier theorem target inside the Hodge endpoint proof and must be reused only after the carrier, branch, and report slots have been instantiated.

\par\noindent\textbf{Hostile reconstruction.} A hostile referee will try to reconstruct this point in the strongest ordinary Hodge-theoretic language: rational Hodge classes are native cohomological objects, Chow groups and cycle-class maps are ordinary algebraic-geometric carriers, and an empty fiber would be ordinary mathematical falsification. The audit therefore never answers by declaring the underlying mathematical content illicit. It answers only by separating raw content, support status, theoremhood, endpoint force, and the fixed cycle-fiber endpoint image.

\par\noindent\textbf{Closure condition.} The claim survives this stress test only if it preserves the same $X$, $p$, and $\alpha$, the same rational Hodge-standing predicate, the same group $\CH^p(X)_\QQ$, the same map $\clQ$, and the same fiber $F_X(\alpha)$. A proposed escape that changes any of these is a domain or carrier shift. A proposed escape that preserves them while retaining endpoint-defeating force enters endpoint-use discipline and must pass through the native empty-fiber exhaustion.

\par\noindent\textbf{Audit result.} Under the fixed endpoint assumptions, the hostile reconstruction leaves only the declared branches: bridge completion, support-only content, bookkeeping, explicit domain shift, or independent same-domain cycle-fiber endpoint discrimination. This packet is therefore not decorative; it is a repeatable closure capsule for the theorem ladder and denial-burden table.


\subsection{Stress test: No post-selection repair (route-exhaustion lens)}

\par\noindent\textbf{Load-bearing claim.} This stress test audits later cycle discovery as fresh bridge completion rather than same-act repair. The claim is not allowed to function by implication, atmosphere, or familiarity with AASC vocabulary. It must be available as a same-carrier theorem target inside the Hodge endpoint proof and must be reused only after the carrier, branch, and report slots have been instantiated.

\par\noindent\textbf{Hostile reconstruction.} A hostile referee will try to reconstruct this point in the strongest ordinary Hodge-theoretic language: rational Hodge classes are native cohomological objects, Chow groups and cycle-class maps are ordinary algebraic-geometric carriers, and an empty fiber would be ordinary mathematical falsification. The audit therefore never answers by declaring the underlying mathematical content illicit. It answers only by separating raw content, support status, theoremhood, endpoint force, and the fixed cycle-fiber endpoint image.

\par\noindent\textbf{Closure condition.} The claim survives this stress test only if it preserves the same $X$, $p$, and $\alpha$, the same rational Hodge-standing predicate, the same group $\CH^p(X)_\QQ$, the same map $\clQ$, and the same fiber $F_X(\alpha)$. A proposed escape that changes any of these is a domain or carrier shift. A proposed escape that preserves them while retaining endpoint-defeating force enters endpoint-use discipline and must pass through the native empty-fiber exhaustion.

\par\noindent\textbf{Audit result.} Under the fixed endpoint assumptions, the hostile reconstruction leaves only the declared branches: bridge completion, support-only content, bookkeeping, explicit domain shift, or independent same-domain cycle-fiber endpoint discrimination. This packet is therefore not decorative; it is a repeatable closure capsule for the theorem ladder and denial-burden table.


\subsection{Stress test: No carrier transfer (route-exhaustion lens)}

\par\noindent\textbf{Load-bearing claim.} This stress test audits known cases and neighboring domains not carrying endpoint standing without target inclusion. The claim is not allowed to function by implication, atmosphere, or familiarity with AASC vocabulary. It must be available as a same-carrier theorem target inside the Hodge endpoint proof and must be reused only after the carrier, branch, and report slots have been instantiated.

\par\noindent\textbf{Hostile reconstruction.} A hostile referee will try to reconstruct this point in the strongest ordinary Hodge-theoretic language: rational Hodge classes are native cohomological objects, Chow groups and cycle-class maps are ordinary algebraic-geometric carriers, and an empty fiber would be ordinary mathematical falsification. The audit therefore never answers by declaring the underlying mathematical content illicit. It answers only by separating raw content, support status, theoremhood, endpoint force, and the fixed cycle-fiber endpoint image.

\par\noindent\textbf{Closure condition.} The claim survives this stress test only if it preserves the same $X$, $p$, and $\alpha$, the same rational Hodge-standing predicate, the same group $\CH^p(X)_\QQ$, the same map $\clQ$, and the same fiber $F_X(\alpha)$. A proposed escape that changes any of these is a domain or carrier shift. A proposed escape that preserves them while retaining endpoint-defeating force enters endpoint-use discipline and must pass through the native empty-fiber exhaustion.

\par\noindent\textbf{Audit result.} Under the fixed endpoint assumptions, the hostile reconstruction leaves only the declared branches: bridge completion, support-only content, bookkeeping, explicit domain shift, or independent same-domain cycle-fiber endpoint discrimination. This packet is therefore not decorative; it is a repeatable closure capsule for the theorem ladder and denial-burden table.


\subsection{Stress test: No generators (route-exhaustion lens)}

\par\noindent\textbf{Load-bearing claim.} This stress test audits Hodge standing not generating Chow-fiber nonemptiness. The claim is not allowed to function by implication, atmosphere, or familiarity with AASC vocabulary. It must be available as a same-carrier theorem target inside the Hodge endpoint proof and must be reused only after the carrier, branch, and report slots have been instantiated.

\par\noindent\textbf{Hostile reconstruction.} A hostile referee will try to reconstruct this point in the strongest ordinary Hodge-theoretic language: rational Hodge classes are native cohomological objects, Chow groups and cycle-class maps are ordinary algebraic-geometric carriers, and an empty fiber would be ordinary mathematical falsification. The audit therefore never answers by declaring the underlying mathematical content illicit. It answers only by separating raw content, support status, theoremhood, endpoint force, and the fixed cycle-fiber endpoint image.

\par\noindent\textbf{Closure condition.} The claim survives this stress test only if it preserves the same $X$, $p$, and $\alpha$, the same rational Hodge-standing predicate, the same group $\CH^p(X)_\QQ$, the same map $\clQ$, and the same fiber $F_X(\alpha)$. A proposed escape that changes any of these is a domain or carrier shift. A proposed escape that preserves them while retaining endpoint-defeating force enters endpoint-use discipline and must pass through the native empty-fiber exhaustion.

\par\noindent\textbf{Audit result.} Under the fixed endpoint assumptions, the hostile reconstruction leaves only the declared branches: bridge completion, support-only content, bookkeeping, explicit domain shift, or independent same-domain cycle-fiber endpoint discrimination. This packet is therefore not decorative; it is a repeatable closure capsule for the theorem ladder and denial-burden table.


\subsection{Stress test: No silent redescription (route-exhaustion lens)}

\par\noindent\textbf{Load-bearing claim.} This stress test audits no replacement of Chow-fiber readout by another carrier under same label. The claim is not allowed to function by implication, atmosphere, or familiarity with AASC vocabulary. It must be available as a same-carrier theorem target inside the Hodge endpoint proof and must be reused only after the carrier, branch, and report slots have been instantiated.

\par\noindent\textbf{Hostile reconstruction.} A hostile referee will try to reconstruct this point in the strongest ordinary Hodge-theoretic language: rational Hodge classes are native cohomological objects, Chow groups and cycle-class maps are ordinary algebraic-geometric carriers, and an empty fiber would be ordinary mathematical falsification. The audit therefore never answers by declaring the underlying mathematical content illicit. It answers only by separating raw content, support status, theoremhood, endpoint force, and the fixed cycle-fiber endpoint image.

\par\noindent\textbf{Closure condition.} The claim survives this stress test only if it preserves the same $X$, $p$, and $\alpha$, the same rational Hodge-standing predicate, the same group $\CH^p(X)_\QQ$, the same map $\clQ$, and the same fiber $F_X(\alpha)$. A proposed escape that changes any of these is a domain or carrier shift. A proposed escape that preserves them while retaining endpoint-defeating force enters endpoint-use discipline and must pass through the native empty-fiber exhaustion.

\par\noindent\textbf{Audit result.} Under the fixed endpoint assumptions, the hostile reconstruction leaves only the declared branches: bridge completion, support-only content, bookkeeping, explicit domain shift, or independent same-domain cycle-fiber endpoint discrimination. This packet is therefore not decorative; it is a repeatable closure capsule for the theorem ladder and denial-burden table.


\subsection{Stress test: Non-factorizability (route-exhaustion lens)}

\par\noindent\textbf{Load-bearing claim.} This stress test audits local bridge fragments not composing into global readout without compatibility. The claim is not allowed to function by implication, atmosphere, or familiarity with AASC vocabulary. It must be available as a same-carrier theorem target inside the Hodge endpoint proof and must be reused only after the carrier, branch, and report slots have been instantiated.

\par\noindent\textbf{Hostile reconstruction.} A hostile referee will try to reconstruct this point in the strongest ordinary Hodge-theoretic language: rational Hodge classes are native cohomological objects, Chow groups and cycle-class maps are ordinary algebraic-geometric carriers, and an empty fiber would be ordinary mathematical falsification. The audit therefore never answers by declaring the underlying mathematical content illicit. It answers only by separating raw content, support status, theoremhood, endpoint force, and the fixed cycle-fiber endpoint image.

\par\noindent\textbf{Closure condition.} The claim survives this stress test only if it preserves the same $X$, $p$, and $\alpha$, the same rational Hodge-standing predicate, the same group $\CH^p(X)_\QQ$, the same map $\clQ$, and the same fiber $F_X(\alpha)$. A proposed escape that changes any of these is a domain or carrier shift. A proposed escape that preserves them while retaining endpoint-defeating force enters endpoint-use discipline and must pass through the native empty-fiber exhaustion.

\par\noindent\textbf{Audit result.} Under the fixed endpoint assumptions, the hostile reconstruction leaves only the declared branches: bridge completion, support-only content, bookkeeping, explicit domain shift, or independent same-domain cycle-fiber endpoint discrimination. This packet is therefore not decorative; it is a repeatable closure capsule for the theorem ladder and denial-burden table.


\subsection{Stress test: Role-occupancy closure (route-exhaustion lens)}

\par\noindent\textbf{Load-bearing claim.} This stress test audits nonempty fiber as role occupation rather than unique representative selection. The claim is not allowed to function by implication, atmosphere, or familiarity with AASC vocabulary. It must be available as a same-carrier theorem target inside the Hodge endpoint proof and must be reused only after the carrier, branch, and report slots have been instantiated.

\par\noindent\textbf{Hostile reconstruction.} A hostile referee will try to reconstruct this point in the strongest ordinary Hodge-theoretic language: rational Hodge classes are native cohomological objects, Chow groups and cycle-class maps are ordinary algebraic-geometric carriers, and an empty fiber would be ordinary mathematical falsification. The audit therefore never answers by declaring the underlying mathematical content illicit. It answers only by separating raw content, support status, theoremhood, endpoint force, and the fixed cycle-fiber endpoint image.

\par\noindent\textbf{Closure condition.} The claim survives this stress test only if it preserves the same $X$, $p$, and $\alpha$, the same rational Hodge-standing predicate, the same group $\CH^p(X)_\QQ$, the same map $\clQ$, and the same fiber $F_X(\alpha)$. A proposed escape that changes any of these is a domain or carrier shift. A proposed escape that preserves them while retaining endpoint-defeating force enters endpoint-use discipline and must pass through the native empty-fiber exhaustion.

\par\noindent\textbf{Audit result.} Under the fixed endpoint assumptions, the hostile reconstruction leaves only the declared branches: bridge completion, support-only content, bookkeeping, explicit domain shift, or independent same-domain cycle-fiber endpoint discrimination. This packet is therefore not decorative; it is a repeatable closure capsule for the theorem ladder and denial-burden table.


\subsection{Stress test: Cycle-fiber single interface (route-exhaustion lens)}

\par\noindent\textbf{Load-bearing claim.} This stress test audits the single endpoint-image interface for algebraic readout. The claim is not allowed to function by implication, atmosphere, or familiarity with AASC vocabulary. It must be available as a same-carrier theorem target inside the Hodge endpoint proof and must be reused only after the carrier, branch, and report slots have been instantiated.

\par\noindent\textbf{Hostile reconstruction.} A hostile referee will try to reconstruct this point in the strongest ordinary Hodge-theoretic language: rational Hodge classes are native cohomological objects, Chow groups and cycle-class maps are ordinary algebraic-geometric carriers, and an empty fiber would be ordinary mathematical falsification. The audit therefore never answers by declaring the underlying mathematical content illicit. It answers only by separating raw content, support status, theoremhood, endpoint force, and the fixed cycle-fiber endpoint image.

\par\noindent\textbf{Closure condition.} The claim survives this stress test only if it preserves the same $X$, $p$, and $\alpha$, the same rational Hodge-standing predicate, the same group $\CH^p(X)_\QQ$, the same map $\clQ$, and the same fiber $F_X(\alpha)$. A proposed escape that changes any of these is a domain or carrier shift. A proposed escape that preserves them while retaining endpoint-defeating force enters endpoint-use discipline and must pass through the native empty-fiber exhaustion.

\par\noindent\textbf{Audit result.} Under the fixed endpoint assumptions, the hostile reconstruction leaves only the declared branches: bridge completion, support-only content, bookkeeping, explicit domain shift, or independent same-domain cycle-fiber endpoint discrimination. This packet is therefore not decorative; it is a repeatable closure capsule for the theorem ladder and denial-burden table.


\subsection{Stress test: Coequal carrier distinction (route-exhaustion lens)}

\par\noindent\textbf{Load-bearing claim.} This stress test audits separating broad outcome carrier from cycle-fiber image carrier. The claim is not allowed to function by implication, atmosphere, or familiarity with AASC vocabulary. It must be available as a same-carrier theorem target inside the Hodge endpoint proof and must be reused only after the carrier, branch, and report slots have been instantiated.

\par\noindent\textbf{Hostile reconstruction.} A hostile referee will try to reconstruct this point in the strongest ordinary Hodge-theoretic language: rational Hodge classes are native cohomological objects, Chow groups and cycle-class maps are ordinary algebraic-geometric carriers, and an empty fiber would be ordinary mathematical falsification. The audit therefore never answers by declaring the underlying mathematical content illicit. It answers only by separating raw content, support status, theoremhood, endpoint force, and the fixed cycle-fiber endpoint image.

\par\noindent\textbf{Closure condition.} The claim survives this stress test only if it preserves the same $X$, $p$, and $\alpha$, the same rational Hodge-standing predicate, the same group $\CH^p(X)_\QQ$, the same map $\clQ$, and the same fiber $F_X(\alpha)$. A proposed escape that changes any of these is a domain or carrier shift. A proposed escape that preserves them while retaining endpoint-defeating force enters endpoint-use discipline and must pass through the native empty-fiber exhaustion.

\par\noindent\textbf{Audit result.} Under the fixed endpoint assumptions, the hostile reconstruction leaves only the declared branches: bridge completion, support-only content, bookkeeping, explicit domain shift, or independent same-domain cycle-fiber endpoint discrimination. This packet is therefore not decorative; it is a repeatable closure capsule for the theorem ladder and denial-burden table.


\subsection{Stress test: Native empty-fiber exhaustion (route-exhaustion lens)}

\par\noindent\textbf{Load-bearing claim.} This stress test audits bridge-neutralized, support-only, bookkeeping, domain shift, or discriminator. The claim is not allowed to function by implication, atmosphere, or familiarity with AASC vocabulary. It must be available as a same-carrier theorem target inside the Hodge endpoint proof and must be reused only after the carrier, branch, and report slots have been instantiated.

\par\noindent\textbf{Hostile reconstruction.} A hostile referee will try to reconstruct this point in the strongest ordinary Hodge-theoretic language: rational Hodge classes are native cohomological objects, Chow groups and cycle-class maps are ordinary algebraic-geometric carriers, and an empty fiber would be ordinary mathematical falsification. The audit therefore never answers by declaring the underlying mathematical content illicit. It answers only by separating raw content, support status, theoremhood, endpoint force, and the fixed cycle-fiber endpoint image.

\par\noindent\textbf{Closure condition.} The claim survives this stress test only if it preserves the same $X$, $p$, and $\alpha$, the same rational Hodge-standing predicate, the same group $\CH^p(X)_\QQ$, the same map $\clQ$, and the same fiber $F_X(\alpha)$. A proposed escape that changes any of these is a domain or carrier shift. A proposed escape that preserves them while retaining endpoint-defeating force enters endpoint-use discipline and must pass through the native empty-fiber exhaustion.

\par\noindent\textbf{Audit result.} Under the fixed endpoint assumptions, the hostile reconstruction leaves only the declared branches: bridge completion, support-only content, bookkeeping, explicit domain shift, or independent same-domain cycle-fiber endpoint discrimination. This packet is therefore not decorative; it is a repeatable closure capsule for the theorem ladder and denial-burden table.


\subsection{Stress test: Bridge-neutralized branch (route-exhaustion lens)}

\par\noindent\textbf{Load-bearing claim.} This stress test audits why a bridge-neutralized negative branch is not live endpoint counterforce. The claim is not allowed to function by implication, atmosphere, or familiarity with AASC vocabulary. It must be available as a same-carrier theorem target inside the Hodge endpoint proof and must be reused only after the carrier, branch, and report slots have been instantiated.

\par\noindent\textbf{Hostile reconstruction.} A hostile referee will try to reconstruct this point in the strongest ordinary Hodge-theoretic language: rational Hodge classes are native cohomological objects, Chow groups and cycle-class maps are ordinary algebraic-geometric carriers, and an empty fiber would be ordinary mathematical falsification. The audit therefore never answers by declaring the underlying mathematical content illicit. It answers only by separating raw content, support status, theoremhood, endpoint force, and the fixed cycle-fiber endpoint image.

\par\noindent\textbf{Closure condition.} The claim survives this stress test only if it preserves the same $X$, $p$, and $\alpha$, the same rational Hodge-standing predicate, the same group $\CH^p(X)_\QQ$, the same map $\clQ$, and the same fiber $F_X(\alpha)$. A proposed escape that changes any of these is a domain or carrier shift. A proposed escape that preserves them while retaining endpoint-defeating force enters endpoint-use discipline and must pass through the native empty-fiber exhaustion.

\par\noindent\textbf{Audit result.} Under the fixed endpoint assumptions, the hostile reconstruction leaves only the declared branches: bridge completion, support-only content, bookkeeping, explicit domain shift, or independent same-domain cycle-fiber endpoint discrimination. This packet is therefore not decorative; it is a repeatable closure capsule for the theorem ladder and denial-burden table.


\subsection{Stress test: Support-only branch (route-exhaustion lens)}

\par\noindent\textbf{Load-bearing claim.} This stress test audits why support status cannot defeat endpoint closure. The claim is not allowed to function by implication, atmosphere, or familiarity with AASC vocabulary. It must be available as a same-carrier theorem target inside the Hodge endpoint proof and must be reused only after the carrier, branch, and report slots have been instantiated.

\par\noindent\textbf{Hostile reconstruction.} A hostile referee will try to reconstruct this point in the strongest ordinary Hodge-theoretic language: rational Hodge classes are native cohomological objects, Chow groups and cycle-class maps are ordinary algebraic-geometric carriers, and an empty fiber would be ordinary mathematical falsification. The audit therefore never answers by declaring the underlying mathematical content illicit. It answers only by separating raw content, support status, theoremhood, endpoint force, and the fixed cycle-fiber endpoint image.

\par\noindent\textbf{Closure condition.} The claim survives this stress test only if it preserves the same $X$, $p$, and $\alpha$, the same rational Hodge-standing predicate, the same group $\CH^p(X)_\QQ$, the same map $\clQ$, and the same fiber $F_X(\alpha)$. A proposed escape that changes any of these is a domain or carrier shift. A proposed escape that preserves them while retaining endpoint-defeating force enters endpoint-use discipline and must pass through the native empty-fiber exhaustion.

\par\noindent\textbf{Audit result.} Under the fixed endpoint assumptions, the hostile reconstruction leaves only the declared branches: bridge completion, support-only content, bookkeeping, explicit domain shift, or independent same-domain cycle-fiber endpoint discrimination. This packet is therefore not decorative; it is a repeatable closure capsule for the theorem ladder and denial-burden table.


\subsection{Stress test: Bookkeeping branch (route-exhaustion lens)}

\par\noindent\textbf{Load-bearing claim.} This stress test audits why labels and ledgers do not determine endpoint status. The claim is not allowed to function by implication, atmosphere, or familiarity with AASC vocabulary. It must be available as a same-carrier theorem target inside the Hodge endpoint proof and must be reused only after the carrier, branch, and report slots have been instantiated.

\par\noindent\textbf{Hostile reconstruction.} A hostile referee will try to reconstruct this point in the strongest ordinary Hodge-theoretic language: rational Hodge classes are native cohomological objects, Chow groups and cycle-class maps are ordinary algebraic-geometric carriers, and an empty fiber would be ordinary mathematical falsification. The audit therefore never answers by declaring the underlying mathematical content illicit. It answers only by separating raw content, support status, theoremhood, endpoint force, and the fixed cycle-fiber endpoint image.

\par\noindent\textbf{Closure condition.} The claim survives this stress test only if it preserves the same $X$, $p$, and $\alpha$, the same rational Hodge-standing predicate, the same group $\CH^p(X)_\QQ$, the same map $\clQ$, and the same fiber $F_X(\alpha)$. A proposed escape that changes any of these is a domain or carrier shift. A proposed escape that preserves them while retaining endpoint-defeating force enters endpoint-use discipline and must pass through the native empty-fiber exhaustion.

\par\noindent\textbf{Audit result.} Under the fixed endpoint assumptions, the hostile reconstruction leaves only the declared branches: bridge completion, support-only content, bookkeeping, explicit domain shift, or independent same-domain cycle-fiber endpoint discrimination. This packet is therefore not decorative; it is a repeatable closure capsule for the theorem ladder and denial-burden table.


\subsection{Stress test: Domain-shift branch (route-exhaustion lens)}

\par\noindent\textbf{Load-bearing claim.} This stress test audits why changing X,p,alpha, coefficients, equivalence, or carrier exits the fixed endpoint. The claim is not allowed to function by implication, atmosphere, or familiarity with AASC vocabulary. It must be available as a same-carrier theorem target inside the Hodge endpoint proof and must be reused only after the carrier, branch, and report slots have been instantiated.

\par\noindent\textbf{Hostile reconstruction.} A hostile referee will try to reconstruct this point in the strongest ordinary Hodge-theoretic language: rational Hodge classes are native cohomological objects, Chow groups and cycle-class maps are ordinary algebraic-geometric carriers, and an empty fiber would be ordinary mathematical falsification. The audit therefore never answers by declaring the underlying mathematical content illicit. It answers only by separating raw content, support status, theoremhood, endpoint force, and the fixed cycle-fiber endpoint image.

\par\noindent\textbf{Closure condition.} The claim survives this stress test only if it preserves the same $X$, $p$, and $\alpha$, the same rational Hodge-standing predicate, the same group $\CH^p(X)_\QQ$, the same map $\clQ$, and the same fiber $F_X(\alpha)$. A proposed escape that changes any of these is a domain or carrier shift. A proposed escape that preserves them while retaining endpoint-defeating force enters endpoint-use discipline and must pass through the native empty-fiber exhaustion.

\par\noindent\textbf{Audit result.} Under the fixed endpoint assumptions, the hostile reconstruction leaves only the declared branches: bridge completion, support-only content, bookkeeping, explicit domain shift, or independent same-domain cycle-fiber endpoint discrimination. This packet is therefore not decorative; it is a repeatable closure capsule for the theorem ladder and denial-burden table.


\subsection{Stress test: Cycle-fiber discriminator (route-exhaustion lens)}

\par\noindent\textbf{Load-bearing claim.} This stress test audits the independent same-domain discriminator induced by endpoint-standing empty-fiber force. The claim is not allowed to function by implication, atmosphere, or familiarity with AASC vocabulary. It must be available as a same-carrier theorem target inside the Hodge endpoint proof and must be reused only after the carrier, branch, and report slots have been instantiated.

\par\noindent\textbf{Hostile reconstruction.} A hostile referee will try to reconstruct this point in the strongest ordinary Hodge-theoretic language: rational Hodge classes are native cohomological objects, Chow groups and cycle-class maps are ordinary algebraic-geometric carriers, and an empty fiber would be ordinary mathematical falsification. The audit therefore never answers by declaring the underlying mathematical content illicit. It answers only by separating raw content, support status, theoremhood, endpoint force, and the fixed cycle-fiber endpoint image.

\par\noindent\textbf{Closure condition.} The claim survives this stress test only if it preserves the same $X$, $p$, and $\alpha$, the same rational Hodge-standing predicate, the same group $\CH^p(X)_\QQ$, the same map $\clQ$, and the same fiber $F_X(\alpha)$. A proposed escape that changes any of these is a domain or carrier shift. A proposed escape that preserves them while retaining endpoint-defeating force enters endpoint-use discipline and must pass through the native empty-fiber exhaustion.

\par\noindent\textbf{Audit result.} Under the fixed endpoint assumptions, the hostile reconstruction leaves only the declared branches: bridge completion, support-only content, bookkeeping, explicit domain shift, or independent same-domain cycle-fiber endpoint discrimination. This packet is therefore not decorative; it is a repeatable closure capsule for the theorem ladder and denial-burden table.


\subsection{Stress test: No-independent closure (route-exhaustion lens)}

\par\noindent\textbf{Load-bearing claim.} This stress test audits the A+ consequence excluding the independent discriminator. The claim is not allowed to function by implication, atmosphere, or familiarity with AASC vocabulary. It must be available as a same-carrier theorem target inside the Hodge endpoint proof and must be reused only after the carrier, branch, and report slots have been instantiated.

\par\noindent\textbf{Hostile reconstruction.} A hostile referee will try to reconstruct this point in the strongest ordinary Hodge-theoretic language: rational Hodge classes are native cohomological objects, Chow groups and cycle-class maps are ordinary algebraic-geometric carriers, and an empty fiber would be ordinary mathematical falsification. The audit therefore never answers by declaring the underlying mathematical content illicit. It answers only by separating raw content, support status, theoremhood, endpoint force, and the fixed cycle-fiber endpoint image.

\par\noindent\textbf{Closure condition.} The claim survives this stress test only if it preserves the same $X$, $p$, and $\alpha$, the same rational Hodge-standing predicate, the same group $\CH^p(X)_\QQ$, the same map $\clQ$, and the same fiber $F_X(\alpha)$. A proposed escape that changes any of these is a domain or carrier shift. A proposed escape that preserves them while retaining endpoint-defeating force enters endpoint-use discipline and must pass through the native empty-fiber exhaustion.

\par\noindent\textbf{Audit result.} Under the fixed endpoint assumptions, the hostile reconstruction leaves only the declared branches: bridge completion, support-only content, bookkeeping, explicit domain shift, or independent same-domain cycle-fiber endpoint discrimination. This packet is therefore not decorative; it is a repeatable closure capsule for the theorem ladder and denial-burden table.


\subsection{Stress test: Local separator occupation (route-exhaustion lens)}

\par\noindent\textbf{Load-bearing claim.} This stress test audits the local empty-fiber countercase as endpoint-defeating non-occupation. The claim is not allowed to function by implication, atmosphere, or familiarity with AASC vocabulary. It must be available as a same-carrier theorem target inside the Hodge endpoint proof and must be reused only after the carrier, branch, and report slots have been instantiated.

\par\noindent\textbf{Hostile reconstruction.} A hostile referee will try to reconstruct this point in the strongest ordinary Hodge-theoretic language: rational Hodge classes are native cohomological objects, Chow groups and cycle-class maps are ordinary algebraic-geometric carriers, and an empty fiber would be ordinary mathematical falsification. The audit therefore never answers by declaring the underlying mathematical content illicit. It answers only by separating raw content, support status, theoremhood, endpoint force, and the fixed cycle-fiber endpoint image.

\par\noindent\textbf{Closure condition.} The claim survives this stress test only if it preserves the same $X$, $p$, and $\alpha$, the same rational Hodge-standing predicate, the same group $\CH^p(X)_\QQ$, the same map $\clQ$, and the same fiber $F_X(\alpha)$. A proposed escape that changes any of these is a domain or carrier shift. A proposed escape that preserves them while retaining endpoint-defeating force enters endpoint-use discipline and must pass through the native empty-fiber exhaustion.

\par\noindent\textbf{Audit result.} Under the fixed endpoint assumptions, the hostile reconstruction leaves only the declared branches: bridge completion, support-only content, bookkeeping, explicit domain shift, or independent same-domain cycle-fiber endpoint discrimination. This packet is therefore not decorative; it is a repeatable closure capsule for the theorem ladder and denial-burden table.


\subsection{Stress test: Reductio annotation (route-exhaustion lens)}

\par\noindent\textbf{Load-bearing claim.} This stress test audits proof-act annotation not object-level premise. The claim is not allowed to function by implication, atmosphere, or familiarity with AASC vocabulary. It must be available as a same-carrier theorem target inside the Hodge endpoint proof and must be reused only after the carrier, branch, and report slots have been instantiated.

\par\noindent\textbf{Hostile reconstruction.} A hostile referee will try to reconstruct this point in the strongest ordinary Hodge-theoretic language: rational Hodge classes are native cohomological objects, Chow groups and cycle-class maps are ordinary algebraic-geometric carriers, and an empty fiber would be ordinary mathematical falsification. The audit therefore never answers by declaring the underlying mathematical content illicit. It answers only by separating raw content, support status, theoremhood, endpoint force, and the fixed cycle-fiber endpoint image.

\par\noindent\textbf{Closure condition.} The claim survives this stress test only if it preserves the same $X$, $p$, and $\alpha$, the same rational Hodge-standing predicate, the same group $\CH^p(X)_\QQ$, the same map $\clQ$, and the same fiber $F_X(\alpha)$. A proposed escape that changes any of these is a domain or carrier shift. A proposed escape that preserves them while retaining endpoint-defeating force enters endpoint-use discipline and must pass through the native empty-fiber exhaustion.

\par\noindent\textbf{Audit result.} Under the fixed endpoint assumptions, the hostile reconstruction leaves only the declared branches: bridge completion, support-only content, bookkeeping, explicit domain shift, or independent same-domain cycle-fiber endpoint discrimination. This packet is therefore not decorative; it is a repeatable closure capsule for the theorem ladder and denial-burden table.


\subsection{Stress test: Annotation discharge (route-exhaustion lens)}

\par\noindent\textbf{Load-bearing claim.} This stress test audits the original empty-fiber assumption is discharged. The claim is not allowed to function by implication, atmosphere, or familiarity with AASC vocabulary. It must be available as a same-carrier theorem target inside the Hodge endpoint proof and must be reused only after the carrier, branch, and report slots have been instantiated.

\par\noindent\textbf{Hostile reconstruction.} A hostile referee will try to reconstruct this point in the strongest ordinary Hodge-theoretic language: rational Hodge classes are native cohomological objects, Chow groups and cycle-class maps are ordinary algebraic-geometric carriers, and an empty fiber would be ordinary mathematical falsification. The audit therefore never answers by declaring the underlying mathematical content illicit. It answers only by separating raw content, support status, theoremhood, endpoint force, and the fixed cycle-fiber endpoint image.

\par\noindent\textbf{Closure condition.} The claim survives this stress test only if it preserves the same $X$, $p$, and $\alpha$, the same rational Hodge-standing predicate, the same group $\CH^p(X)_\QQ$, the same map $\clQ$, and the same fiber $F_X(\alpha)$. A proposed escape that changes any of these is a domain or carrier shift. A proposed escape that preserves them while retaining endpoint-defeating force enters endpoint-use discipline and must pass through the native empty-fiber exhaustion.

\par\noindent\textbf{Audit result.} Under the fixed endpoint assumptions, the hostile reconstruction leaves only the declared branches: bridge completion, support-only content, bookkeeping, explicit domain shift, or independent same-domain cycle-fiber endpoint discrimination. This packet is therefore not decorative; it is a repeatable closure capsule for the theorem ladder and denial-burden table.


\subsection{Stress test: Ordinary theoremhood protection (route-exhaustion lens)}

\par\noindent\textbf{Load-bearing claim.} This stress test audits negative theoremhood is allowed until used as official endpoint resolution. The claim is not allowed to function by implication, atmosphere, or familiarity with AASC vocabulary. It must be available as a same-carrier theorem target inside the Hodge endpoint proof and must be reused only after the carrier, branch, and report slots have been instantiated.

\par\noindent\textbf{Hostile reconstruction.} A hostile referee will try to reconstruct this point in the strongest ordinary Hodge-theoretic language: rational Hodge classes are native cohomological objects, Chow groups and cycle-class maps are ordinary algebraic-geometric carriers, and an empty fiber would be ordinary mathematical falsification. The audit therefore never answers by declaring the underlying mathematical content illicit. It answers only by separating raw content, support status, theoremhood, endpoint force, and the fixed cycle-fiber endpoint image.

\par\noindent\textbf{Closure condition.} The claim survives this stress test only if it preserves the same $X$, $p$, and $\alpha$, the same rational Hodge-standing predicate, the same group $\CH^p(X)_\QQ$, the same map $\clQ$, and the same fiber $F_X(\alpha)$. A proposed escape that changes any of these is a domain or carrier shift. A proposed escape that preserves them while retaining endpoint-defeating force enters endpoint-use discipline and must pass through the native empty-fiber exhaustion.

\par\noindent\textbf{Audit result.} Under the fixed endpoint assumptions, the hostile reconstruction leaves only the declared branches: bridge completion, support-only content, bookkeeping, explicit domain shift, or independent same-domain cycle-fiber endpoint discrimination. This packet is therefore not decorative; it is a repeatable closure capsule for the theorem ladder and denial-burden table.


\subsection{Stress test: Incompleteness firewall (route-exhaustion lens)}

\par\noindent\textbf{Load-bearing claim.} This stress test audits same-domain independence matters only if endpoint-standing relevant. The claim is not allowed to function by implication, atmosphere, or familiarity with AASC vocabulary. It must be available as a same-carrier theorem target inside the Hodge endpoint proof and must be reused only after the carrier, branch, and report slots have been instantiated.

\par\noindent\textbf{Hostile reconstruction.} A hostile referee will try to reconstruct this point in the strongest ordinary Hodge-theoretic language: rational Hodge classes are native cohomological objects, Chow groups and cycle-class maps are ordinary algebraic-geometric carriers, and an empty fiber would be ordinary mathematical falsification. The audit therefore never answers by declaring the underlying mathematical content illicit. It answers only by separating raw content, support status, theoremhood, endpoint force, and the fixed cycle-fiber endpoint image.

\par\noindent\textbf{Closure condition.} The claim survives this stress test only if it preserves the same $X$, $p$, and $\alpha$, the same rational Hodge-standing predicate, the same group $\CH^p(X)_\QQ$, the same map $\clQ$, and the same fiber $F_X(\alpha)$. A proposed escape that changes any of these is a domain or carrier shift. A proposed escape that preserves them while retaining endpoint-defeating force enters endpoint-use discipline and must pass through the native empty-fiber exhaustion.

\par\noindent\textbf{Audit result.} Under the fixed endpoint assumptions, the hostile reconstruction leaves only the declared branches: bridge completion, support-only content, bookkeeping, explicit domain shift, or independent same-domain cycle-fiber endpoint discrimination. This packet is therefore not decorative; it is a repeatable closure capsule for the theorem ladder and denial-burden table.


\subsection{Stress test: Route-locus exhaustion (route-exhaustion lens)}

\par\noindent\textbf{Load-bearing claim.} This stress test audits status, codomain, existence, carrier, temporal, locality, selector, and second-classifier loci. The claim is not allowed to function by implication, atmosphere, or familiarity with AASC vocabulary. It must be available as a same-carrier theorem target inside the Hodge endpoint proof and must be reused only after the carrier, branch, and report slots have been instantiated.

\par\noindent\textbf{Hostile reconstruction.} A hostile referee will try to reconstruct this point in the strongest ordinary Hodge-theoretic language: rational Hodge classes are native cohomological objects, Chow groups and cycle-class maps are ordinary algebraic-geometric carriers, and an empty fiber would be ordinary mathematical falsification. The audit therefore never answers by declaring the underlying mathematical content illicit. It answers only by separating raw content, support status, theoremhood, endpoint force, and the fixed cycle-fiber endpoint image.

\par\noindent\textbf{Closure condition.} The claim survives this stress test only if it preserves the same $X$, $p$, and $\alpha$, the same rational Hodge-standing predicate, the same group $\CH^p(X)_\QQ$, the same map $\clQ$, and the same fiber $F_X(\alpha)$. A proposed escape that changes any of these is a domain or carrier shift. A proposed escape that preserves them while retaining endpoint-defeating force enters endpoint-use discipline and must pass through the native empty-fiber exhaustion.

\par\noindent\textbf{Audit result.} Under the fixed endpoint assumptions, the hostile reconstruction leaves only the declared branches: bridge completion, support-only content, bookkeeping, explicit domain shift, or independent same-domain cycle-fiber endpoint discrimination. This packet is therefore not decorative; it is a repeatable closure capsule for the theorem ladder and denial-burden table.


\subsection{Stress test: Weakening resistance (route-exhaustion lens)}

\par\noindent\textbf{Load-bearing claim.} This stress test audits weaker regimes preserve behavior, shift domain, lower status, or reintroduce forbidden governance. The claim is not allowed to function by implication, atmosphere, or familiarity with AASC vocabulary. It must be available as a same-carrier theorem target inside the Hodge endpoint proof and must be reused only after the carrier, branch, and report slots have been instantiated.

\par\noindent\textbf{Hostile reconstruction.} A hostile referee will try to reconstruct this point in the strongest ordinary Hodge-theoretic language: rational Hodge classes are native cohomological objects, Chow groups and cycle-class maps are ordinary algebraic-geometric carriers, and an empty fiber would be ordinary mathematical falsification. The audit therefore never answers by declaring the underlying mathematical content illicit. It answers only by separating raw content, support status, theoremhood, endpoint force, and the fixed cycle-fiber endpoint image.

\par\noindent\textbf{Closure condition.} The claim survives this stress test only if it preserves the same $X$, $p$, and $\alpha$, the same rational Hodge-standing predicate, the same group $\CH^p(X)_\QQ$, the same map $\clQ$, and the same fiber $F_X(\alpha)$. A proposed escape that changes any of these is a domain or carrier shift. A proposed escape that preserves them while retaining endpoint-defeating force enters endpoint-use discipline and must pass through the native empty-fiber exhaustion.

\par\noindent\textbf{Audit result.} Under the fixed endpoint assumptions, the hostile reconstruction leaves only the declared branches: bridge completion, support-only content, bookkeeping, explicit domain shift, or independent same-domain cycle-fiber endpoint discrimination. This packet is therefore not decorative; it is a repeatable closure capsule for the theorem ladder and denial-burden table.


\subsection{Stress test: Official endpoint identity (report-preservation lens)}

\par\noindent\textbf{Load-bearing claim.} This stress test audits the identification of the Clay-facing endpoint with universal nonempty rational cycle-class fibers. The claim is not allowed to function by implication, atmosphere, or familiarity with AASC vocabulary. It must be available as a same-carrier theorem target inside the Hodge endpoint proof and must be reused only after the carrier, branch, and report slots have been instantiated.

\par\noindent\textbf{Hostile reconstruction.} A hostile referee will try to reconstruct this point in the strongest ordinary Hodge-theoretic language: rational Hodge classes are native cohomological objects, Chow groups and cycle-class maps are ordinary algebraic-geometric carriers, and an empty fiber would be ordinary mathematical falsification. The audit therefore never answers by declaring the underlying mathematical content illicit. It answers only by separating raw content, support status, theoremhood, endpoint force, and the fixed cycle-fiber endpoint image.

\par\noindent\textbf{Closure condition.} The claim survives this stress test only if it preserves the same $X$, $p$, and $\alpha$, the same rational Hodge-standing predicate, the same group $\CH^p(X)_\QQ$, the same map $\clQ$, and the same fiber $F_X(\alpha)$. A proposed escape that changes any of these is a domain or carrier shift. A proposed escape that preserves them while retaining endpoint-defeating force enters endpoint-use discipline and must pass through the native empty-fiber exhaustion.

\par\noindent\textbf{Audit result.} Under the fixed endpoint assumptions, the hostile reconstruction leaves only the declared branches: bridge completion, support-only content, bookkeeping, explicit domain shift, or independent same-domain cycle-fiber endpoint discrimination. This packet is therefore not decorative; it is a repeatable closure capsule for the theorem ladder and denial-burden table.


\subsection{Stress test: Context non-smuggling (report-preservation lens)}

\par\noindent\textbf{Load-bearing claim.} This stress test audits the claim that the Hodge endpoint context fixes the carrier and slots without asserting nonempty fiber. The claim is not allowed to function by implication, atmosphere, or familiarity with AASC vocabulary. It must be available as a same-carrier theorem target inside the Hodge endpoint proof and must be reused only after the carrier, branch, and report slots have been instantiated.

\par\noindent\textbf{Hostile reconstruction.} A hostile referee will try to reconstruct this point in the strongest ordinary Hodge-theoretic language: rational Hodge classes are native cohomological objects, Chow groups and cycle-class maps are ordinary algebraic-geometric carriers, and an empty fiber would be ordinary mathematical falsification. The audit therefore never answers by declaring the underlying mathematical content illicit. It answers only by separating raw content, support status, theoremhood, endpoint force, and the fixed cycle-fiber endpoint image.

\par\noindent\textbf{Closure condition.} The claim survives this stress test only if it preserves the same $X$, $p$, and $\alpha$, the same rational Hodge-standing predicate, the same group $\CH^p(X)_\QQ$, the same map $\clQ$, and the same fiber $F_X(\alpha)$. A proposed escape that changes any of these is a domain or carrier shift. A proposed escape that preserves them while retaining endpoint-defeating force enters endpoint-use discipline and must pass through the native empty-fiber exhaustion.

\par\noindent\textbf{Audit result.} Under the fixed endpoint assumptions, the hostile reconstruction leaves only the declared branches: bridge completion, support-only content, bookkeeping, explicit domain shift, or independent same-domain cycle-fiber endpoint discrimination. This packet is therefore not decorative; it is a repeatable closure capsule for the theorem ladder and denial-burden table.


\subsection{Stress test: Hodge-standing exactness (report-preservation lens)}

\par\noindent\textbf{Load-bearing claim.} This stress test audits the representation of rational Hodge standing as cohomological type status. The claim is not allowed to function by implication, atmosphere, or familiarity with AASC vocabulary. It must be available as a same-carrier theorem target inside the Hodge endpoint proof and must be reused only after the carrier, branch, and report slots have been instantiated.

\par\noindent\textbf{Hostile reconstruction.} A hostile referee will try to reconstruct this point in the strongest ordinary Hodge-theoretic language: rational Hodge classes are native cohomological objects, Chow groups and cycle-class maps are ordinary algebraic-geometric carriers, and an empty fiber would be ordinary mathematical falsification. The audit therefore never answers by declaring the underlying mathematical content illicit. It answers only by separating raw content, support status, theoremhood, endpoint force, and the fixed cycle-fiber endpoint image.

\par\noindent\textbf{Closure condition.} The claim survives this stress test only if it preserves the same $X$, $p$, and $\alpha$, the same rational Hodge-standing predicate, the same group $\CH^p(X)_\QQ$, the same map $\clQ$, and the same fiber $F_X(\alpha)$. A proposed escape that changes any of these is a domain or carrier shift. A proposed escape that preserves them while retaining endpoint-defeating force enters endpoint-use discipline and must pass through the native empty-fiber exhaustion.

\par\noindent\textbf{Audit result.} Under the fixed endpoint assumptions, the hostile reconstruction leaves only the declared branches: bridge completion, support-only content, bookkeeping, explicit domain shift, or independent same-domain cycle-fiber endpoint discrimination. This packet is therefore not decorative; it is a repeatable closure capsule for the theorem ladder and denial-burden table.


\subsection{Stress test: Cycle-fiber exactness (report-preservation lens)}

\par\noindent\textbf{Load-bearing claim.} This stress test audits the representation of algebraic readout by the fiber $F_X(\alpha)$. The claim is not allowed to function by implication, atmosphere, or familiarity with AASC vocabulary. It must be available as a same-carrier theorem target inside the Hodge endpoint proof and must be reused only after the carrier, branch, and report slots have been instantiated.

\par\noindent\textbf{Hostile reconstruction.} A hostile referee will try to reconstruct this point in the strongest ordinary Hodge-theoretic language: rational Hodge classes are native cohomological objects, Chow groups and cycle-class maps are ordinary algebraic-geometric carriers, and an empty fiber would be ordinary mathematical falsification. The audit therefore never answers by declaring the underlying mathematical content illicit. It answers only by separating raw content, support status, theoremhood, endpoint force, and the fixed cycle-fiber endpoint image.

\par\noindent\textbf{Closure condition.} The claim survives this stress test only if it preserves the same $X$, $p$, and $\alpha$, the same rational Hodge-standing predicate, the same group $\CH^p(X)_\QQ$, the same map $\clQ$, and the same fiber $F_X(\alpha)$. A proposed escape that changes any of these is a domain or carrier shift. A proposed escape that preserves them while retaining endpoint-defeating force enters endpoint-use discipline and must pass through the native empty-fiber exhaustion.

\par\noindent\textbf{Audit result.} Under the fixed endpoint assumptions, the hostile reconstruction leaves only the declared branches: bridge completion, support-only content, bookkeeping, explicit domain shift, or independent same-domain cycle-fiber endpoint discrimination. This packet is therefore not decorative; it is a repeatable closure capsule for the theorem ladder and denial-burden table.


\subsection{Stress test: Empty-fiber exactness (report-preservation lens)}

\par\noindent\textbf{Load-bearing claim.} This stress test audits the negative local branch as fiber non-occupation. The claim is not allowed to function by implication, atmosphere, or familiarity with AASC vocabulary. It must be available as a same-carrier theorem target inside the Hodge endpoint proof and must be reused only after the carrier, branch, and report slots have been instantiated.

\par\noindent\textbf{Hostile reconstruction.} A hostile referee will try to reconstruct this point in the strongest ordinary Hodge-theoretic language: rational Hodge classes are native cohomological objects, Chow groups and cycle-class maps are ordinary algebraic-geometric carriers, and an empty fiber would be ordinary mathematical falsification. The audit therefore never answers by declaring the underlying mathematical content illicit. It answers only by separating raw content, support status, theoremhood, endpoint force, and the fixed cycle-fiber endpoint image.

\par\noindent\textbf{Closure condition.} The claim survives this stress test only if it preserves the same $X$, $p$, and $\alpha$, the same rational Hodge-standing predicate, the same group $\CH^p(X)_\QQ$, the same map $\clQ$, and the same fiber $F_X(\alpha)$. A proposed escape that changes any of these is a domain or carrier shift. A proposed escape that preserves them while retaining endpoint-defeating force enters endpoint-use discipline and must pass through the native empty-fiber exhaustion.

\par\noindent\textbf{Audit result.} Under the fixed endpoint assumptions, the hostile reconstruction leaves only the declared branches: bridge completion, support-only content, bookkeeping, explicit domain shift, or independent same-domain cycle-fiber endpoint discrimination. This packet is therefore not decorative; it is a repeatable closure capsule for the theorem ladder and denial-burden table.


\subsection{Stress test: Counterexample object/use distinction (report-preservation lens)}

\par\noindent\textbf{Load-bearing claim.} This stress test audits the separation between an object satisfying the negative branch and the act of offering it as endpoint defeat. The claim is not allowed to function by implication, atmosphere, or familiarity with AASC vocabulary. It must be available as a same-carrier theorem target inside the Hodge endpoint proof and must be reused only after the carrier, branch, and report slots have been instantiated.

\par\noindent\textbf{Hostile reconstruction.} A hostile referee will try to reconstruct this point in the strongest ordinary Hodge-theoretic language: rational Hodge classes are native cohomological objects, Chow groups and cycle-class maps are ordinary algebraic-geometric carriers, and an empty fiber would be ordinary mathematical falsification. The audit therefore never answers by declaring the underlying mathematical content illicit. It answers only by separating raw content, support status, theoremhood, endpoint force, and the fixed cycle-fiber endpoint image.

\par\noindent\textbf{Closure condition.} The claim survives this stress test only if it preserves the same $X$, $p$, and $\alpha$, the same rational Hodge-standing predicate, the same group $\CH^p(X)_\QQ$, the same map $\clQ$, and the same fiber $F_X(\alpha)$. A proposed escape that changes any of these is a domain or carrier shift. A proposed escape that preserves them while retaining endpoint-defeating force enters endpoint-use discipline and must pass through the native empty-fiber exhaustion.

\par\noindent\textbf{Audit result.} Under the fixed endpoint assumptions, the hostile reconstruction leaves only the declared branches: bridge completion, support-only content, bookkeeping, explicit domain shift, or independent same-domain cycle-fiber endpoint discrimination. This packet is therefore not decorative; it is a repeatable closure capsule for the theorem ladder and denial-burden table.


\subsection{Stress test: Official counterexample force (report-preservation lens)}

\par\noindent\textbf{Load-bearing claim.} This stress test audits the nonclassical dependency identifying endpoint-defeating counterexample force with theorem-bearing endpoint use. The claim is not allowed to function by implication, atmosphere, or familiarity with AASC vocabulary. It must be available as a same-carrier theorem target inside the Hodge endpoint proof and must be reused only after the carrier, branch, and report slots have been instantiated.

\par\noindent\textbf{Hostile reconstruction.} A hostile referee will try to reconstruct this point in the strongest ordinary Hodge-theoretic language: rational Hodge classes are native cohomological objects, Chow groups and cycle-class maps are ordinary algebraic-geometric carriers, and an empty fiber would be ordinary mathematical falsification. The audit therefore never answers by declaring the underlying mathematical content illicit. It answers only by separating raw content, support status, theoremhood, endpoint force, and the fixed cycle-fiber endpoint image.

\par\noindent\textbf{Closure condition.} The claim survives this stress test only if it preserves the same $X$, $p$, and $\alpha$, the same rational Hodge-standing predicate, the same group $\CH^p(X)_\QQ$, the same map $\clQ$, and the same fiber $F_X(\alpha)$. A proposed escape that changes any of these is a domain or carrier shift. A proposed escape that preserves them while retaining endpoint-defeating force enters endpoint-use discipline and must pass through the native empty-fiber exhaustion.

\par\noindent\textbf{Audit result.} Under the fixed endpoint assumptions, the hostile reconstruction leaves only the declared branches: bridge completion, support-only content, bookkeeping, explicit domain shift, or independent same-domain cycle-fiber endpoint discrimination. This packet is therefore not decorative; it is a repeatable closure capsule for the theorem ladder and denial-burden table.


\subsection{Stress test: No autonomous counterexample endpoint role (report-preservation lens)}

\par\noindent\textbf{Load-bearing claim.} This stress test audits the denial of a fifth role outside endpoint use for official endpoint defeat. The claim is not allowed to function by implication, atmosphere, or familiarity with AASC vocabulary. It must be available as a same-carrier theorem target inside the Hodge endpoint proof and must be reused only after the carrier, branch, and report slots have been instantiated.

\par\noindent\textbf{Hostile reconstruction.} A hostile referee will try to reconstruct this point in the strongest ordinary Hodge-theoretic language: rational Hodge classes are native cohomological objects, Chow groups and cycle-class maps are ordinary algebraic-geometric carriers, and an empty fiber would be ordinary mathematical falsification. The audit therefore never answers by declaring the underlying mathematical content illicit. It answers only by separating raw content, support status, theoremhood, endpoint force, and the fixed cycle-fiber endpoint image.

\par\noindent\textbf{Closure condition.} The claim survives this stress test only if it preserves the same $X$, $p$, and $\alpha$, the same rational Hodge-standing predicate, the same group $\CH^p(X)_\QQ$, the same map $\clQ$, and the same fiber $F_X(\alpha)$. A proposed escape that changes any of these is a domain or carrier shift. A proposed escape that preserves them while retaining endpoint-defeating force enters endpoint-use discipline and must pass through the native empty-fiber exhaustion.

\par\noindent\textbf{Audit result.} Under the fixed endpoint assumptions, the hostile reconstruction leaves only the declared branches: bridge completion, support-only content, bookkeeping, explicit domain shift, or independent same-domain cycle-fiber endpoint discrimination. This packet is therefore not decorative; it is a repeatable closure capsule for the theorem ladder and denial-burden table.


\subsection{Stress test: Kernel forcing (report-preservation lens)}

\par\noindent\textbf{Load-bearing claim.} This stress test audits the derivation of reference, standing, admissibility, and irreversibility from target adequacy. The claim is not allowed to function by implication, atmosphere, or familiarity with AASC vocabulary. It must be available as a same-carrier theorem target inside the Hodge endpoint proof and must be reused only after the carrier, branch, and report slots have been instantiated.

\par\noindent\textbf{Hostile reconstruction.} A hostile referee will try to reconstruct this point in the strongest ordinary Hodge-theoretic language: rational Hodge classes are native cohomological objects, Chow groups and cycle-class maps are ordinary algebraic-geometric carriers, and an empty fiber would be ordinary mathematical falsification. The audit therefore never answers by declaring the underlying mathematical content illicit. It answers only by separating raw content, support status, theoremhood, endpoint force, and the fixed cycle-fiber endpoint image.

\par\noindent\textbf{Closure condition.} The claim survives this stress test only if it preserves the same $X$, $p$, and $\alpha$, the same rational Hodge-standing predicate, the same group $\CH^p(X)_\QQ$, the same map $\clQ$, and the same fiber $F_X(\alpha)$. A proposed escape that changes any of these is a domain or carrier shift. A proposed escape that preserves them while retaining endpoint-defeating force enters endpoint-use discipline and must pass through the native empty-fiber exhaustion.

\par\noindent\textbf{Audit result.} Under the fixed endpoint assumptions, the hostile reconstruction leaves only the declared branches: bridge completion, support-only content, bookkeeping, explicit domain shift, or independent same-domain cycle-fiber endpoint discrimination. This packet is therefore not decorative; it is a repeatable closure capsule for the theorem ladder and denial-burden table.


\subsection{Stress test: A+ consequence layer (report-preservation lens)}

\par\noindent\textbf{Load-bearing claim.} This stress test audits the fixed-domain consequence package binding endpoint use. The claim is not allowed to function by implication, atmosphere, or familiarity with AASC vocabulary. It must be available as a same-carrier theorem target inside the Hodge endpoint proof and must be reused only after the carrier, branch, and report slots have been instantiated.

\par\noindent\textbf{Hostile reconstruction.} A hostile referee will try to reconstruct this point in the strongest ordinary Hodge-theoretic language: rational Hodge classes are native cohomological objects, Chow groups and cycle-class maps are ordinary algebraic-geometric carriers, and an empty fiber would be ordinary mathematical falsification. The audit therefore never answers by declaring the underlying mathematical content illicit. It answers only by separating raw content, support status, theoremhood, endpoint force, and the fixed cycle-fiber endpoint image.

\par\noindent\textbf{Closure condition.} The claim survives this stress test only if it preserves the same $X$, $p$, and $\alpha$, the same rational Hodge-standing predicate, the same group $\CH^p(X)_\QQ$, the same map $\clQ$, and the same fiber $F_X(\alpha)$. A proposed escape that changes any of these is a domain or carrier shift. A proposed escape that preserves them while retaining endpoint-defeating force enters endpoint-use discipline and must pass through the native empty-fiber exhaustion.

\par\noindent\textbf{Audit result.} Under the fixed endpoint assumptions, the hostile reconstruction leaves only the declared branches: bridge completion, support-only content, bookkeeping, explicit domain shift, or independent same-domain cycle-fiber endpoint discrimination. This packet is therefore not decorative; it is a repeatable closure capsule for the theorem ladder and denial-burden table.


\subsection{Stress test: ATS locking (report-preservation lens)}

\par\noindent\textbf{Load-bearing claim.} This stress test audits same target, carrier, branch, readout trajectory, and no second classifier. The claim is not allowed to function by implication, atmosphere, or familiarity with AASC vocabulary. It must be available as a same-carrier theorem target inside the Hodge endpoint proof and must be reused only after the carrier, branch, and report slots have been instantiated.

\par\noindent\textbf{Hostile reconstruction.} A hostile referee will try to reconstruct this point in the strongest ordinary Hodge-theoretic language: rational Hodge classes are native cohomological objects, Chow groups and cycle-class maps are ordinary algebraic-geometric carriers, and an empty fiber would be ordinary mathematical falsification. The audit therefore never answers by declaring the underlying mathematical content illicit. It answers only by separating raw content, support status, theoremhood, endpoint force, and the fixed cycle-fiber endpoint image.

\par\noindent\textbf{Closure condition.} The claim survives this stress test only if it preserves the same $X$, $p$, and $\alpha$, the same rational Hodge-standing predicate, the same group $\CH^p(X)_\QQ$, the same map $\clQ$, and the same fiber $F_X(\alpha)$. A proposed escape that changes any of these is a domain or carrier shift. A proposed escape that preserves them while retaining endpoint-defeating force enters endpoint-use discipline and must pass through the native empty-fiber exhaustion.

\par\noindent\textbf{Audit result.} Under the fixed endpoint assumptions, the hostile reconstruction leaves only the declared branches: bridge completion, support-only content, bookkeeping, explicit domain shift, or independent same-domain cycle-fiber endpoint discrimination. This packet is therefore not decorative; it is a repeatable closure capsule for the theorem ladder and denial-burden table.


\subsection{Stress test: UEAP report discipline (report-preservation lens)}

\par\noindent\textbf{Load-bearing claim.} This stress test audits no report exceeding carrier, bridge, and support status. The claim is not allowed to function by implication, atmosphere, or familiarity with AASC vocabulary. It must be available as a same-carrier theorem target inside the Hodge endpoint proof and must be reused only after the carrier, branch, and report slots have been instantiated.

\par\noindent\textbf{Hostile reconstruction.} A hostile referee will try to reconstruct this point in the strongest ordinary Hodge-theoretic language: rational Hodge classes are native cohomological objects, Chow groups and cycle-class maps are ordinary algebraic-geometric carriers, and an empty fiber would be ordinary mathematical falsification. The audit therefore never answers by declaring the underlying mathematical content illicit. It answers only by separating raw content, support status, theoremhood, endpoint force, and the fixed cycle-fiber endpoint image.

\par\noindent\textbf{Closure condition.} The claim survives this stress test only if it preserves the same $X$, $p$, and $\alpha$, the same rational Hodge-standing predicate, the same group $\CH^p(X)_\QQ$, the same map $\clQ$, and the same fiber $F_X(\alpha)$. A proposed escape that changes any of these is a domain or carrier shift. A proposed escape that preserves them while retaining endpoint-defeating force enters endpoint-use discipline and must pass through the native empty-fiber exhaustion.

\par\noindent\textbf{Audit result.} Under the fixed endpoint assumptions, the hostile reconstruction leaves only the declared branches: bridge completion, support-only content, bookkeeping, explicit domain shift, or independent same-domain cycle-fiber endpoint discrimination. This packet is therefore not decorative; it is a repeatable closure capsule for the theorem ladder and denial-burden table.


\subsection{Stress test: AMetric no-selector boundary (report-preservation lens)}

\par\noindent\textbf{Load-bearing claim.} This stress test audits no preferred representative or scalar ranking as pre-bridge authority. The claim is not allowed to function by implication, atmosphere, or familiarity with AASC vocabulary. It must be available as a same-carrier theorem target inside the Hodge endpoint proof and must be reused only after the carrier, branch, and report slots have been instantiated.

\par\noindent\textbf{Hostile reconstruction.} A hostile referee will try to reconstruct this point in the strongest ordinary Hodge-theoretic language: rational Hodge classes are native cohomological objects, Chow groups and cycle-class maps are ordinary algebraic-geometric carriers, and an empty fiber would be ordinary mathematical falsification. The audit therefore never answers by declaring the underlying mathematical content illicit. It answers only by separating raw content, support status, theoremhood, endpoint force, and the fixed cycle-fiber endpoint image.

\par\noindent\textbf{Closure condition.} The claim survives this stress test only if it preserves the same $X$, $p$, and $\alpha$, the same rational Hodge-standing predicate, the same group $\CH^p(X)_\QQ$, the same map $\clQ$, and the same fiber $F_X(\alpha)$. A proposed escape that changes any of these is a domain or carrier shift. A proposed escape that preserves them while retaining endpoint-defeating force enters endpoint-use discipline and must pass through the native empty-fiber exhaustion.

\par\noindent\textbf{Audit result.} Under the fixed endpoint assumptions, the hostile reconstruction leaves only the declared branches: bridge completion, support-only content, bookkeeping, explicit domain shift, or independent same-domain cycle-fiber endpoint discrimination. This packet is therefore not decorative; it is a repeatable closure capsule for the theorem ladder and denial-burden table.


\subsection{Stress test: Surrogate carrier separation (report-preservation lens)}

\par\noindent\textbf{Load-bearing claim.} This stress test audits motivic, variational, or obstruction objects as bridge support rather than readout. The claim is not allowed to function by implication, atmosphere, or familiarity with AASC vocabulary. It must be available as a same-carrier theorem target inside the Hodge endpoint proof and must be reused only after the carrier, branch, and report slots have been instantiated.

\par\noindent\textbf{Hostile reconstruction.} A hostile referee will try to reconstruct this point in the strongest ordinary Hodge-theoretic language: rational Hodge classes are native cohomological objects, Chow groups and cycle-class maps are ordinary algebraic-geometric carriers, and an empty fiber would be ordinary mathematical falsification. The audit therefore never answers by declaring the underlying mathematical content illicit. It answers only by separating raw content, support status, theoremhood, endpoint force, and the fixed cycle-fiber endpoint image.

\par\noindent\textbf{Closure condition.} The claim survives this stress test only if it preserves the same $X$, $p$, and $\alpha$, the same rational Hodge-standing predicate, the same group $\CH^p(X)_\QQ$, the same map $\clQ$, and the same fiber $F_X(\alpha)$. A proposed escape that changes any of these is a domain or carrier shift. A proposed escape that preserves them while retaining endpoint-defeating force enters endpoint-use discipline and must pass through the native empty-fiber exhaustion.

\par\noindent\textbf{Audit result.} Under the fixed endpoint assumptions, the hostile reconstruction leaves only the declared branches: bridge completion, support-only content, bookkeeping, explicit domain shift, or independent same-domain cycle-fiber endpoint discrimination. This packet is therefore not decorative; it is a repeatable closure capsule for the theorem ladder and denial-burden table.


\subsection{Stress test: No post-selection repair (report-preservation lens)}

\par\noindent\textbf{Load-bearing claim.} This stress test audits later cycle discovery as fresh bridge completion rather than same-act repair. The claim is not allowed to function by implication, atmosphere, or familiarity with AASC vocabulary. It must be available as a same-carrier theorem target inside the Hodge endpoint proof and must be reused only after the carrier, branch, and report slots have been instantiated.

\par\noindent\textbf{Hostile reconstruction.} A hostile referee will try to reconstruct this point in the strongest ordinary Hodge-theoretic language: rational Hodge classes are native cohomological objects, Chow groups and cycle-class maps are ordinary algebraic-geometric carriers, and an empty fiber would be ordinary mathematical falsification. The audit therefore never answers by declaring the underlying mathematical content illicit. It answers only by separating raw content, support status, theoremhood, endpoint force, and the fixed cycle-fiber endpoint image.

\par\noindent\textbf{Closure condition.} The claim survives this stress test only if it preserves the same $X$, $p$, and $\alpha$, the same rational Hodge-standing predicate, the same group $\CH^p(X)_\QQ$, the same map $\clQ$, and the same fiber $F_X(\alpha)$. A proposed escape that changes any of these is a domain or carrier shift. A proposed escape that preserves them while retaining endpoint-defeating force enters endpoint-use discipline and must pass through the native empty-fiber exhaustion.

\par\noindent\textbf{Audit result.} Under the fixed endpoint assumptions, the hostile reconstruction leaves only the declared branches: bridge completion, support-only content, bookkeeping, explicit domain shift, or independent same-domain cycle-fiber endpoint discrimination. This packet is therefore not decorative; it is a repeatable closure capsule for the theorem ladder and denial-burden table.


\subsection{Stress test: No carrier transfer (report-preservation lens)}

\par\noindent\textbf{Load-bearing claim.} This stress test audits known cases and neighboring domains not carrying endpoint standing without target inclusion. The claim is not allowed to function by implication, atmosphere, or familiarity with AASC vocabulary. It must be available as a same-carrier theorem target inside the Hodge endpoint proof and must be reused only after the carrier, branch, and report slots have been instantiated.

\par\noindent\textbf{Hostile reconstruction.} A hostile referee will try to reconstruct this point in the strongest ordinary Hodge-theoretic language: rational Hodge classes are native cohomological objects, Chow groups and cycle-class maps are ordinary algebraic-geometric carriers, and an empty fiber would be ordinary mathematical falsification. The audit therefore never answers by declaring the underlying mathematical content illicit. It answers only by separating raw content, support status, theoremhood, endpoint force, and the fixed cycle-fiber endpoint image.

\par\noindent\textbf{Closure condition.} The claim survives this stress test only if it preserves the same $X$, $p$, and $\alpha$, the same rational Hodge-standing predicate, the same group $\CH^p(X)_\QQ$, the same map $\clQ$, and the same fiber $F_X(\alpha)$. A proposed escape that changes any of these is a domain or carrier shift. A proposed escape that preserves them while retaining endpoint-defeating force enters endpoint-use discipline and must pass through the native empty-fiber exhaustion.

\par\noindent\textbf{Audit result.} Under the fixed endpoint assumptions, the hostile reconstruction leaves only the declared branches: bridge completion, support-only content, bookkeeping, explicit domain shift, or independent same-domain cycle-fiber endpoint discrimination. This packet is therefore not decorative; it is a repeatable closure capsule for the theorem ladder and denial-burden table.


\subsection{Stress test: No generators (report-preservation lens)}

\par\noindent\textbf{Load-bearing claim.} This stress test audits Hodge standing not generating Chow-fiber nonemptiness. The claim is not allowed to function by implication, atmosphere, or familiarity with AASC vocabulary. It must be available as a same-carrier theorem target inside the Hodge endpoint proof and must be reused only after the carrier, branch, and report slots have been instantiated.

\par\noindent\textbf{Hostile reconstruction.} A hostile referee will try to reconstruct this point in the strongest ordinary Hodge-theoretic language: rational Hodge classes are native cohomological objects, Chow groups and cycle-class maps are ordinary algebraic-geometric carriers, and an empty fiber would be ordinary mathematical falsification. The audit therefore never answers by declaring the underlying mathematical content illicit. It answers only by separating raw content, support status, theoremhood, endpoint force, and the fixed cycle-fiber endpoint image.

\par\noindent\textbf{Closure condition.} The claim survives this stress test only if it preserves the same $X$, $p$, and $\alpha$, the same rational Hodge-standing predicate, the same group $\CH^p(X)_\QQ$, the same map $\clQ$, and the same fiber $F_X(\alpha)$. A proposed escape that changes any of these is a domain or carrier shift. A proposed escape that preserves them while retaining endpoint-defeating force enters endpoint-use discipline and must pass through the native empty-fiber exhaustion.

\par\noindent\textbf{Audit result.} Under the fixed endpoint assumptions, the hostile reconstruction leaves only the declared branches: bridge completion, support-only content, bookkeeping, explicit domain shift, or independent same-domain cycle-fiber endpoint discrimination. This packet is therefore not decorative; it is a repeatable closure capsule for the theorem ladder and denial-burden table.


\subsection{Stress test: No silent redescription (report-preservation lens)}

\par\noindent\textbf{Load-bearing claim.} This stress test audits no replacement of Chow-fiber readout by another carrier under same label. The claim is not allowed to function by implication, atmosphere, or familiarity with AASC vocabulary. It must be available as a same-carrier theorem target inside the Hodge endpoint proof and must be reused only after the carrier, branch, and report slots have been instantiated.

\par\noindent\textbf{Hostile reconstruction.} A hostile referee will try to reconstruct this point in the strongest ordinary Hodge-theoretic language: rational Hodge classes are native cohomological objects, Chow groups and cycle-class maps are ordinary algebraic-geometric carriers, and an empty fiber would be ordinary mathematical falsification. The audit therefore never answers by declaring the underlying mathematical content illicit. It answers only by separating raw content, support status, theoremhood, endpoint force, and the fixed cycle-fiber endpoint image.

\par\noindent\textbf{Closure condition.} The claim survives this stress test only if it preserves the same $X$, $p$, and $\alpha$, the same rational Hodge-standing predicate, the same group $\CH^p(X)_\QQ$, the same map $\clQ$, and the same fiber $F_X(\alpha)$. A proposed escape that changes any of these is a domain or carrier shift. A proposed escape that preserves them while retaining endpoint-defeating force enters endpoint-use discipline and must pass through the native empty-fiber exhaustion.

\par\noindent\textbf{Audit result.} Under the fixed endpoint assumptions, the hostile reconstruction leaves only the declared branches: bridge completion, support-only content, bookkeeping, explicit domain shift, or independent same-domain cycle-fiber endpoint discrimination. This packet is therefore not decorative; it is a repeatable closure capsule for the theorem ladder and denial-burden table.


\subsection{Stress test: Non-factorizability (report-preservation lens)}

\par\noindent\textbf{Load-bearing claim.} This stress test audits local bridge fragments not composing into global readout without compatibility. The claim is not allowed to function by implication, atmosphere, or familiarity with AASC vocabulary. It must be available as a same-carrier theorem target inside the Hodge endpoint proof and must be reused only after the carrier, branch, and report slots have been instantiated.

\par\noindent\textbf{Hostile reconstruction.} A hostile referee will try to reconstruct this point in the strongest ordinary Hodge-theoretic language: rational Hodge classes are native cohomological objects, Chow groups and cycle-class maps are ordinary algebraic-geometric carriers, and an empty fiber would be ordinary mathematical falsification. The audit therefore never answers by declaring the underlying mathematical content illicit. It answers only by separating raw content, support status, theoremhood, endpoint force, and the fixed cycle-fiber endpoint image.

\par\noindent\textbf{Closure condition.} The claim survives this stress test only if it preserves the same $X$, $p$, and $\alpha$, the same rational Hodge-standing predicate, the same group $\CH^p(X)_\QQ$, the same map $\clQ$, and the same fiber $F_X(\alpha)$. A proposed escape that changes any of these is a domain or carrier shift. A proposed escape that preserves them while retaining endpoint-defeating force enters endpoint-use discipline and must pass through the native empty-fiber exhaustion.

\par\noindent\textbf{Audit result.} Under the fixed endpoint assumptions, the hostile reconstruction leaves only the declared branches: bridge completion, support-only content, bookkeeping, explicit domain shift, or independent same-domain cycle-fiber endpoint discrimination. This packet is therefore not decorative; it is a repeatable closure capsule for the theorem ladder and denial-burden table.


\subsection{Stress test: Role-occupancy closure (report-preservation lens)}

\par\noindent\textbf{Load-bearing claim.} This stress test audits nonempty fiber as role occupation rather than unique representative selection. The claim is not allowed to function by implication, atmosphere, or familiarity with AASC vocabulary. It must be available as a same-carrier theorem target inside the Hodge endpoint proof and must be reused only after the carrier, branch, and report slots have been instantiated.

\par\noindent\textbf{Hostile reconstruction.} A hostile referee will try to reconstruct this point in the strongest ordinary Hodge-theoretic language: rational Hodge classes are native cohomological objects, Chow groups and cycle-class maps are ordinary algebraic-geometric carriers, and an empty fiber would be ordinary mathematical falsification. The audit therefore never answers by declaring the underlying mathematical content illicit. It answers only by separating raw content, support status, theoremhood, endpoint force, and the fixed cycle-fiber endpoint image.

\par\noindent\textbf{Closure condition.} The claim survives this stress test only if it preserves the same $X$, $p$, and $\alpha$, the same rational Hodge-standing predicate, the same group $\CH^p(X)_\QQ$, the same map $\clQ$, and the same fiber $F_X(\alpha)$. A proposed escape that changes any of these is a domain or carrier shift. A proposed escape that preserves them while retaining endpoint-defeating force enters endpoint-use discipline and must pass through the native empty-fiber exhaustion.

\par\noindent\textbf{Audit result.} Under the fixed endpoint assumptions, the hostile reconstruction leaves only the declared branches: bridge completion, support-only content, bookkeeping, explicit domain shift, or independent same-domain cycle-fiber endpoint discrimination. This packet is therefore not decorative; it is a repeatable closure capsule for the theorem ladder and denial-burden table.


\subsection{Stress test: Cycle-fiber single interface (report-preservation lens)}

\par\noindent\textbf{Load-bearing claim.} This stress test audits the single endpoint-image interface for algebraic readout. The claim is not allowed to function by implication, atmosphere, or familiarity with AASC vocabulary. It must be available as a same-carrier theorem target inside the Hodge endpoint proof and must be reused only after the carrier, branch, and report slots have been instantiated.

\par\noindent\textbf{Hostile reconstruction.} A hostile referee will try to reconstruct this point in the strongest ordinary Hodge-theoretic language: rational Hodge classes are native cohomological objects, Chow groups and cycle-class maps are ordinary algebraic-geometric carriers, and an empty fiber would be ordinary mathematical falsification. The audit therefore never answers by declaring the underlying mathematical content illicit. It answers only by separating raw content, support status, theoremhood, endpoint force, and the fixed cycle-fiber endpoint image.

\par\noindent\textbf{Closure condition.} The claim survives this stress test only if it preserves the same $X$, $p$, and $\alpha$, the same rational Hodge-standing predicate, the same group $\CH^p(X)_\QQ$, the same map $\clQ$, and the same fiber $F_X(\alpha)$. A proposed escape that changes any of these is a domain or carrier shift. A proposed escape that preserves them while retaining endpoint-defeating force enters endpoint-use discipline and must pass through the native empty-fiber exhaustion.

\par\noindent\textbf{Audit result.} Under the fixed endpoint assumptions, the hostile reconstruction leaves only the declared branches: bridge completion, support-only content, bookkeeping, explicit domain shift, or independent same-domain cycle-fiber endpoint discrimination. This packet is therefore not decorative; it is a repeatable closure capsule for the theorem ladder and denial-burden table.


\subsection{Stress test: Coequal carrier distinction (report-preservation lens)}

\par\noindent\textbf{Load-bearing claim.} This stress test audits separating broad outcome carrier from cycle-fiber image carrier. The claim is not allowed to function by implication, atmosphere, or familiarity with AASC vocabulary. It must be available as a same-carrier theorem target inside the Hodge endpoint proof and must be reused only after the carrier, branch, and report slots have been instantiated.

\par\noindent\textbf{Hostile reconstruction.} A hostile referee will try to reconstruct this point in the strongest ordinary Hodge-theoretic language: rational Hodge classes are native cohomological objects, Chow groups and cycle-class maps are ordinary algebraic-geometric carriers, and an empty fiber would be ordinary mathematical falsification. The audit therefore never answers by declaring the underlying mathematical content illicit. It answers only by separating raw content, support status, theoremhood, endpoint force, and the fixed cycle-fiber endpoint image.

\par\noindent\textbf{Closure condition.} The claim survives this stress test only if it preserves the same $X$, $p$, and $\alpha$, the same rational Hodge-standing predicate, the same group $\CH^p(X)_\QQ$, the same map $\clQ$, and the same fiber $F_X(\alpha)$. A proposed escape that changes any of these is a domain or carrier shift. A proposed escape that preserves them while retaining endpoint-defeating force enters endpoint-use discipline and must pass through the native empty-fiber exhaustion.

\par\noindent\textbf{Audit result.} Under the fixed endpoint assumptions, the hostile reconstruction leaves only the declared branches: bridge completion, support-only content, bookkeeping, explicit domain shift, or independent same-domain cycle-fiber endpoint discrimination. This packet is therefore not decorative; it is a repeatable closure capsule for the theorem ladder and denial-burden table.


\subsection{Stress test: Native empty-fiber exhaustion (report-preservation lens)}

\par\noindent\textbf{Load-bearing claim.} This stress test audits bridge-neutralized, support-only, bookkeeping, domain shift, or discriminator. The claim is not allowed to function by implication, atmosphere, or familiarity with AASC vocabulary. It must be available as a same-carrier theorem target inside the Hodge endpoint proof and must be reused only after the carrier, branch, and report slots have been instantiated.

\par\noindent\textbf{Hostile reconstruction.} A hostile referee will try to reconstruct this point in the strongest ordinary Hodge-theoretic language: rational Hodge classes are native cohomological objects, Chow groups and cycle-class maps are ordinary algebraic-geometric carriers, and an empty fiber would be ordinary mathematical falsification. The audit therefore never answers by declaring the underlying mathematical content illicit. It answers only by separating raw content, support status, theoremhood, endpoint force, and the fixed cycle-fiber endpoint image.

\par\noindent\textbf{Closure condition.} The claim survives this stress test only if it preserves the same $X$, $p$, and $\alpha$, the same rational Hodge-standing predicate, the same group $\CH^p(X)_\QQ$, the same map $\clQ$, and the same fiber $F_X(\alpha)$. A proposed escape that changes any of these is a domain or carrier shift. A proposed escape that preserves them while retaining endpoint-defeating force enters endpoint-use discipline and must pass through the native empty-fiber exhaustion.

\par\noindent\textbf{Audit result.} Under the fixed endpoint assumptions, the hostile reconstruction leaves only the declared branches: bridge completion, support-only content, bookkeeping, explicit domain shift, or independent same-domain cycle-fiber endpoint discrimination. This packet is therefore not decorative; it is a repeatable closure capsule for the theorem ladder and denial-burden table.


\subsection{Stress test: Bridge-neutralized branch (report-preservation lens)}

\par\noindent\textbf{Load-bearing claim.} This stress test audits why a bridge-neutralized negative branch is not live endpoint counterforce. The claim is not allowed to function by implication, atmosphere, or familiarity with AASC vocabulary. It must be available as a same-carrier theorem target inside the Hodge endpoint proof and must be reused only after the carrier, branch, and report slots have been instantiated.

\par\noindent\textbf{Hostile reconstruction.} A hostile referee will try to reconstruct this point in the strongest ordinary Hodge-theoretic language: rational Hodge classes are native cohomological objects, Chow groups and cycle-class maps are ordinary algebraic-geometric carriers, and an empty fiber would be ordinary mathematical falsification. The audit therefore never answers by declaring the underlying mathematical content illicit. It answers only by separating raw content, support status, theoremhood, endpoint force, and the fixed cycle-fiber endpoint image.

\par\noindent\textbf{Closure condition.} The claim survives this stress test only if it preserves the same $X$, $p$, and $\alpha$, the same rational Hodge-standing predicate, the same group $\CH^p(X)_\QQ$, the same map $\clQ$, and the same fiber $F_X(\alpha)$. A proposed escape that changes any of these is a domain or carrier shift. A proposed escape that preserves them while retaining endpoint-defeating force enters endpoint-use discipline and must pass through the native empty-fiber exhaustion.

\par\noindent\textbf{Audit result.} Under the fixed endpoint assumptions, the hostile reconstruction leaves only the declared branches: bridge completion, support-only content, bookkeeping, explicit domain shift, or independent same-domain cycle-fiber endpoint discrimination. This packet is therefore not decorative; it is a repeatable closure capsule for the theorem ladder and denial-burden table.


\subsection{Stress test: Support-only branch (report-preservation lens)}

\par\noindent\textbf{Load-bearing claim.} This stress test audits why support status cannot defeat endpoint closure. The claim is not allowed to function by implication, atmosphere, or familiarity with AASC vocabulary. It must be available as a same-carrier theorem target inside the Hodge endpoint proof and must be reused only after the carrier, branch, and report slots have been instantiated.

\par\noindent\textbf{Hostile reconstruction.} A hostile referee will try to reconstruct this point in the strongest ordinary Hodge-theoretic language: rational Hodge classes are native cohomological objects, Chow groups and cycle-class maps are ordinary algebraic-geometric carriers, and an empty fiber would be ordinary mathematical falsification. The audit therefore never answers by declaring the underlying mathematical content illicit. It answers only by separating raw content, support status, theoremhood, endpoint force, and the fixed cycle-fiber endpoint image.

\par\noindent\textbf{Closure condition.} The claim survives this stress test only if it preserves the same $X$, $p$, and $\alpha$, the same rational Hodge-standing predicate, the same group $\CH^p(X)_\QQ$, the same map $\clQ$, and the same fiber $F_X(\alpha)$. A proposed escape that changes any of these is a domain or carrier shift. A proposed escape that preserves them while retaining endpoint-defeating force enters endpoint-use discipline and must pass through the native empty-fiber exhaustion.

\par\noindent\textbf{Audit result.} Under the fixed endpoint assumptions, the hostile reconstruction leaves only the declared branches: bridge completion, support-only content, bookkeeping, explicit domain shift, or independent same-domain cycle-fiber endpoint discrimination. This packet is therefore not decorative; it is a repeatable closure capsule for the theorem ladder and denial-burden table.


\subsection{Stress test: Bookkeeping branch (report-preservation lens)}

\par\noindent\textbf{Load-bearing claim.} This stress test audits why labels and ledgers do not determine endpoint status. The claim is not allowed to function by implication, atmosphere, or familiarity with AASC vocabulary. It must be available as a same-carrier theorem target inside the Hodge endpoint proof and must be reused only after the carrier, branch, and report slots have been instantiated.

\par\noindent\textbf{Hostile reconstruction.} A hostile referee will try to reconstruct this point in the strongest ordinary Hodge-theoretic language: rational Hodge classes are native cohomological objects, Chow groups and cycle-class maps are ordinary algebraic-geometric carriers, and an empty fiber would be ordinary mathematical falsification. The audit therefore never answers by declaring the underlying mathematical content illicit. It answers only by separating raw content, support status, theoremhood, endpoint force, and the fixed cycle-fiber endpoint image.

\par\noindent\textbf{Closure condition.} The claim survives this stress test only if it preserves the same $X$, $p$, and $\alpha$, the same rational Hodge-standing predicate, the same group $\CH^p(X)_\QQ$, the same map $\clQ$, and the same fiber $F_X(\alpha)$. A proposed escape that changes any of these is a domain or carrier shift. A proposed escape that preserves them while retaining endpoint-defeating force enters endpoint-use discipline and must pass through the native empty-fiber exhaustion.

\par\noindent\textbf{Audit result.} Under the fixed endpoint assumptions, the hostile reconstruction leaves only the declared branches: bridge completion, support-only content, bookkeeping, explicit domain shift, or independent same-domain cycle-fiber endpoint discrimination. This packet is therefore not decorative; it is a repeatable closure capsule for the theorem ladder and denial-burden table.


\subsection{Stress test: Domain-shift branch (report-preservation lens)}

\par\noindent\textbf{Load-bearing claim.} This stress test audits why changing X,p,alpha, coefficients, equivalence, or carrier exits the fixed endpoint. The claim is not allowed to function by implication, atmosphere, or familiarity with AASC vocabulary. It must be available as a same-carrier theorem target inside the Hodge endpoint proof and must be reused only after the carrier, branch, and report slots have been instantiated.

\par\noindent\textbf{Hostile reconstruction.} A hostile referee will try to reconstruct this point in the strongest ordinary Hodge-theoretic language: rational Hodge classes are native cohomological objects, Chow groups and cycle-class maps are ordinary algebraic-geometric carriers, and an empty fiber would be ordinary mathematical falsification. The audit therefore never answers by declaring the underlying mathematical content illicit. It answers only by separating raw content, support status, theoremhood, endpoint force, and the fixed cycle-fiber endpoint image.

\par\noindent\textbf{Closure condition.} The claim survives this stress test only if it preserves the same $X$, $p$, and $\alpha$, the same rational Hodge-standing predicate, the same group $\CH^p(X)_\QQ$, the same map $\clQ$, and the same fiber $F_X(\alpha)$. A proposed escape that changes any of these is a domain or carrier shift. A proposed escape that preserves them while retaining endpoint-defeating force enters endpoint-use discipline and must pass through the native empty-fiber exhaustion.

\par\noindent\textbf{Audit result.} Under the fixed endpoint assumptions, the hostile reconstruction leaves only the declared branches: bridge completion, support-only content, bookkeeping, explicit domain shift, or independent same-domain cycle-fiber endpoint discrimination. This packet is therefore not decorative; it is a repeatable closure capsule for the theorem ladder and denial-burden table.


\subsection{Stress test: Cycle-fiber discriminator (report-preservation lens)}

\par\noindent\textbf{Load-bearing claim.} This stress test audits the independent same-domain discriminator induced by endpoint-standing empty-fiber force. The claim is not allowed to function by implication, atmosphere, or familiarity with AASC vocabulary. It must be available as a same-carrier theorem target inside the Hodge endpoint proof and must be reused only after the carrier, branch, and report slots have been instantiated.

\par\noindent\textbf{Hostile reconstruction.} A hostile referee will try to reconstruct this point in the strongest ordinary Hodge-theoretic language: rational Hodge classes are native cohomological objects, Chow groups and cycle-class maps are ordinary algebraic-geometric carriers, and an empty fiber would be ordinary mathematical falsification. The audit therefore never answers by declaring the underlying mathematical content illicit. It answers only by separating raw content, support status, theoremhood, endpoint force, and the fixed cycle-fiber endpoint image.

\par\noindent\textbf{Closure condition.} The claim survives this stress test only if it preserves the same $X$, $p$, and $\alpha$, the same rational Hodge-standing predicate, the same group $\CH^p(X)_\QQ$, the same map $\clQ$, and the same fiber $F_X(\alpha)$. A proposed escape that changes any of these is a domain or carrier shift. A proposed escape that preserves them while retaining endpoint-defeating force enters endpoint-use discipline and must pass through the native empty-fiber exhaustion.

\par\noindent\textbf{Audit result.} Under the fixed endpoint assumptions, the hostile reconstruction leaves only the declared branches: bridge completion, support-only content, bookkeeping, explicit domain shift, or independent same-domain cycle-fiber endpoint discrimination. This packet is therefore not decorative; it is a repeatable closure capsule for the theorem ladder and denial-burden table.


\subsection{Stress test: No-independent closure (report-preservation lens)}

\par\noindent\textbf{Load-bearing claim.} This stress test audits the A+ consequence excluding the independent discriminator. The claim is not allowed to function by implication, atmosphere, or familiarity with AASC vocabulary. It must be available as a same-carrier theorem target inside the Hodge endpoint proof and must be reused only after the carrier, branch, and report slots have been instantiated.

\par\noindent\textbf{Hostile reconstruction.} A hostile referee will try to reconstruct this point in the strongest ordinary Hodge-theoretic language: rational Hodge classes are native cohomological objects, Chow groups and cycle-class maps are ordinary algebraic-geometric carriers, and an empty fiber would be ordinary mathematical falsification. The audit therefore never answers by declaring the underlying mathematical content illicit. It answers only by separating raw content, support status, theoremhood, endpoint force, and the fixed cycle-fiber endpoint image.

\par\noindent\textbf{Closure condition.} The claim survives this stress test only if it preserves the same $X$, $p$, and $\alpha$, the same rational Hodge-standing predicate, the same group $\CH^p(X)_\QQ$, the same map $\clQ$, and the same fiber $F_X(\alpha)$. A proposed escape that changes any of these is a domain or carrier shift. A proposed escape that preserves them while retaining endpoint-defeating force enters endpoint-use discipline and must pass through the native empty-fiber exhaustion.

\par\noindent\textbf{Audit result.} Under the fixed endpoint assumptions, the hostile reconstruction leaves only the declared branches: bridge completion, support-only content, bookkeeping, explicit domain shift, or independent same-domain cycle-fiber endpoint discrimination. This packet is therefore not decorative; it is a repeatable closure capsule for the theorem ladder and denial-burden table.


\subsection{Stress test: Local separator occupation (report-preservation lens)}

\par\noindent\textbf{Load-bearing claim.} This stress test audits the local empty-fiber countercase as endpoint-defeating non-occupation. The claim is not allowed to function by implication, atmosphere, or familiarity with AASC vocabulary. It must be available as a same-carrier theorem target inside the Hodge endpoint proof and must be reused only after the carrier, branch, and report slots have been instantiated.

\par\noindent\textbf{Hostile reconstruction.} A hostile referee will try to reconstruct this point in the strongest ordinary Hodge-theoretic language: rational Hodge classes are native cohomological objects, Chow groups and cycle-class maps are ordinary algebraic-geometric carriers, and an empty fiber would be ordinary mathematical falsification. The audit therefore never answers by declaring the underlying mathematical content illicit. It answers only by separating raw content, support status, theoremhood, endpoint force, and the fixed cycle-fiber endpoint image.

\par\noindent\textbf{Closure condition.} The claim survives this stress test only if it preserves the same $X$, $p$, and $\alpha$, the same rational Hodge-standing predicate, the same group $\CH^p(X)_\QQ$, the same map $\clQ$, and the same fiber $F_X(\alpha)$. A proposed escape that changes any of these is a domain or carrier shift. A proposed escape that preserves them while retaining endpoint-defeating force enters endpoint-use discipline and must pass through the native empty-fiber exhaustion.

\par\noindent\textbf{Audit result.} Under the fixed endpoint assumptions, the hostile reconstruction leaves only the declared branches: bridge completion, support-only content, bookkeeping, explicit domain shift, or independent same-domain cycle-fiber endpoint discrimination. This packet is therefore not decorative; it is a repeatable closure capsule for the theorem ladder and denial-burden table.


\subsection{Stress test: Reductio annotation (report-preservation lens)}

\par\noindent\textbf{Load-bearing claim.} This stress test audits proof-act annotation not object-level premise. The claim is not allowed to function by implication, atmosphere, or familiarity with AASC vocabulary. It must be available as a same-carrier theorem target inside the Hodge endpoint proof and must be reused only after the carrier, branch, and report slots have been instantiated.

\par\noindent\textbf{Hostile reconstruction.} A hostile referee will try to reconstruct this point in the strongest ordinary Hodge-theoretic language: rational Hodge classes are native cohomological objects, Chow groups and cycle-class maps are ordinary algebraic-geometric carriers, and an empty fiber would be ordinary mathematical falsification. The audit therefore never answers by declaring the underlying mathematical content illicit. It answers only by separating raw content, support status, theoremhood, endpoint force, and the fixed cycle-fiber endpoint image.

\par\noindent\textbf{Closure condition.} The claim survives this stress test only if it preserves the same $X$, $p$, and $\alpha$, the same rational Hodge-standing predicate, the same group $\CH^p(X)_\QQ$, the same map $\clQ$, and the same fiber $F_X(\alpha)$. A proposed escape that changes any of these is a domain or carrier shift. A proposed escape that preserves them while retaining endpoint-defeating force enters endpoint-use discipline and must pass through the native empty-fiber exhaustion.

\par\noindent\textbf{Audit result.} Under the fixed endpoint assumptions, the hostile reconstruction leaves only the declared branches: bridge completion, support-only content, bookkeeping, explicit domain shift, or independent same-domain cycle-fiber endpoint discrimination. This packet is therefore not decorative; it is a repeatable closure capsule for the theorem ladder and denial-burden table.


\subsection{Stress test: Annotation discharge (report-preservation lens)}

\par\noindent\textbf{Load-bearing claim.} This stress test audits the original empty-fiber assumption is discharged. The claim is not allowed to function by implication, atmosphere, or familiarity with AASC vocabulary. It must be available as a same-carrier theorem target inside the Hodge endpoint proof and must be reused only after the carrier, branch, and report slots have been instantiated.

\par\noindent\textbf{Hostile reconstruction.} A hostile referee will try to reconstruct this point in the strongest ordinary Hodge-theoretic language: rational Hodge classes are native cohomological objects, Chow groups and cycle-class maps are ordinary algebraic-geometric carriers, and an empty fiber would be ordinary mathematical falsification. The audit therefore never answers by declaring the underlying mathematical content illicit. It answers only by separating raw content, support status, theoremhood, endpoint force, and the fixed cycle-fiber endpoint image.

\par\noindent\textbf{Closure condition.} The claim survives this stress test only if it preserves the same $X$, $p$, and $\alpha$, the same rational Hodge-standing predicate, the same group $\CH^p(X)_\QQ$, the same map $\clQ$, and the same fiber $F_X(\alpha)$. A proposed escape that changes any of these is a domain or carrier shift. A proposed escape that preserves them while retaining endpoint-defeating force enters endpoint-use discipline and must pass through the native empty-fiber exhaustion.

\par\noindent\textbf{Audit result.} Under the fixed endpoint assumptions, the hostile reconstruction leaves only the declared branches: bridge completion, support-only content, bookkeeping, explicit domain shift, or independent same-domain cycle-fiber endpoint discrimination. This packet is therefore not decorative; it is a repeatable closure capsule for the theorem ladder and denial-burden table.


\subsection{Stress test: Ordinary theoremhood protection (report-preservation lens)}

\par\noindent\textbf{Load-bearing claim.} This stress test audits negative theoremhood is allowed until used as official endpoint resolution. The claim is not allowed to function by implication, atmosphere, or familiarity with AASC vocabulary. It must be available as a same-carrier theorem target inside the Hodge endpoint proof and must be reused only after the carrier, branch, and report slots have been instantiated.

\par\noindent\textbf{Hostile reconstruction.} A hostile referee will try to reconstruct this point in the strongest ordinary Hodge-theoretic language: rational Hodge classes are native cohomological objects, Chow groups and cycle-class maps are ordinary algebraic-geometric carriers, and an empty fiber would be ordinary mathematical falsification. The audit therefore never answers by declaring the underlying mathematical content illicit. It answers only by separating raw content, support status, theoremhood, endpoint force, and the fixed cycle-fiber endpoint image.

\par\noindent\textbf{Closure condition.} The claim survives this stress test only if it preserves the same $X$, $p$, and $\alpha$, the same rational Hodge-standing predicate, the same group $\CH^p(X)_\QQ$, the same map $\clQ$, and the same fiber $F_X(\alpha)$. A proposed escape that changes any of these is a domain or carrier shift. A proposed escape that preserves them while retaining endpoint-defeating force enters endpoint-use discipline and must pass through the native empty-fiber exhaustion.

\par\noindent\textbf{Audit result.} Under the fixed endpoint assumptions, the hostile reconstruction leaves only the declared branches: bridge completion, support-only content, bookkeeping, explicit domain shift, or independent same-domain cycle-fiber endpoint discrimination. This packet is therefore not decorative; it is a repeatable closure capsule for the theorem ladder and denial-burden table.


\subsection{Stress test: Incompleteness firewall (report-preservation lens)}

\par\noindent\textbf{Load-bearing claim.} This stress test audits same-domain independence matters only if endpoint-standing relevant. The claim is not allowed to function by implication, atmosphere, or familiarity with AASC vocabulary. It must be available as a same-carrier theorem target inside the Hodge endpoint proof and must be reused only after the carrier, branch, and report slots have been instantiated.

\par\noindent\textbf{Hostile reconstruction.} A hostile referee will try to reconstruct this point in the strongest ordinary Hodge-theoretic language: rational Hodge classes are native cohomological objects, Chow groups and cycle-class maps are ordinary algebraic-geometric carriers, and an empty fiber would be ordinary mathematical falsification. The audit therefore never answers by declaring the underlying mathematical content illicit. It answers only by separating raw content, support status, theoremhood, endpoint force, and the fixed cycle-fiber endpoint image.

\par\noindent\textbf{Closure condition.} The claim survives this stress test only if it preserves the same $X$, $p$, and $\alpha$, the same rational Hodge-standing predicate, the same group $\CH^p(X)_\QQ$, the same map $\clQ$, and the same fiber $F_X(\alpha)$. A proposed escape that changes any of these is a domain or carrier shift. A proposed escape that preserves them while retaining endpoint-defeating force enters endpoint-use discipline and must pass through the native empty-fiber exhaustion.

\par\noindent\textbf{Audit result.} Under the fixed endpoint assumptions, the hostile reconstruction leaves only the declared branches: bridge completion, support-only content, bookkeeping, explicit domain shift, or independent same-domain cycle-fiber endpoint discrimination. This packet is therefore not decorative; it is a repeatable closure capsule for the theorem ladder and denial-burden table.


\subsection{Stress test: Route-locus exhaustion (report-preservation lens)}

\par\noindent\textbf{Load-bearing claim.} This stress test audits status, codomain, existence, carrier, temporal, locality, selector, and second-classifier loci. The claim is not allowed to function by implication, atmosphere, or familiarity with AASC vocabulary. It must be available as a same-carrier theorem target inside the Hodge endpoint proof and must be reused only after the carrier, branch, and report slots have been instantiated.

\par\noindent\textbf{Hostile reconstruction.} A hostile referee will try to reconstruct this point in the strongest ordinary Hodge-theoretic language: rational Hodge classes are native cohomological objects, Chow groups and cycle-class maps are ordinary algebraic-geometric carriers, and an empty fiber would be ordinary mathematical falsification. The audit therefore never answers by declaring the underlying mathematical content illicit. It answers only by separating raw content, support status, theoremhood, endpoint force, and the fixed cycle-fiber endpoint image.

\par\noindent\textbf{Closure condition.} The claim survives this stress test only if it preserves the same $X$, $p$, and $\alpha$, the same rational Hodge-standing predicate, the same group $\CH^p(X)_\QQ$, the same map $\clQ$, and the same fiber $F_X(\alpha)$. A proposed escape that changes any of these is a domain or carrier shift. A proposed escape that preserves them while retaining endpoint-defeating force enters endpoint-use discipline and must pass through the native empty-fiber exhaustion.

\par\noindent\textbf{Audit result.} Under the fixed endpoint assumptions, the hostile reconstruction leaves only the declared branches: bridge completion, support-only content, bookkeeping, explicit domain shift, or independent same-domain cycle-fiber endpoint discrimination. This packet is therefore not decorative; it is a repeatable closure capsule for the theorem ladder and denial-burden table.


\subsection{Stress test: Weakening resistance (report-preservation lens)}

\par\noindent\textbf{Load-bearing claim.} This stress test audits weaker regimes preserve behavior, shift domain, lower status, or reintroduce forbidden governance. The claim is not allowed to function by implication, atmosphere, or familiarity with AASC vocabulary. It must be available as a same-carrier theorem target inside the Hodge endpoint proof and must be reused only after the carrier, branch, and report slots have been instantiated.

\par\noindent\textbf{Hostile reconstruction.} A hostile referee will try to reconstruct this point in the strongest ordinary Hodge-theoretic language: rational Hodge classes are native cohomological objects, Chow groups and cycle-class maps are ordinary algebraic-geometric carriers, and an empty fiber would be ordinary mathematical falsification. The audit therefore never answers by declaring the underlying mathematical content illicit. It answers only by separating raw content, support status, theoremhood, endpoint force, and the fixed cycle-fiber endpoint image.

\par\noindent\textbf{Closure condition.} The claim survives this stress test only if it preserves the same $X$, $p$, and $\alpha$, the same rational Hodge-standing predicate, the same group $\CH^p(X)_\QQ$, the same map $\clQ$, and the same fiber $F_X(\alpha)$. A proposed escape that changes any of these is a domain or carrier shift. A proposed escape that preserves them while retaining endpoint-defeating force enters endpoint-use discipline and must pass through the native empty-fiber exhaustion.

\par\noindent\textbf{Audit result.} Under the fixed endpoint assumptions, the hostile reconstruction leaves only the declared branches: bridge completion, support-only content, bookkeeping, explicit domain shift, or independent same-domain cycle-fiber endpoint discrimination. This packet is therefore not decorative; it is a repeatable closure capsule for the theorem ladder and denial-burden table.

\section{Conclusion}



The manuscript has one nonclassical dependency: official endpoint-defeating counterexample force is theorem-bearing endpoint use. Once that is accepted, the proof is target instantiation plus route exhaustion. The official Hodge endpoint is preserved: every rational Hodge class must have nonempty rational cycle-class fiber. Empty-fiber content is lawful as raw content, descriptive nonmembership, support, obstruction, or ordinary negative theoremhood. It becomes endpoint-governed only when used with official endpoint-defeating force on the same fixed carrier.

The live exact-complement empty-fiber branch cannot occupy an autonomous ordinary-negative endpoint role. Bridge-neutralized, support-only, bookkeeping, and domain-shift routes are unavailable in the local endpoint subproof. The remaining branch is $D_{\mathrm{fib}}(X,p,\alpha)$. Endpoint use supplies $N_{\mathrm{fib}}(X,p,\alpha)$ and hence $\neg D_{\mathrm{fib}}(X,p,\alpha)$. The contradiction discharges the local assumption $F_X(\alpha)=\varnothing$. Thus every rational Hodge class has nonempty rational cycle-class fiber:
\[
\forall X,p,\alpha\in\Hdg^p(X),\quad F_X(\alpha)\neq\varnothing.
\]
Equivalently, every rational Hodge class is a rational linear combination of classes of algebraic cycles. This is the Hodge Conjecture.

\section{References}



\begin{thebibliography}{99}
\bibitem{ClayHodge} Clay Mathematics Institute, \emph{Hodge Conjecture}, official Millennium Problem statement.
\bibitem{Voisin1} Claire Voisin, \emph{Hodge Theory and Complex Algebraic Geometry I}, Cambridge University Press, 2002.
\bibitem{Voisin2} Claire Voisin, \emph{Hodge Theory and Complex Algebraic Geometry II}, Cambridge University Press, 2003.
\bibitem{GH} Phillip Griffiths and Joseph Harris, \emph{Principles of Algebraic Geometry}, Wiley, 1978.
\bibitem{Fulton} William Fulton, \emph{Intersection Theory}, second edition, Springer, 1998.
\bibitem{Lewis} James D. Lewis, \emph{A Survey of the Hodge Conjecture}, American Mathematical Society, 1999.
\bibitem{CDK} Eduardo Cattani, Pierre Deligne, and Aroldo Kaplan, \emph{On the locus of Hodge classes}, Journal of the American Mathematical Society 8 (1995), 483--506.
\bibitem{Andre} Yves Andre, \emph{Une introduction aux motifs}, Societe Mathematique de France, 2004.
\bibitem{Deligne} Pierre Deligne, \emph{Hodge cycles on abelian varieties}, Lecture Notes in Mathematics 900, Springer, 1982.
\bibitem{Kleiman} Steven L. Kleiman, \emph{The standard conjectures}, Proceedings of Symposia in Pure Mathematics 55, American Mathematical Society, 1994.
\bibitem{AASC-Kernel} Amos Jay Maley, \emph{Non-Degenerate Construction and the Kernel of Admissibility}, 2026.
\bibitem{AASC-UEAP} Amos Jay Maley, \emph{Claim Standing and Legitimacy under UEAP}, 2026.
\bibitem{AASC-ATS} Amos Jay Maley, \emph{Anchor, Tensor, and Skin}, 2026.
\bibitem{AASC-AMetric} Amos Jay Maley, \emph{The AMetric Boundary}, 2026.
\bibitem{AASC-Role} Amos Jay Maley, \emph{Role-Occupancy Closure under Admissibility}, 2026.
\bibitem{AASC-Godel} Amos Jay Maley, \emph{Admissible Construction and the Non-Instantiability of Godel Architectures}, 2026.
\end{thebibliography}

\end{document}
